% 本文件由 LaTeX 排版 Agent 根据有效 translation.json 自动生成。
% author=Augustine of Hippo; work_id=1301; title=论天主圣三

\chapter{论天主圣三}

% structure: p0001 source=h1 target=omitted reason=page-title-already-rendered-by-parent

\Emph{第一卷}\newline{}\Emph{第二卷}\newline{}\Emph{第三卷}\newline{}\Emph{第四卷}\newline{}\Emph{第五卷}\newline{}\Emph{第六卷}\newline{}\Emph{第七卷}\newline{}\Emph{第八卷}\newline{}\Emph{第九卷}\newline{}\Emph{第十卷}\newline{}\Emph{第十一卷}\newline{}\Emph{第十二卷}\newline{}\Emph{第十三卷}\newline{}\Emph{第十四卷}\newline{}\Emph{第十五卷}\par

\SourceRecord{Translated by Arthur West Haddan.；Nicene and Post-Nicene Fathers, First Series；Vol. 3.；Edited by Philip Schaff.；Buffalo, NY: Christian Literature Publishing Co.,；1887.；<http://www.newadvent.org/fathers/1301.htm>.}

% structure: p0001 source=h1 target=section reason=page-title

\section{论天主圣三（卷一）}

% structure: p0002 source=h2 target=subsection reason=relative-heading-rank; base=h2

\subsection{绪论}

在\Quote{\WorkTitle{订正录}}（ii. 15）中，\Person{奥斯定}论及本书如下：—\par

\Quote{我用了几年时间写了关于天主圣三（即天主）的十五卷书。然而，当我尚未完成第十三卷时，一些极为渴望得到这部著作的人等得超过了他们的忍耐限度，此书便以不够精确的状态被从我这里偷走，其精确程度既不及它可能达到的，也不及我本来打算出版时的样子。我一旦发现此事，由于另有副本，最初决定不亲自出版，而是在其他著作中提及此事；但在一些我无法拒绝的弟兄们的迫切请求下，我尽我所能地对其进行了修正，并完成和出版，并在开头附加了一封我写给可敬的迦太基主教\Person{奥勒留}的信，在其中以序言的方式陈述了发生的事，我原本打算自己做什么，以及弟兄们的爱迫使我做了什么。}\par

他所提及的信函如下：—\par

致至为有福的主，他以最诚挚的爱所敬爱的那一位，致他的圣洁兄弟兼同司铎，\Person{奥勒留}主教，\Person{奥斯定}在主内问安。\par

\Quote{我很年轻时开始，在老年时出版了一些关于天主圣三（即至高而真实的天主）的书。事实上，我之所以放下这项工作，是因为我发现它在我完成之前，甚至在我按原意修订和润色之前，就被过早地、或者更确切地说，被暗中偷走。因为我原不打算逐卷出版，而是一起出版，因为探究的进展使我后来加上的卷次与前文相连。因此，当这些人妨碍了我计划的完成（因为他们中有些人在我打算之前就获得了手稿），我便放弃了口授，打算在我可能写的其他著作中公开抱怨，以便任何人都知道这些书并非由我亲自出版，而是在我认为尚未适合出版之前就被从我手中夺走。然而，由于许多弟兄的热切要求，尤其是您的命令，我费尽心力，靠着天主的助佑完成了这部艰巨的著作；并经过修正（不是如我所愿，而是尽我所能，以免这些书与那些暗中流入人们手中的书相差太大），我已将之托付给我至亲爱的儿子和同执事，呈交给您的主教尊驾，并允许任何人乐意听、抄写和阅读。毫无疑问，如果我能实现我最初的意图，尽管它们会包含同样的思想，但它们会得到更彻底、更清晰的阐述，只要如此深奥的主题的难度和我自己的能力允许。然而，有些人拥有最初四卷或五卷（没有序言）以及第十二卷（缺少其后部分的大量章节）。但如果他们愿意且能够，通过熟悉本版，他们将修正全书。无论如何，我恳求您命令将这封信单独附加在书卷的开头。再见。请为我祈祷。}\par

% structure: p0008 source=h2 target=subsection reason=relative-heading-rank; base=h2

\subsection{第一章——本书是为驳斥那些滥用理性、以诡辩攻击天主圣三信德的人而写。那些关于天主辩论的人，从三重原因上犯错。圣经除去虚假，逐步引导我们达于神圣之事。何谓真正的不朽。我们由信德滋养，以便能够领会神圣之事。}

1. 以下关于天主圣三的论文，正如读者应当知道的，是为了防范那些不屑于从信德开始、又被粗野乖僻的理性之爱所欺骗之人的诡辩而写。这类人中有一种试图将他们从身体感官的经验、或从自然的人类才智和勤勉的敏锐、或借助技艺从有形之物所形成的观念，转移到无形和属神的事物中；以便用前者来度量并构想后者。另一些人则根据人类心灵的本性或情感来形成他们关于天主的任何情感；并通过这一错误，在辩论天主时，用扭曲且虚妄的规则来支配他们的言论。还有第三种人，他们确实努力超越整个受造界——它无疑是可变的——以便将他们的思想提升到不变的本体，即天主；但因必死性的重负所压，当他们既想假装知道他们所不知道的、又不能知道他们所想知道的时，就因过份肯定自己武断的判断而将自己关在理解的正途之外；宁愿不纠正自己乖僻的意见，也不愿改变他们曾经捍卫的东西。事实上，这是我所提到的三种人共通的病症——即\Emph{viz.}，那些依据有形之物来构想天主的人，以及那些依据属神的受造物（如灵魂）来构想的人，以及那些既不参照身体也不参照属神受造物、却仍对天主作出错误判断的人；而且他们离真理更远，以至于无论是在身体中、在被造或被造的精神体中、还是在造物主本身中，都找不到与他们的观念相符之物。例如，认为天主是白色或红色的人，虽在错误中；然而这些颜色可见于身体。同样，认为天主时而遗忘时而记起，或有任何类似情况的人，也并未少犯错误；然而这些情况可见于心灵。但认为天主具有如此大能以至于自生的人，则错误更大，因为不仅天主并非如此存在，而且连精神体和身体受造物也并非如此；因为没有任何东西可以自生其存在。\par

2. 因此，为使人的心灵摆脱这类虚妄，圣经俯就婴孩，从一切真实存在的事物中汲取词语，借此如同滋养一般，使我们的理解力逐渐上升至神圣和超越之物。因为，在谈论天主时，圣经既使用了取自形体的词语，例如说：\BibleQuote{请将我隐藏在你的翅膀荫下}；又从精神受造物借用许多词语，以表达那并非如此、却必须如此言说的事：例如，\BibleQuote{我，上主，你的天主，是忌邪的天主}；以及\BibleQuote{我后悔造了人}。但圣经从未从根本不存的事物中抽取任何词语，用以构成修辞或谜语。因此，那些被第三种错误隔绝于真理之外的人，比同类更加有害而虚妄；因为他们对天主的揣测，既不能在天主自身内找到，也不能在任何受造物中找到。因为圣经惯于从受造物中寻得的事物，如同为儿童设下诱饵；以便按他们的尺度，仿佛逐步地，使软弱者的情感被激发去寻求上面的事，而离开下面的事。但同一圣经却很少使用那些专指天主、不在任何受造物中的词语；例如，对梅瑟所说的：\BibleQuote{我是自有者}；以及\BibleQuote{那自有者派遣我到你们这里来}。因为既然身体和灵魂在某种意义上也被称为\Emph{是}，圣经当然不会这样表达，除非它有意在某种特殊含义下被理解。同样，宗徒所说的：\BibleQuote{唯独祂拥有不朽}。既然灵魂在某种方式上也被称为不朽且确实是不朽的，圣经不会说\BibleQuote{唯独拥有}，除非因为真正的不朽就是不变性；这是任何受造物都不能拥有的，因为它唯独属于造物主。同样，雅各伯说：\BibleQuote{一切美好的赠与和完美的恩赐，都是从上头来的，从光明之父那里降下来的，在他那里没有变化或转动的影儿}。同样，达味说：\BibleQuote{你要改变它们，它们就被改变；但你是同一的}。\par

3. 此外，默想并完全认识天主的本体是困难的；祂塑造可变化之物，而自身毫无变化；创造暂时之物，而自身毫无时间性的变动。因此，我们必须净化我们的心灵，以便能够不可言喻地看见那不可言喻者；既然尚未达到，我们应由信德滋养，被引导到更适合我们能力的途径，好使我们变得适宜并能理解它。因此，宗徒说：\BibleQuote{在基督内确实蕴藏着一切智慧和知识的宝藏}；却又向我们推荐祂，如同向在基督内的婴孩，这些人虽然已由祂的恩宠重生，却仍是属血肉和属魂的，并非凭那使祂与圣父平等的天主性德能，而是凭祂被钉十字架的人性软弱。因为他（宗徒）说：\BibleQuote{我曾决定，在你们中间什么也不知道，只知道耶稣基督，并被钉在十字架上的耶稣基督}；接着他继续说：\BibleQuote{我在你们那里，又软弱，又恐惧，又战兢}。稍后他对他们说：\BibleQuote{弟兄们，我不能对你们说话，如同对属神的人，只能如同对属血肉的人，如同对基督内的婴孩。我用奶喂养你们，没有用饭：因为那时你们还不能接受，就是现在也不能}。有些人恼恨这类言语，认为这是对他们的轻视，而且大多宁愿相信那些对他们这样说的人无话可说，也不愿承认自己不能理解他们所说的话。有时，我们确实向他们陈述——当然不是他们请求中关于天主的解释，因为他们自己不能接受，或许我们也无法领会或表达——而是向他们指出，他们多么无能，完全不适于理解他们向我们要求的东西。但他们呢，因为听不见自己所渴望的，便以为我们或是为了掩饰自己的无知而欺骗他们，或是因恶意而说话，因为不愿他们得到知识；于是愤愤不平、心烦意乱地离去。\par

% structure: p0012 source=h2 target=subsection reason=relative-heading-rank; base=h2

\subsection{第二章——本书将如何谈论天主圣三}

4. 因此，依靠我们的主天主的助佑，我们将尽力做出他们如此迫切要求的说明：即天主圣三是唯一、独一、真实的天主，以及圣父、圣子、圣神如何被正确地言说、相信、理解为本体或本质相同的一体。如此，他们不会自以为被我方的推诿所戏弄，而是通过实际体验发现：至善正是被最纯净的心智所辨识的，正因为此，他们自己无法辨识或理解它，因为人类心智的眼睛是软弱的，在如此超然的光照中目眩，除非它被信德之义的滋养所加强。然而，首先我们必须根据圣经的权威来证明信德是否如此。然后，如果天主愿意并帮助我们，我们或许至少能为这些多嘴的争论者服务——他们自大胜过能力，因而受着更危险的疾病之苦——使他们能够找到一些他们无法怀疑的东西，这样，在那些他们找不到类似之物的情况下，他们会被引导将过错归于自己的心智，而非归于真理本身或我们的推理；并且，如果他们心中尚有对天主的爱或畏惧，他们就会回头，并按照适当的次序从信德开始：最终认识到在圣教会中为信友预备了何等有益的良药，借此，谨慎的虔诚医治了心智的软弱，使其能够领悟不变的真理，并防止它因无秩序的鲁莽而一头栽入有害的虚假意见。我自己也不会在任何存疑之处退缩而不探究；也不会在任何错误之处耻于学习。\par

% structure: p0014 source=h2 target=subsection reason=relative-heading-rank; base=h2

\subsection{第三章 \Person{奥古斯丁}向读者提出的请求。理解力迟钝的读者的错误不应归咎于作者。}

5. 此外，让我请求我的读者：无论何处，凡他与我同样确定的地方，就与我一同前进；凡他与我同样犹豫的地方，就与我一同探究；凡他自认犯错的地方，就转向我；凡他认出我犯错的地方，就唤我回头：好使我们一同踏上爱德之路，并向那被说\Quote{你们要寻求祂的面容，永不休止}者迈进。我在我们的主天主面前，与所有阅读我著作的人订立这个虔诚而安全的盟约，其他一切情况都如此，尤其是那些探讨圣父、圣子和圣神之圣三合一的作品；因为没有任何其他主题的错误更危险、探究更艰辛或发现真理更有益。如果任何读者说：\InnerQuote{这说得不好，因为我不明白。}这样的人指责的是我的语言，而非我的信德：也许事实上本可以表达得更清楚；但没有人能说话让所有人在所有事上都明白。因此，若有人对我的论述提出这样的指责，让他看看，当他不理解我时，是否能理解其他处理过类似主题和问题的人：如果能，就让他放下我的书，甚至，如果他愿意，把它扔掉；并让他将精力和时间花在他能理解的那些人身上。然而，他不应当因此认为我本该保持沉默，因为我未能像他所理解的那些人那样优美清楚地向他表达自己。因为并非所有人写的一切东西都落入所有人的手中。而且，可能有些能够理解我著作的人，反而找不到那些更清晰的作品，只遇到了我的作品。因此，许多人就同一问题写许多书，风格不同但信德相同，是有益的，以便该事理本身能以不同方式触及最多的人——一些人经由这种方式，另一些人经由另一种方式。但如果那个抱怨他不理解这些事的人，从未能理解任何关于此类主题的仔细而精确的推理，那么在这种情况下，让他通过决心和研习来对付自己，以便更好地明白；而不是通过争吵和争执来对付我，让我闭口。再者，当他读我的书时说：\InnerQuote{我当然明白所说的内容，但它不是真的。}如果他愿意，让他陈述自己的意见，并在可能时驳斥我的意见。如果他以爱德和真理这样做，并费心让我知道（如果我还活着），那么我将从我这劳苦中得到最丰盛的果实。如果他不能通知我本人，我非常愿意并高兴他能通知那些他能通知的人。至于我，我\Quote{默思主的律法，}纵使不是\InnerQuote{昼夜，}至少在我能做的短暂时间里；我将我的默想记下来，免得因遗忘而丢失；希望因天主的仁慈，祂会使我坚定地持守所有我确定无疑的真理；\Quote{但我若在什么事上心存异见，愿祂亲自向我启示这一点，}无论是通过隐秘的启示和劝诫，还是通过祂明确的言语，或通过我弟兄们的推理。我祈求这一点，并将我的信赖和渴望交托给祂，祂足以保守祂所赐予我的一切，并实现祂所应许的一切。\par

6. 诚然，我预料有些领悟较为迟钝的人，会以为我在本书的某些部分持守了我实际并未持守的见解，或者未曾持守我实际持守的见解。但他们的错误——无人可不知——不应归于我，倘若他们因追随我的脚步却不理解我而误入虚假教义；同时，我被迫在艰涩隐晦的论题中摸索道路。因为任何人都不能以任何方式正当地将异端者的众多不同错误归于圣经的神圣见证本身；尽管他们全都试图从同一圣经中为自己的虚假错误意见辩护。基督的律法，即爱德，清楚地训诫我，并以甘甜的约束命令我：当人们以为我在书中持守了某些虚假的见解（实际我并未持守），而同一虚假令一人不悦、令另一人喜悦时，我宁愿被那指责虚假的人责备，也不愿被那称赞虚假的人称赞。因为虽然我（从未持守那错误的）被前者责备并不正确，但错误本身却受到正确的谴责；而后者称赞我（以为我持守了真理所谴责的）也不正确，且那见解本身（真理也谴责它）也是如此。因此，让我们在主的名下，开始我们已经着手的工作。\par

% structure: p0017 source=h2 target=subsection reason=relative-heading-rank; base=h2

\subsection{第四章——论天主圣三的公教信仰教义是什么。}

7. 所有那些我能阅读的、在我之前论述天主圣三（祂就是天主）的圣经（包括旧约和新约）的公教诠释者，都立意按照圣经教导这一教义：圣父、圣子、圣神在不可分的同等中显示同一实体之神性合一；因此他们不是三位神，而是一位天主：虽然圣父生了圣子，所以为父者并非圣子；圣子由圣父所生，所以为子者并非圣父；圣神既非圣父也非圣子，而只是圣父和圣子的神，祂自己与圣父和圣子同等，并属于天主圣三的合一。然而，并非这位天主圣三由童贞玛利亚所生，在般雀比拉多治下被钉十字架，被埋葬，第三日复活，升天，而只是圣子。同样，也不是这位天主圣三在耶稣受洗时以鸽子的形状降临；也不是在五旬节主升天之后，当\BibleQuote{从天上有响声下来，好像一阵大风吹过}，同一圣三\BibleQuote{有舌头如火焰显现出来，分开落在他们各人头上}，而只是圣神。也不是这位圣三从天上传话：\BibleQuote{你是我的儿子}，无论是当祂受若翰洗礼时，或当三位门徒与祂在山上的时候，或当有声音说：\BibleQuote{我已经荣耀了它，还要再荣耀它}；而那只是圣父对圣子所说的话；虽然圣父、圣子、圣神，正如他们不可分，也以不可分的方式工作。这也是我的信德，因为这是公教信德。\par

% structure: p0019 source=h2 target=subsection reason=relative-heading-rank; base=h2

\subsection{第五章——论关于天主圣三的难题：三者如何是一位天主，以及虽不可分地工作，却如何分别行某些事。}

8. 然而，有些人在这种信德中遇到困难：他们听说圣父是天主，圣子是天主，圣神是天主，而这一圣三不是三位神，而是一位天主；他们询问如何理解这一点，尤其当说圣三在一切天主所做的工中不可分地工作，然而某一话语是圣父说的，而非圣子的声音；并且除了圣子以外，没有别者降生为人、受苦、复活、升天；并且除了圣神以外，没有别者以鸽子形状降临。他们想要理解：圣三如何发出了那只是圣父的话语；同一圣三如何创造了那肉身，其中只有圣子由童贞女所生；同一圣三本身如何造作了那鸽子的形状，其中只有圣神显现。否则，圣三就不是不可分地工作，而是圣父做某些事，圣子做另一些事，圣神又做别的事；或者，如果他们一同做某些事，又分别做另一些事，那么圣三就不是不可分的。另一个困难是：圣神在圣三中是怎样的情况？祂既非圣父所生，也非圣子所生，也非两者所生，虽然祂是圣父和圣子的神。既然人们以这样的问题使我们疲倦，就让我们尽己所能，将天主之恩赐予我们软弱者的任何智慧，就此主题向他们展开；也不可\BibleQuote{怀着嫉妒的心前行}。若我们说我们不习惯思考这类事情，那并非实话；但若我们承认这类主题常萦绕我们思绪（因我们被探索真理的爱所牵引），那么他们就以爱德之律法要求我们将我们在此所能发现的告诉他们。\BibleQuote{这不是说我已经达到，或已经成为成全的人}（因为，如果宗徒\Person{保禄}尚且如此，何况我远在他脚下，岂敢自以为已经得着了！）但按我的尺度，\BibleQuote{我忘记背后，努力面前，向着标竿直跑，要得天主在基督耶稣里从天上召我来得的奖赏}，我请求揭示我已走过的路程，以及我所到达的点，从那里还有路程引我至终点。那请求我的人，正是慷慨的爱德促使我为他们服务。也必须如此，天主要恩赐我在供他们阅读的材料中，自己也受益；在寻求答复他们疑问的过程中，我自己也会找到我所寻索的。因此，我承担此任务，奉我主天主的命令和帮助，与其说是以权威论述我已知道的事，不如说是以虔诚的论述学习那些事。\par

% structure: p0021 source=h2 target=subsection reason=relative-heading-rank; base=h2

\subsection{第六章　论圣子是真天主，与圣父同一性体。不仅圣父，而且天主圣三被宣告为不朽的。万有并非仅来自圣父，也来自圣子。论圣神是真天主，与圣父和圣子同等。}

9. 凡是说我们的主耶稣基督不是天主，或不是真天主，或不是与圣父同是独一天主，或因其可变而非真正不朽的，都被最明确而一致的神圣见证之声驳倒了。例如：\BibleQuote{在起初已有圣言，圣言与天主同在，圣言就是天主。}显然，我们应将天主的圣言理解为天主的独生子，论及他，后来又说：\BibleQuote{圣言成了血肉，居住在我们中间，}这是指他降生成人的诞生，是在童贞女内按时成就的。但这里不仅宣告他是天主，而且他与圣父同一性体；因为说了\BibleQuote{圣言就是天主}之后，又说：\BibleQuote{这位圣言在起初就与天主同在：万物是藉着他而造成的；凡受造的，没有一样不是由他而造成的。}不仅仅是\Quote{万物}，而是所有\Emph{受造的}万物，即整个受造界。由此清楚可见，那藉以造成万物的他本身不是受造的。如果他不是受造的，那么他就不是受造物；如果他不是受造物，那么他就与圣父同一性体。因为凡不是天主的性体就是受造物；凡不是受造物的就是天主。如果圣子不与圣父同一性体，那么他就是受造的性体；如果他是受造的性体，那么万物就不是藉他造成的；但\BibleQuote{万物是藉着他造成的，}因此他与圣父是同一性体。所以他不仅是天主，而且是真天主。同一位\Person{若望}在他的书信中最明确地肯定了这一点：\BibleQuote{我们知道圣子来了，并且赐给了我们理智，叫我们认识那真实者；我们确实是在那真实者内，即在他的圣子耶稣基督内，他是真天主和永生。}\par

10. 因此也必然得出，宗徒\Person{保禄}说\BibleQuote{惟独祂具有不死不灭}，并不是仅仅指圣父，而是指独一的天主，即天主圣三本身。因为那本身就是永生的，按其任何可变性都不是可朽的；因此圣子，因为\BibleQuote{他就是永生}，也被理解为与圣父同在，经文说\BibleQuote{惟独祂具有不死不灭}时也是如此。因为我们也被分施这永生，并在自己的限度内成为不朽的。但我们所分施的永生本身是一回事；我们这些因分施它而将永远生活的人是另一回事。因为如果他说：\BibleQuote{到了适当的时期，天父要显示他，那真福者和独有权能者，万王之王，万主之主；惟独祂具有不死不灭}；即使如此，也未必理解为圣子被排除在外。因为圣子也没有将圣父与自己分开，因为他自己，以智慧的声音（他就是天主的智慧）在另一处说：\BibleQuote{我独自绕行天穹。}因此，更不必将\BibleQuote{惟独具有不死不灭}理解为单指圣父而排除圣子，因为经文是这样说的：\BibleQuote{你要遵守这诫命，无玷污、无可指摘，直到我们的主耶稣基督的显现。在预定的时期，那真福者和独有权能者，万王之王，万主之主要显现他；惟独他具有不死不灭，住于不可接近的光中，没有人见过他，也不能看见他。愿尊崇和永远的威能归于他！阿们。}在这些话中，既没有特别提到圣父，也没有提到圣子或圣神；而是提到真福者和独有权能者，万王之王，万主之主；即独一、唯一、真实的天主，即天主圣三本身。\par

11. 但可能接下来的话会干扰这含义，因为经上说\BibleQuote{没有人见过他，也不能看见他}；虽然这也可以指基督按其神性而言，犹太人从未见过他的神性，他们却在肉身上看见并钉死了他；而他的神性无论如何不能为肉眼所见，只能为那种眼光所见，具有那种眼光的人不再是凡人，而是超越凡人。因此，正确地理解，天主圣三本身就是\BibleQuote{真福者和独有权能者}，他\BibleQuote{将按自己的时期显示我们的主耶稣基督的来临}。因为说话的方式是：\BibleQuote{惟独他具有不死不灭}和\BibleQuote{惟独他行奇事}。我很想知道他们认为这些话是指谁说的。如果仅指圣父，那么圣子亲自说的话如何正确：\BibleQuote{凡父所做的，子也照样做}？在奇妙作为中，还有比复活并赐生命给死者更奇妙的吗？然而同一位圣子说：\BibleQuote{父怎样复活死者并赐给他们生命，子也照样随自己的意愿赐生命。}那么，圣父怎能\BibleQuote{惟独行奇事}呢？这些话让我们理解，既不是仅指圣父，也不是仅指圣子，而确实是指独一真实的天主，即圣父、圣子、圣神。\par

12. 此外，当同一宗徒说：\BibleQuote{但我们只有一位天主，就是圣父，万物都出于祂，我们也归于祂；并只有一位主耶稣基督，万物都藉着祂，我们也藉着祂，}谁能怀疑他是指一切受造之物呢？正如\Person{若望}所说：\BibleQuote{万物是藉着祂造成的。}因此我问，他在另一处说的是谁？\BibleQuote{因为万物都出于祂，藉着祂，并在祂内：愿光荣归于祂，直到永远。阿们。}因为如果这是指圣父、圣子和圣神，将每个短语分别归给每一位：出于祂，即出于圣父；藉着祂，即藉着圣子；在祂内，即在圣神内——那么显然，圣父、圣子和圣神是一位天主，因为语句持续使用单数：\BibleQuote{愿光荣归于祂，直到永远。}因为在经文的开头，他没有说：\BibleQuote{啊，智慧和知识的富藏多么深奥}，指的是圣父、圣子或圣神，而是说：\BibleQuote{天主的智慧和知识！}\BibleQuote{祂的判断何其难测，祂的道路何其难寻！谁曾知道主的心意？谁曾作过祂的参谋？谁先给了祂，使得祂偿还呢？因为万物都出于祂，藉着祂，并在祂内：愿光荣归于祂，直到永远。阿们。}但如果他们坚持认为这只是指圣父，那么万物如何出于圣父，如这里所说；万物又如何藉着圣子，如给格林多人所说的\BibleQuote{并只有一位主耶稣基督，万物藉着祂}，以及\BibleBook{若望}{福音}所说的\BibleQuote{万物是藉着祂造成的？}因为如果有些事物是圣父造的，有些是圣子造的，那么万物就不是都出于圣父，也不是都出于圣子；但如果万物都是圣父造的，万物也都是圣子造的，那么相同的事物就是由圣父和圣子造的。因此，圣子与圣父平等，圣父与圣子的工作是不可分割的。因为如果圣父创造了圣子（圣子当然没有创造自己），那么万物就不是都藉着圣子；但万物都是藉着圣子，所以祂自己并不是受造的，为能与圣父一同创造一切受造之物。宗徒也没有避讳使用这个词本身，而是最明确地说：\BibleQuote{祂虽具有天主的形体，却没有将与自己同天主平等看作强取之事}，这里特别用\Quote{天主}一词指圣父；如同另一处所说：\BibleQuote{但基督的元首是天主。}\par

13. 关于圣神，也收集了类似的证据，在我们之前讨论过这个主题的人已经充分利用了这些证据，表明祂也是天主，而不是受造物。但如果祂不是受造物，那么祂不仅是神（因为人也被称为神），而且是真天主；因此与圣父和圣子绝对平等，并在天主圣三的合一中，与圣父圣子同体、同永恒。但圣神不是受造物，这一点在上述经文中尤为明显，那里命令我们不要事奉受造物，而要事奉造物主；这并非指我们因爱而彼此\Quote{事奉}的那种含义，在希腊文是\OriginalTerm{服务（奴仆式）}{\GreekText{δουλεύειν}}，而是指唯独事奉天主的那种含义，在希腊文是\OriginalTerm{敬拜（钦崇式）}{\GreekText{λατρεύειν}}。因此，那些将应归于天主的敬拜献给偶像的人被称为偶像崇拜者。因为正是这种敬拜，经上说：\BibleQuote{你要朝拜主你的天主，唯独事奉祂。}这在希腊文圣经中更为明确，写作\OriginalTerm{你要事奉（钦崇式）}{\GreekText{λατρεύσεις}}。既然我们被禁止用这种敬拜事奉受造物，如经上记载：\BibleQuote{你要朝拜主你的天主，唯独事奉祂}（因此，宗徒也弃绝那些敬拜事奉受造物甚于造物主的人），那么圣神肯定不是受造物，因为众圣徒都将这种敬拜献给祂；正如宗徒所说：\BibleQuote{我们才是受割损的，我们以圣神事奉天主}，在希腊文是\OriginalTerm{事奉（钦崇式）}{\GreekText{λατρεύοντες}}。因为甚至大多数拉丁文抄本也是这样：\BibleQuote{我们事奉天主圣神}；但所有的希腊文抄本，或几乎所有的，都是这样。不过在有些拉丁文抄本中，我们发现不是\BibleQuote{我们朝拜天主圣神}，而是\BibleQuote{我们在圣神内朝拜天主}。但那些在这方面犯错，且不肯屈服于更权威之证据的人，请告诉他们，他们是否发现这经文在抄本中也有异文：\BibleQuote{你们不知道你们的身体是圣神的宫殿，这圣神住在你们内，你们是从天主领受的吗？}还有什么比这更愚昧或更亵渎的呢？竟然有人敢说基督的肢体是某位的宫殿，而他们认为这位是次于基督的受造物？因为宗徒在另一处说：\BibleQuote{你们的身体是基督的肢体。}但如果基督的肢体也是圣神的宫殿，那么圣神就不是受造物；因为我们的身体既然是祂的宫殿，我们必须将唯独归于天主的敬拜归于祂，这在希腊文称为\OriginalTerm{敬拜（钦崇式）}{\GreekText{λατρεία}}。因此，宗徒说：\BibleQuote{所以你们要在你们的身体上光荣天主。}\par

% structure: p0027 source=h2 target=subsection reason=relative-heading-rank; base=h2

\subsection{第7章——圣子如何小于圣父，且小于祂自己。}

14. 在这些以及类似的神圣经文的证词中，正如我所说，我们的前辈自由地运用这些证词驳斥了异端者的诡辩或错误，\OriginalTerm{天主圣三}{Trinity}的合一与平等被启示给我们的信德。但正因为为了完成我们的救恩，天主\OriginalTerm{圣言}{Word}\OriginalTerm{降生成人}{Incarnation}，使基督耶稣这个人成为天主与人之间的中保，所以在圣经中有许多这样的说法，表明甚至最明确地宣告圣父比圣子更大；人们因缺乏对圣经整体脉络的仔细查考或思考而犯错，试图将那些按照血肉论及耶稣基督的事，转移到祂在降生成人之前就永恒存在、且永远存在的那一实体上。例如，他们说圣子小于圣父，因为经上记载主亲自说：\BibleQuote{我的父比我大}。但真理表明，按同一意义，圣子也小于祂自己；因为那\BibleQuote{空虚了自己，取了仆人的形式}的，祂怎么会不比自己小呢？因为祂取了仆人的形式，并没有因此而失去\OriginalTerm{天主的形式}{form of God}，在那种形式中祂与圣父平等。那么，既然仆人的形式被取了，而天主的形式并未失去，因为无论是在\OriginalTerm{仆人的形式}{form of a servant}中还是在天主的形式中，祂本身就是圣父的同一独生子，在天主的形式中与圣父平等，在仆人的形式中作为天主与人之间的中保——基督耶稣这个人；有谁不能看出，祂自己在天主的形式中也大于祂自己，然而同样在仆人的形式中小于祂自己呢？因此，圣经说这两者——圣子与圣父平等，以及圣父比圣子大——并非没有缘故。因为当前者理解为因天主的形式，后者理解为因仆人的形式时，就没有混乱。事实上，贯穿所有圣经以澄清这一问题的规则，在宗徒保禄的一封书信的一章中已清楚地向我们推荐了这种区别。因为他说：\BibleQuote{祂虽具有天主的形式，却没有以自己与天主同等为应当把持不舍的，反而空虚了自己，取了仆人的形式，成了人的模样；且以人的形态出现。}所以，天主子在天主性上与天主圣父相等，但在\Quote{形态}上较小。因为在祂所取的仆人的形式中，祂小于父；但在天主的形式中，在祂取仆人的形式之前祂也如此存在，祂与父平等。在天主的形式中，祂是圣言，\BibleQuote{万物是藉着祂造的}；但在仆人的形式中，祂\BibleQuote{生于女人，生于法律之下，为要救赎那些在法律之下的人}。同样，在天主的形式中，祂造了人；在仆人的形式中，祂成了人。因为如果只有父没有子造了人，就不会记载：\BibleQuote{让我们按\Emph{我们的}肖像，照\Emph{我们的}模样造人}。因此，因为天主的形式取了仆人的形式，所以祂既是天主又是人；但说祂是天主，是因为那取者的天主性；说祂是人，是因为那被取的人性。因为藉着那取，并没有使两者中的一方转化为另一方：天主性不转变为受造物而不再是天主性，受造物也不转变为天主性而不再是受造物。\par

% structure: p0029 source=h2 target=subsection reason=relative-heading-rank; base=h2

\subsection{第八章——圣经中关于圣子服从圣父的经文解释（这些经文曾被误解）。基督不会将王国交给父以至于从自身夺走。瞻仰祂是所有行动的应许终点。圣神与父同等足以使我们得享真福。}

15. 关于宗徒所说的：\BibleQuote{凡事都服从于祂以后，那时，子自己也要屈服于那使万物服从于祂的}，这经文或如此措辞，是免得有人以为基督按人性受取的\OriginalTerm{形态}{fashion}将来会转化为神性，或更确切地说，转化为神格本身，祂不是受造物，而是天主圣三的合一，是无形体、不变、\OriginalTerm{同质}{consubstantial}、与自身同永恒的本性；或若有人如某些人所想，主张经文\BibleQuote{那时，子自己也要屈服于那使万物服从于祂的}如此措辞，是为了使人相信那\OriginalTerm{屈服}{subjection}本身是受造物未来转变为造物主实体或本质的改变和转化，即受造物的实体变成造物主的实体——这种人至少承认，当主说\BibleQuote{我的父比我大}时，这事尚未发生，这是无可置疑的。因为祂说这话不仅是在升天以前，也是在受难和从死者中复活以前。但那些以为祂内的人性将改变和转化为神性实体，并以为如此说\BibleQuote{那时，子自己也要屈服于那使万物服从于祂的}——如同说，那时人子自己、以及天主圣言所取的的人性，将改变为那使万物服从于祂者的本性——也必须认为这将在审判之日以后发生，当\BibleQuote{祂把王国交给天主父}之后。因此，即使按照这种意见，父仍大于那取自重贞女的\OriginalTerm{仆人的形体}{form of a servant}。但若有人进一步肯定，说基督耶稣的人已经改变为天主的实体，至少他们不能否认，当祂在受难前说\BibleQuote{因为我的父比我大}时，人性仍然存在；因此毫无疑问，这话的意思是说，父大于仆人的形体，而在\OriginalTerm{天主的形体}{form of God}中，子是与父相等的。谁也不要听宗徒说\BibleQuote{但当祂说万物都服从于祂时，显然那使万物服从于祂的除外}，就以为万物服从于子这话是指父，从而认为子自己并没有使万物服从于自己。因为宗徒在致斐理伯人书中明确宣告：\BibleQuote{我们的家乡在天上，我们也等待救主，主耶稣基督：祂将改变我们卑微的身体，使之相似祂光荣的身体，因为祂有能力使万物服从于自己}。因为父和子的作为是不可分的。否则，父自己并没有使万物服从于自己，而是子使万物服从于祂，子将王国交给父，并废除一切率领者、一切掌权者和一切异能者。因为这些话是针对子说的：\BibleQuote{当祂把王国交给天主父时}，宗徒说，\BibleQuote{当祂废除一切率领者、一切掌权者和一切异能者时}。因为那废除的，也使它服从。\par

16. 我们也不可认为基督将如此把王国交给天主父，以致祂从自己手中夺走。因为有些虚妄的谈论者竟这样想。因为当经文说\BibleQuote{祂把王国交给天主父}时，祂自己并未被排除；因为祂同父是一个天主。但\Quote{\OriginalTerm{直到}{until}}这个词欺骗了那些粗心阅读圣经而热衷争论的人。因为经文继续说：\BibleQuote{因为祂必须为王，直到把一切仇敌放在祂脚下}；好像当祂如此放置后，祂就不再为王了。他们也不明白这话与另一处经文相同：\BibleQuote{他的心坚定，不畏惧，直到他看见仇敌遭报}。因为当他看见时，他就不再畏惧了。那么，\BibleQuote{当祂把王国交给天主父}是什么意思，好像天主父现在没有王国？而是因为祂将来要把所有义人——就是现在因信德生活、天主与人之间的中保、人基督耶稣所统治的——带到那看见，也就是宗徒所称的\Quote{\OriginalTerm{面对面}{face to face}}；因此，\BibleQuote{当祂把王国交给天主父}这话等于说：当祂把信众带到对天主父的默观时。因为祂说：\BibleQuote{我父将一切交给了我；除了父，没有人认识子；除了子和子愿意启示的人，也没有人认识父}。那时，子将启示父，\BibleQuote{当祂废除一切率领者、一切掌权者和一切异能者}；也就是说，不再需要任何藉着天使的率领者、掌权者和异能者的\OriginalTerm{比喻的经纶}{economy of similitudes}。关于他们，在\BibleBook{}{雅歌}中对新娘所说的话很合适：\BibleQuote{王坐在席上，我们要为你编金辫，镶上银钉}；也就是说，只要基督在祂的隐密处：因为\BibleQuote{你们的生命与基督一同藏在天主内；当基督，我们的生命，显现时，你们也要与祂一同在光荣中显现}。在此以前，\BibleQuote{我们现在是藉着镜子观看，模糊不清}，即藉着比喻，\BibleQuote{那时，就要面对面了}。\par

17. 这种默观被赐予我们，作为一切行动的目标和永远圆满的喜乐。因为\BibleQuote{我们是天主的子女；将来如何，还没有显明；但我们知道：当祂显现时，我们必要相似祂，因为我们要看见祂实在怎样。}因为祂对祂的仆人梅瑟所说的话：\BibleQuote{我是自有者；你要这样对以色列子民说：那自有者派遣我到你们这里来。}这就是我们当永远生活时所将默观的。因为经上也是这样说的：\BibleQuote{永生就是：认识祢，唯一的真天主，和祢所派遣的耶稣基督。}这将在主来临，并且\BibleQuote{揭露了黑暗中的隐密事物}的时候；那时现世死亡与腐朽的黑暗将过去。那时将是我们的清晨，圣咏上论及这清晨说：\BibleQuote{早晨我要向祢献上祈祷，并要默观祢。}我理解，关于这一默观，经上所说的\BibleQuote{当祂把王国交给天主，即圣父的时候}，就是指：当祂将那些义人——如今因信德而生活、在天主与人之间的中保、基督耶稣为人者所统治的——带到对天主，即圣父的默观中。如果我在此愚昧，愿那更有知识者纠正我；至少在我看来，情况正如同我所说的。因为当我们来到对祂的默观时，我们将不再寻求其他事物。但是，只要我们的喜乐仍在希望中，那默观就尚未实现。因为\BibleQuote{所见到的希望不是希望：人若看见，为什么还希望呢？但我们若希望那没有看见的，就安心等待。} \Emph{viz.} \BibleQuote{君王坐席的时候。}那时，经上所写的将会实现：\BibleQuote{在祢面前有圆满的喜乐。}那喜乐之外，不再需要别的；因为可能需要的，也不会多于那喜乐。因为圣父将显示给我们，而这就足够我们了。斐理伯对此理解得很好，所以他向主说：\BibleQuote{请将父显示给我们，我们就满足了。}但他尚未明白，他自己也能这样说出同样的事：主，请将祢自己显示给我们，我们就满足了。因为，为使他明白这事，主答复他说：\BibleQuote{斐理伯！这么长久的时间，我和你们在一起，你还不认识我吗？谁看见了我，就是看见了父。}但因为祂愿意他在这事能看见之前，先因信德而生活，祂继续说：\BibleQuote{你不信我在父内，父在我内吗？}因为\BibleQuote{我们住在肉身内时，是与主远离的：因为我们凭着信德行走，并非凭着眼见。}因为默观是信德的报酬，为此我们借信德洁净心灵；正如经上所记：\BibleQuote{借信德洁净了他们的心。}而我们的心为了这默观必须被洁净，这尤其由这段经文所证明：\BibleQuote{心里洁净的人是有福的，因为他们要看见天主。}而这就是永生，天主在圣咏中说：\BibleQuote{我要以长寿满足他，并将我的救恩显示给他。}因此，不论我们听见\Quote{请显示子给我们}，或听见\Quote{请显示父给我们}，其实都一样，因为两者不能分开而显示。因为祂们原是一体，正如祂自己也说：\BibleQuote{我与父原是一体。}最后，正因这不可分割性，有时只提及圣父，或只提及圣子，就已足够，因为以后祂的面容将充满我们的喜乐。\par

18. 同样，父和圣子的圣神也不可排除，祂被特别称为\BibleQuote{真理之神，世界不能接受祂。}因为享受天主圣三本身——我们是按祂的肖像受造的——确实是我们喜乐的圆满，再没有比这更大的喜乐了。因此，有时圣神被描述为好像唯独祂就足以使我们幸福：而祂确实独自就够了，因为祂不能与父和子分离；正如父独自就够了，因为祂不能与子和圣神分离；子独自就够了，因为祂不能与父和圣神分离。祂说：\BibleQuote{如果你们爱我，就要遵守我的命令；我要请求父，祂就必赐给你们另一位护慰者，使祂永远与你们同在；祂就是真理之神，世界不能接受祂，}就是说，世界的爱好者不能接受祂？因为\BibleQuote{属血气的人不能领受天主圣神的事。}但也许进一步看来，似乎\BibleQuote{我要请求父，祂就必赐给你们另一位护慰者}这句话，好像是说子独自不够。而另一处谈到圣神时，好像祂独自完全足够：\BibleQuote{当祂，真理之神，来到时，祂将引导你们进入一切真理。}请问，难道这里把子排除在外，好像祂没有教导一切真理，或者好像圣神要填补子未能完全教导的吗？那么，若他们愿意，就让他们说圣神大于子，而他们惯常称子较小吧。或者，难道诚然因为经文没有说\Quote{唯独祂}，或\Quote{除祂以外没有别人}，会引导你们进入一切真理，所以他们可以相信子也与祂一同教导？既然如此，宗徒排除了子知道天主的事，因为他说：\BibleQuote{同样，天主的事，除了天主的神，也没有人知道。}于是这些乖谬的人便可借此说，唯独圣神教导关于天主的事，甚至子也受祂教导，如同大的教导小的；而子自己却极其尊重圣神，以致说：\BibleQuote{只因我给你们讲了这些事，你们心中就充满了忧愁。然而，我将真情告诉你们：我去为你们更有益，因为我若不去，护慰者就不会到你们这里来。}\par

% structure: p0034 source=h2 target=subsection reason=relative-heading-rank; base=h2

\subsection{第九章——有时一切皆在一个位格中被领悟。}

但这样说，并非因为天主的圣言和圣神有任何不平等，而是好像人子与他们同在会阻碍那位的来临，那位并非更小，因为祂没有像圣子那样\BibleQuote{空虚了自己，取了奴仆的形体}。因此，必须将那奴仆的形体从他们眼前除去，因为他们凝视它时，以为他们所见到的就是基督。因此也有这句话：\BibleQuote{如果你们爱我，你们就该欢喜，因为我说了：\InnerQuote{我往父那里去；因为父比我大}}。也就是说，因此我必须到父那里去，因为你们这样看见我时，你们因所见到的而以为我比父小；因而，你们专注在受造物和我所取的样式上，就看不到我与父的平等。同样也有这话：\BibleQuote{不要摸我，因为我还没有升到父那里}。因为触摸仿佛限制了他们的概念，主因此不愿意人内心的思想因转向祂而受限制，以至于祂被看作只是祂所显现的那样。但升到父那里的意思是，这样显现为与父平等，以至于我们足以看见的界限可以在那里达到。有时也单独论及圣子，说祂自己就足够了，我们慈爱和渴慕的全部赏赐就在看见祂时展示。因为如此说：\BibleQuote{那拥有我的诫命并遵守它们的，就是爱我的；爱我的必蒙我父爱他，我也要爱他，并将我自己显示给他}。请问，因为祂在这里没有说\Quote{我也要把父显示给他}，难道祂就排除了父吗？相反，因为\BibleQuote{我与父原是一体}是真实的，当父被显示时，在祂内的圣子也被显示；当圣子被显示时，在祂内的父也被显示。因此，正如当说\BibleQuote{我将我自己显示给他}时，理解为祂也显示父；同样，在说\BibleQuote{当祂把王国交给天主，即父}时，理解为祂并没有从自己那里夺去什么；因为当祂带领信徒默观天主，即父时，无疑祂将带领他们默观祂自己，祂曾说过：\BibleQuote{我将我自己显示给他}。因此，当犹达斯对祂说：\BibleQuote{主，祢为什么将祢自己显示给我们，而不显示给世界呢？}耶稣回答他说：\BibleQuote{人若爱我，必遵守我的话；我父也必爱他，我们要到他那里去，并在他那里作我们的住所}。看，祂不仅将祂自己显示给那爱祂的人，因为祂与父一同来到他那里，并与他同住。\par

19. 或许有人会想，当父与子在爱他们者那里作住所时，圣神是否被排除于那住所之外？那么，关于圣神以上所说的是什么呢：\BibleQuote{世界不能接受祂，因为看不见祂；你们却认识祂，因为祂与你们同在，并在你们内}。因此，那被说为\BibleQuote{祂与你们同在，并在你们内}的圣神，并不被排除于那住所之外；除非有人愚蠢到以为，当父与子来在爱他们者那里作住所时，圣神会离开那里，仿佛让位给更大的那两位。但圣经本身回应了这种属肉体的想法；因为上面一点说：\BibleQuote{我要求父，祂必赐给你们另一位护慰者，使祂永远与你们同在}。因此，当父与子来时，祂不会离开，而是与它们永远住在同一住所；因为祂不会不带他们而来，他们也不会不带祂而来。但为了暗示天主圣三，有些事是分别肯定的，每位也分别被命名；然而不应理解为其他位格被排除，因为同一圣三、父及子及圣神同一实体和天主性的合一。\par

% structure: p0037 source=h2 target=subsection reason=relative-heading-rank; base=h2

\subsection{第十章 — 基督以何种方式将王国交给天主，即父。王国已交给天主，即父后，基督将不再为我们代祷。}

20. 因此，我们的主耶稣基督将这样把王国交给天主，即圣父，而祂自己并不因此被排除，圣神也不被排除，那时祂将带领信徒进入对天主的默观，这默观就是一切善行的终极目标，永久的安息，以及永不会被夺去的喜乐。因为祂在以下的话中预示了这点：\BibleQuote{我要再见到你们，那时你们心里要喜乐，并且你们的喜乐谁也不能从你们夺去。} \BibleRef{若 16:22} 玛利亚坐在主脚前，专心聆听祂的话，预示了这种喜乐的肖像；她安息于一切事务之外，专注于真理，按照今生所能的方式，但仍是那永恒之事的预像。因为当她的姊妹玛尔塔为了必要的职务忙碌时——这些职务虽美好有用，但在安息来临时将要过去——玛利亚自己却安息于主的话语中。因此，当玛尔塔抱怨她姊妹不帮助她时，主回答说：\BibleQuote{玛利亚选择了更好的一份，是不能从她夺去的。} \BibleRef{路 10:42} 祂并不是说玛尔塔做得不好，而是说\BibleQuote{最好的一份是不能被夺去的。}因为那忙于服务需要的一份，在需要本身过去之后，必被夺去。因为那将要过去的善功的赏报，是不会过去的安息。因此，在那默观中，天主将成为万物中的万有；因为除了祂自己以外，将不再需要任何事物，而被祂光照并享有祂就足够了。因此，那位在他内\BibleQuote{圣神以无可言喻的叹息代我们转求} \BibleRef{罗 8:26} 的人说：\BibleQuote{我有一事祈求主，我要恳切寻求此事：使我一生岁月，常居住在主殿里，欣赏主的甘饴慈祥，瞻仰主的圣殿。} \BibleRef{咏 27:4} 因为我们那时将默观天主圣父、圣子和圣神，就是当天主与人之间的中保，那人基督耶稣，把王国交给天主，即圣父的时候，祂不再作为我们的中保和司祭、天主子与人子为我们转求；而是祂自己，就祂为我们取了奴仆的形象而成为司祭而言，也将被置于那使万有服从祂（圣子）、且祂（圣子）也将万有置于其下的那一位（圣父）之下。所以，就祂是天主而言，祂将同那位一起使我们服从祂自己；就祂是司祭而言，祂将同我们一起被置于那位之下。因此，正如降生成人的圣子既是天主又是人，更恰当的说法是：圣子内的人性是另一性体（不同于圣子），而不是圣子在圣父内是另一性体（不同于圣父）；正如我灵魂的肉性相对于我灵魂本身而言，虽在同一人身内，却比另一个人的灵魂相对于我的灵魂更具另一性体。\par

21. 因此，当祂\BibleQuote{将王国交付给天主，即圣父}——即当祂带领那些因信德而相信并生活的人（祂如今作为中保为他们转祷）达到那凝思，我们为此叹息呻吟，而在劳苦和呻吟逝去后——那时，既然王国已交付给天主，即圣父，祂将不再为我们转祷。祂在说以下的话时预示了这一点：\BibleQuote{这些事我用比喻对你们说了；时候将到，我不再用比喻对你们说，而要明明地将圣父传报给你们。}即那时不再有\BibleQuote{比喻}，当看见成为\BibleQuote{面对面}时。这就是祂所说的：\BibleQuote{但我将明明地将圣父传报给你们}；好像祂说：我将明明地给你们显示圣父。因为祂说，我将\BibleQuote{传报}给你们，因为祂是祂的\OriginalTerm{圣言}{word}。祂接着说：\BibleQuote{到那日，你们要因我的名祈求；我对你们不说，我要为你们求圣父，因为圣父自己爱你们，因为你们爱了我，且相信我是从天主出来的。我从圣父出来，到了世界；我又离开世界，往圣父那里去。}\BibleQuote{我从圣父出来}是什么意思？除非是：我没有以与圣父相等的形像显现，而是以另种方式，即小于圣父的、我所取受造物之形像？而\BibleQuote{我到了世界}是什么意思？除非是：我向甚至爱此世界的罪人显明了我所取的仆人形像，虚己？而\BibleQuote{我又离开世界}是什么意思？除非是：我从爱世界者的眼前取走他们所见的？而\BibleQuote{我往圣父那里去}是什么意思？除非是：我教导那些忠于我的人，在我与圣父相等的存在中理解我？相信这事的人将被认为配得由信德被带到看见，即带到那看见，在带领他们到此看见时，祂被称为\BibleQuote{将王国交付给天主，即圣父}。因为祂的忠诚者，祂用血所救赎的，被称为祂的王国，祂如今为他们转祷；但那时，使他们安息在祂自己里面，在那里祂与圣父相等，祂将不再为他们在圣父面前转祷。因为祂说：\BibleQuote{圣父自己爱你们}。因为祂\BibleQuote{转祷}，是就祂小于圣父而言；但就祂与圣父相等而言，祂与圣父同赐予。因此，祂当然不把自己排除在祂所说的\BibleQuote{圣父自己爱你们}之外；而是指我上面已谈论并充分暗示的方式理解——即大多时候，天主圣三的每位位格如此被称呼，以致其他位格也可被理解。因此，\BibleQuote{因为圣父自己爱你们}如此说，以至于可以理解圣子和圣神也在内：并非祂现在不爱我们，祂没有怜惜自己的儿子，反而为众人把祂交出；但天主爱我们，爱我们将成为的样式，而非我们现在的样式，因为祂所爱的人是什么样，祂永远保守的也是什么样；这将在那时成就，当祂如今为我们转祷的那位\BibleQuote{将王国交付给天主，即圣父}时，以致不再求圣父，因为圣父自己爱我们。但除了信德，我们凭什么功绩相信那应许呢？我们凭此信德将到达看见；以致祂可以爱我们，爱我们成为的样式，正如祂爱我们以使我们成为那样；而非爱我们现在的样式，祂恨我们因为我们是这样，并劝勉并使我们有能力愿不再如此。\par

% structure: p0040 source=h2 target=subsection reason=relative-heading-rank; base=h2

\subsection{第十一章——依据圣经中的何种规则理解圣子现在相等，现在较小。}

22. 因此，掌握了这一解释有关天主子圣经的规则，即我们必须区分其中哪些涉及天主的形状，在此形状下他与圣父平等，哪些涉及他所取的仆人的形状，在此形状下他小于圣父；我们就不会被圣书表面上相反且互相矛盾的说法所困扰。因为圣子和圣神，按照天主的形状，与圣父平等，因为他们都不是受造物，正如我们已经指出的；但是按照仆人的形状，他小于圣父，因为他自己说过：\BibleQuote{父大于我}；他也小于自己，因为论到他说，他空虚了自己；他也小于圣神，因为他自己说：\BibleQuote{凡说话干犯人子的，可得赦免；但说话干犯圣神的，必不得赦免}。他也借着圣神行奇迹，说：\BibleQuote{我若靠着圣神驱魔，天主的国就临到你们了}。在\BibleBook{依撒意亚}{书}中他说——在他自己在会堂中宣读的那段经文里，并且毫无疑惑地表明是关于他自己应验的——\BibleQuote{主圣神在我身上，因为他用油傅了我，叫我传福音给温良的人，他派遣我向俘虏宣告自由}等等：他因此宣称自己是\Quote{被派遣的}，因为圣神在他身上。按照天主的形状，万物是他造的；按照仆人的形状，他自己是由女人生的，生于法律之下。按照天主的形状，他与父是一体；按照仆人的形状，他不是来行自己的旨意，而是行那派遣他者的旨意。按照天主的形状，\BibleQuote{怎样父自己有生命，也赐给子在自己内有生命}；按照仆人的形状，他的\BibleQuote{心灵忧闷得要死}；并且他说：\BibleQuote{父啊！若是可能，就让这杯离开我}。按照天主的形状，\BibleQuote{他是真天主，是永生}；按照仆人的形状，\BibleQuote{他服从至死，且死在十字架上}。——23. 按照天主的形状，凡父所有的，都是他的，他说：\BibleQuote{我的一切都是你的，你的一切都是我的}；按照仆人的形状，那道理不是他自己的，而是那派遣他者的。\par

% structure: p0042 source=h2 target=subsection reason=relative-heading-rank; base=h2

\subsection{第十二章——圣子如何被称为不知道父所知道的日子和时辰。论基督有些话是按照天主的形状说的，另一些是按照仆人的形状说的。基督如何赐予王国，如何不赐予。基督将审判，也将不审判。}

再次，\BibleQuote{至于那日子和那时辰，没有人知道，连天上的天使也不知道，圣子也不知道，只有父知道}。因为他对此是\Quote{不知}的，是\Emph{使别人不知}；也就是说，他当时没有那样知道以便告诉门徒：正如对\Person{亚巴郎}所说：\BibleQuote{如今我知道你敬畏天主了}，也就是说，如今我使你知道了；因为他在那试探中受过考验，就认识了自己。因为他当然会在适当的时候告诉门徒同样的事；讲论那尚未来临的事如同已成过去，他说：\BibleQuote{今后我不再称你们为仆人，因为仆人不知道他主人所做的事；我称你们为朋友，因为凡我从我父所听见的，我都已告诉了你们}；这他还没有做，却好像已经做了，因为他一定会做。因为他自己对门徒说：\BibleQuote{我还有好多事要告诉你们，但你们现在不能承担}。在这些事中也要包括\Quote{那日子和时辰}。因为宗徒也说：\BibleQuote{我决定在你们中什么也不知道，只知道耶稣基督，并他被钉在十字架上}；因为他在对那些人说话，他们不能领受关于基督天主性的更高深事理。稍后他对他们说：\BibleQuote{我不能对你们说话如同对属神的人，只能如同对属血肉的人}。因此，他在他们中间是\Quote{不知}的，是他们不能从他那里知道的。他只说他知道了那些他们适宜从他那里知道的事。总之，在成熟的人中间他知道，而在婴孩中间他不知道；因为他那里说：\BibleQuote{我们在成熟的人中间讲智慧}。因为一个人被称为不知道他所隐藏的东西，按照那种说法，如同一条沟渠被称为盲点，因为它是隐藏的。因为圣经不使用不同于人在使用中的任何说法，因为圣经是对人说话。\par

24. 按天主的形象，经上说：\BibleQuote{在一切山丘之先，祂生了我}，即是在一切受造物的崇高之先；又说：\BibleQuote{在晨曦之先，我生了祢}，即是在一切时间和暂时事物之先。但按仆人的形象，经上说：\BibleQuote{主在祂道路的起始就造了我}。因为按天主的形象，祂说：\BibleQuote{我是真理}；按仆人的形象，说：\BibleQuote{我是道路}。因为祂自己是死者的首生者，为祂的教会开辟了一条通往天主的国、通往永生的通道；祂是教会的元首，也使身体成为不朽的。因此，在祂的工作中，祂在\BibleQuote{天主道路的起始}被造。因为按天主的形象，祂就是元始，那也向我们讲话的；在这个\BibleQuote{元始}内，天主创造了天地。但按仆人的形象，\BibleQuote{祂如同新郎走出洞房}。按天主的形象，\BibleQuote{祂是万物的首生者，祂先万有而存在，万有靠祂而存在}；按仆人的形象，\BibleQuote{祂是身体——教会的元首}。按天主的形象，\BibleQuote{祂是光荣的主}。由此明显可见，祂亲自光荣祂的圣人们。因为经上说：\BibleQuote{祂所预定的人，也召叫了；所召叫的人，也成义了；所成义的人，也光荣了。}因此，论到祂说：祂使不虔敬的人成义；论到祂说：祂是公义的，并使他人成义。所以，如果祂也使那些祂所成义的人受了光荣，那使他人成义的，祂自己也光荣他们；祂就是我所说的光荣的主。然而，按仆人的形象，当门徒们询问自己的光荣时，祂回答说：\BibleQuote{坐在我的右边和左边，不是我可以给的，而是[必要给他们的]那些我父为他们预备的人。}\par

25. 但那由圣父预备的，也是由圣子自己预备的，因为祂与圣父原是一体。因为我们已经在圣经中许多表达方式中显出，在这天主圣三内，对每一位所说的，也是对全体说的，这是由于同一实体的不可分割的运作。正如祂论及圣神时也说：\BibleQuote{若我去，我就要派遣祂到你们这里来。}祂没有说：\Emph{我们}要派遣；但似乎只是圣子将派遣祂，而非圣父；然而在另一处祂说：\BibleQuote{我还与你们在一起的时候，给你们讲论了这些事；但那护慰者，就是圣神，就是圣父因我的名所要派遣的，祂要教训你们一切。}这里再次如此说，似乎圣子也不会派遣祂，而只有圣父。因此，正如在这些经文里，同样在祂说：\BibleQuote{而是给我父为谁预备的人，}祂的意思是让人明白，祂自己与圣父一同为祂愿意的人预备光荣的座位。但有人会说：在那里，当祂论及圣神时，祂这样说祂自己要派遣祂，而并不否认圣父也会派遣祂；在另一处，祂这样说圣父要派遣祂，而并不否认祂自己也会这样做；但在这里，祂明确说：\BibleQuote{不是我可以给的，}接着又说这些是由圣父预备的。但这正是我们已经确定的：这是按仆人的形象说的。\Emph{即}，我们当这样理解\BibleQuote{不是我可以给的}，就好像说：这不是在人能力之内可以给的；这样，祂便可以在祂与圣父同等的天主性内被理解为赐予。\BibleQuote{不是我的，}祂说，\BibleQuote{可以给的；}也就是说，我不是凭人的能力给这些事物，而是\BibleQuote{给我父为谁预备的人。}但随后你也要注意理解，如果\BibleQuote{凡父所有的，都是我的，}那么这当然也是我的，并且我与圣父一起预备了这些事物。\par

26. 因为我又问，这话是怎样说的：\BibleQuote{若有人不听我的话，我不审判他？} 因为或许祂在这里说\BibleQuote{我不审判他}，与那里说\BibleQuote{不是我可以给的}意义相同。但这里接着说？祂说：\BibleQuote{我来不是为审判世界，乃是为拯救世界}；然后祂又加上：\BibleQuote{那拒绝我、不接受我话的，有审判他的。} 那么，这里我们应该理解为圣父，除非祂又加上：\BibleQuote{我所讲的话，在末日要审判他。} 那么，圣子既不审判，圣父也不审判，而是圣子所说的话？不对，且听下文：\BibleQuote{因为我没有凭自己说话，而是那派遣我的圣父给了我诫命，该说什么，该讲什么；我知道祂的诫命是永生：所以我所说的，正如圣父对我说的，我就如此说。} 因此，如果圣子不审判，而是\BibleQuote{圣子所说的话}，而圣子所说的话所以审判，是因为圣子\BibleQuote{没有凭自己说话，而是那派遣他的圣父给了他诫命，该说什么，该讲什么}：那么，审判的肯定是圣父，因为圣子所说的话是圣父的话；而且圣子自己就是\OriginalTerm{圣父的圣言}{Word of the Father}。因为圣父的诫命并非一样东西，圣父的话又是另一样；因为祂称它既为话，也为诫命。因此，我们来看，或许当祂说\BibleQuote{我没有凭自己说话}时，祂的意思是——我不是从自己生的。因为如果祂说圣父的话，那么祂就是说自己，因为祂自己就是圣父的圣言。因为祂通常说：\BibleQuote{圣父给了我}；祂的意思是让人明白圣父生了祂：不是给已有的、尚未拥有的祂什么，而是\Quote{给祂使祂拥有}的真正意义是\Quote{生了祂使祂存在}。因为在降生成人、取了我们肉躯之前，天主子——\OriginalTerm{独生子}{Only-begotten}，万物由祂而造——并不像受造物那样，祂的\Emph{是}是一回事，\Emph{有}是另一回事；而是祂\Emph{是}的方式就是\Emph{有}祂所\Emph{是}的。这一点，若人能够领悟，在祂说\BibleQuote{因为圣父在自己内有生命，同样也给了圣子在自己内有生命}的地方说得更明白。因为祂不是给已经存在但没有生命的祂，使祂在自己内有生命；而是就祂\Emph{是}而言，祂是生命。因此\BibleQuote{给了圣子在自己内有生命}的意思是，祂生了圣子成为不变的生命，即永生。因此，既然天主的圣言就是天主子，天主子就是\BibleQuote{真天主和永生}，正如\Person{若望}在他的书信中所说的；那么在这里，当主说\BibleQuote{我所讲的话，在末日要审判他}，并称那话为圣父的话和圣父的诫命，而那诫命就是永生时，我们还能承认别的什么吗？祂说：\BibleQuote{我知道祂的诫命是永生。}\par

27. 因此，我问，我们该如何理解：\BibleQuote{我不审判他；但我所说的话要审判他：}从下文看来，这话似乎是说，我不审判；但圣父的圣言要审判。但圣父的圣言就是天主子本身。是否应该这样理解：我不审判，但我要审判？这怎么可能为真呢，除非这样：\Emph{viz.}，我不凭人的能力审判，因为我是人子；但我要凭圣言的能力审判，因为我是天主子？或者，若说\BibleQuote{我不审判，但我要审判}似乎仍然矛盾和不一致，那么，对于祂所说的这段经文，我们该说什么：\BibleQuote{我的教训不是我的？}怎么是\BibleQuote{我的，}却又\BibleQuote{不是我的？}因为祂不是说：\Emph{这个}教训不是我的，而是说：\BibleQuote{\Emph{我的}教训不是我的：}凡祂称为自己的，祂也称作不是自己的。这怎么可能为真呢，除非祂在某一方面称为自己的，在另一方面称为不是自己的？按照天主的形象，是祂自己的；按照仆人的形象，不是祂自己的。因为当祂说：\BibleQuote{不是我的，而是那派遣我者的，}祂使我们回到圣言本身。因为圣父的教训就是圣父的圣言，那就是独生子。还有，这句话又是什么意思：\BibleQuote{信我的，不是信我？}怎么信祂，又不信祂？如此对立且不一致的事如何能理解——祂说：\BibleQuote{凡信我的，}祂说，\BibleQuote{不是信我，而是信那派遣我者；}——除非你这样理解：凡信我的，不是信他所看见的那一位，免得我们的希望在于受造物，而是信那取了受造物的一位，祂借此才能在人的眼中显现，并藉着信德洁净我们的心，好瞻仰祂与圣父同等。因此，当祂将信徒的注意力转向圣父，并说：\BibleQuote{不是信我，而是信那派遣我者，}祂当然不是指祂自己与圣父分离，即与那派遣祂者分离；而是使人能这样信祂，如同他们信圣父，祂与圣父同等。祂在另一处明确地说：\BibleQuote{你们信天主，也要信我：}也就是说，如同你们信天主一样，也要信我；因为我和父是同一的天主。因此，正如在这里，祂好像将人的信德从祂自己身上撤回，并转移到圣父身上，藉着说：\BibleQuote{不是信我，而是信那派遣我者，}然而祂当然没有与那派遣者分离；同样，当祂说：\BibleQuote{给的不是我，而是【将要给那些】我父为他们所预备的人，}我认为在何种意义上两者都应被理解是清楚的。因为另一处也属于同类：\BibleQuote{我不审判；}然而祂自己要审判活人死人。但祂不会凭人的能力这样做，因此，回转到天主性，祂将人的心提升向上；而为提升人心，祂自己降了下来。\par

% structure: p0048 source=h2 target=subsection reason=relative-heading-rank; base=h2

\subsection{第十三章——关于同一基督，因同一一位格（神人）的两种本性而说不同的事。为什么说父不审判，而将审判权交给了子。}

28. 然而，除非那同一者，因祂所取仆人的形体而为人子，而祂因祂所在的天主的形体而为天主之子，宗徒保禄就不会论及这世界的掌权者说：\BibleQuote{因为他们若认识，就不至于将光荣的主钉在十字架上了。}因为祂是按仆人的形体被钉的，然而\BibleQuote{光荣的主}竟被钉了。因为那个\Quote{摄取}是这样的，它使天主成为人，使人成为天主。然而，什么是因什么而说的，什么是按照什么而说的，由勤勉而虔诚的读者在主帮助下自己辨别。例如，我们曾说祂光荣自己的人，作为天主，当然也作为光荣的主；然而光荣的主被钉了，因为天主被说成被钉，不是按神性的能力，而是按肉身的软弱，这是合理的：正如我们说，祂作为天主审判，即按神圣的能力，而非按人的能力；然而那人自己将审判，正如光荣的主被钉：因为祂明说：\BibleQuote{当人子在祂的光荣中，与祂一同前来的众圣天使时，那时祂要坐在光荣的宝座上，万民都要聚集在祂面前；}以及那处所述的关于未来审判的其余部分，直到最后判决。而犹太人，因他们坚持恶行将在那审判中受罚，正如别处所写：\BibleQuote{他们要仰望他们所刺透的。}因为无论善人恶人都将看见审判生者死者，无疑恶人不能看见祂，除非按祂是人子的形体；然而是在祂审判的光荣中，而非在祂被审判的卑微中。但不敬天主的人无疑将不能看见祂那与圣父同等的天主的形体。因为他们心不纯洁；而\BibleQuote{心里洁净的人是有福的，因为他们要看见天主。}那看见是面对面，正是应许给义人作为最高奖赏的看见，并且那时将发生，当祂\BibleQuote{把王国交给天主，即圣父}；在这\BibleQuote{王国}中，祂也要理解为看见祂自己的形体，整个受造物都服从天主，包括其中天主之子成为人子的那一位。因为，按照这受造物，\BibleQuote{圣子自己也要服从那使万有服从于祂的，为叫天主成为万物之中的万有。}否则，如果天主之子以祂与圣父同等的形体审判，在审判时也向不敬天主的人显现；那么祂应许给爱祂的人作为大事的话又成了什么？祂说：\BibleQuote{我要爱他，并要将我自己显示给他？}因此，祂将作为人子审判，但不是凭人的能力，而是凭祂作为天主之子的能力；另一方面，祂将作为天主之子审判，却不以祂与圣父同等的[非降生的]形体显现，而是以祂作为人子的[降生的]形体显现。\par

29. 因此，两种说法都可以使用：人子将审判，以及，人子将不审判；因为人子将审判，为使那经文得以成就：\BibleQuote{当人子来临时，那时万民都要聚集在祂面前；}并且人子将不审判，为使那经文得以成就：\BibleQuote{我不审判他；}以及\BibleQuote{我不寻求自己的光荣：有一位寻求且审判的。}因为就这一点而言，在审判中将要显现的不是天主的形体，而是人子的形体，圣父自己将不会审判；因为按此所说：\BibleQuote{圣父不审判任何人，但将全部审判交给了圣子。}无论这是按照我们上面提到的那种说话方式而说的，其中说：\BibleQuote{祂给了圣子在自己内有生命，}这表示祂如此生了圣子；或是按照宗徒所说之方式：\BibleQuote{因此，天主也极其高举祂，赐给祂一个超乎万名之名；}——（这是关于人子说的，因为天主之子是就那人子而从死者中复活的；既然祂，以与圣父同等之天主的形体，从那里祂\Quote{空虚了自己}而取了仆人的形体，就在那同一仆人的形体中行动、受苦和接受，正如宗徒接着提到的：\BibleQuote{祂贬抑自己，听命至死，且死在十字架上。因此，天主也极其高举祂，赐给祂一个超乎万名之名，使天上、地上和地下的一切，因耶稣的名都屈膝，并且众口都承认耶稣基督是主，以光荣天主圣父。}——那么，这些话：\BibleQuote{祂将全部审判交给了圣子，}是按照这一种或那一种的说话方式而说的，从这节经文已充分看出，如果它们是按照那种方式说的，正如\BibleQuote{祂给了圣子在自己内有生命}，那么\BibleQuote{圣父不审判任何人}肯定就不会说了。因为就圣父生了与自己相等的圣子而言，祂是与圣子一起审判的。因此，就这一点而言，所说是：在审判中将要显现的不是天主的形体，而是人子的形体。不是因为那不审判的人将把全部审判交给圣子，因为圣子论及祂说：\BibleQuote{有一位寻求且审判的；}而是如此说：\BibleQuote{圣父不审判任何人，但将全部审判交给了圣子，}好比说：在审判生者死者时，没有人会看见圣父，但众人都将看见圣子；因为祂也是人子，所以连不敬天主的人也能看见祂，因为他们也要看见他们所刺透的那一位。\par

30. 然而，为避免我们似乎只是猜测而非明确证明这一点，让我们引用主自己的一句确定而清晰的说话，以此表明祂说\BibleQuote{圣父不审判任何人，但将一切审判全交给了圣子}的原因，\Emph{即}：因为祂要以人子的形状显现为审判者，这形状并非圣父的形状，而是圣子的形状；但也不是圣子与圣父相等的那个形状，而是祂小于圣父的那个形状；如此，在审判中，祂可能同时被善人和恶人看见。因为不久之后祂说：\BibleQuote{我实实在在告诉你们：那听我的话，并相信派遣我者，便有永生，不会进入审判，而是从死亡进入了生命。}这永生就是一种不属于恶人的看见。接着祂说：\BibleQuote{我实实在在告诉你们：时候将至，现在就是，死者要听见天主子的声音，听见的人将生活。}这专属于虔诚的人，他们如此听从祂的道成肉身，以至于相信祂是天主子，即他们如此接受祂——为了他们而被造得小于圣父、具有仆人的形状——也相信祂具有天主的形状、与圣父相等。于是祂接着强调这一点说：\BibleQuote{因为父怎样在自己内有生命，也照样赐给子在自己内有生命。}然后祂谈到祂自己的光荣的显现，祂将以此显现来审判；这显现将是恶人和义人都能看见的。因为祂继续说：\BibleQuote{并赐给祂权柄去执行审判，因为祂是人子。}我想没有什么比这更清楚的了。因为天主子既然与圣父相等，祂并非接受这执行审判的权柄，而是与圣父在隐秘中共同拥有它；但祂接受它，是为了让善人和恶人都能看见祂审判，因为祂是人子。因为人子的显现也将显示给恶人：因为天主形状的显现只显示给心地纯洁的人，因为他们将看见天主；即只显示给虔诚的人，祂向他们的爱情许诺这件事本身，即祂将显示自己给他们。因此，请看接下来的话：\BibleQuote{不要惊讶于此，}祂说。为什么祂禁止我们惊讶？除非实际上，每个不理解的人都会惊讶，即祂说圣父也赐给祂执行审判的权柄，是因为祂是人子；而原本人们可能预料祂会说，因为祂是天主子？但是，因为恶人不能看见天主子如同祂在圣子的形状中与圣父相等，然而又必须使义人和恶人都看见那位审判生者死者的审判者，当他们将在祂面前受审判时；祂说：\BibleQuote{不要惊讶于此，因为时候将至，凡在坟墓里的都要听见祂的声音，并要出来：行过善的，复活进入生命；行过恶的，复活进入定罪。}为此，祂必须因此接受那权柄，因为祂是人子，以便所有在复活时都能看见祂，在那所有人都能看见的形状中——但有些人是为了定罪，有些人是为了永生。永生是什么？不就是那不赐予恶人的看见吗？祂说：\BibleQuote{他们认识祢，唯一的真天主，并认识祢所派遣的耶稣基督。}他们又如何认识耶稣基督自己呢？除非作为那唯一的真天主，祂将要显示给他们；而不是像祂将要显示给那也要受惩罚的人那样，以人子的形状显示。\par

31. 按照那一种看见，即天主向心里洁净的人所显示的看见，祂是\BibleQuote{善}的；因为\BibleQuote{天主实在善待以色列，善待那些心地纯洁的人。}但是当恶人看见审判者时，祂在他们眼中不再是善的；因为他们不会在心中欢喜地看见祂，反而\BibleQuote{地上的一切种族都要因祂而哀号}，即被视为与所有恶人和不信者同类。为此，祂对那位称祂为善师、求问如何获得永生的人回答说：\BibleQuote{你为什么问我关于善的事？只有一位是善的，就是天主。}然而主自己在另一处称人是善的：\BibleQuote{善人，}祂说，\BibleQuote{从心里所存的善就发出善来；恶人，从心里所存的恶就发出恶来。}但是因为那人寻求永生，而永生在于那种默观，在其中天主被看见，不是为了惩罚，而是为了永远的喜乐；又因为他不明白与谁说话，以为祂只是人子，所以祂说：你为什么问我关于善的事？那就是，对于你所看见的这形态，你为什么问我关于善，并按照你所看见的称我为善师？这是人子的形态，是被取的形态，将在审判中显现的形态，不仅对义人，也对不虔敬的人；这形态的看见对恶人将不是善的。但还有我那个形态的看见，当我原先存在时，我以自己与天主同等不为抢夺，但为了取这形态，我空虚了自己。因此，那唯一的天主，圣父、圣子及圣神，除了为了义人不能被夺去的喜乐之外，不会显现；为了这未来的喜乐，那叹息者说：\BibleQuote{我有一件事祈求上主，我要寻求这事：使我一生住在上主的殿中，瞻仰上主的美善。}所以，那唯一的天主，祂自己，我说，唯独是善的，理由是没有人在忧伤和哀号中看见祂，而只在救恩和真喜乐中。如果你按后一种形态理解我，那么我是善的；但如果只按前一种，那么你为什么问我关于善的事？如果你属于那些\BibleQuote{要仰望他们所刺透的那一位}的人，那看见本身对他们将是恶的，因为它将是惩罚性的。那么，按这个意思，主说\BibleQuote{你为什么问我关于善的事？没有一个是善的，唯独天主}是可信的，据我所举的证据，因为那种对天主的看见，我们借以默观天主那不变且人眼不能看见的实体（这只应许给圣人；宗徒\Person{保禄}提到它，说是\BibleQuote{面对面；}宗徒\Person{若望}说：\BibleQuote{我们将相似祂，因为我们将看见祂实在怎样；}经上又说：\BibleQuote{我有一件事祈求上主，使我瞻仰上主的美善，}主自己也说：\BibleQuote{我要爱他，并要向他显示我自己；}并且唯独为此我们因信德洁净我们的心，使我们成为那些\BibleQuote{心里洁净的人，他们是有福的，因为他们要看见天主：}以及所有其他关于那看见的说法；凡以爱的眼目寻求它的人，可以在全部圣经中找到极其丰富的论述）——我说，唯有那看见是我们的至善，我们为获得它而被引导去正当地做一切事。但是那预言的人子的看见，当万民都聚集在祂面前，向祂说：\BibleQuote{主，我们几时看见祢饿了，或渴了，等等？}对不虔敬的人不是善的，他们将被投入永火，对义人也不是至善。因为祂继续召唤这些义人进入从创世以来为他们预备的国度。因为正如祂将对那些人说：\BibleQuote{离开我，进入永火里去}，同样对这些人说：\BibleQuote{我父所祝福的，你们来吧，承受为你们预备的国度。}那些人将进入永远焚烧，义人将进入永生。但永生是什么？除非\BibleQuote{他们认识祢，}祂说，\BibleQuote{唯一的真天主，和祢所派遣的耶稣基督？}但要认识祂，现在是在那荣耀中，祂对父说：\BibleQuote{在世界未有以前，我同祢所有的荣耀。}那时祂要把国度交给天主，即圣父，使良善的仆人进入他主人的喜乐，并使祂隐藏那些天主所保护的人，在祂圣容的隐密处，使他们免受人们的混乱，即那些将因听到这判决而惊慌的人；对这恶的听闻，\BibleQuote{义人必不恐惧}，只要他被保护在\BibleQuote{帐幕}里，即在天主教会真正的信德中，远离\BibleQuote{口舌之争}，即异端者的诡辩。但如果对主的话\BibleQuote{你为什么问我关于善的事？没有一个是善的，唯独天主}有其他解释，只要不因此认为圣父的实体比圣子的实体有更大的善，圣子就是圣言，万物借祂而造成；并且若这解释与纯正道理无抵触，我们可以安全地使用它，并且不仅一种解释，而是我们所能找到的许多解释。因为异端者被证明错误的程度越有力，就有越多的出路可以避开他们的陷阱。但现在让我们重新开始，致力于思考那仍然存留的部分。\par

\SourceRecord{Translated by Arthur West Haddan.；Nicene and Post-Nicene Fathers, First Series；Vol. 3.；Edited by Philip Schaff.；Buffalo, NY: Christian Literature Publishing Co.,；1887.；<http://www.newadvent.org/fathers/130101.htm>.}

% structure: p0001 source=h1 target=section reason=page-title

\section{\WorkTitle{论天主圣三（第二卷）}}

\EditorialNote{奥思定继续辩护天主圣三的平等；在论述圣子和圣神的被派遣以及天主的各种显现时，他证明被派遣者并不因此小于派遣者，因为一位派遣，另一位被派遣；但圣三在一切事上平等，在其本性上同样不变、不可见、无所不在，在每个派遣或显现中不可分割地工作。}\par

% structure: p0003 source=h2 target=subsection reason=relative-heading-rank; base=h2

\subsection{序言}

当人寻求认识天主，并按其人性软弱的限度将心思转向理解天主圣三时，他们必然从经验中体会到这任务的艰辛困难：无论是从心灵本身努力凝视不可接近之光，还是从圣经中所使用的多种多样表达方式（在我看来，这些表达方式无非是粗略地操练心灵，好使它在被基督的恩宠荣耀后能得尝甘甜）；——我说，这样的人，当他们驱散一切模糊之处，获得某种确定性之后，应当比任何人都更容易宽容那些在如此深奥奥秘的探索中犯错的人。但在犯错者身上，有两件事最难容忍：在真理未显明之前轻率地断言，以及在真理已显明之后固执地维护那轻率断言而成的谬误。但愿天主，如我所祈求并希望的，以祂善意之盾和仁慈恩宠的护佑，保护我远离这两种敌对发现真理和诠释神圣圣经的过失，我就不会懈怠于通过祂的圣经或通过受造物来探寻天主的本质。因为这两者都陈列于我们眼前，供我们默观，为使我们寻求祂、爱慕祂——祂默感了圣经，也创造了受造物。我也不会害怕发表我的意见，我宁愿被正直者审查，也不怕被乖戾者挑剔。因为爱德最为卓越而谦逊，它乐意接纳鸽眼；但对犬牙，则只能要么以最谨慎的谦卑避开，要么以最坚实的真理挫钝它；我宁愿被任何人责难，也不愿被犯错者或奉承者赞扬。因为爱真理者无需惧怕任何人的指责。因为指责者要么是敌人，要么是朋友。若是敌人辱骂，必须忍耐；若是朋友犯错，必须教导；若他教诲，则当聆听。但若犯错者赞扬你，那是肯定你的错误；若奉承者赞扬你，那是引诱你陷入错误。因此，\Quote{让义人，}\Quote{击打我，这乃是恩慈；让他责斥我，但罪人的油不可膏我的头。}\par

% structure: p0005 source=h2 target=subsection reason=relative-heading-rank; base=h2

\subsection{第一章——理解圣经中关于天主圣子的表达方式有两重规则。这些表达方式有三重种类。}

2. 因此，尽管我们最坚定地持守——关于我们的主耶稣基督，所谓正典规则——正如它通过圣经传播，并由博学而公教的圣经诠释者所证明的那样——即：按照祂所是的天主性形象，圣子被理解为与圣父同等；而按照祂所取的奴仆形象，祂小于圣父；在这个形象中，祂不仅小于圣父，也小于圣神；不仅如此，甚至小于祂自己——不是小于祂所曾是的那位，而是小于祂所是的那位；因为取了奴仆形象，祂并未失去天主性形象，正如我们在前一卷书中引用过的圣经见证所教导我们的：然而，圣经中有一些说法表述得如此含糊，以至于难以确定它们应按哪一规则来理解：是按我们理解圣子因取了受造物而较小的规则，还是按我们理解圣子并非较小而是与圣父同等，但祂是由圣父而来，即出自天主的天主、出自光的光的规则呢？因为我们称圣子是出自天主的\Emph{of}天主；而圣父只是天主，并非\Emph{of}天主。由此可见，圣子有另一位作为祂\Emph{of}的那位，并且祂是那位之子；但圣父没有一个圣子作为祂\Emph{of}的那位，而只有一位作为祂的父亲。因为每个儿子之所以是儿子，是\Emph{of}他的父亲，且是相对于他的父亲而言的；但没有一位父亲之所以是父亲，是\Emph{of}他的儿子，而是相对于他的儿子而言的。\par

3. 因此，在圣经中，有些关于圣父和圣子的说法是为了表明他们本质的合一与平等；例如：\BibleQuote{我与圣父原是一体}，以及\BibleQuote{他虽具有天主的形体，却没有以自己与天主同等为应当把持不舍的}，以及其他这类经文。另有一些说法则是为了显示圣子因取了奴仆的形体，即祂所取的可变及人性受造物的本质，而处于较低的地位；例如那句话说：\BibleQuote{因为我的圣父比我大}，以及\BibleQuote{圣父不审判任何人，但把一切审判权交给了圣子}。因为稍后祂接着说：\BibleQuote{并且赐给祂行审判的权柄，因为祂是人子}。此外，还有些说法是表明那时祂既非较低，也非同等，而只是指示祂是出于圣父；例如那句话说：\BibleQuote{因为圣父怎样在自己内有生命，也赐给圣子在自己内有生命}；以及另一处：\BibleQuote{圣子不能由自己做什么，唯有看见圣父做什么，才能做什么}。因为如果我们把这理解为因此圣子在被取受造形像中是较低者，那么必然得出圣父也必须行在水面上，或者用泥土和唾沫开启某个生来瞎眼者的眼睛，并做了圣子肉身显现时在人间所做的其他事，在圣子做之前；这样祂才能做那些事，因为圣子说祂不能由自己做什么，除非看见圣父做什么。然而，谁，即使是疯子，会这样想呢？因此，这些经文之所以如此表述，乃是因为圣子的生命与圣父的生命一样是不变的，然而祂又是出于圣父；圣父与圣子的工作是不可分的，然而这工作是由圣父赐给圣子，圣子从祂所是的圣父那里领受；圣子如此看见圣父，以至于在看见祂的同时，祂就是圣子。因为出于圣父，即由圣父所生，对祂来说，无非就是看见圣父；而看见圣父工作，无非就是与祂一同工作：但因此不是出于自己，因为祂不是出于自己。因此，那些\BibleQuote{祂看见圣父所做的事，圣子也同样做}的事，是因为祂是出于圣父。因为祂既不像是画家按照他人所画的画来画别的画那样做别的事；也不像是身体表达思想所想的相同字母那样以不同的方式做相同的事；而是\BibleQuote{无论什么事}，祂说，\BibleQuote{圣父所做的，圣子也同样做}。祂说了\BibleQuote{同样的事}，也说了\BibleQuote{照样}；因此圣父与圣子的工作是不可分且平等的，但它是从圣父到圣子的。因此，圣子不能由自己做什么，除非看见圣父做什么。因此，根据这一规则，圣经如此说话的意思，并非要表明一个比另一个小，而只是要表明哪一个是从哪一个来的，有些人却从中得出仿佛圣子被说成是较小的意思。在我们中间，那些更不学无术、最不熟悉这些事情的人，企图按照奴仆的形体来理解这些经文，从而误解它们，就感到困扰。为了避免这一点，必须遵守这一规则：圣子并非较小，而只是简单地指示祂是出于圣父，在这些话中，所宣告的不是祂的不平等，而是祂的出生。\par

% structure: p0008 source=h2 target=subsection reason=relative-heading-rank; base=h2

\subsection{第二章——某些关于圣子的说法应按任一规则理解。}

4. 因此，正如我开始所说的，在圣书中有些说法是如此表述的，以致于我们不清楚该把它们归入哪一规则：是归入那条因圣子取了受造物而地位较低的规则，还是归入那条表明虽然平等，但祂仍是出于圣父的规则。依我之见，如果这种情况如此不确定，以致于无法解释或辨别其真实含义，那么这些段落可以安全地按照任一规则来理解，例如：\BibleQuote{我的教训不是我的，而是那派遣我者的}。因为这句话既可以按照奴仆的形体来理解，如同我们在前一卷书中已经讨论过的；也可以按照天主的形体来理解，在那形体中，祂与圣父平等，然而出于圣父。因为按照天主的形体，正如圣子不是一位而祂的生命是另一位，而是生命本身就是圣子；同样，圣子不是一位而祂的教训是另一位，而是教训本身就是圣子。因此，正如经文\BibleQuote{祂赐给了圣子生命}只能被理解为：祂生了圣子，圣子是生命；同样，当说到祂赐给了圣子教训时，可以正确地理解为：祂生了圣子，圣子是教训。所以，当说到\BibleQuote{我的教训不是我的，而是那派遣我者的}，就要这样理解，如同说：我不是出于自己，而是出于那派遣我者。\par

% structure: p0010 source=h2 target=subsection reason=relative-heading-rank; base=h2

\subsection{第三章——某些关于圣神的事只能按一条规则理解。}

5. 因为关于圣神，并没有说\Quote{祂使自己空虚，取了奴仆的形状}；然而主自己却说：\Quote{但当那真理之神来到时，祂要引导你们进入一切真理。因为祂不是凭自己说话，而是把祂所听到的说出来，并要把将来的事告诉你们。祂要光荣我，因为祂要把由我所领受的，传告给你们。}除非祂紧接着又说：\Quote{凡父所有的一切，都是我的；为此我说：祂要把由我所领受的，传告给你们。}否则，也许有人会相信圣神是由基督所生，正如基督由父所生一样。因为祂论自己说过：\Quote{我的教训不是我的，而是那派遣我者的。}但论圣神却说：\Quote{因为祂不是凭自己说话，而是把祂所听到的说出来。}又说：\Quote{祂要光荣我，因为祂要把由我所领受的，传告给你们。}但是，因为祂已经说明了祂说\Quote{祂要由我所领受}的理由（祂说：\Quote{凡父所有的一切，都是我的；为此我说：祂要把由我所领受的}），所以，我们必须明白圣神拥有属于父的东西，如同子也拥有一样。然而，除非依据我们上面所说的：\Quote{但当护慰者，就是我从父那里要派遣给你们的，那从父出发的真理之神来到时，祂必要为我作证。}这怎么可能呢？因此，论到祂说\Quote{不是凭自己说话}，是因为祂从父出发；正如子说\Quote{子不能凭自己做什么，只有看到父做什么，才能做什么}并不意味子较小（因为祂说这话不是按奴仆的形状，而是按天主的形状，如我们已指出的，这些话并不表示祂较小，而是表示祂出自父），同样，论到圣神说\Quote{因为祂不是凭自己说话，而是把祂所听到的说出来}，也不意味着圣神较小；因为这话属于祂作为从父出发者。但是，既然子出自父，圣神也从父出发，为什么两者都不被称为子，两者都不说是受生，而只有前者被称为独生子，后者，即圣神，既不是子也不是受生（因为如果受生，则一定是子），我们将在别处讨论，如果天主恩准，并按其恩赐的程度。\par

% structure: p0012 source=h2 target=subsection reason=relative-heading-rank; base=h2

\subsection{第四章——子由父受光荣并不证明不平等。}

6. 但在这里，如果那些认为这也是支持他们主张的见证，以证明父大于子，因为子曾说\Quote{父啊，光荣我}的人，他们能够醒悟的话，就请醒悟吧。怎么？圣神也光荣祂。请问：难道圣神也大于祂吗？再者，如果圣神光荣子是因为祂将领受属于子的，并因此领受属于子的，因为凡父所有的一切也都是子的，那么显然，当圣神光荣子时，父也光荣子。由此可以觉察，凡父所有的一切，不仅属于子，也属于圣神，因为圣神能够光荣子，而父也光荣子。但如果光荣者大于被光荣者，那么请他们承认那些互相光荣的者是平等的。但圣经也记载子光荣父，因为祂说：\Quote{我在世上已光荣了祢。}的确，要提防的是，不要以为圣神比两者都大，因为祂光荣子，而父也光荣子，却没有记载祂自己被父或子所光荣。\par

% structure: p0014 source=h2 target=subsection reason=relative-heading-rank; base=h2

\subsection{第五章——子和圣神并不因为被派遣而较小。子也自己派遣自己。论圣神的派遣。}

7. 但人在此被证明错误后，转而说：派遣者大于被派遣者；因此圣父大于圣子，因为圣子常称自己被圣父所派遣；圣父也大于圣神，因为耶稣论及圣神曾说：\BibleQuote{圣父因我的名所要派遣的}；而圣神小于二者，因为圣父派遣祂，如我们所说，圣子也说：\BibleQuote{我若离去，我必派遣祂到你们这里来}。那么，我在这探究中首先问：圣子是从何处被派遣，往何处去？祂说：\BibleQuote{我从圣父那里来，来到了世界}。因此，被派遣就是从圣父那里来，并来到世界。那么，同一福音作者论及祂所说的话是什么意思呢？\BibleQuote{祂在世界，世界藉祂而造，世界却不认识祂}；接着又说：\BibleQuote{祂来到自己的地方}？当然，祂被派往祂所到之处；但如果祂因从圣父那里来而进入世界，那么祂既来到世界，又早已在世界。因此，祂是被派往祂已在那里之处。因为也请思量先知所记载的，天主说：\BibleQuote{我岂不充满天地吗}？如果这是指圣子而言（因为有些人认为圣子本人曾藉先知或是在先知内发言），那么祂除了被派往祂已在那里之处，还能往哪里去呢？因为说\BibleQuote{我充满天地}的那位，无所不在。但如果这是指圣父而言，那么没有祂的圣言和祂的智慧，祂能在哪里呢？这智慧\BibleQuote{大能地从地极到地极治理万物，和谐地安排一切}。但祂不可能没有祂的圣神而存在于任何地方。因此，如果天主无所不在，祂的圣神也无所不在。那么圣神也是被派往祂已在那里之处。因为那找不到地方可以躲避天主面容的人，他说：\BibleQuote{我何处能离开祢的神？或我何处能逃避祢的面容？}；他愿意使人体会天主处处临在，在上一节称祂的圣神；因为他说：\BibleQuote{我若升到天上，祢在那里；我若下到阴间，祢也在那里。}\par

8. 因此，如果圣子和圣神都被派往他们已在那里之处，那么我们必须探究，圣子或圣神的派遣该如何理解；因为唯独圣父，我们在任何地方都从未读到祂被派遣。关于圣子，宗徒这样写道：\BibleQuote{但时期一满，天主派遣了自己的儿子，生于女人，生于法律之下，为把在法律之下的人赎出来。} \BibleQuote{祂派遣了}，他说，\BibleQuote{自己的儿子，生于女人。} 藉着\Quote{女人}这个用词，哪位天主教徒不知道他并非意指失去童贞，而是按照\OriginalTerm{希伯来语用法}{Hebraism}，指性别差异？因此，当他说\BibleQuote{天主派遣了自己的儿子，生于女人}时，他充分表明圣子是因\Quote{生于女人}而被\Quote{派遣}的。因此，就祂由天主所生而言，祂在世界；就祂由玛利亚所生而言，祂被派遣而来到世界。此外，祂不能没有圣神而为圣父所派遣，不仅因为圣父在派遣祂时，即当祂使祂生于女人时，显然没有祂自己的圣神就不会这样做；而且因为福音中童贞玛利亚问天使时，得到的答复非常明确而清楚地说了：\BibleQuote{这事怎能成就呢？} \BibleQuote{圣神要临于你，至高者的能力要庇荫你。} \Person{玛窦}也说：\BibleQuote{她被查出因圣神怀孕。} 此外，在\Person{依撒意亚}先知书中，基督自己也被认为论及自己将来的来临说：\BibleQuote{现今主上主和祂的神派遣了我。}\par

9. 或许有人会逼我们说，圣子也是由自己派遣的，因为玛利亚的孕育和生产是天主圣三的工程，万有都是由祂们的创造行动所造的。他会接着说，如果圣子派遣了自己，那么圣父又怎能派遣祂呢？我首先回答他，请他告诉我，若他能的话，圣父是如何圣化祂的，如果祂圣化了自己？因为同一主说两者都要说：祂说：\BibleQuote{你们说那被圣父所圣化并派遣到世界上来的那个人，因为我说了\InnerQuote{我是天主子}，你们就亵渎。}而在另一处祂说：\BibleQuote{我为他们的缘故，圣化我自己。}我也问，圣父是如何交出祂的，如果祂交出了自己？因为宗徒保禄说了两者：他说：\BibleQuote{祂没有怜惜自己的儿子，反而为我们众人把祂交出来。}而在另一处，关于救主本人，他说：\BibleQuote{祂爱我，为我交出自己。}我想他会回答，如果他对这些事有正确的理解，因为圣父和圣子的旨意是一，它们的工作不可分。那么同样地，让他理解童贞女的降生成人和诞生，其中圣子被认为是派遣的，是由圣父和圣子不可分地同一运作而完成的；圣神当然不被排除在外，论到祂曾明确地说：\BibleQuote{她被发现因圣神怀了孕。}也许如果我们问天主是如何派遣祂的圣子的，我们的意思会更清楚地展开。祂命令祂来，祂遵命来了。那么，祂是要求了，还是仅仅暗示了？但无论是哪一种，肯定是由一个言语完成的，而天主的圣言就是天主子自己。因此，既然圣父藉一个言语派遣了祂，祂的被派遣就是圣父和祂的圣言两者的工作；所以同一圣子是由圣父和圣子派遣的，因为圣子本人就是圣父的圣言。因为谁会接受如此不敬的意见，以为圣父在时间内发出一个言语，以便永恒圣子因此被派遣，并在时期圆满时在肉身中显现？但确实，正是在那从起初就与天主同在，且是天主的圣言本身，即天主的智慧本身，超脱时间，在那智慧必须在肉身中显现的时刻。因此，既然无时间的开始，圣言在起初，圣言与天主同在，圣言是天主，所以是在圣言本身中无时间地，在那圣言将要成为血肉并居住在我们中间的时刻。当这时期圆满来到时，\BibleQuote{天主派遣了祂的儿子，由女人所生}，也就是说，在时间内被造，以便降生成人的圣言显现给人类；而同时在圣言本身中，超脱时间，有着这将要成就的时刻；因为时间的秩序是在天主永恒智慧中无时间的。既然圣子要在肉身中显现是由圣父和圣子两者所成就的，那么适宜地说，那在肉身中显现的一位是被派遣的，而那在肉身中不显现的一位派遣了祂；因为那些在外部经肉体的眼所完成的事，其存在来自灵性本性的内蕴结构（\OriginalTerm{内蕴结构}{apparatu}），因此适宜地被称为派遣。此外，祂所取的人性形式是圣子的位格，而非圣父的；为此，无形可见的圣父，连同与圣父一样无形可见的圣子，被说成派遣了同一圣子，藉着使祂成为可见。但如果祂成为可见的方式是使祂不再与圣父同属无形可见，即如果无形可见的圣言的实体通过变化和转变成了可见的受造物，那么圣子将被理解为只是被圣父派遣，而不与圣父一同派遣。但由于祂取了仆人的形体，而天主的不可改变的形式仍然保留，显然那在圣子中成为显明的，是由不显明的圣父和圣子所作成的；即，由无形可见的圣父，连同无形可见的圣子，同一圣子自己被打发成为可见的。那么，祂为什么说：\BibleQuote{我不是由自己来的？}我们现在可以说，这是按照仆人的形体说的，正如祂说：\BibleQuote{我不判断任何人。}\par

10. 因此，如果祂被说成是受派遣，是因为祂在身体受造物中外在地显现，而祂内在地在祂的灵性本性中总是对凡人的眼隐藏的，那么现在也容易理解圣神为何也被说成是受派遣的。因为在一个适当的时候，一个受造物的外在显现被做成了，在其中圣神能以可见的方式显示；无论是当祂以鸽子的形体降在主身上时，还是在祂升天后十天，在五旬节日，忽然从天上来了一阵响声，好像暴风刮来，又有分开的舌头如同火焰，显现出来，停在每个人身上。这个操作，被可见地展示，并呈现给凡人的眼，就被称为圣神的派遣；不是祂的实体本身显现了——祂自己也是无形可见、不可改变的，如同圣父和圣子一样——而是让人的心，被外在可见的事物触动，可能会从祂来临的时间性显现转向祂永远临在的隐藏永恒。\par

% structure: p0019 source=h2 target=subsection reason=relative-heading-rank; base=h2

\subsection{第六章——受造物被圣神所取，不如血肉被圣言所取那样。}

11. 因此，圣经中绝无记载父大于圣神，或圣神小于天主父，因为圣神将显现于其中的受造物，其取用方式不同于人子被取的方式——圣言本身的位格所应呈现的形体，并非为使祂拥有圣言（如其他圣善智慧之人拥有圣言），而是\Quote{超越祂的同伴}；当然不是指祂比他们拥有圣言更多，以致智慧远超其余，而是因为祂就是圣言本身。因为血肉中的圣言是一回事，成为血肉的圣言是另一回事；\Emph{即}人中的圣言是一回事，而成为人的圣言是另一回事。因为\Quote{血肉}在此代表\Quote{人}，正如经上说：\BibleQuote{圣言成了血肉}；又说：\BibleQuote{凡有血肉的，都要看见天主的救恩}。因为这里不是指没有灵魂和理智的肉体，而是\Quote{凡有血肉的}等同于\Quote{每一个人}。因此，圣神将显现于其中的受造物，其取用方式不同于从童贞玛利亚所取的血肉和人体。因为圣神并没有使鸽子、风或火受享真福，并将它们与自身永远合一于同一性体和\BibleQuote{样式}中。圣神的性体也不是易变无常的；以致这些事物并非由受造物所造，而是圣神自身先变成一种形态，后又变为另一种，如同水变为冰一样。但这些事物在应当显现的时候出现，受造物事奉造物主，并在祂的命令下被改变和转化——祂自身恒久不变——为要以适合的方式向必死之人指示和彰显祂。因此，虽然那只鸽子被称为神；论到那火，经上说：\BibleQuote{有散开好像火舌的舌头，出现给他们}并停在各人头上，他们就充满圣神，说起别国的话来，正如圣神赐给他们口才。这是要显示，圣神如同藉鸽子般，也藉那火被彰显。然而我们不能像称圣子既是天主又是人那样，称圣神既是天主又是鸽子，或既是天主又是火；也不能像我们称圣子为天主的羔羊那样，称圣神为鸽子或火。这羔羊不仅若翰洗者说\BibleQuote{请看，天主的羔羊}；若望福音作者也在默示录中看见被杀的羔羊。因为那先知性的神视，并非藉有形体之物显示给肉眼，而是在神魂中藉有形之物的属灵形象。但看见那鸽子和那火的人，都是亲眼所见。虽然关于那火，或许会有争议：究竟是用肉眼看见还是在神魂中看见，因为句子的形式。因为经文并未说\Quote{他们看见散开好像火舌的舌头}，而是说\BibleQuote{有……出现给他们}。但我们一般不把\Quote{有……出现给我}和\Quote{我看见}说成同一件事。在那些涉及形体的属灵神视中，通常两种说法都会用：\Quote{有……出现给我}和\Quote{我看见}；但在那些藉具体形体显示给眼睛的事物中，通常不说\Quote{有……出现给我}，而说\Quote{我看见}。因此，关于那火是如何被看见的，可能提出问题：是在神魂中如同在外面那样，还是真正在外面用肉眼看见？但对于那鸽子，据说它以形体降下，无人怀疑它是被肉眼看见的。再者，我们不能像称圣子为磐石那样——经上记着：\BibleQuote{那磐石就是基督}——也称圣神为鸽子或火。因为那磐石原已是受造之物，随其作用方式被称为基督之名，基督是其预象；正如雅各伯放在头下并傅油的石头，他用来象征主；又如依撒格背着献祭的木柴时，预象基督。一种特定的表意动作被附加于那些已经存在的事物上；但这些事物并非如鸽子和火那样，突然出现只为单纯地表示象征。因为鸽子和火，在我看来，更类似于向梅瑟显现的荆棘火焰，或百姓在旷野中跟随的云柱，或法律在西乃山颁赐时出现的雷轰闪电。因为这些事物的形体，正是为了表征某物而出现，随即消逝。\par

% structure: p0021 source=h2 target=subsection reason=relative-heading-rank; base=h2

\subsection{第七章——关于神性显现的一个疑问。}

12. 因此，圣神也被说成被派遣，是基于这些随时间出现的形体，为向人的感官指示和彰显祂（正如必须彰显祂那样）；然而，祂并不像圣子因处于奴仆形体而被说成小于父；因为那奴仆形体持续存在于圣子位格的合一中，但那些形体暂时出现，为要显示必须显示的事，然后便不复存在。那么，为什么父不被说成是因那些形体——荆棘火焰、云柱或火柱、山上的闪电，以及那时出现的其他同类事物——而被派遣呢？那时（我们从圣经见证得知）祂曾与父祖们面对面说话，如果祂便是藉那些受造物的样式和形体，以有形可见的方式向人的眼目显明自身的话？但如果圣子是藉它们被彰显，为什么祂在其后这么久，当祂生于女人时，才如宗徒所说\BibleQuote{但时期一满，天主就派遣了自己的儿子，生于女人}？既然祂以前也藉那些可变的受造物形体向父祖们显现，祂早已被派遣了。或者，除非圣言成了血肉，祂才应被正确地称为被派遣；那么为什么圣神被说成被派遣呢？圣神从未如此降生成人。但如果藉法律和先知中摆在我们眼前的那些可见事物，所彰显的既不是父也不是子，而是圣神，那么为什么现在祂被说成被派遣？祂以前也曾以这些样式被派遣过。\par

13. 在此困惑的探讨中，主助佑我们，我们须先询问：是圣父、或圣子、或圣神；或者是有时圣父、有时圣子、有时圣神；或是毫无位格区分，正如那独一真天主被论及的方式，即：天主圣三本身藉那些受造物的形态显现于众圣祖？其次，无论这些选项中哪一项被发现或被认为真实，即：是否只为这一目的而塑造了受造物，好让天主按祂当时认为适当的方式，显示于人的视线中；或是现存的天使被如此派遣，以致他们藉受造的形体从受造物那里取了一个有形的形态，以天主位格发言，为完成他们的职分，照各人所需；或者，依照造物主所赐予的能力，他们改变并转化自己的身体本身——他们并不受其支配，而是治理它如属下——使之成为他们想要的各种外貌，适合并适宜于他们的各种行动。最后，我们将分辨我们原要询问的问题，\Emph{即}：圣子和圣神以前是否也曾被派遣；若曾被派遣，那派遣与我们在福音中所读到的那一位有何区别；或者，事实上两者皆未被派遣，除非是当圣子由童贞玛利亚所生，或圣神以可见的形式显现时——不论是在鸽子上还是在火舌中。\par

% structure: p0024 source=h2 target=subsection reason=relative-heading-rank; base=h2

\subsection{第八章——整个天主圣三是不可见的。}

14. 因此，我们暂且不提那些过于肉性思想的人，他们认为天主的圣言与智慧——那\Quote{她虽居留在自己中，却使万物焕然一新}的存在，我们称之为天主的独生子——不仅是可变的，而且是可见的。因为这些人以其冒失胜过虔诚，带着极其迟钝的心智来探讨神圣的事物。因为灵魂是一种精神实体，且本身也是受造的，但只能由那创造万有、没有祂则一无所成的那位所造，它虽可变，却不可见；他们竟相信圣言本身和那创造灵魂的天主智慧本身也是如此；然而这智慧不仅如灵魂般不可见，且同样永恒不变，而灵魂则非如此。正是这智慧的不变性，当经上说\Quote{她虽居留在自己中，却使万物焕然一新}时所指的。然而这些人，试图用神圣圣经的见证来支撑他们摇摇欲坠的谬误，引用了宗徒保禄的话；并将那论及独一真天主（其中包含天主圣三本身）的话，理解为仅指圣父，而非圣子或圣神：\Quote{惟愿尊贵和荣耀归于永世的君王，那不朽坏、不可见、独一智慧的天主，直到永远。}以及另一处：\Quote{那蒙福的、独一的掌权者，万王之王，万主之主，独享不死不灭，住于不可接近的光中，没有人见过，也不能看见。}这些经文应如何理解，我想我们已经讲论得足够充分了。\par

% structure: p0026 source=h2 target=subsection reason=relative-heading-rank; base=h2

\subsection{第九章——反驳那些相信只有圣父是不朽且不可见的人。真理应以平和的研究来寻求。}

15. 然而，那些坚持认为这些经文只能理解为论及圣父，而非圣子或圣神的人，宣称圣子是可见的，并非由于祂取了童贞女的肉躯，而是早在先前祂本身就是可见的。因为他们说，祂本人曾向祖先们的眼睛显现。如果你对他们说：那么，无论圣子本身以何种方式可见，祂本身也就是以那种方式有死的；这样就显然得出，你也会将这句经文理解为仅指圣父，即\BibleQuote{唯独具有不朽性的；}因为如果圣子是由于取了我们的肉躯而成为有死的，那么你就得承认，祂也正是因为这肉躯而成为可见的：他们回答说，他们之所以说圣子是有死的，并非因为这肉躯；而是正如祂先前就是可见的，同样祂先前也是有死的。因为如果他们声称圣子是由于取了我们的肉躯而成为有死的，那么就不是唯独圣父没有圣子而具有不朽性，因为祂的圣言——万物藉祂而造——也具有不朽性。祂并没有因为取了有死的肉躯而失去祂的不朽性，因为就连人的灵魂也不可能与身体同死，主亲自说：\BibleQuote{不要害怕那些杀害身体却不能杀害灵魂的。} 或者，诚然，圣神也取了肉躯？关于这一点，他们无疑会困惑地说——如果圣子由于取了我们的肉躯而成为有死的，那么他们如何理解唯独圣父没有圣子和圣神而具有不朽性？既然圣神确实没有取我们的肉躯；如果祂没有不朽性，那么圣子就不是由于取了我们的肉躯而成为有死的；但如果圣神具有不朽性，那么\BibleQuote{唯独具有不朽性的}就不是仅指圣父而言。因此，他们认为自己能够证明圣子在降生成人之前本身也是有死的，因为可变性本身并非不当地被称为有死性，灵魂也依此说会死；不是因为灵魂改变并转变为身体或其他不同于自身的本质，而是因为，凡在其同一本质中现在与先前不同的，在它已不再是它原来的样子这一点上，就被发现是有死的。因为他们说，那时，在圣子由童贞玛利亚诞生之前，祂本人曾向我们的祖先显现，并非只以同一形式，而是以多种形式；先以一种形式，然后以另一种形式；祂本身既是可见的（因为当祂尚未取我们的肉躯时，祂的本质对于有死的眼睛是可见的），又是有死的，因为祂是可变的。同样，圣神也曾一度显现为鸽子，另一度显现为火。因此，他们说，以下经文不属于天主圣三，而是单独且恰当地仅指圣父：\BibleQuote{愿尊崇、荣耀归于永生的、不朽的、不可见的君王，独一的全智天主；} 以及\BibleQuote{唯独具有不朽性的，住在人不能接近的光中，没有人见过，也没有人能见到。}\par

16. 因此，撇开这些推理者，他们连灵魂的不可见的本体也无法认识，因而确实远不能认识独一天主的本体，即圣父、圣子和圣神，这本体始终不仅是不可见的，而且是不变的，从而拥有真实而真正的不朽性；我们否认天主（无论是圣父、圣子还是圣神）曾以肉体眼睛看见，除非是通过那服从祂自己权柄的受造物；我们，我说——准备接受纠正，如果我们受到兄弟般的正直精神的责备，也准备如此，即使被敌人挑剔，只要他说真话——在公教会的和平中，以平心静气的研究探询：在天主在肉身内来临之前，天主是否不加区别地向我们的祖先显现？或者是天主圣三中的某一位？或者轮流地，如同交替显现。\par

% structure: p0029 source=h2 target=subsection reason=relative-heading-rank; base=h2

\subsection{第十章——是否天主圣三不加区别地向祖先显现，或仅以天主圣三中的一位显现。天主向亚当显现。论同一次显现。向亚巴郎的显现。}

17. 首先，在\BibleBook{}{创世纪}所记载的、天主与祂用尘土所造的人说话的事上：如果我们撇开比喻意义，仅按字面相信叙述，那么，天主当时似乎是藉着人的形象与人说话的。这固然没有在书中明确说明，但阅读全篇的脉络皆有此意，尤其是在记载\Person{亚当}听见上主天主的声音，就是祂行走在晚凉中，\Person{亚当}就藏在园中的树木间；当天主说：\BibleQuote{亚当，你在哪里？}他回答说：\BibleQuote{我听见你的声音，就害怕，因为我赤身露体，便藏起来躲避你的面。}一段之中。因为我看不出这种天主的行走和交谈如何能按字面理解，除非祂显现为人。因为当天主被说成行走时，不能仅说祂发出声音，也不能说在某个地方行走的祂是不可见的；而\Person{亚当}也说，他躲避了天主的面。那么，祂是谁？是圣父，或是圣子，或是圣神？还是天主圣三本身毫无区别地以人的形态向人说话？其实，圣经上下文似乎没有一处表明位格之间的转换；而祂似乎仍在向那说\BibleQuote{有光}和\BibleQuote{有穹苍}以及诸日各事的第一人说话；我们通常将这位视为天主圣父，祂藉着圣言造了祂所愿造的一切。因为祂藉着祂的圣言造了万物，而这圣言，按正确的信德准则，我们知道祂是祂的独生子。因此，如果天主圣父曾向第一人说话，祂自己在晚凉中行走于乐园，而罪人从祂面前藏在园中的树木间，那么，我们为何不进而理解：也正是祂，通过那从属于祂的、可变而可见的受造物，向\Person{亚巴郎}和\Person{梅瑟}显现，并随祂所愿、按祂所愿的方式显现，而祂自己却留在祂自己内，留在祂的实体中，在其中祂是不可变、不可见的呢？但是，可能圣经是隐秘地转换了位格，在记载圣父说了\BibleQuote{有光}以及其余祂藉圣言所行之事后，转而暗示圣子向第一人说话；并未公开说明，而是让能理解的人领会。\par

18. 因此，凡有能力以心灵之眼洞察这奥秘，以致能清楚看出：或圣父亦然，或唯有圣子和圣神能通过可见的受造物向人眼显现的人，就让他，我说，去考察这些事——若他能做到，甚或以言辞表达和处理它们；但就此事本身而言，就圣经中天主与人说话的这段证词而论，依我之见，是无法发现的，因为连\Person{亚当}是否通常以肉眼看见天主，也不明显；尤其是当他们尝了禁果后，什么眼睛被开启了，这是一个大问题；因为尝之前，这些眼睛是闭着的。然而，我不敢冒昧断言，即使那段经文暗示乐园是个物质的地方，天主在那里行走也只能以某种形体方式。因为可以说，只有话语被造出来让人听见，而未见任何形象。也不因为记载说\BibleQuote{亚当躲避了天主的面}，就立刻推论他通常看见祂的面。因为，假如他自己看不见，但害怕被那声音的主人看见，并在祂行走时感受到祂的同在，那又如何呢？同样，\Person{加音}也对天主说：\BibleQuote{我要躲避你的面}；但我们并不因此被迫承认他惯常以肉眼在可见形象中看见天主的面，虽然他听过天主的声音，讯问并与他谈论他的罪。然而，天主当时用什么话语向人的外在耳朵说话，尤其是在向第一人说话时，既难发现，我们也不打算在此言论中讨论。但是，如果仅凭话语和声音被造出来，以便为那些最初的人带来天主某些感官上的临在，那么我不知道为何我不能在那里理解为天主圣父的位格，因为祂的位格也在那声音中显明，就是当\Person{耶稣}在山上三位门徒前荣耀显现时；以及当祂在受洗时鸽子降下时；以及当祂为祂自己的光荣向圣父呼求，并得到回答：\BibleQuote{我已光荣了，并要再光荣}时。并非声音可以不经圣子和圣神的工程而造出（因为天主圣三不可分割地工作），而是这样造出的声音只显明了圣父的位格；正如天主圣三从童贞\Person{玛利亚}造出了那人形，但那是圣子独有的位格；因为不可见的天主圣三造出了圣子独一的可见位格。也没有任何事禁止我们，不仅将那些对\Person{亚当}说的话理解为天主圣三所说，而且也将其视为显明天主圣三的位格。因为我们被迫只就圣父理解那句：\BibleQuote{这是我的爱子。}因为\Person{耶稣}既不能相信也不能被理解为圣神的儿子，甚至祂自己的儿子。而在发出\BibleQuote{我已光荣了，并要再光荣}那声音之处，我们承认那只是圣父的位格；因为这是对主的那句话的回应，主曾说过：\BibleQuote{父啊，光荣你的子}，这句话祂只能对天主圣父说，不能也对圣神说，因为祂不是圣神的儿子。但在这里，经上记载\BibleQuote{上主天主对亚当说}，没有理由不应理解为天主圣三本身。\par

19. 同样，在经上所记载的：\BibleQuote{上主曾对亚巴郎说：\InnerQuote{离开你的故乡、你的家族和你的父家，}}我们不清楚是仅有声音传到亚巴郎耳中，还是有某种事物也出现在他眼前。但稍后，经上说得较为清楚：\BibleQuote{上主显现给亚巴郎，说：\InnerQuote{我要将这地方赐给你的后裔。}}但那里也没有明确说出天主以什么形状向他显现，或是圣父、圣子、还是圣神向他显现。也许有人会认为向亚巴郎显现的是圣子，因为经上不是写着\Quote{天主显现给他}，而是写着\BibleQuote{上主显现给他}。因为圣子似乎被称为主，好像这名号是特属于祂的；例如宗徒说：\BibleQuote{虽有称为神的，或在天上，或在地上，就如那许多的神，许多的主，可是我们只有一个天主，就是圣父，万物都出于祂，我们也归于祂；也只有一个主耶稣基督，万物藉祂而有，我们也藉祂而有。}但是，既然发现在许多地方天主圣父也被称为主——例如：\BibleQuote{上主对我说：\InnerQuote{你是我的儿子，我今日生了你。}}以及：\BibleQuote{上主对我主说：\InnerQuote{你坐在我右边。}}而且，圣神也被发现被称为主，如宗徒所说：\BibleQuote{主就是那圣神；}接着，恐怕有人以为所指的是圣子，并因祂的非物质实体而被称为神，又继续说：\BibleQuote{主的圣神在那里，那里就有自由。}没有人怀疑主的圣神就是圣神。因此，这里也并不明显表明是圣三的哪一位向亚巴郎显现，或是天主圣三本身，这唯一的天主，经上说：\BibleQuote{你要敬畏上主你的天主，只事奉祂。}但在玛默勒橡树旁，他看见三个人，便邀请他们，热情款待，在他们用餐时伺候他们。然而，圣经在叙述开头并没有说三个人向他显现，而是说：\BibleQuote{上主向他显现。}然后，按次序说明上主如何向他显现之后，又加上了三个人的叙述，亚巴郎以复数邀请他们，后来又以单数如同对一个人说话；且如同一个人应许他从撒辣得一个儿子，\Emph{即}圣经称为上主的那一位，正如同一叙述开头所说：\BibleQuote{上主显现给亚巴郎。}于是他邀请他们，洗他们的脚，在他们离开时送他们，好像他们是凡人；但他说话如同与上主天主说话，无论是当儿子被应许给他时，或是当显示给他那即将临于索多玛的毁灭时。\par

% structure: p0033 source=h2 target=subsection reason=relative-heading-rank; base=h2

\subsection{第十一章 论同一显现}

20. 那段圣经既不容轻忽也不容草率对待。因为如果当时只有一个人显现，那些主张圣子在降孕于童贞女之前已用祂的本体显现的人，岂不立刻呼喊说，那就是祂自己？因为他们说，关于圣父，经上写着：\BibleQuote{独一的不朽天主。}然而我仍可以追问：在祂尚未取得我们的血肉之前，祂如何\BibleQuote{以人的形状显现}？既然祂的脚被洗，并且吃了地上的食物？这如何可能？那时祂还是\BibleQuote{以天主的形状存在，不以自己与天主同等为僭越}呢？请问，祂是否已经\BibleQuote{空虚自己，取了奴仆的形状，成了人的模样，以人的形状显现}？我们知道祂是在何时通过童贞女的诞生做到这一点的。那么，在祂尚未做到这一点之前，祂如何以一个单独的人向亚巴郎显现？或者，那形状不是真实的？如果向亚巴郎显现的是一个人，且那人被认为是天主子，我本来可以提出这些问题。但既然有三个人显现，且没有人在形状、年龄或能力上被说成比其余更大，我们为何不在这里理解为由可见的受造物可见地暗示了圣三的平等，以及三个位格中同一且相同的实体呢？\par

21. 因为，恐怕有人会认为三人中的这一个被暗示为更大，并应理解为是主、天主子，而其他两人是祂的天使；因为，虽然有三人显现，亚巴郎却向其中一人说话如同向主。圣经并未忘记通过一个反例来预先阻止这些将来的思考和意见：稍后它记载，有两个天使来到罗特那里，而那位义人——他理应被从索多玛的焚烧中救出——也向其中一人说话如同向主。因为圣经继续写道：\BibleQuote{上主在与亚巴郎说完话后便离开了；亚巴郎也回到了自己的地方。}\par

% structure: p0036 source=h2 target=subsection reason=relative-heading-rank; base=h2

\subsection{第十二章 考察向罗特显现}

\BibleQuote{但是黄昏时，有两个天使来到\Place{索多玛}。} 这里，我们必须更仔细地思考我开始阐述的内容。当然，\Person{亚巴郎}正在与三位说话，并用单数称呼那一位为\OriginalTerm{主}{Lord}。也许有人会说，他认出三位中的一位是主，而另外两位是他的天使。那么，圣经接着说下去的话是什么意思呢：\BibleQuote{主与\Person{亚巴郎}说完话就走了；\Person{亚巴郎}也回到自己的地方。黄昏时，有两个天使来到\Place{索多玛}吗？} 我们是否应当认为，三位中被认出是主的那位已经离去，并派与他同在的两个天使去毁灭\Place{索多玛}？那么，让我们看看接下来发生的事。据说，\BibleQuote{黄昏时，有两个天使来到\Place{索多玛}；\Person{罗特}正坐在\Place{索多玛}城门口。\Person{罗特}看见他们，就起身迎接他们，俯伏在地，说：\InnerQuote{看，诸位主上，请你们转身，进入你们仆人的家里。}} 这里清楚表明，一方面有两个天使，另一方面他们以复数受到款待的邀请，并被尊称为主，尽管他们可能被当作是人。\par

22. 然而，又有人反对说，除非他们被认出是天主的使者，否则\Person{罗特}不会俯伏在地。那么，为什么又向他们提供款待和食物，好像他们需要这样人的援助呢？但无论这里隐藏着什么，现在让我们继续我们已着手的工作。有两位显现；两者都被称为天使；他们受到复数的邀请；他复数地同两人说话，直到离开\Place{索多玛}。然后圣经接着说：\BibleQuote{当那两人领他们到外面时，他们说：\InnerQuote{逃命吧！不要回头看，也不要留在整个平原上；逃到山上，在那里你们将得救，免得被消灭。} \Person{罗特}对他们说：\InnerQuote{哦！不，我主；看，你的仆人在你眼前蒙恩。}}等等。他对他们说\BibleQuote{哦！不，我主}是什么意思——如果那位主已经离去并派了天使？为什么说\BibleQuote{哦！不，我主}而不说\BibleQuote{哦！不，我的主上们}？或者，如果他希望只对其中一位说话，为什么圣经说：\BibleQuote{\Person{罗特}对他们说：\InnerQuote{哦！不，我主；看，你的仆人在你眼前蒙恩。}}等等？我们在这里是否也可理解为复数形式中的两位，但当这两位被当作一位称呼时，就是一位本质同一的主\OriginalTerm{天主}{God}？但在这里我们理解为哪两位呢？——\OriginalTerm{圣父}{Father}和\OriginalTerm{圣子}{Son}，或\OriginalTerm{圣父}{Father}和\OriginalTerm{圣神}{Holy Spirit}，或\OriginalTerm{圣子}{Son}和\OriginalTerm{圣神}{Holy Spirit}？最后一种或许更合适；因为他们自称被派遣，这正如我们所说的圣子和圣神。因为我们在圣经中从未找到圣父被派遣。\par

% structure: p0039 source=h2 target=subsection reason=relative-heading-rank; base=h2

\subsection{第十三章——荆棘丛中的显现}

23. 但當\Person{梅瑟}被派遣去帶領以色列子民出離埃及時，經上記載主這樣向他顯現：\BibleQuote{現在\Person{梅瑟}牧養他岳父\Person{耶特洛}（\Place{米德楊}的司祭）的羊群；他領羊群往曠野的後邊，來到天主的山，即\Place{曷勒布}。上主的使者從荊棘叢的火焰中顯現給他；他觀看，見荊棘被火燒著，卻沒有燒毀。\Person{梅瑟}說：\InnerQuote{我要轉過去，看看這個大異象，為什麼荊棘沒有燒掉。}上主見他轉過去觀看，天主便從荊棘叢中呼叫他說：\InnerQuote{我是你父親的天主，\Person{亞巴郎}的天主，\Person{依撒格}的天主，\Person{雅各伯}的天主。}}此處也首先稱祂為上主的使者，然後稱之為天主。那麼，使者是\Person{亞巴郎}、\Person{依撒格}和\Person{雅各伯}的天主嗎？因此，祂可以正確地被理解為救主自己，關於祂，宗徒說：\BibleQuote{他們的祖先，並且按血統說，\Person{基督}也是從他們來的，祂是在萬有之上，永受讚美的天主。}所以，\BibleQuote{在萬有之上，永受讚美的天主}在此也被合理地理解為\Person{亞巴郎}的天主、\Person{依撒格}的天主和\Person{雅各伯}的天主。但為何祂先前在荊棘火焰中顯現時被稱為上主的使者？是因為祂是眾多天使之一，因著一種\OriginalTerm{安排}{economy}（或arrangement）而負起祂主子的\OriginalTerm{位格}{person}？還是祂取用了某種受造物，以便為當前的任務產生可見的顯現，並從中可聽見話語，藉著服從的受造物，以適宜的方式向人肉體感官顯示主的臨在？因為如果祂是天使之一，誰能輕易斷言，祂被賦予宣告聖子、聖神、天主聖父，還是整個天主聖三本身（祂是獨一無二的天主）的位格，以便祂可以說：\BibleQuote{我是\Person{亞巴郎}的天主，\Person{依撒格}的天主，\Person{雅各伯}的天主}？因為我們不能說天主之子是\Person{亞巴郎}的天主、\Person{依撒格}的天主、\Person{雅各伯}的天主，而聖父卻不是；也沒有人敢否認聖神或天主聖三本身（我們相信並理解為獨一真神）是那些祖先的天主。因為那不是天主的，就不是那些祖先的天主。此外，不僅聖父是天主（連異端者都承認），而且聖子也是（無論他們願意與否，都不得不承認，因為宗徒說：\BibleQuote{祂是在萬有之上，永受讚美的天主}），聖神也是（因為同一宗徒說：\BibleQuote{所以你們要在你們身上光榮天主}，他之前曾說：\BibleQuote{你們不知道你們的身體是聖神的宮殿，這聖神在你們內，是你們從天主那裡領受的嗎？}）這三者是一個天主，正如公教的正統信仰所相信的那樣；那麼，如果那位天使是其他天使之一，祂所負的位格是哪一位，或者不是某一位，而是整個天主聖三本身，就不夠明顯了。但如果受造物是為當前的任務而被取用，以便向人眼顯現、向人耳發聲，並被稱為上主的使者、主和天主，那麼此處的天主就不能理解為聖父，而是聖子或聖神。雖然我想不起聖神在任何其他場合被稱為使者，但從祂的工作可以理解；因為經上論到祂說：\BibleQuote{祂要告訴你們將來的事}；而希臘文的\Quote{天使}當然等同於拉丁文的\Quote{使者}；但我們在先知書中最清楚地讀到主\Person{耶穌基督}被稱為\BibleQuote{偉大的謀士天使}，而聖神和天主之子都是天使的天主和主。\par

% structure: p0041 source=h2 target=subsection reason=relative-heading-rank; base=h2

\subsection{第十四章——論雲柱和火柱中的顯現。}

24. 在以色列子民出離埃及時也記載著：\BibleQuote{上主在他們前面行，白天用雲柱為他們領路，夜間用火柱。白天沒有把雲柱從百姓面前撤去，夜間也沒有把火柱撤去。}此處，誰會懷疑天主是藉服從祂的受造的形體向可朽的人眼顯現，而不是藉祂自己的本體呢？但是，是聖父、聖子、聖神，還是整個天主聖三這獨一真神，並不清楚。在我看來，那裏也沒有加以區分，正如經上記載：\BibleQuote{上主的榮耀在雲中顯現，上主對梅瑟說：\InnerQuote{我聽到了以色列子民的怨言。}}等等。\par

% structure: p0043 source=h2 target=subsection reason=relative-heading-rank; base=h2

\subsection{第十五章——論西乃山上的顯現。在那次顯現中說話的是天主聖三，還是某一位特別的位格？}

25. 但如今论到西乃山上的云彩、声音、闪电、号角和烟雾，当时记载说：\BibleQuote{西乃山全笼罩在烟雾中，因为主在火中降临其上，烟上腾，如炉烟；营中众百姓都战栗；角声越长越大，\Person{梅瑟}说话，天主用声音回答他。}稍后，当律法在十诫中颁布以后，经文接着记载：\BibleQuote{众百姓看见雷霆、闪电、角声、冒烟的山。}再稍后：\BibleQuote{［百姓一见］就后退，远远站着；\Person{梅瑟}走近天主所在的浓云中，主对\Person{梅瑟}说}等等。对此我该说什么呢？除非没有一个人会疯狂到相信烟、火、云雾、黑暗以及任何这类的东西，就是天主的圣言和智慧（即基督）或圣神的实体。连亚略派异端者也从不敢说它们是圣父天主的实体。所以，这一切都是通过服役于造物主的受造物而造成的，并以适宜的\OriginalTerm{经济}{dispensatio}呈现给人的感官；除非，也许因为经上说：\BibleQuote{\Person{梅瑟}走近天主所在的云彩，}属血肉的念头就必定以为那云彩确实被百姓看见了，但云彩之内\Person{梅瑟}用肉眼看见了天主子，痴心妄想的异端者硬要说祂在祂自己的实体中被看见。的确，如果不仅天主的智慧即基督，就连你随意选择的任何一个人，无论多么聪明，都能用肉眼看见的话，那么\Person{梅瑟}可能用肉眼看见了祂；或者，因为关于以色列的长老们记载说：\BibleQuote{他们看见了以色列天主所立的地方，}又说：\BibleQuote{在祂脚下仿佛有蓝宝石铺成的平地，好像天空本身那样明净，}因此我们就得相信，天主的圣言和智慧在祂自己的实体中占据了一个尘世地方的空间，而祂\BibleQuote{刚强地伸展到地极，温和地管理万物；}并且天主的圣言（万物都是借着祂造的）是如此可变的，时而收缩，时而扩张；（愿主洁净祂忠实信徒的心，远离如此思想！）但所有这些可见可感的事物，正如我们常说，是通过受造的受造物呈现出来，以指示那不可见、不可理解的天主，不仅圣父，还有圣子和圣神，\BibleQuote{万物出于祂，万物通过祂，万物在祂内；}虽然\BibleQuote{天主那看不见的事，即祂永远的大能和神性，从创世以来，借着所造之物，就可以明白看见。}\par

26. 但就我们目前的工作而言，在\Place{西乃山}上，透过那些可怖地展示给必死之人感官的一切事物，我并未看出显现的是天主圣三，还是圣父、圣子或圣神各自说话。但若可以不经轻率断言，而对这段经文作审慎和犹豫的推测，若可以理解为涉及圣三中的一位，那我们为何不更愿意理解为圣神呢？因为在那里颁布的律法本身，被说是用天主的手指写在石版上，我们知道在福音中这称呼是指圣神。从宰杀羔羊和庆祝逾越节到\Place{西乃山}上开始发生这些事，算来是五十天；正如在主受难后，从祂复活算起五十天，然后天主子所预许的圣神降临了。而在祂降临的那次（我们在宗徒大事录中读到），有火焰似的舌头显现，分开落在他们每人身上：这与出谷纪吻合，其中写道：\BibleQuote{\Place{西乃山}全笼罩在烟雾中，因为主在火中降临其上；}稍后又说：\BibleQuote{主的荣耀显现，}他说：\BibleQuote{在以色列子民眼中，像烈火在山顶上吞灭一切。}或者，如果这些事之所以成就，是因为没有圣神（律法本身必须由祂书写），圣父或圣子就不能以那种模式临现，那么我们就确知天主在那里显现了，不是凭祂自己的实体（那实体是不可见、不变的），而是凭上面提到的受造物的显现；但至于圣三中某一个特别位格以适当标记显现，就我的理解能力所及，我们并未看到。\par

% structure: p0046 source=h2 target=subsection reason=relative-heading-rank; base=h2

\subsection{第十六章 — \Person{梅瑟}看见天主的方式}

26. 还有另一难题困扰着大多数人，即经上记载：\BibleQuote{主曾与梅瑟面对面地谈话，如同人对朋友说话一样；}而不久之后，同一位梅瑟说：\BibleQuote{所以如今我恳求祢，如果我在祢眼中蒙恩，求祢将祢自己清楚地显示给我，使我能看见祢，并能在祢眼中蒙恩，且使我晓得这民族是祢的人民；}又过了一会儿，梅瑟再次对主说：\BibleQuote{求祢显示祢的光荣给我。}那么这是什么意思呢？在这一切所行的事中，如上所述，人们以为天主凭祂自己的\OriginalTerm{实体}{substance}显现了；因此这些可怜的人相信天主子之所以可见，不是藉着受造物，而是凭祂自己；并且梅瑟进入云中，似乎正是为了这个目的而进入，即一片幽暗云雾确实可向百姓的双眼显示，但梅瑟在里面却能听见天主的话，如同看见祂的面；而且，如经上所载：\BibleQuote{主曾与梅瑟面对面地谈话，如同人对朋友说话一样；}然而，看哪，同一位梅瑟说：\BibleQuote{如果我在祢眼中蒙恩，求祢将祢自己清楚地显示给我？}无疑地，他知道自己看见的是属形体的东西，而他所寻求的是属神地看见天主的真面目。相应地，那藉言语形成的说话方式，已被调整得如同朋友对朋友交谈。然而，谁用肉体的眼睛看见天主圣父呢？那\Quote{在起初}、\Quote{与天主同在}、\Quote{是天主}的圣言，万物藉着祂而造成——谁用肉体的眼睛看见祂呢？智慧之神，又谁用肉体的眼睛看见呢？那么，\BibleQuote{求祢现在将祢自己清楚地显示给我，使我能看见祢}，除非是指\Quote{请将祢的\OriginalTerm{实体}{substance}显示给我}，还有什么意思呢？但如果梅瑟没有说过这话，我们确实只得尽量容忍那些愚昧的人，他们认为天主的\OriginalTerm{实体}{substance}藉着上述诸言诸行，被弄成可被他肉眼看见。但这里最明显地证明，这事并未\Emph{准许}给他，尽管他渴望如此；那么，谁敢说，藉着类似的形式，这些形式也曾有形可见地出现在他眼前，而那显现给一个必朽之人的，不是事奉天主的受造物，而是天主自己本身呢？\par

28. 另外，再加上此后主对梅瑟所说的话：\BibleQuote{你不能看见我的面，因为人见了我不能活。主又说：看，在我这里有个地方，你要站在磐石上：当我的光荣经过时，我把你放在磐石的缝隙中，用我的手遮掩你，直到我过去；然后我收回我的手，你将看见我的\OriginalTerm{背面}{back parts}；但我的面却不得看见。}\par

% structure: p0049 source=h2 target=subsection reason=relative-heading-rank; base=h2

\subsection{第十七章——如何看见了\OriginalTerm{天主的背面}{Back Parts of God}。基督复活的信德。惟独天主教会是从那里看见\OriginalTerm{天主的背面}{Back Parts of God}的地方。以色列人曾看见\OriginalTerm{天主的背面}{Back Parts of God}。认为只有天主圣父从未被圣祖们看见，乃是一种轻率的意见。}

未尝不合理地，普遍认为从我们的主耶稣基督的位格中预表了：祂的\BibleQuote{背}\BibleRef{出 33:23}应被视为祂的肉躯，祂由童贞女所生，受死并复活；或称之为背，是由于死亡的后期性，或是因为几乎在世界末期，即晚期，祂屈尊取了它；但祂的\Quote{面}是天主的本性，祂\BibleQuote{持守与天主同等不是僭越}\BibleRef{斐 2:6}，这当然无人能见而仍活着；或是因为在此生之后，我们离主而居，可朽坏的身体压迫灵魂，我们将看见\BibleQuote{面对面}\BibleRef{格前 13:12}，如宗徒所说——（因为在圣咏中论到此生说：\BibleQuote{凡人实在全是虚幻}\BibleRef{咏 39:6}；又说：\BibleQuote{在祢面前，凡活着的人没有一个可称义}\BibleRef{咏 143:2}；在此生，按若望所说：\BibleQuote{将来如何，还没有显明}\BibleRef{若一 3:2}，他说：\BibleQuote{当祂显现时，我们将相似祂，因为我们将看见祂实在怎样}\BibleRef{若一 3:2}，他当然意指在此生之后，当我们偿还了死亡的债务，并得到了复活的应许）；——或者是因为即使在现在，我们以何种程度属灵地理解创造万有的天主的智慧，也以同样的程度死于肉性的情欲，以致看这个世界于我们已死，我们也向这个世界而死，并说宗徒所说的：\BibleQuote{世界于我已被钉在十字架上，我于世界亦然}\BibleRef{迦 6:14}。因为正是论及这种死亡，他也说：\BibleQuote{所以，如果你们与基督同死，为什么还像活在世俗中，服从那些规条呢？}\BibleRef{哥 2:20}。因此，并非无故，没有人能看见\Quote{面}，即天主的智慧本身显现，而仍活着。正是这显现，每一个全心、全灵、全意爱天主的人所叹息渴慕的；也尽己所能爱邻人如同自己的人所建立的；这两条诫命统摄全部法律和先知。这也在梅瑟身上预示了。当他因对天主特别炽热的爱而说：\BibleQuote{如果我在祢眼中蒙恩，求祢将祢自己明明地指示给我，使我在祢眼中蒙恩}\BibleRef{出 33:13}；他紧接着因爱邻人的缘故说：\BibleQuote{叫我知道这民族是祢的百姓}\BibleRef{出 33:13}。因此，正是那\Quote{显现}吸引每个理性灵魂，越纯洁就越热切；灵魂越纯洁就越上升至灵性事物；越上升至灵性事物就越死于肉性事物。但我们既远离主，凭信德行事，而非凭目睹，就应当凭信德看见基督的\Quote{背}，即祂的肉躯，就是站在信德的坚固根基（磐石所象征的）上，从如此安全的瞭望台（即在天主教会中，关于它曾说：\BibleQuote{在这磐石上，我要建立我的教会}\BibleRef{玛 16:18}）观看。我们越确认基督爱我们有多深，就越确定地爱那热切渴望看见的基督的面。\par

29. 但在肉躯本身，对祂复活的信德拯救我们并使我们成义。因为，他说：\BibleQuote{你若心里相信，天主使祂从死者中复活，你必得救}\BibleRef{罗 10:9}；又说：\BibleQuote{祂被交付，是为了我们的过犯，并且祂复活，是为了我们的成义}\BibleRef{罗 4:25}。如此，我们信德的赏报就是我们的主身体的复活。因为就连祂的仇敌也相信那肉躯在祂苦难的十字架上死了，但不相信它又复活了。我们最坚定地相信这一点，就如从磐石的坚固处凝视它：由此我们以确定的希望等待义子的名分，即我们身体的救赎；因为我们盼望在基督的肢体中，即在我们自身中，那借着纯正的信仰我们承认在祂内（我们的元首内）是完美的。因此，祂不愿意祂的\BibleQuote{背}被看见，除非当祂经过时，好使祂的复活得以相信。因为希伯来语中\OriginalTerm{逾越节}{Pascha}译作\Quote{逾越}。因此若望福音者也说：\BibleQuote{在逾越节庆日前，耶稣知道祂的时辰已到，祂要从这世界过渡到父那里去}\BibleRef{若 13:1}。\par

30. 但那些相信这事，却不是在天主教会内，而是在某种分裂或异端中相信的人，并不能从\Quote{祂旁边的那个地方}看到主的后面。因为主所说的\BibleQuote{看，在我旁边有一个地方，你要站在磐石上}是什么意思呢？有什么尘世的地方是主\Quote{旁边}的，除非那是属神地接触祂的\Quote{祂旁边}？因为哪有一个地方不是主\Quote{旁边}的呢？祂\BibleQuote{以大力从地极达到地极，从容地治理万物}；关于祂曾说过：\BibleQuote{天是我的宝座，地是我的脚凳}；祂又说：\BibleQuote{你们要为我建筑什么样的房屋？哪里是我安息的地方？这一切不是我的手所造的吗？}但显然，天主教会本身就被理解为\Quote{祂旁边的那个地方}，人在其中站在磐石上，有益地看到\OriginalTerm{主的巴斯卦}{Pascha Domini}，即主的\Quote{经过}，以及祂的后面，即祂的身体——那位相信祂复活的人。祂说：\BibleQuote{你要站在磐石上，当我的光荣经过的时候}。事实上，在主的光荣经过以后——即主在光荣中复活并升到父那里，我们立刻坚定地站在磐石上。伯多禄自己那时也坚定了，以致他满怀信心地宣讲祂，而在坚定以前，他曾因恐惧三次否认了祂；虽然，事实上，他早已在预定中被安置在磐石的守望台上，但主的手仍覆在他身上，使他不能看见。因为他要看到祂的后面，而主还没有\Quote{经过}，即从死亡到生命；祂还没有因复活而受光荣。\par

31. 因为关于出谷纪中接着的话：\BibleQuote{我要在你经过时用我的手遮掩你，等我过去后，我要挪开我的手，你将从后面看见我}，许多以色列人——梅瑟当时是他们的预像——在主复活后相信了主，好像祂的手从他们的眼前挪开，他们现在看到了祂的后面。因此，福音作者也提到依撒意亚的预言：\BibleQuote{使这百姓的心变迟钝，使他们的耳朵沉重，并闭上他们的眼睛}。最后，在圣咏中，不无理由地被认为是在他们口中所说的话：\BibleQuote{因为日日夜夜，你的手沉重地压在我身上}。\Quote{白日}大概是指祂显明奇迹时，他们却不认识祂；\Quote{黑夜}是指祂在苦难中死去时，他们更加确信祂像凡人一样被剪除而消灭了。但是，当祂已经经过，以致祂的后面被看见，宗徒伯多禄向他们宣讲基督必须受难并复活，他们心中因悔改的忧伤而被刺透，使受洗者中所发生的事应验了那篇圣咏开头的话：\BibleQuote{罪恶得赦免、过犯被遮掩的人是有福的}。因此，在说了\BibleQuote{你的手沉重地压在我身上}之后，主仿佛经过，如今祂挪开祂的手，祂的后面被看见了，随后就听见一个忧伤、忏悔并因着对主复活的信德而获得罪赦的声音：他说：\BibleQuote{我的精力变成了夏日的干旱。我向祢承认了我的罪，我没有隐瞒我的过犯。我说：我要向上主承认我的过犯，祢就赦免了我罪孽的过犯}。因为我们不该如此被肉身的黑暗所蒙蔽，以致认为天主的容貌固然看不见，但祂的后面却是可见的，因为两者都在仆人的形像中可见地出现了；但决不可让我们对天主本身有任何这类想法；决不可让我们认为天主的圣言和天主的智慧有一面是脸，另一面是背，像人的身体那样，或因其显现或运动而在地方或时间上有所改变。\par

32. 因此，如果在出谷纪中所说的那些话，以及所有那些有形体的显现中，主耶稣基督被显示出来；或者，如果有些情形下基督被显示，正如本段文字的考量所说服我们的，而在其他情形下则是圣神，正如我们上文所提醒的；无论如何，绝不能得出这样的结论，即天主圣父从未以任何这类形式显现给先祖们。因为在那时代发生了许多这样的显现，其中既没有明确指出是圣父，还是圣子，还是圣神；但通过某些非常可能的解释，仍有某些暗示，因此若说天主圣父从未以任何可见的形式显现给先祖或先知们，那就太轻率了。因为那些无法就天主圣三的合一而理解诸如\BibleQuote{愿尊崇和光荣归于永生的君王，那不可见、不朽坏、独一全智的天主}和\BibleQuote{从来没有人看见过祂，也不能看见祂}这些经文的人，便产生了这种见解。这些经文，凭着对那个实体本身——那至高的、至神圣的、不变的天主——的正确信德来理解，也就是圣父、圣子、圣神同是那独一的天主。但是那些显现是通过可变的受造物，服从于不变的天主而实现的，并没有正确地按祂的所是将天主显示出来，而是按照适合各种环境的原因和时期的暗示而显示。\par

% structure: p0055 source=h2 target=subsection reason=relative-heading-rank; base=h2

\subsection{第十八章 达尼尔的神视}

我不明白这些人如何理解，万古常存者曾显现给达尼尔，祂是人子（祂为了我们屈尊成为人子）从祂那里领受王国的；即那位在圣咏中对祂说：\BibleQuote{祢是我的儿子，我今日生了祢；祢向我请求，我就将外邦人赐给祢作产业}，并且\BibleQuote{将万有置于祂的脚下}的那一位。然而，如果圣父赐予王国和圣子领受王国两者均以形体显现给达尼尔，这些人怎能说圣父从未显现给先知们，因此只有那位无人看见过，也无人能看见的才应被理解为不可见的呢？因为达尼尔如此告诉我们：\BibleQuote{我观望，直到宝座设立，万古常存者坐于其上，祂的衣服洁白如雪，头发如纯净的羊毛；祂的宝座有如火焰，其轮毂如同燃烧的火；一道火河从祂面前涌出并流出；成千上万服事于祂，亿万众侍立于祂面前；审判已设立，案卷已展开}等等。稍后他又说：\BibleQuote{我在夜间的异象中观望，看哪，有一位像人子的，乘着天上的云彩而来，来到万古常存者面前，被引到祂跟前。于祂被赐予统权、荣耀和王国，使各民族、邦国和语言都事奉祂；祂的统权是永远的统权，必不消逝，祂的王国也不被毁灭。}看，圣父赐予，圣子领受一个永恒的王国；两者都在预言者的眼前以可见的形体显现。因此，相信天主圣父也曾以那种方式显现给终有一死的人，并非不合宜。\par

除非或许有人会说，圣父之所以不可见，是因为祂显现在一个做梦者的视线中；但圣子和圣神之所以可见，是因为梅瑟在清醒时看见了这一切；仿佛梅瑟真的用血肉的眼睛看见了天主的圣言和智慧，或者甚至那使肉身活跃的人的精神可以被看见，或者甚至那被称为气的有形之物——更何况那超越所有人灵和天使之思想的天主圣神，因着神圣本体的不可言喻的卓越，如何能被看见呢？或者有谁能猛然跌入这样的错误，竟敢说圣子和圣神对清醒的人也是可见的，而圣父只对那些做梦者才可见？那么，他们如何理解唯独对圣父所说的\BibleQuote{无人看见过，也无人能看见}？当人睡觉时，他们就不算人了吗？或者，那位能造出身体的相似物以通过做梦者的异象来显示自己的，难道不能也造出同样的有形之物来向清醒者的眼睛显示自己吗？然而，祂自己的本体——祂就是祂所是的——不能通过任何身体的相似物向睡着的人显示，也不能通过任何身体的外貌向清醒者显示；但这不仅限于圣父，也适用于圣子和圣神。当然，对于那些因清醒者的异象而相信不是圣父，而是只有圣子或圣神曾以有形的样式出现在人的肉眼前之人——暂且不提浩瀚的圣经及其多种解释（没有任何健全理性的人应当断言圣父的位格从未以任何有形的外貌出现在清醒者的眼前）——但正如我所说，暂且不提这一点，他们对我们先祖亚巴郎又作何解释呢？当圣经先声明\BibleQuote{主显现给亚巴郎}后，不是一位、也不是两位，而是三个人显现给他；其中没有一位说是比其他更为突出，没有一位比其他闪耀着更大的荣耀，或行动更有权威。\par

因此，既然在我们那三重划分中，我们决定首先探究：是圣父、或圣子、或圣神？或者是有时圣父、有时圣子、有时圣神？还是如所说的，不分位格的那独一的天主，即天主圣三本身，藉着那些受造物的形象显现给先祖们？既然我们已经查考了圣经中我们所能找到的足够多的章节，那么，对我所能判断的而言，对神圣奥秘的谦虚谨慎的思考引导我们得出这样的结论：除非上下文给叙述附带了某些关于该主题的可能的暗示，否则我们不应鲁莽地断言天主圣三的哪一位以某种身体或身体的相似物显现给这位或那位先祖或先知。因为那本性、或实体、或本质（或无论我们称那构成天主的、无论是什么的东西为哪个名字），本身不能被有形地看见；但我们必须相信，藉着服从于祂的受造物，不仅圣子或圣神，而且圣父也可以藉着有形的形态或相似物向终有一死之人的感官显示自己。既然情况如此，为使这第二本书不致过于冗长，让我们在接下来的书中思考余下的问题。\par

\SourceRecord{Translated by Arthur West Haddan.；Nicene and Post-Nicene Fathers, First Series；Vol. 3.；Edited by Philip Schaff.；Buffalo, NY: Christian Literature Publishing Co.,；1887.；<http://www.newadvent.org/fathers/130102.htm>.}

% structure: p0001 source=h1 target=section reason=page-title

\section{\WorkTitle{论天主圣三（第三卷）}}

\EditorialNote{此處討論的是前卷所述以形體方式出現的天主顯現，問題在於：是否僅有一個受造物被形成，為使天主按祂每時認為合適的方式顯現於人目；還是早已存在的天使被派遣，以天主位格說話，藉由從形體受造物中取用形體外貌，或按造物主賜予他們的能力，將自身軀體變化為適合特定行動的任何形狀；而天主的本體自身從未被直接得見。}\par

% structure: p0003 source=h2 target=subsection reason=relative-heading-rank; base=h2

\subsection{序言——奧斯定為何寫論天主聖三。他要求讀者什麼。前卷已言之事。}

1. 我願那些願意相信的人相信，我寧願在閱讀上勞苦，而不願口述供他人閱讀。但若有人不願相信此事，卻有能力且願意嘗試，那麼，請將從閱讀中可能獲得的任何答覆，無論是針對我自己的疑問，或是針對他人因我在基督內的職務以及我為堅固信德、抵禦屬血氣和屬人性之人的錯謬所懷的熱忱而必須承受的質問，全都賜予我；然後，他們便會看到我是多麼樂於免除這勞苦，甚至多麼歡喜地讓我的筆休假。然而，若我們在這些主題上所讀到的內容，在拉丁文中要麼闡述不足，要麼根本無從覓得，或至少我們不易覓得，而我們對希臘文又不熟悉，不足以勝任閱讀和理解論述此類主題的書籍——按我們已有的少數譯本判斷，此類著作中無疑包含我們所能得益的一切；然而，我無法抗拒我的弟兄們，他們按那使我成為他們僕人的律法要求我，尤其要以我的口和筆——在我內，愛是這兩者的御者——在基督內服務於他們可讚美的勤學；同時我亦自認，藉由寫作，我學會了許多先前不知道的事。若果如此，那麼我的這項勞苦，對任何懶惰或極有學問的讀者而言，都不應顯得多餘；而對許多忙碌和許多無學問的人——後者中包括我自己——則在相當大的程度上是必需的。因此，憑著我們已讀過的其他人在此主題上的著作所給予的極大支持和幫助，我承擔起探究和討論一切關於天主聖三——獨一至高而至善的天主——我認為可以敬虔探究和討論之事；祂親自鼓勵我探究，並幫助我討論此事；如此，若沒有其他同類著作，則那些願意且有能力的讀者可以擁有並閱讀一些東西；但若已有此類著作，那麼這類著作越多，就越容易找到它們。\par

2. 確實，如同在我所有的著作中，我不僅希望有虔誠的讀者，也希望有自由的糾正者；在這次特別重要的探究中，我更如此期望，巴不得質問者有多少，探究者也有多少。但正如我不願我的讀者被束縛於我，我也不願我的糾正者被束縛於自己。願前者愛我不超過愛公教信仰，願後者愛自己不超過愛公教真理。我對前者說：不要情願將我的著作等同於正典聖經；但在正典聖經中，即使你發現了先前所不信的，也當毫不猶豫地相信；而在我的著作中，除非你確切理解了你先前未確信之事，否則不要堅持持守。我對後者說：不要情願以你自己的意見或辯論來修改我的著作，而要根據天主聖言或無可反駁的理性。你若在其中領悟了任何真理，它在其中並不使它成為我的；你通過理解和愛它，讓它既成為你的，也成為我的；但若你指證了任何錯誤，即使它曾經是我的（因我犯錯），如今因你避免它，讓它既不屬於你，也不屬於我。\par

3. 那么，让这第三卷从第二卷达到的地方开始。因为我们达到这一点后，我们想要表明：圣子并不因此而小于圣父，因为圣父派遣，圣子被派遣；圣神也不因此而小于两者，因为我们在福音中读到祂由两者派遣；于是我们着手探究：既然圣子被派遣到祂已所在的地方（因为祂来到世界，并且\BibleQuote{曾在这世界内}），既然圣神也被派遣到祂已所在的地方（因为\BibleQuote{上主的神充满了世界，且包罗万有的通晓声音}），那么，主是否因此而被说成是\Quote{被派遣的}，因为祂生于肉身而不再隐藏，并且，可以说，从圣父怀中出来，以\OriginalTerm{仆人的形体}{forma servi}出现在人眼前；圣神是否也因此而被说成是\Quote{被派遣的}，因为祂也以鸽子的形体可见，并以分开的火舌显现；从而，当说到祂们时，被派遣意味着从精神的隐藏处，以某种有形的形式出现于凡人眼前；因为圣父没有这样做，所以祂只被称为派遣者，而不被称为被派遣者。我们接着探究：如果圣父自己也曾通过那些出现在古人眼前的形体形式显现，为何祂有时不被说成是被派遣的？但如果圣子在这些时候显现，为什么祂要在很久之后、时期满了、应由女人所生时，才被说成是\Quote{被派遣的}？既然，实际上，祂以前也曾在那些形式中形体显现时被派遣过？或者，如果只有在\OriginalTerm{圣言成了血肉}{Verbum caro factum}时祂才被正确地称为\Quote{被派遣的}，那么为什么圣神也被记载为\Quote{被派遣的}，而祂从未有过这样的降生成人？但如果这些古代显现中显现的既不是圣父，也不是圣子，而是圣神，那么为什么现在祂也被说成是\Quote{被派遣的}，而以前祂也以这些不同方式被派遣过？接着，我们将主题细分，以便最仔细地处理，我们提出了三重问题，其中一部分已在第二卷中解释，还有两个部分剩下，接下来我将讨论它们。因为我们已经探究并确定：在那些古代的有形体形式和异象中显现的，不仅不是圣父单独，也不是圣子单独，也不是圣神单独，而是要么无区别地显示主天主（即理解为此圣三本身），要么显示圣三的一位，无论叙述的文字通过上下文暗示而指出的那一位。\par

% structure: p0007 source=h2 target=subsection reason=relative-heading-rank; base=h2

\subsection{第一章——对此需言之事。}

4. 那么，让我们现在按顺序继续探究。因为在那个划分中的第二个标题下，出现了问题：是否受造物是特为那项工作而形成的，在其中天主以祂当时认为合适的方式，可能向人眼显现；还是已存在的天使被这样派遣，以致以天主的名义说话，从有形的受造物中为他们的职务而取用有形的外貌；或者他们改变并转变自己的身体本身（那身体不是他们受其支配的，而是他们如同支配属己之物般管理它）为任何他们愿意的、适合并适宜于他们行动的形式，按照造物主赐给他们的能力。当这一部分问题被探究完，在天主许可的情况下，最后，我们将要看我们开始时提出的问题，\Emph{即}，圣子和圣神以前是否也被\Quote{派遣过}；如果是，那么那种派遣与我们在福音中读到的那一种有什么区别；或者是否两者都没有被派遣，除非当圣子由童贞玛利亚所生，或者当圣神以可见的形式显现，无论是作为鸽子还是火舌之时。\par

5. 然而，我承认，探究以下问题超出了我的目的所能承担：天使们是否暗中藉他们身上仍存留的他们身体的精神品质，从低级的、更具物质性的元素中取用一些东西，这些东西经他们调整后，他们可以像衣服一样改变并转变为任何他们愿意的有形体显现，而且这些显现本身也是真实的，就像主将真实的水变为真实的酒一样；或者他们是否将自己的身体本身转变为适于特定行动的任何形态。但无论哪一种，对当前问题都无关紧要。虽然我不能通过实际经验理解这些事，因为我是人，而天使们做这些事并比我更了解它们，\Emph{即}，我的身体能被我的意志运作改变到何种程度，无论这是通过我对自身的经验，还是我从他人那里收集的；然而，在这里没有必要说明我应当相信这些替代方案中的哪一个，依据神圣圣经的权威，以免我被迫证明它，从而我的论述在关于当前问题无关的题目上变得太长。\par

6. 那么我们目前的探究是：天使们当时是否既是将那些有形体显现显于人眼的中介，也是将那些话送入他们耳中的中介，当可感觉的受造物本身，服从造物主的命令，当时被转变为任何需要的东西；正如在 \BibleBook{}{智慧篇} 中所记载的：\BibleQuote{因为受造物服事祢，祢是造物主，为惩罚不义者而增强它的力量，为信赖祢的人而减弱它的力量。因此，甚至那时它也改变成各种样式，并顺从祢的恩宠，这恩宠按渴望祢者的愿望滋养万有。}因为天主旨意的能力通过精神的受造物甚至达到有形的受造物可见可感的效果。因为全能天主的智慧——它\BibleQuote{从一端强有力地抵达另一端，并温和地治理万有}——哪里不运作祂愿意的事呢？\par

% structure: p0011 source=h2 target=subsection reason=relative-heading-rank; base=h2

\subsection{第二章——天主的旨意是一切有形变化的高级原因。这由例子说明。}

但是，在物体的变化与变动中，有一种自然秩序。这一秩序本身也服从天主的命令，但由于它持续不断，已不再令人惊奇；例如，在天上、地上或海中，那些在极短或至少不长的时段内发生变化的事物，诸如升起、落下或外观的随时而变；而另一些事物虽出于同一秩序，但因间隔时间较长而不那么为人熟知。这些事物尽管众多愚人惊奇，但为探究现世的人所理解，且随世代推移，它们因重复次数越多、为更多人所知而越不令人惊奇。诸如日蚀、月蚀、某些罕见出现的星辰、地震、反常的生物诞生及其他类似事物；其中没有一件是未经天主旨意而发生的；然而，大多数人并不清楚事实如此。因此，哲学家们的虚荣便有了借口，将这些事物也归因于其他原因——也许是真实的原因，但只是近似的原因——而他们完全无法看见那高于一切的原因，即天主的旨意；或者又归因于错误的原因，这些原因甚至不是出于对有形事物与运动的认真探究，而是出于他们自己的猜测与谬误。\par

我将举一个例子，若可能，使这一点更清楚。我们知道，在人身上，有某种肉体的体积、外在形状、肢体的安排与分布、健康的体质；又有一个灵魂被吹入其中，管理这身体，且这灵魂是理性的；因此，它虽可变，却能分享那不变的智慧，以\Quote{分享那自在自为者}，正如圣咏上论及诸圣人时所说的，他们如同活石，建造成那为众人之母的耶路撒冷，在天上永远常存。因为如此歌唱：\Quote{耶路撒冷被建造如同一座城，分享那自在自为者。}因为在那里，\Quote{自在自为}被理解为那至高的、不变的美善，即天主，以及祂自己的智慧与旨意。在另一处，向祂歌唱说：\Quote{祢要改变它们，它们就被改变；但祢始终如一。}\par

% structure: p0014 source=h2 target=subsection reason=relative-heading-rank; base=h2

\subsection{第三章　同一论题。}

那么，我们取一位智者的例子：他的理性灵魂已分享那不变而永恒的真理，以致他关于自己一切行动都要咨询这真理，凡他所做的，无不是藉这真理知道该做的，为因服从它、听命它而行事正直。现在假设，此人经与神圣正义的最高理性商议——这理性他在心中秘密地以心灵之耳聆听——并奉其命令，在某慈悲职分上劳累身体，以致染病；然后咨询医生，一个告诉他病因是身体过度干燥，另一个则说是过度潮湿；二者中无疑有一个会指出真正原因，另一个则错误，但二人都是就近似的原因，即形体的原因发言。但假如要探究那干燥的原因，发现是自我强加的劳累，那么我们就到达了一个更高层的原因，它来自灵魂，从而影响灵魂所管理的身体。然而这还不是第一因，因为那第一因无疑更高，在于那不变的智慧本身；智者的灵魂因爱而事奉它、服从它那不可言说的命令，才承担了自我强加的劳累；因此，除了天主的旨意，找不出那疾病的真正第一因。但假如在那虔敬劳累的职分中，这位智者曾借助他人的帮助来合作行善，但这些人并不像他一样怀有同一旨意事奉天主，而是或欲获得自己肉欲的报酬，或仅逃避肉体的不适；——再假设，他使用了役畜，如果完成工作需此类供应，这些役畜当然是非理性的动物，因而不会因想到那善工而移动肢体驮负重物，而是出于自己喜好的自然欲望和避免不适；——最后假设，他使用了本身缺乏一切感觉的形体之物，但对于那工作是必需的，例如\Emph{谷物、酒、油、衣服、货币、书卷}或诸如此类的东西；——当然，所有这些用在此工作中的形体之物，不管是有生命的还是无生命的，其中所发生的任何移动、磨损、修复、毁坏、更新或以其他方式的变化，随地点和时间而受影响；请问，我能说这些可见又可变之事的任何其他原因，除了那无形不变的天主旨意，它使用所有这些——包括恶的与非理性的灵魂，以及最后那些被这些灵魂所激发并赋予生命的形体，或那缺乏一切感觉的形体——藉那义人的灵魂作为祂智慧的宝座，因为那善意圣洁的灵魂本身首先使用了它们，而祂的智慧已使那灵魂在虔敬与宗教的服从下顺服于自己？\par

% structure: p0016 source=h2 target=subsection reason=relative-heading-rank; base=h2

\subsection{第四章　天主随意使用一切受造物，并造有形之物以显示自己。}

9. 那么，我们刚才举了一个智者的例子，尽管他还带着可朽的身体，并且只如雾里看花，但若人间事务的最高管理和治理由智者、由那些虔诚而完全服从天主的人掌管，这个例子就可以合理地延伸到一个家庭——那里有这样一个团体——或一个城市，甚至整个世界；但由于目前情况并非如此（因为我们首先必须在此尘世旅途中按可朽的方式接受训练，并借着温和与耐心的力量接受惩戒），让我们将思绪转向那至高天上的家乡，我们在此是旅客。因为在那里，天主的旨意\BibleQuote{祂使祂的天使成为神风，使祂的仆役成为火焰}，临在于那些在完美的平安与友谊中结合、并由爱德的神火（如同属灵之火）合为一体的诸神之间，仿佛在一个崇高、神圣而隐秘的宝座上，如同在自己的殿堂和殿宇中，然后借着受造物最完美有序的运动——先是属灵的，后是物质的——将自己散布于万有，并根据其永不改变的旨意之喜悦来使用一切，无论是无形的事物还是有形的事物，无论是理性还是非理性的神灵，无论是因祂恩宠而良善的，还是因自身意志而邪恶的。然而，正如较粗重低等的形体由较精细有力的形体按适当秩序所治理，同样，一切形体都由有生命的灵所治理；而无理性的活灵由理性的活灵所治理；而犯错犯罪理性的活灵由虔诚正义的活灵所治理；而后者又由天主亲自治理，如此，整个受造物由其创造者治理，从祂、借着祂、在祂内受造并得以建立。因此，天主的旨意是一切有形现象和运动的首要且最高原因。因为没有任何可见或可感的事物发生，除非是出于最高统治者在无形可知的内心殿堂中的命令或允许，并按照那不可言喻的奖惩和恩惠报应的正义，在遍及整个受造物的广阔无边的国度中。\par

10. 因此，如果宗徒保禄，虽然他仍带着被腐朽所压制并压重灵魂的身体，虽仍只如雾里看花、只凭隐喻看见，渴望离世与基督同在，内心叹息，等待义子的名分，即身体的救赎，但他竟能有效地宣讲主耶稣基督，时而言语，时而去信，时而藉着祂体血的圣事（因为我们绝不称宗徒的舌头、羊皮纸、墨水、他口中发出的有意义的声音或写在皮上的字母是基督的体血；而只是那我们从地里收获、用神秘的祈祷祝圣，然后为了我们的属灵健康而合宜地领受的，以纪念主为我们所受的苦难：且这圣事，虽由人手带到那可见的形态，却除非透过天主圣神无形的工作，不能成为如此伟大的圣事；因为天主在那工作的进行中，透过物质的运动，首先推动祂仆人的不可见事物，无论是人的灵魂还是服从祂的隐秘神灵的服役），那么，如果天主在天地、海洋和空气的受造物中，按照祂自己所知恰当的方式，为在其中彰显和显示自己，而不显明祂的本质（即祂所是的），那完全不变而又比所有祂所创造的神灵更内在地、更隐秘地崇高的本性，又有什么奇怪呢？\par

% structure: p0019 source=h2 target=subsection reason=relative-heading-rank; base=h2

\subsection{第五章 为何奇迹并非寻常作为}

11. 因为既然天主的权能掌管整个属灵和有形的受造物，每年某些日子里海水被召聚并倾注在地面上。但当这事因圣厄里亚的祈祷而发生时，因为之前持续了那么久长时间的晴天，人们都饥饿了；且在那时，天主仆人在祈祷的那时刻，空气本身没有显出任何潮湿的迹象预示降雨；随后的大雨迅速来临就显明了天主的权能，并由此赐予并分施了那个奇迹。同样，天主通常藉着雷和闪电工作：但在西乃山上，这些是以异常方式进行的，那些声音不是以混乱的噪音发出的，而是以最确定的证据表明它们传递了某些信息，因此它们是奇迹。谁能使葡萄树通过根部将汁液吸到葡萄串，并酿造葡萄酒，除非是天主，当人种植和浇灌时，祂自己赐予生长？但当主下令时，水以超常的速度变为酒，神圣权能便彰显出来，甚至愚人也承认。谁通常用树叶和花朵装饰树木，除非是天主？然而，当亚伦司铎的棍杖开花时，天主性就以某种方式与疑惑的人性对话。再者，泥土元素既为生产并形成各种木材和一切动物的肉身服务：谁做了这些，但只有那说地要生出活物的那一位？谁又藉着同一位圣言治理和引导祂所创造的一切？然而，当祂将同样的泥土从梅瑟的棍杖变为蛇的肉身时，立刻而迅速地，那个异常的变化，虽然对象是可变的，却是一个奇迹。但谁在出生时赋予一切生物生命，除非是祂也在那一刻按需要赐予那条蛇生命？\par

% structure: p0021 source=h2 target=subsection reason=relative-heading-rank; base=h2

\subsection{第六章 独特性造成奇迹}

那使死者复活时，将各自的灵魂归还给尸体的，难道不是那位在母腹中赐生命给肉体，使那些将要死去的人得以存在的那一位吗？当这些事情像一条连续不断的河流，按照常例和既定的途径，从隐藏到可见，又从可见到隐藏地发生时，它们便被称为自然的；而当它们为了劝诫人，以不寻常的变化被插入时，便被称为奇迹。\par

% structure: p0023 source=h2 target=subsection reason=relative-heading-rank; base=h2

\subsection{第七章  藉巫术所行的大奇迹。}

12. 我在这里看到一件可能发生在软弱判断中的事，即：为何此类奇迹也藉巫术而行；因为法老的贤士也曾造出蛇，并行了其他类似的事。然而，更令人惊奇的是，那些术士的能力能够造出蛇，但当遇到极小的苍蝇时，却完全失败了。因为虱子——埃及骄傲的人民所受的第三灾——是极短命的小飞虫；然而在那里，术士确实失败了，说：\Quote{这是天主的指头。} 因此，我们得以明白，即使是那些犯罪的天空中的天使和权势，他们从崇高的以太纯洁的居所被推入最深的黑暗，如同进入特殊的监狱，巫术藉着他们拥有任何能力，也不能做任何事，除非由上而来的权力赐予。那权力或被赐予来欺骗那些欺诈者，正如赐予来对付埃及人，也对付术士们自己，以便在引诱那些神灵时，他们可能被那些施行者视为可惊奇的，但被天主的真理所定罪；或为劝诫信友，免得他们渴望行任何此类的事，以为是什么大事，为此它们也藉着圣经的权威传给了我们；或最后，为操练、证明和显示义人的坚忍。因为约伯失去他所拥有的一切，包括他的儿女和身体的健康本身，并非藉着可见奇迹的小小能力。\par

% structure: p0025 source=h2 target=subsection reason=relative-heading-rank; base=h2

\subsection{第八章  唯独天主创造那些被巫术改变的事物。}

13. 然而，不可因此认为有形之物受那些邪恶天使的命令支配；而是受天主的命令支配，是由祂赐予这能力的，恰如祂——那位不变的——在祂崇高而精神的居所中所决定赐予的。因为水、火和土地甚至服从那些被判处下矿的恶人，为使他们能随意使用它们，但只限于被允许的范围。事实上，也不应将那些邪恶天使称为创造者，因为藉着他们，术士们——与天主的仆人对抗——造出了青蛙和蛇；并不是他们创造了它们。但事实上，所有有形可见地诞生之物的某些隐藏的种子，都隐藏在这世界的有形元素中。因为那些现在从果实和生物中我们眼睛可见的种子，与那些先前种子的隐藏种子截然不同；从这些种子，按照创造者的命令，水产生了第一代水生动物和飞禽，大地产生了最初各从其类的芽和最初各从其类的生物。因为那时那些种子并没有如此被释放出来成为各种产品，以致生产能力在这些产品中耗尽；相反，常常缺少合适的条件组合，使它们能够爆发出来并完成其物种。因为，试想，最微小的芽也是一颗种子；如果被适当地置于土中，它就会长成一棵树。但在这芽中，在同一种类的某些谷粒中还有更精细的种子，这甚至是我们可见的。但在这谷粒中还有进一步的种子，虽然我们无法用眼睛看见，却可以用理性推断其存在；因为除非在那些元素中有某种这样的力量，否则就不会如此频繁地从地里生出没有播种的东西；也不会有如此多的动物，没有事先的雌雄交配，无论在地面上还是水中，它们却生长，并通过交配生出其他动物，而它们自己却是在没有任何父母结合的情况下生出来的。当然，蜜蜂并不是通过交配来孕育其幼崽的种子，而是用嘴从地上收集散落的种子。因为这些无形种子的创造者就是万物的创造者自己；因为凡是藉着诞生出现在我们眼前的一切，都从隐藏的种子开始其进程，并按照这些原始规则依次获得其适当的大小和独特的形式。因此，正如我们不称父母为人的创造者，也不称农民为谷物的创造者——尽管通过他们行动的外在应用，天主的能力内在运作以创造这些东西——同样地，也不应认为不仅恶天使，甚至善天使是创造者，如果他们通过其知觉和身体的精微性，知道那些对我们来说更隐藏的事物种子，并秘密地通过元素的适当调配将它们散布，从而提供产生事物的机会，并加速其增长。但善天使也不做这些事，除非天主命令；恶天使也不公义地做它们，除非祂公义地允许。因为恶者的恶意使他的意志成为不义的；但他行事的权力，他正当地领受，或为惩罚他自己，或就别人而言，为惩罚恶人，或为赞美善人。\par

14. 相应地，宗徒保禄区别了天主内在的创造与形成，和受造物从外在施加的操作，并从农业取比喻说：\BibleQuote{我栽种了，阿颇罗浇灌了，但使之生长的是天主。}\BibleRef{格前 3:6}因此，正如在属灵生命本身，除了天主，没有人能在我们内完成义德，然而人也能以外在方式宣讲福音——不仅善人真诚地，恶人也虚伪地；同样在可见事物的创造中，是天主从内在工作；而外在的操作，无论是善或恶的天使或人，甚至任何动物，都按照祂自己绝对的能力、祂所赋予的官能分配和对愉悦之物的各种欲望，由祂应用于祂在其中创造万物的本性，正如农业应用于土壤。因此，我不能称被法术召唤的恶天使为青蛙和蛇的创造者，正如我不能说恶人是谷物收成的创造者——我看见谷物通过他们的劳动生长出来。\par

15. 同样，雅各伯也不是羊群中毛色的创造者，因为他把各种颜色的棍子放在各母羊面前，让它们怀孕时观看。然而，牲畜自身也不是它们后代变色的创造者，因为通过眼睛看到各种棍子而印入灵魂的杂色形象，附着在灵魂上，但只能通过这种混合的同情作用影响被如此影响的灵魂所赋予生命的身体，在幼体柔软的开端上染上颜色。因为它们从自身受到这样的影响，无论是灵魂影响身体，还是身体影响灵魂，实际上都源于适当的理由，这些理由不变地存在于天主自身最高的智慧中，这智慧不为任何空间所包含；它本身不变，却不离弃任何一个变化的事物，因为没有一样不是由它创造的。正是天主智慧的不变和不可见的理由，一切受造之物都由它创造，才使得从牲畜生出牲畜，而不是从棍子生出牲畜；但怀孕牲畜的颜色受到各种棍子的影响，是因为怀孕牲畜的灵魂通过眼睛从外在受到影响，然后按其自身尺度，从它自己创造者的内在能力中接受了形成规则，并将其内在化。然而，灵魂在影响和改变物质实体方面有多大能力（尽管它当然不能被称为身体的创造者，因为每一个可变和可感知实体的原因及其全部尺度、数目和重量——正是通过这些，它得以存在并具有这样那样的性质——都来自可理解且不变的生命，这生命高于万有，并延伸到最遥远和最属世的事物），这是一个非常广泛的话题，目前不必要讨论。但我认为应当注意雅各伯关于牲畜的行为，\Emph{即}是为了使人明白：如果放这些棍子的人不能被称为羔羊和山羊羔毛色的创造者；甚至母羊自己的灵魂（它通过身体的眼睛所看到的杂色形象，按照自然所允许的程度，使肉身中受孕的种子着色）也不能被称为创造者；那么，就更不能说法郎的术士通过恶天使所制造的青蛙和蛇的创造者是恶天使了。\par

% structure: p0029 source=h2 target=subsection reason=relative-heading-rank; base=h2

\subsection{第九章。—— 万物的原始原因是来自天主。}

16. 因为这是一回事——从内在和最高的因果枢轴创造和管理受造物，这唯独作为创造者的天主所做的；而另一回事是根据祂所赋予每个受造物的能力和官能，从外在应用某种操作，以便受造物在此时或彼时、以这样或那样的方式产生出来。因为所有这些事物，就本源和起始而言，早已在一种元素的织体中受造，但它们遇到时机才出现。正如母亲怀胎，世界本身也怀有将要诞生之物的原因；这些原因不是在世界内受造，除非从那个最高本质而来——在那里没有任何事物出生或死亡，开始存在或停止存在。而从外在应用附加的原因（虽然它们不是本性的，却仍需按照本性来应用），以便那些被包含和隐藏在本性秘密怀抱中的事物，通过展开它们从祂那里秘密接受的适当尺度、数目和重量而迸发出来，并在某种程度上外在受造——\BibleQuote{按尺度、数目和重量安排万物的那位}\BibleRef{智 11:21}——这不仅在恶天使的能力中，也在恶人的能力中，如我以上用农业的例子所表明的。\par

17. 但为避免有人因动物的境遇略为不同而困扰，因为它们有生命的呼吸，并具有按本性欲求那些合乎本性的事物、避免那些相反本性之物的感觉；我们也必须思考：有多少人知道从何种草药或肉、从何种汁液或液体（无论何种类型），无论是如此放置、如此掩埋、如此捣碎或如此混合，通常能生出这种或那种动物；然而，谁能愚蠢到胆敢称自己为这些动物的创造者？那么，如果正如最卑贱之人也能知道如此这般的蠕虫和苍蝇从何而生，恶天使按其感知的敏锐，在元素更隐秘的种子中辨明青蛙和蛇从何而生，并通过某些确定且已知的适时配合，以隐秘的运动应用这些种子，使它们被创造，但并非它们自己创造，这又有什么可惊奇的呢？只是人们不惊奇那些通常由人所做的事。但若有人偶然惊异于那些生长之迅速，因为这些活物如此迅速地生成，就让他思考人按其人类能力的尺度如何也能做到这一点。因为何以同一躯体在夏天比冬天、或在较热之地比较冷之地更快地生出蠕虫呢？只是人们应用这些物质时越困难，因为他们的属土和迟钝的肢体缺乏感知的敏锐和身体运动的迅速。由此产生的结果是：对于任何种类的天使，按其更容易从元素中抽取近因的程度，他们在这类工作上的迅速就越是奇妙。\par

18. 但只有那一位是创造者，即这些事物首要的塑造者。除非是那位万有存在的度量、数目和重量根本与之相系的那一位，无人能成为这创造者；祂就是天主，唯一的创造者，凭其不可言喻的大能，也使得这些天使若被允许所能做的事，因不被允许而无法做到。因为别无其他原因，为何那些制造了青蛙和蛇的天使却不能制造最微小的苍蝇，除非是天主更强大的力量通过圣神禁止了它们；连术士们也承认了这一点，说：\Quote{这是天主的手指。}但他们凭本性所能做、却因被禁止而无法做的事，以及他们本性的条件本身不允许他们做的事，除非通过宗徒所说的天主之恩赐，即\Quote{另有一人能辨别神恩}，否则人难以甚至无法探究。因为我们知道人能够行走，但若不被允许就不能行走；但他不能飞行，即使被允许也不行。同样，那些天使，若被更强大的天使按天主的至高命令允许，便能做某些事；但某些其他事，即使被那些天使允许，也不能做，因为赐予他们如此这般自然能力尺度的那一位不允许；祂甚至通过祂的天使，通常不允许那些祂已赐予能力去做的事。\par

19. 因此，除了那些在完全惯常的时间系列中按自然秩序发生的形体事物，例如星辰的升起和落下、动物的生育和死亡、种子和萌芽的无数差异、蒸汽和云、雪和雨、闪电和雷、霹雳和冰雹、风和火、冷和热，以及一切类似之物；此外，还有在同一自然秩序中罕见发生的事物，如月蚀、星辰的异常显现、怪物、地震等等——我再说，所有这些都要除外，其首要和主要的原因实为天主的旨意；由此，在圣咏中，当提及此类事物时，即\BibleQuote{火与冰雹，雪和蒸气，暴风}，以免有人以为这些是出于偶然或仅来自形体原因，或甚至来自精神原因而脱离天主的旨意，便立刻加上：\BibleQuote{成就祂的圣言。}\par

% structure: p0034 source=h2 target=subsection reason=relative-heading-rank; base=h2

\subsection{第十章　受造物以标记的方式被领受的多种方式。圣体圣事。}

因此，正如我刚才所说，除了所有这些事物之外，还有另一类事物；它们虽然来自同一个物质实体，却通过我们的感官被呈现出来，用以宣告来自天主的信息，这些被称为奇迹与征兆；然而，凡是由主天主向我们宣告的一切，并非都取用了天主自己的位格。不过，当这一位格被取用时，有时以天使的形象显现；有时以一种形式呈现，这种形式本身并非天使，但却由天使命令和服事。再者，当它以并非天使本身的形式被取用时，有时是一个已经存在的身体，经过某种变化后取用，以显明那信息；有时是为特定目的而生成，事成之后便被抛弃。正如人们作为使者时，有时以自身位格说天主的话，如开头说\Quote{主说}，或\Quote{主这样说}，或其他类似语句，但有时不加任何这样的前缀，便亲自取用天主的位格，例如：\Quote{我必训导你，指教你当行的路}；同样，不仅在言语上，也在行动上，先知身上被赋予了天主的位格之象征，以便他在预言服务中担当那位格；就像某人那样，他担当了那位将外衣撕成十二块，将十块赐给撒罗满王的仆人、即未来的以色列王的位格。有时，一件本身并非先知、且已存在于世间的事物，也被取用来象征这一点；如雅各伯在梦见天梯醒来后，用他睡觉时枕在头下的石头所做的。有时，一件东西只是为了特定目的而被制造；它要么存在一段时间，如同在旷野中被举起的铜蛇一样，以及书面记录也能如此；要么在执行完其服务后便消逝，如同为圣事而制备的饼在领受中被消耗一样。\par

20. 然而，由于这些事为人所知，且是由人完成的，它们固然值得因神圣而敬畏，却不能因奇迹而令人惊叹。因此，由天使所行的事在我们看来更为奇妙，因为它们更困难、也更知名；但对天使而言，它们是已知且容易的，因为是它们自身的行动。一位天使以天主的位格向人说话，说：\Quote{我是亚巴郎的天主，依撒格的天主，雅各伯的天主}；而圣经在前文已说：\Quote{主的天使向他显现。}一个人也以天主的位格说：\Quote{请听，我的百姓，我要向你作证：我是主你的天主。}一根棍子被取来作为标记，并藉天使之力变成了一条蛇；尽管人缺乏那种能力，但人也曾取用一块石头作类似的标记。天使的作为与人的作为之间有很大差异。前者既令人惊叹又需要理解，后者则只需理解。从两者中所理解的东西，也许是一样的；但借以理解的那些事物，却是不同的。就像天主的名号用金子和墨水书写一样；前者更贵重，后者更低贱；但两者所表达的东西却是同一的。尽管从梅瑟的棍子变来的蛇与雅各伯的石头象征同样的事，但雅各伯的石头所象征的比术士的蛇更好。因为用油傅抹石头象征了在肉身中的基督，他被喜油傅抹，胜过他的同伴；同样，梅瑟的棍子变为蛇，象征了基督自己服从至死，且死在十字架上。因此经上说：\Quote{梅瑟在旷野里怎样举起了蛇，人子也必照样被举起来，使凡信他的人不至丧亡，反而获得永生}；正如仰望旷野中被举起的蛇，他们就不死于蛇咬。因为\Quote{我们的旧人已与他同钉十字架，使罪恶的身体失效}。因为蛇象征了死亡，这死亡是由乐园中的蛇造成的，修辞上以原因表达结果。因此，棍子变为蛇，基督进入死亡；蛇又变回棍子，整个基督及其身体进入复活；这身体就是教会；这将在时期的终结发生，由梅瑟握住的尾巴所象征，好让它变回棍子。而术士的蛇，如同在世上死亡的人，除非借由相信基督而被他的身体所吞并、进入他的身体，就不能在他内复活。因此，雅各伯的石头，如我所说，象征了比术士的蛇更好的东西；但术士的作为却更为奇妙。不过，这些事并不妨碍对事情的理解；就像一个人的名字用金子写，而天主的名字用墨水写一样。\par

21. 有哪一个人，又知道天使们是如何制造或取用了那些云彩和火焰，以表示他们所传递的信息呢？即使我们假定主或圣神是以那些有形体的形式显现的，也是如此。正如婴孩不知道那放在祭台上、在举行神圣庆典后所消耗的东西，究竟从何而来、如何制成、或从何处取来用于宗教用途。如果他们从未从自己的经验或他人的经验中学习，也从未在施行圣事时（当它被献上并赐予时）见过那类东西；如果有最具权威的人告诉他们，那是谁的身体和血；那么他们除了相信主确实以这形像向凡人的眼睛显现，并且那液体实际上是从被刺透的肋旁流出的、与这相似之外，其他什么也不会相信。但对我自己来说，这确实是一个有益的警戒，使我记住我自己的能力有多大，并劝诫我的弟兄们也记住他们自己的能力有多大，免得人的软弱超过安全的界限。因为天使们如何做这些事，或者更确切地说，天主如何藉祂的天使们做这些事，以及祂愿意这些事甚至由恶天使施行到何等程度——是出于许可、命令还是强迫，从祂自己至高权能的隐密宝座——这些我既不能用眼睛看透，也不能用理性的确证讲明，更不能被提升以领悟，以致能像我是天使、先知或宗徒那样，确实地回答关于这些事可能提出的所有问题。\BibleQuote{因为凡人的思想是悲惨的，我们的计谋是不确定的。因为可朽坏的身体压迫灵魂，地上的帐幕压服那默想许多事的心思。我们难以正确猜测地上的事，劳苦才找到面前的事物；但天上的事，谁能探究呢？} 但因接着又说：\BibleQuote{谁知道你的计谋，除非你赐予智慧，并从高处派遣你的圣神？} 因此，我们固然不去探究天上的事——这包括按他们各自尊严的天使的身体，以及这些身体的任何有形体的行动；然而，按照从高处派遣给我们的天主圣神，以及祂赐予我们心灵的恩宠，我敢自信地说：天主圣父、祂的圣言、祂的圣神——即唯一的天主——在祂所是的那一方面，以及藉此而成为祂所是的那一方面，绝无变化；更不会在这一方面是可见的。因为无疑有些事物是可变的，却不可见，如我们的思想、记忆和意志，以及整个无形的受造物；但没有可见的事物不是可变的。\par

% structure: p0038 source=h2 target=subsection reason=relative-heading-rank; base=h2

\subsection{第11章 天主的本质从未其自身显现。藉天使的职务而给祖先们的神视。从一个说话方式引出的异议被消除。天主向亚巴郎本人的显现，如同向梅瑟的显现一样，是藉天使进行的。同样的事由法律藉天使赐给梅瑟而得到证明。本书已论及的和下本书将论及的。}

因此，天主的本体——或者更好说是本质——我们在其中按我们的尺度、哪怕在最小的程度上理解圣父、圣子和圣神，既然它绝无变化，就绝不可能在其自身中为可见的。\par

22. 因此，很明显，所有那些向祖先们显现的事，当天主按祂自己适合时节的经纶呈现给他们时，都是藉着受造物施行的。即使我们不能辨别祂如何藉天使的职务施行它们，我们仍说它们是藉天使施行的；但这并非出于我们自己的辨别力，免得被人看为过于我们当有的智慧，而我们却是以清醒的头脑为有智慧，照着天主分给我们信德的度量；我们信，所以才说话。因为圣经的权威是确实的，我们的理性不该偏离它，也不该离开天主话语的坚固支撑，而坠入自己猜测的悬崖——在那些既无身体知觉主宰，也无真理的清晰理性照耀的事上。现在，在\WorkTitle{希伯来书}中，当新约的经纶要与旧约的经纶区分时，按时代和时期的适宜，确实写得最清楚：不仅是那些可见的事物，连那话本身也是藉天使施行的。因为这样记着说：\BibleQuote{但天主曾对哪一个天使说过：\InnerQuote{你坐在我右边，等我使你的仇敌作你的脚凳？}众天使不都是服役的神，奉差遣为那些将要承受救恩的人服务吗？} 由此可见，那些事不仅是由天使施行的，也是为我们的缘故施行的，即为天主的子民，就是那些应许承受永生的人。正如也写给\WorkTitle{格林多前书}说：\BibleQuote{这一切事发生在他们身上，都是预像；并且记载下来，是为了劝诫我们这些到了万世终结的人。} 然后藉着明确的推论指明，正如当时那话是由天使说的，如今却是藉着圣子说的：\BibleQuote{因此，我们必须越发重视所听见的事，免得我们失去它们。因为如果藉天使所说的话是确定的，凡过犯和悖逆都得了应得的报应；那么，如果我们忽略了这么大的救恩，怎能逃罪呢？} 然后，好像你问：什么救恩？——为要表明他现在是在讲新约，即那话不是由天使说的，而是由主说的——他说：\BibleQuote{这救恩开始时是由主宣讲的，后来由听见的人给我们证实了；天主又以神迹、奇事、各种异能和圣神的恩赐，按自己的旨意，一同作证。}\par

23. 但有人可以说：那么，为何经上写着：\Quote{主对\Person{梅瑟}说；}而不更恰当地说：天使对\Person{梅瑟}说？因为，当传令员宣告审判官的话语时，在记录中通常不会写着：某某传令员说，而是某某审判官说。同样，当圣先知说话时，虽然我们说先知说了，但我们所理解的不过是主说了；如果我们说主说了，我们并不是撇开先知，而只是指出是谁藉着他说话。确实，这些圣经屡屡显示天使就是主，论到他的说话，时常说：\Quote{主说，}正如我们已经指明的。但是，因为有些人，鉴于圣经在那个地方指明是一位天使，便要理解为天主圣子本身及在他本身内，因为祂被先知称为天使，是宣布祂圣父和祂自己的旨意；所以我认为应该从那封书信中提出一个更清楚的证据，那里不是通过一位天使说，而是\Quote{通过众位天使。}\par

24. 此外，\Person{斯德望}在\BibleBook{宗徒}{大事录}中也以同样的方式叙述这些事情，正如它们在\WorkTitle{旧约}中所写的：\Quote{诸位仁人弟兄，诸位父老，请听！}他说；\Quote{荣耀的天主曾显现给我们的祖先\Person{亚巴郎}，当他还在\Place{美索不达米亚}的时候。}但是，免得有人以为荣耀的天主那时曾以祂本身的样子显现给任何凡人的眼睛，他接着说，有一位天使显现给了\Person{梅瑟}。\Quote{那时\Person{梅瑟}听了这话就逃走了，}他说，\Quote{并在\Place{米德杨}地作了客旅，在那里生了两个儿子。过了四十年，在\Place{西乃山}的旷野里，有一位主的天使在荆棘火焰中向他显现。\Person{梅瑟}看见了，就惊讶于这个景象；当他走近观看时，主的声音向他来说：我是你父辈的天主，\Person{亚巴郎}的天主，\Person{依撒格}的天主，\Person{雅各伯}的天主。\Person{梅瑟}战栗起来，不敢注视。然后主对他说：脱下你脚上的鞋，}等等。这里，他当然既说到天使，也说到主；而且说到同一位就是\Person{亚巴郎}的天主、\Person{依撒格}的天主、\Person{雅各伯}的天主；正如在\BibleBook{}{创世纪}中所写的。\par

25. 岂有人会说：主是藉一位天使向梅瑟显现，却亲自向亚巴郎显现？我们不要从斯德望的言论来回答这个问题，而要从那本书本身，即斯德望取材的出处来回答。因为，请问，既然记载着：\BibleQuote{上主天主对亚巴郎说}；不久之后又说：\BibleQuote{上主天主向亚巴郎显现}，难道这些事因此就不是藉着天使完成的吗？\newline{}然而在另一处也类似地记载着：\BibleQuote{上主在玛默勒平原向他显现，那时他坐在帐幕门口，在白天最热的时候}；并且又立即加上：\BibleQuote{他举目观看，见有三个人站在他面前}——关于这些人我们已经讨论过了。因为那些既不肯从字面上升到意义，又轻易从意义跌落到字面的人，他们怎么能够解释天主在三个人中被看见，除非他们承认那些人是天使，正如后面所表明的那样？既然没有说\InnerQuote{一位天使向他说话或显现}，他们就敢说：向梅瑟赐下的异象和声音是藉着天使完成的（因为是这样记载的），但天主是以自己的实体向亚巴郎显现和说话，因为那里没有提到天使吗？\newline{}那么，就连亚巴郎的事例中，天使不也是没有被遗漏吗？因为当他的儿子被命令献为祭时，我们读到这样一段话：\BibleQuote{这些事以后，天主试探亚巴郎，对他说：亚巴郎！他答说：我在这里。祂说：带你心爱的独生子依撒格，往摩黎雅地去，在我所要指示你的山上，将他献为全燔祭。}这里明确提到的是天主，而不是天使。但稍后圣经如此说：\BibleQuote{亚巴郎就伸手拿刀要宰杀他的儿子。上主的天使从天上呼唤他说：亚巴郎！亚巴郎！他答说：我在这里。天使说：不可在这孩子身上下手，不要伤害他。}对此能回答什么呢？他们会说：天主命令杀死依撒格，而天使禁止了？更进一步，父亲本人违背了天主命令杀子的旨意，却听从了天使叫他饶恕的命令？这种解释荒谬之极，应予弃绝。然而，即使这样的解释粗鄙不堪，圣经也没有留下任何余地，因为它立即加上：\BibleQuote{现在我确知你敬畏天主，因为你为了我，没有顾惜你的独生子。}\Quote{为了我}是什么意思，难道不是为了那命令他杀子的那一位吗？那么，亚巴郎的天主是否与这位天使相同？或者，不是天主藉着天使行事吗？再看看下面的话。这里，无疑已经非常清楚地提到了天使；但注意上下文：\BibleQuote{亚巴郎举目观看，见在他后面有一只公绵羊，两角缠在灌木中；亚巴郎就去取了那公绵羊，代替自己的儿子献为全燔祭。亚巴郎给那地方起名叫上主看见，直到今日人还说：在山上，上主被看见了。}正如之前天主藉天使所说：\BibleQuote{现在我确知你敬畏天主}，这并不是要理解为天主那时才知道，而是祂促成亚巴郎藉着天主自己知道他的内心有多大的力量去服从天主，甚至牺牲自己的独生子。这是一种表达方式，其中的效果被施动者所代表——正如寒冷被说成是迟钝的，因为它使人迟钝。同样，天主之所以被说成\Quote{知道}，是因为祂使亚巴郎自己知道，而亚巴郎若未经如此考验，本可能无法辨别自己信德的坚定。同样，这里亚巴郎给那地方起名叫\Quote{上主看见}，即让自己被看见。因为他立即接着说：\BibleQuote{直到今日人还说：在山上，上主被看见了。}你看，这里同一位天使被称为\Quote{主}；那是什么原因呢，岂不是因为主藉着天使说话？但如果我们继续往下看，天使完全像先知一样说话，并明确表示天主藉着天使说话。\BibleQuote{上主的天使第二次从天上呼唤亚巴郎说：我指着自己起誓——上主说：因为你做了这事，没有顾惜你的独生子，为了我，等等。}的确，这些话，\Emph{viz.}即由主藉以说话的那一位应说\BibleQuote{上主这样说}，也是先知们常用的。难道天主子对圣父说\BibleQuote{上主说}吗？而祂自己就是圣父的那位天使？那么怎样呢？难道他们看不见，当先前已经说过\BibleQuote{上主向他显现？}时，关于向亚巴郎显现的那三个人，他们处于何等困境？难道因为他们被称为人，就不是天使吗？让他们读读达尼尔的话：\BibleQuote{看，那人加俾额尔}。\par

26. 我们为何还要迟疑，不藉另一最清晰、最有力的证据来堵住他们的口呢？此处所提到的既非单数的天使，也非复数的人，而是纯粹的天使；藉由他们，并非某种个别的言语被实行，而是法律本身最明确地被宣告为所赐下的；任何信友都不会怀疑，这法律是天主赐给梅瑟，为管理以色列子民的，然而它也是藉天使所赐下的。为此，斯德望说：\BibleQuote{你们这硬着颈项、心耳未受割损的人，时常抗拒圣神；你们的祖先怎样，你们也怎样。哪一个先知你们的祖先没有迫害过？他们杀害了那些预言那义者来临的人，现在你们竟成了那义者的出卖者和凶手。你们藉天使所颁赐的法律，却未遵守。}\BibleRef{宗 7:51-53} 还有什么比这更清楚的呢？还有什么比这权威更强有力的呢？法律确实是藉天使的传递赐给那百姓的，但我们的主耶稣基督的来临却是由此预备和预言的；而祂自己，作为\OriginalTerm{圣言}{Word of God}，以某种奇妙而不可言喻的方式存在于那些天使中，法律本身正是藉天使的传递而赐下的。因此祂在福音中说：\BibleQuote{若是你们相信梅瑟，必会相信我，因为他是指着我而写的。}\BibleRef{若 5:46} 可见主是藉天使发言的；而天主子，祂将成为天主与人之间的\OriginalTerm{中保}{Mediator}，出自亚巴郎的后裔，正藉天使预备自己的来临，为要找到那些因法律未能成全而成为罪人、且认罪的人，以便接收祂。因此宗徒也对迦拉达人这样说：\BibleQuote{那么，法律为什么有了呢？它是为彰显过犯而添设的，直到那预许的后裔来到，这法律是由天使经中保之手颁布的}\BibleRef{迦 3:19} ——即藉天使并在\Emph{祂自己的}手中颁布。因为祂不是在限制中，而是在大能中诞生的。但你在另一处得知，他所说的中保并非指任何一位天使，而是指主耶稣基督自己，就祂屈尊降生成人而言：\BibleQuote{因为天主只有一个，在天主与人之间的中保也只有一个，就是降生成人的基督耶稣。}\BibleRef{弟前 2:5} 因此，那逾越节宰杀羔羊的礼仪，以及所有在法律中以预像表达基督将按肉身来临、受难并复活的那些事——这法律是由天使的传递所赐下的——在其中，圣父、圣子和圣神确实存在于那些天使中；并且，有时圣父、有时圣子、有时圣神，有时天主，不分位格，藉它们以形象的方式被指示出来，尽管它们以可见可感的形态显现，却是藉自己的受造物，而非凭自己的本体；为能看见这本体，人心需藉那些眼所见、耳所闻的一切事得以洁净。\par

27. 但现在，我认为，我们在这本书中本打算要说明的，已按我们的能力充分讨论和论证了；并且已确立——既凭概然性推理（按一个人，或更确切说，按我所能的），也凭权威的力量（按圣经中的神圣声明已阐明）——那些在救主降生成人以前，赐给我们这些先祖的言语和形体显现，当时被称为天主显现的，都是由天使所行的：无论是他们以天主的位格说话或做事（如我们已表明先知们也惯常做的），还是从受造物中取用其自身所不是的形状，在其中天主能以形象向人显示；这种显示方式，圣经也以许多例子教导我们，先知们也并未忽略。因此，现在我们必须考虑——既然在主由童贞女所生、圣神以鸽子般的形体降临、以及五旬节主升天后从天上有响声如火舌显现的那些事中，不是圣言本身凭其本体（祂在本体上与圣父同等且永恒），也不是圣神凭其本体（祂自身也与两者同等且共永恒），而确确实实是一个受造物，能够以这些形态形成并存在，向可朽的身体感官显现——我再说，必须考虑这些显现与那些虽由可见受造物所行、却专属圣子和圣神的显现之间有何区别；这个问题我们将在另一卷书中更方便地开始讨论。\par

\SourceRecord{Translated by Arthur West Haddan.；Nicene and Post-Nicene Fathers, First Series；Vol. 3.；Edited by Philip Schaff.；Buffalo, NY: Christian Literature Publishing Co.,；1887.；<http://www.newadvent.org/fathers/130103.htm>.}

% structure: p0001 source=h1 target=section reason=page-title

\section{论天主圣三（卷四）}

\EditorialNote{解释为什么圣子被派遣：即，为使藉基督为罪人而死，我们得以确信天主对我们的爱是多么伟大，也得以认识我们所爱的是怎样的人。圣言降生成人的目的，也是使我们得以被洁净，以便默观并依附于天主。我们的双重死亡被祂单一的死所废除。此处讨论救主的单数与我们的双数如何和谐地指向救恩；并详细论述六数（senary number）的完美，单与双的比例本身可还原为六数。藉着生命唯一的中保基督，众多人被聚集为一；唯有藉着祂，灵魂才能得到真正的洁净。此外还证明：圣子虽因取了奴仆形体而被派遣成为较小的，但按天主的形体而言，祂并不小于圣父，因为祂是由自己所派遣；对圣神的派遣也应作同样的理解。}\par

% structure: p0003 source=h2 target=subsection reason=relative-heading-rank; base=h2

\subsection{序言——认识天主当求之于天主}

1. 世人通常极为看重对地上和天上事物的认识。然而，那些宁肯认识自己而不认识这些的人，无疑判断得更好；那认识自身软弱的心灵，比那不顾此点、探究并终于知道星体运行途径，或已掌握此知识却不知自己如何进入自身健全与力量之途径的心灵，更值得称赞。但若有人已向着天主觉醒，被圣神的热焰点燃，因爱天主而在自己眼中变为卑贱；并因渴望进入祂内却无力做到，藉着祂赐下的光，留意自己，发现自己，认识到自己的污秽不能与祂的纯洁相混；感到流泪恳求祂一次又一次怜悯是甘甜的，直到除去一切悲惨；并因已藉着唯一救主和人类启蒙者白白领受救恩的保证而充满自信地祈祷——这样一个人，如此行事、如此哀叹，知识不会使他自大，因为爱德在建造；他偏好知识胜过知识，他偏好认识自己的软弱，胜过认识世界的围墙、大地的根基和诸天的顶点。藉着获得这种知识，他也获得了忧愁；但忧愁是因为偏离了到达自己本乡及其创造者、自己蒙福的天主的渴望。如果在这些人中间，在祢基督的家中，上主我的天主，我同祢的穷人们一起呻吟，求祢从祢的食粮中赐给我，使我得以答复那些不饥不渴慕义、反而饱足有余的人。但使他们饱足的是那些事物的虚像，而非祢的真理；他们拒绝并回避祢的真理，因而陷入自己的虚妄。我确实知道人心产生多少幻象。我的心不也是人心吗？但我祈求我心的天主，使我在这些著作中不吐出（\Emph{eructuem}）任何这些幻象当作坚实的真理，而是让祂真理的气息所吹入我的，才能进入其中；尽管我被逐出祂双眼的视线，从远处努力经由祂独生子的天主性所藉着祂的人性开辟的道路返回。我虽是可变的，我这样啜饮此真理，直到在其中看不见任何可变之物：既不在空间和时间中——如物体那样；也不单独在时间中并在某种感觉的空间中——如我们自己的精神思想那样；也不单独在时间中甚至不在任何空间的相似中——如我们心灵的一些推理那样。因为天主的本质，祂所是的，绝无任何可变之物，无论在永恒、真理还是意志上；因为那里真理是永恒的，爱是永恒的；那里爱是真实的，永恒是真实的；那里永恒被爱，真理被爱。\par

% structure: p0005 source=h2 target=subsection reason=relative-heading-rank; base=h2

\subsection{第一章——承认我们的软弱使我们成全。降生成人的圣言驱除我们的黑暗}

2. 但既然我们被放逐离开那不变的喜乐，却又没有完全被切断或撕离，以致在那些可变和暂时的事物中不寻求永恒、真理、福乐（因为我们不愿死、不愿受骗、不愿受苦）；因此，从天上赐给了我们适合我们朝圣状况的异象，以提醒我们：我们所寻求的并不在此，而是必须从这朝圣之旅回到那里去，除非我们源自那里，我们不会在此寻求这些事。首先，我们必须被说服：天主多么爱我们，免得我们因绝望而不敢仰望祂。我们也需要被显示：我们所爱的我们是怎样的人，免得我们骄傲，仿佛出于自己的功德，从而更远离祂，在自己的力量中更失败。因此祂如此对待我们，使我们藉祂的力量得益，从而在谦卑的软弱中使爱德的美德得以成全。这在圣咏中有所暗示，那里说：\Quote{天主，祢降下甘霖，使祢疲惫的产业得以成全。}所谓\Quote{甘霖}，无非就是恩宠，不是回报功德，而是白白赐予，故称为恩宠；因为祂赐予，不是因为我们配得，而是因为祂愿意。知道这一点，我们便不自信；这便成为\Quote{软弱}。但祂自己使我们成全，祂也对宗徒保禄说：\Quote{我的恩宠为你足够了，因为我的能力在软弱中得以成全。}因此，人必须被说服：天主多么爱我们，以及我们所爱的我们是怎样的人；前者免得我们绝望，后者免得我们骄傲。宗徒如此解释这最必要的主题：\Quote{但是基督在我们还是罪人的时候，为我们死了，这证明了天主对我们的爱。} \BibleRef{罗 5:8} \Quote{既然现在我们因祂的血而成义，我们更要藉着祂脱离天主的义怒。因为我们作仇敌的时候，尚且藉着祂圣子的死与天主和好；那么，和好之后，我们更要藉着祂的生命得救了。}又在另一处说：\Quote{对此我们还有什么可说的呢？如果天主偕同我们，谁能反对我们呢？祂既然没有怜惜自己的儿子，反而为我们众人把祂交出来，岂不也把一切与祂一同白白赐给我们吗？}那已为我们成就的事，也曾向古时的义人显示为将要成就的事；为使他们也藉着同一的信德而谦卑，成为软弱；并使他们成为软弱，从而得以成全。\par

3. 因此，天主的圣言是唯一的，万物都是藉祂而受造，祂是不变的真理，万物在祂内同时存在，是潜在地且不变地存在；不仅是现在整个受造界中的事物，也包括那些曾经有过的和将来要有的。在祂内，它们既未有过，也未曾有，而只是\Emph{是}；并且万物都是生命，万物都是一体；更确切地说，祂是一体且一个生命。因为万物都是由祂所造，以致凡在它们内受造的，并不是在祂内受造的，而是在祂内是生命。因为\BibleQuote{在起初}，圣言不是受造的，而是\BibleQuote{圣言与天主同在，圣言就是天主，万物都是藉祂而造成的}；除非祂本身在万有之先并且不是受造的，万物也不会由祂而造成。但在那些由祂所造的事物中，即使是身体——它并不是生命——也不会由祂造成，除非它在受造之前在祂内已经是生命。因为\BibleQuote{那受造的，在祂内已经是生命}；并且不是任何一种生命：因为灵魂也是身体的生命，但灵魂也是受造的，因为它是可变的；那么它是藉什么造成的，除非藉着不变的天主圣言？因为\BibleQuote{万物是藉着祂造的；凡受造的，没有一样不是藉着祂造的}。\BibleQuote{所以，那受造的，在祂内已经是生命}；并且不是任何一种生命，而是\BibleQuote{那生命是人的光}；这光无疑是理性心智的光，人藉此与野兽区别开来，因此才成为人。因此，那不是肉体的光——即肉体的光，无论它是从天上照耀，还是由地上的火点燃；也不只是人类肉体的光，还包括野兽的，甚至直到最微小的虫子的光。因为所有这些事物都看见那光：但那生命是人的光；祂离我们每个人并不远，因为在祂内\BibleQuote{我们生活、行动、存在}。\par

% structure: p0008 source=h2 target=subsection reason=relative-heading-rank; base=h2

\subsection{第二章 我们如何通过降生成人的圣言而适于感知真理。}

4. 但是\BibleQuote{光在黑暗中照耀，黑暗却没有接纳它}。这里的\Quote{黑暗}是指人愚昧的心智，因邪恶的欲望和不信而变得盲目。为使万物藉着祂而造的圣言眷顾并医治这些，\BibleQuote{圣言成了血肉，居住在我们中间}。因为我们的光照就是分享圣言，即那生命是人的光。但对于这种分享，我们因罪的污秽而完全不配，并且有所欠缺。因此我们需要被洁净。进一步说，不义者和骄傲者唯一的洁净就是义人的血，以及天主自身的谦卑；好使我们藉着祂得到洁净，祂使自己成为我们本性所是的，以及因罪我们不是的，好使我们能默观天主，这本是我们本性所不是的。因本性我们不是天主；因本性我们是人，因罪我们不是义人。因此，天主成为义人，在天主面前为罪人代祷。因为罪人与义人不相符，但人与人相符。因此，藉着将祂人性的肖像结合于我们，祂消除了我们不义的差异；藉着分享我们的必死性，祂使我们分享了祂的天主性。因为罪人之死出自定罪的必然性，理当由义人之死所废除，义人之死出自祂慈悲的自由选择，而祂单次的死亡与复活对应我们的双重死亡与复活。这种适应、适合、和谐、协和，或用任何更合适的词，使一者与二者结合，在受造物的一切组合或更佳的调协中具有重大意义。因为（正如我刚才想到的），我所说的正是希腊人称为\OriginalTerm{和谐}{\GreekText{ἁρμονία}}的那种调协。然而这里不是阐述那种一与二协和的力量的地方，这种力量特别存在于我们内，并且是自然铭刻于我们内的（除了创造我们的那一位，还有谁呢？），以致连无知的人，无论自己唱歌还是听别人唱，都无法不察觉。因为正是藉此，高音和低音才和谐一致，因此任何人在音调上偏离它，都会极大地冒犯，不仅冒犯训练有素的技巧（大部分人都缺乏），而且冒犯听觉本身。要证明这一点无疑需要长篇论述；但知道它的人可以通过调好的一弦琴使耳朵明了。\par

% structure: p0010 source=h2 target=subsection reason=relative-heading-rank; base=h2

\subsection{第三章 基督身体的一次死亡与复活，与我们的身体和灵魂的双重死亡与复活协调一致，以达致救恩。基督的一次死亡如何施予我们的双重死亡。}

5. 但是，为了我们现在的需要，我们必须讨论——按天主赐给我们的能力——我们的主和救主耶稣基督的单一性如何适应、并可以说与我们的双重性相协调，以达致救恩的效果。我们当然——正如任何基督徒都不怀疑的——在灵魂和身体上都是死的：在灵魂上，因罪而死；在身体上，因罪的惩罚而死，并且因此也在身体上因罪而死。并且对我们自身的这两部分，即灵魂和身体，都需要一味药和一次复活，使已变得较坏的能被更新为较好的。灵魂的死亡是不虔诚，身体的死亡是可朽坏性，由此也发生了灵魂与身体的分离。因为正如灵魂在天主离开它时就死去，同样，身体在灵魂离开它时就死去；因此前者变为愚妄，后者变为无生命。因为灵魂藉悔改再次被举起，而在仍会死的身体中，生命的更新藉信德开始，人藉此信德相信那使不敬虔者成义的天主；这生命的更新藉良好的习惯日复一日地增长和加强，正如内心的人被日益更新。但身体，作为外在的人，这生命延续得越久，就越是朽坏，或由于年龄，或由于疾病，或由于各种苦难，直到它达到那所有人都称为死亡的最终苦难。它的复活被延迟到终结；那时我们的成义本身也将不可言喻地得以完成。因为那时我们将肖似祂，因为我们将看见祂实在怎样。但现在，只要这可朽坏的身体压迫灵魂，而且地上的人生全是试探，在祂面前没有一个活人能成义，与我们将被造成与天使相等的正义以及那将显示在我们身上的光荣相比。但何必再提更多关于灵魂之死与身体之死区别的证明呢？当主在福音中的一句话里已经使任何人都能轻易地将这两种死亡彼此区分，祂说：\BibleQuote{让死人埋葬他们的死人}？因为埋葬是对死尸的适当处理。但祂所指那些要埋葬它的人，是因不信的不虔诚而在灵魂上死了的人，即当说到：\BibleQuote{你这睡眠的，醒起来罢！从死者中起来罢！基督必要光照你。}时被唤醒的人。还有一种死亡是宗徒所谴责的，论及寡妇说：\BibleQuote{但那好宴乐的寡妇，虽生犹死。}因此，先前不敬虔而现今敬虔的灵魂，因信德之正义，被称为已从死者中复活而活着。但身体不仅因将要发生的灵魂离去而被说成将要死亡；而且因血肉的重大软弱，在圣经的某处甚至被说成现在已经死了，即宗徒说：\BibleQuote{身体固然因罪恶而死亡，但神魂却因正义而生活。}而这生命是藉信德实行的，因为\BibleQuote{义人必因信德而生活。}但接着呢？\BibleQuote{再者，如果那使耶稣从死者中复活的圣神住在你们内，那么，那使基督从死者中复活的，也必要藉那住在你们内的圣神，使你们有死的身体复活。}\par

6. 因此，关于我们的这双重死亡，我们的救主献上了祂自己的单一死亡；并且为了促成我们的两种复活，祂预先安排并以奥秘和预像展示了祂自己的一次复活。因为祂不是罪人或不虔敬的人，以至于仿佛在精神上死了，需要在内里被更新，并通过悔改被唤回到义的生命；但祂穿上了可朽的血肉，并且唯独在那血肉中死去，唯独在那血肉中复活，唯独在那血肉中祂替我们应对了双重；因为在血肉中，对于内里的人，成就了一个奥秘，对于外在的人，则成就了一个预像。因为对于我们的内里的人而言，这是一个奥秘，以表示我们灵魂的死亡，所以那些话不仅是在圣咏中，也在十字架上说出：\BibleQuote{我的天主，我的天主，袮为什么舍弃了我？}宗徒同意这些话，说：\BibleQuote{我们知道，我们的旧人已与祂同钉在十字架上，使罪身消逝，以致我们不再作罪恶的奴隶}因为内里的人被钉十字架，理解到悔改的痛苦和一种有益的克己挣扎，通过这死亡，不虔敬的死亡被消灭，在这死亡中，天主舍弃了我们。这样，罪身通过这样的十字架被消灭，以至于我们不再将我们的肢体作为不义的工具献给罪恶。因为，即使内里的人确实日益更新，但无疑，在更新之前它是旧的。因为那是在内里发生的，同一位宗徒说：\BibleQuote{脱去旧人，穿上新人}他接着解释说：\BibleQuote{因此，你们应弃绝谎言，各人说实话}但是，哪里弃绝了谎言，除非是内里，以至于那从心里说实话的人才能居住在天主的圣山上呢？然而，主的身体的复活被显示属于我们内里复活的奥秘，在祂复活后，祂对那女人说：\BibleQuote{不要摸我，因为我还没有升到父那里去}宗徒的话与这个奥秘一致，他说：\BibleQuote{你们既然与基督一同复活了，就该追求天上的事，那里有基督坐在天主的右边；你们要思念天上的事}因为不触摸基督，除非当祂升到父那里，意思是不要以属肉体的方式思想基督。同样，我们主肉体之死包含我们外在人之死的预像，因为正是通过这样的苦难，祂尤其劝诫祂的仆人们不要害怕那些杀死身体却不能杀死灵魂的人。因此宗徒说：\BibleQuote{为基督的身体，就是为教会，要在我的肉身上补充基督苦难的欠缺}并且我们发现主身体的复活包含我们外在人复活的预像，因为祂对门徒说：\BibleQuote{你们摸摸我，看看！因为鬼神没有骨肉，如同你们看我一样}还有一个门徒，摸着祂的伤痕，喊道：\BibleQuote{我的主！我的天主！}而且，既然那肉体的完整无缺是明显的，这显示在祂劝诫门徒时所说的话：\BibleQuote{连你们的一根头发也不会失落}因为怎么会先说：\BibleQuote{不要摸我，因为我还没有升到父那里去}；而祂在升到父之前，实际上被门徒触摸呢？除非因为在前者暗示了内里的人的奥秘，在后者则给出了外在的人的一个预像？难道有人能如此无知，如此背离真理，竟敢说祂在升天前被男人触摸，而升天后被女人触摸吗？正是由于这个在主内预先存在的、关于我们将来身体复活的预像，宗徒说：\BibleQuote{基督是初果，然后是属于基督的人}因为这里指的是身体的复活，为此他也说：\BibleQuote{祂将改变我们卑贱的身体，使之相似祂光荣的身体}因此，我们救主的一次死亡为我们的双重死亡带来了救恩，祂的一次复活为我们成就了两种复活；因为祂的身体在两者中，即在祂的死亡和复活中，通过一种有益的适应，作为内里之人的奥秘和外在之人的预像，为我们服务。\par

% structure: p0013 source=h2 target=subsection reason=relative-heading-rank; base=h2

\subsection{第四章 — 单单的与双倍的比率来自六数之完美。圣经中赞扬六数之完美。年岁因六数而丰盈。}

7. 现在，单一与双倍的比率无疑来自于三元数，因为一加二得三；但这些数组成的整体达到六元数，因为一、二、三得六。这数因此被称为完数，因为它由自身的部分完成：它有这三个部分：六分之一、三分之一、二分之一；在它之内找不到其他我们可称为约数的部分。它的六分之一是一，三分之一是二，二分之一是三。而一、二、三完成同一个六。圣经向我们称赞这个数的完美，尤其在此：天主在六天内完成了祂的工程，第六天，人按天主的肖像被造。圣子来临并成了人子，为按天主的肖像重新创造我们，正是在人类历史的第六时代。因为现在是当前的时代，无论每时代被分配一千年，还是我们在圣经中追溯可纪念和显著的时代或时间的转折点，使得第一时代从亚当直到诺厄，第二时代从那里直到亚巴郎，然后按福音作者玛窦的划分，从亚巴郎到达味，从达味到被掳往巴比伦，从那里到童贞玛利亚的产痛；这三个时代加上另两个共五个。因此，主的诞生开始了第六时代，这时代现在继续直到时间隐藏的终结。我们也在这六元数中认识了一种时间的形象，在那三分的方式中，我们计算一部分时间在法律之前，第二部分在法律之下，第三部分在恩宠之下。在这最后的时间，我们领受了更新的圣事，以便我们也在时间的终结，在每一部分，借着肉身的复活而得以更新，并因此能从我们整个的软弱——不仅是灵魂，也是身体——得以痊愈。因此那妇女被理解为教会的预象，她被主治愈并直立，之前她被撒殚的捆绑所弯曲而软弱。因为\BibleBook{}{圣咏}的那些话哀叹这样的隐藏敌人：\Quote{他们压弯了我的灵魂。}这妇女的软弱持续了十八年，即三个六。十八年的月数在数目上被确认为六的立方，\Emph{即}六乘六乘六。几乎在\BibleBook{}{福音}同一地点，那棵无花果树也被第三年的可怜不结果而定罪。但有人为它代祷，使它那一年被留下，那一年，如果它结果，就好；否则，它将被砍下。因为三年都属于同样的三分，而三年的月数构成六的平方，即六乘六。\par

8. 若将整个十二个月都计算在内，每月由三十天构成（因自古以来以月的运行为准的月），则一年充满六这个数。因为在数的第一序中，六是什么，第二位中六十就是什么，第一位由单位到十，第二位由十到百。因此六十天是一年的六分之一。进一步，若将第二序的第六数乘以第一序的第六数，则我们得六乘六十，即三百六十天，这正是整个十二个月。但由于月的运行决定人的月，同样年的运行由太阳的运行标志；并有五天四分之一天剩余，以便太阳完成其运行结束一年；因四个四分之一天构成一天，每四年必须闰一日，他们称之为\OriginalTerm{闰年}{bissextile}，以免时间秩序被打乱：如果我们考虑这五天四分之一天本身，六这个数也在其中盛行。首先，因为习惯上从部分计算整体，我们不应称其为五天，而应称六天，将四分之一天算作一天。其次，因为五天本身是一个月的六分之一；而四分之一天包含六小时。因整个一天（\Emph{即}包括其夜间）为二十四小时，其四分之一，即一日的四分之一，被确认为六小时。在一年中，第六数如此盛行。\par

% structure: p0016 source=h2 target=subsection reason=relative-heading-rank; base=h2

\subsection{第五章——六这个数在建立基督的身体及耶路撒冷圣殿中也受到称赞。}

9. 并非毫无理由，六这个数被理解为代表一年，用于建立主的身体；作为此身体的预象，祂曾说祂要在三天内重建被犹太人毁坏的圣殿。因为他们说：\Quote{这圣殿建造了四十六年。}六乘四十六得二百七十六。这数目的天数完成九个月零六天，这被算作妇女生产的十个月；并非所有人在第九个月后都到第六天，而是因为主身体的完美本身被发现在这么多天达到出生，正如教会的权威根据长老的传统所维护。因为据信祂在三月二十五日受孕，祂也在同一日受难；因此童贞玛利亚的子宫（祂在其中受孕，没有凡人被生在其中）对应祂被埋葬的新坟墓，其中从未有人被放置，以前或以后都没有。但据传统，祂生于十二月二十五日。因此，你若从那天算到这天，你发现二百七十六天，即四十六乘六。圣殿也在这些年数中被建造，因为在那些六中，主的身体被成全；这身体被死亡的苦难摧毁后，祂在第三日复活了。因为\Quote{祂这话是指祂身体的圣殿说的}，正如\BibleBook{}{福音}最清楚、最坚实的见证所宣告的；在那里祂说：\Quote{就如约纳曾在大鱼腹中三天三夜，同样人子也要在地心里三天三夜。}\par

% structure: p0018 source=h2 target=subsection reason=relative-heading-rank; base=h2

\subsection{第六章——复活的三天，其中也显明了单一与双倍的比率。}

10. 圣经又见证，这三天的时间本身并非完整，而是第一天从其最后部分算为完整，第三天从其开始部分算为完整；但中间那一天，即\Emph{即}第二天，则完全是完整的二十四小时，白天十二小时，夜晚十二小时。因为祂首先在第三时辰被犹太人的喊声钉死，那时是一周的第六天。然后祂在第六时辰被悬挂在十字架上，在第九时辰交出了祂的灵魂。但祂被埋葬时，正如福音作者所说：\BibleQuote{现在傍晚已至}，意思是一天的结束。因此，无论你从何处开始——即使可以有别的解释，以便不与\BibleBook{若望}{福音}矛盾，而是理解祂在第三时辰被悬挂在十字架上——你仍然不能使第一天成为完整的一天。它将被从其最后部分算为完整的一天，第三天则从其开始部分算为完整。因为从夜晚到黎明，当主的复活被宣告时，属于第三天；因为天主（祂命令光从黑暗中照耀，好藉着新约的恩宠和分享基督的复活，那话对我们说出：\BibleQuote{因为你们从前是黑暗，但现在在主内却是光明}）以某种方式向我们暗示，日从夜开始。因为正如最初的日子都是从光到暗来计算的，为预表人的将来堕落；同样，这些日子，为预表人的恢复，则是从暗到光来计算的。所以，从祂死亡的时刻到复活的黎明，共四十小时，也包括第九时辰本身。这个数字也与祂复活后在世上四十天的生命相吻合。这个数字在圣经中经常被用来表达四重世界中的完美奥秘。因为数字十有某种完美，乘以四得四十。但从埋葬的傍晚到复活的黎明，是三十六小时，即六的平方。这关系到单一与双重的比例，其中有着最大的协调一致。因为十二加二十四符合单一加双重的比例，得三十六：即一整夜、一整天和另一整夜，这并非没有我上面注意到的奥秘。因为我们并非不合适地将精神比作白天，将身体比作黑夜。因为主在死亡和复活中的身体是我们精神的预象和身体的预表。这样，在三十六小时中，当十二加到二十四时，单一与双重的比例也显明了。至于为什么这些数字在圣经中如此安排，其他人可能找出别的原因，要么我所给出的原因优于它们，要么与我的原因同样可能，甚至更可能；但确实没有一个如此愚蠢或荒谬的人会争辩说，它们在圣经中的安排毫无目的，也没有神秘的理由说明为什么这些数字在那里被提及。但我在此所给出的理由，要么是从教会的权威、根据我们祖先的传统收集而来，要么是从圣经的见证，要么是从数字本身和相似性的性质而来。没有一个清醒的人会违背理性，没有一个基督徒会违背圣经，没有一个爱好和平的人会违背教会。\par

% structure: p0020 source=h2 target=subsection reason=relative-heading-rank; base=h2

\subsection{第七章——我们如何通过唯一的中保从众多中被聚集为一。}

11. 这奥秘、这祭献、这司铎、这天主，在祂被派遣和来临、由女人所生之前，所有那些以神圣和神秘的方式藉着天使的奇迹向我们的祖先显现的事，或由祖先自己所行的事，都是预象；为要使每个受造物藉其行为以某种方式论及那将要来的那一位，在祂内将要有救恩，即从死亡中恢复一切。因为因着不虔敬的邪恶，我们已退避并从独一真实至高的天主那里因失和而堕落，在许多事上变为虚妄，被许多事物分心并紧抓着许多事物；因此，因天主慈爱的命令和吩咐，那些同一的许多事物必须联合宣告那将要来的那一位，而那一位必须由这些许多事物如此宣告而来，并且这些许多事物必须联合见证这一位已经来了；从而，我们从这些许多事物的重担中释放出来，应来到那一位面前，我们灵魂因许多罪而死，注定因罪在肉身中死去，我们应爱那位无罪的、为我们在肉身中死去的那一位；并藉着相信祂现已复活，并藉着信德在精神上与祂一同复活，我们应因在唯一正义者中成为一而成义；并且我们不应绝望于自己肉身的复活，当我们考虑元首已先于我们众多肢体前行时；在祂内，我们现已藉着信德被洁净，然后藉着显现得以更新，并藉着祂作为中保与天主和好，我们应依附于那一位，享受那一位，保持合一。\par

% structure: p0022 source=h2 target=subsection reason=relative-heading-rank; base=h2

\subsection{第八章——基督怎样愿意所有人在祂内合而为一。}

12. 因此，天主子自己、圣言、本身也是天主与人之间的中保、人子，因天主性的合一与圣父平等，又因摄取人性而与我们同有份，在祂作为人时为我们在圣父前代祷，却不隐藏祂是天主、与圣父为一，在其他事中这样说道：\BibleQuote{我不但为他们祈求，}祂说，\BibleQuote{也为那些因他们的话而信我的人，使他们合而为一，如同祢，父，在我内，我在祢内，使他们也在我们内合而为一，为叫世界相信是祢派遣了我。祢赐给我的光荣，我已赐给了他们，为叫他们合而为一，如同我们原是一体。}\par

% structure: p0024 source=h2 target=subsection reason=relative-heading-rank; base=h2

\subsection{第九章——同一论证继续。}

祂没有说：我与他们是一体；虽然，祂是教会——祂的身体——的头，祂本可说：他们与我是同一（不只是一体，而是一个位格），因为头与身体是一个基督；但是为了显示祂自己的天主性与圣父\OriginalTerm{同体}{consubstantial}（为此祂在另一处说：\BibleQuote{我与父原是一体}），在祂自己的方式中，即在同一本性的同体平等中，祂愿意祂的（信徒）在祂内成为一体；因为他们自身不能如此，由于他们被各种快乐、欲望和罪恶的不洁彼此分离；因此他们通过中保得以洁净，使他们在祂内成为一体，不仅是通过同一本性——在其中所有人都从必死的人变为与天使平等，而且是通过同一意志，最和谐地趋向同一真福，并藉着爱德之火以某种方式融为一灵。因为祂接下来的话正是指此：\BibleQuote{使他们合而为一，如同我们合而为一}；即，正如父与子不仅在本体平等上是一体，也在意志上是一体，同样那些人也可能成为一体——在他们与天主之间，圣子是中保——不仅在于他们拥有同一本性，也在于同一的爱之结合。然后祂继续这样启示真理本身，即祂是那中保，藉着祂我们与天主和好，说：\BibleQuote{我在他们内，祢在我内，使他们完全合而为一}。\par

% structure: p0026 source=h2 target=subsection reason=relative-heading-rank; base=h2

\subsection{第十章——正如基督是生命的中保，同样魔鬼是死亡的中保。}

13. 在此中我们与造物主有真正的和平与坚固的结合之纽带，即我们应藉着生命的中保被净化、和好，正如我们曾藉着死亡的中保被污染、疏离，并因此而离弃了祂。因为正如魔鬼藉着骄傲，通过骄傲引领人走向死亡；同样基督藉着谦卑，通过服从将人带回生命。既然，一方因高升而坠落，也把同意他的人坠落；另一方因受贬抑而被高举，也把相信他的人高举。因为魔鬼自己并未到达他所引领的道路（因为他的确带着不敬虔的灵魂的死亡，但并未经历肉身的死亡，因为他并未披上肉身的覆盖），他在人面前显出是魔鬼军团中一位大首领，藉着他行使欺骗统治；从而更加膨胀人——那些渴求权力甚于正义的人——通过高傲的骄傲，或虚伪的哲学；或用亵渎的仪式缠绕他，在其中藉着欺骗和幻象将心灵中更好奇和更骄傲的人抛入深渊，他也以巫术的把戏俘虏他；并承诺藉着那些他们称为\OriginalTerm{入教仪式}{\GreekText{τελεταί}}的仪式净化灵魂，藉着将自己变形为光明之天使，通过种种在虚假征兆和奇迹中的阴谋。\par

% structure: p0028 source=h2 target=subsection reason=relative-heading-rank; base=h2

\subsection{第十一章——由魔鬼所行的奇迹当被唾弃。}

14. 因为最不值一提的恶灵藉着气状身体能做出许多事，足以使被属地身体所压制的灵魂感到惊奇，即使这些灵魂倾向于更好的一方。因为如果属地身体本身，经过某种技巧和实践的训练，在戏剧表演中向人展示如此伟大的奇迹，以至于从未见过的人几乎难以相信所听闻的；那么，为何魔鬼及其使者难以从物质元素中，藉着他们自己的气状身体，造出让肉身惊异的事物；或甚至藉着隐秘的感动，制造幻象以欺骗人的感官，从而欺骗他们（无论醒着或睡着），或将他们驱入疯狂呢？但正如一个在生活中和品格上比他们更好的人，可能注视着最不值一提的人，或走钢丝，或以各种身体动作做出许多难以置信的事，而他可能完全不想做这些事，也不认为那些人因此而比他更优越；同样，虔敬而信德的灵魂，不仅看到，甚至若因肉身的软弱而战栗于魔鬼的奇迹，也断不会因此哀叹自己无力做这类事，或因此评判他们比自己更好；尤其因为它与圣者们为伴，无论他们是人还是善天使，他们藉着天主——万物都服从于祂——的权能，成就远为伟大且绝不欺人的奇迹。\par

% structure: p0030 source=h2 target=subsection reason=relative-heading-rank; base=h2

\subsection{第十二章——魔鬼是死亡的中保，基督是生命的中保。}

15. 因此，灵魂绝不能通过亵圣的仿效、不虔的巫术或魔法咒语得到洁净并与天主和好；因为假的中保并不将他们提升到更高的事物，反而借着情感（这些情感越骄傲就越恶毒），堵塞并切断通往那里的道路，他将这些情感灌输给属于他团体的人；这些情感不能滋育美德之翼以向上飞翔，反而堆积恶行之重以向下压沉；因为灵魂越是自以为被高举，就越沉重地坠落。因此，正如那些蒙天主警告的贤士，被星引导去朝拜主卑微的处境；同样，我们也应当归回本乡，不是走我们来的那条路，而是走那位卑微的君王所教导的另一条路，那条路是骄傲的君王（即那位卑微君王的仇敌）所无法堵塞的。因为，为了叫我们也能朝拜卑微的基督，\BibleQuote{诸天宣布了天主的荣耀，当它们的声音传遍普世，它们的话语到达地极。}罪在亚当内为我们开辟了一条通向死亡的路。因为，\BibleQuote{藉着一个人，罪进入了世界，死亡藉着罪也进来了，于是死亡就临到了众人，因为众人都犯了罪。}在这条路上，魔鬼是那中保，是劝诱人犯罪、将人投入死亡的那一位。因为他也用他自己的一个死亡来成就我们的双重死亡。他固然在灵性上因不虔而死了，但确实没有在肉身上死；然而他既劝诱我们行不虔，从而使我们理应陷入肉身的死亡。因此，我们因邪恶的劝诱而渴望了那一个，另一个则因公义的惩罚而跟随了我们；所以经上记载：\BibleQuote{天主并未制造死亡，}因为天主自己不是死亡的原因；然而死亡是因着祂最公义的报应而加在罪人身上的。正如法官对罪犯施加惩罚；然而惩罚的原因不是法官的公义，而是罪行的应得。那么，死亡的中保使我们陷入的终点（即肉身的死亡），他自己并没有到达，而我们的主天主却为我们在那里引入了矫正的良药——祂所不配得的——借着一种隐秘且极其神秘的、神圣而深奥公义的谕令。因此，为要如同因一人而来死亡，也因一人而带来死人的复活；因为人们更竭力躲避他们所不能躲避的，即肉身的死亡，超过灵性的死亡，也就是更躲避惩罚而非惩罚的应得（因为对于不犯罪，人们要么不关心，要么太少关心；而不死虽不可企及，却被热切追求）；生命的中保使人明白，死亡（因人类现状已无法逃脱）并不可怕，而可怕的是不虔，这是可以借着信心防备的。祂在我们所到达的终点迎接我们，但不是走我们所来的路。因为我们确实是通过罪来到死亡；祂则是通过义；因此，正如我们的死亡是罪的惩罚，祂的死亡则被献为赎罪的祭献。\par

% structure: p0032 source=h2 target=subsection reason=relative-heading-rank; base=h2

\subsection{第十三章——基督死亡的主动性。生命的中保如何征服死亡的中保。魔鬼如何引导其属下轻视基督的死亡。}

16. 既然灵应优先于身体，灵的死意味着天主离开了它，而身体的死则意味着灵离开了它；并且身体的死之惩罚就在于：灵违背自己的意愿离开身体，因为它曾自愿离开天主；所以，灵因乐意而离开天主，却虽不愿意而离开身体；而且它不能在自己意愿的时刻离开，除非它施暴于自身，从而使身体被杀：中保的灵表明，祂来到肉身的死绝非出于罪的惩罚，因为祂不是违背自己的意愿离开身体，而是因祂意愿，在祂意愿的时刻，照祂所意愿的那样。因为祂被天主的圣言如此结合[于肉身]而成为一体，祂说：\BibleQuote{我有权舍弃我的性命，也有权再取回它。没有人夺去我的性命，而是我自愿舍弃，为能再取回它。}而且，正如福音告诉我们的，在场的人对此非常惊讶，即在祂说出那最后的话（其中祂预表了我们的罪）之后，祂便立即交付了灵魂。因为被钉在十字架上的人通常要受长时间的折磨。因此，窃贼的腿被打断，以便他们立刻死去，并在安息日之前从十字架上取下。而发现祂已经死了，这引起了惊奇。这也是，我们读到，比拉多所惊奇的，当时有人向他请求主的遗体安葬。\par

17. 因为那个欺骗者——他本是引人死亡的中保，虚伪地打着通过渎神的仪式和祭献而得洁净的旗号（骄傲者正被此引入歧途）把自己显现为生命的中保——既不能分享我们的死亡，也不能从自己的死亡中复活：他的确能将他单一的死亡加于我们双重的死亡；但他绝不可能给予单一的复活，这复活应同时是我们更新的奥迹和终末苏醒的预像。然后，那在神内是活着的，并使自己的死肉复活的，乃是真实生命的中保，祂将那位在神内是死的、乃是死亡的中保者，从那些信靠祂的人的灵魂中驱逐出去，使他不致在里面为王，只可从外面攻击，却不能得胜。并且，祂也向那欺骗者交出自己，让他试探，为要成为一位中保，不仅借着援助，也借着榜样，去战胜他的诱惑。但是，当魔鬼起初虽然竭力从每一个入口潜入祂的内里，却被驱逐出去，在他于洗礼后在旷野中完成了一切迷人的试探之后；因为他在神内是死的，无法进入那在神内是活者的里面，他便转而以其对人之死的渴望，无论以何种方式，去成就那他能成就的死亡，并且被允许将那死亡加于那永生中保从我们领受的必死元素之上。并且，在他任何能有所作为之处，他都在各方面被征服了；而在他外在得到杀死主在肉身之权的那个方面，他内在用来掌控我们的权能也被杀死了。因为藉由那一位无任何罪过在先的死亡，许多死亡中许多罪过的锁链得以解开。这死亡虽非当受的，主却因而为我们偿还，为使那当受的死亡不致伤害我们。因为祂并非因某权柄的强制而被剥去肉身，而是祂自己脱下了肉身。因为无疑地，那位本可以不死、若祂不愿就不死的那一位，却因祂愿意而死；如此，祂公然羞辱了率领者与掌权者，在自身内凯旋胜过它们。因为藉着祂的死亡，那独一无二且最真实的祭献为我们献上了，无论有什么过犯，使率领者与掌权者有权将我们牢牢扣住以偿付其惩罚，祂都洗净、废除、消灭了；并且藉着祂自己的复活，祂也召叫了祂所预定得新生命的我们；祂所召叫的，也使他们成义；祂所成义的，也使他们获得荣耀。因此，魔鬼在肉身的死亡中丧失了人——他原以绝对的权利占据着人，因人因自己的同意而被诱惑，且他统治着人——他自己虽不被血肉的腐败所妨碍，却因人的必死身体的脆弱而既贫乏又软弱；他自以为又富又强，因而骄傲，统治着那仿佛衣衫褴褛、充满苦难的人。因为他将罪人驱入堕落之处，自己却不跟随，而那里，他因跟随而迫使救赎者降下。因此，天主子屈尊进入死亡的共融中成为我们的朋友，因为那仇敌未曾进入死亡，便自以为比我们更好、更大。因为我们的救赎者说：\BibleQuote{人若为自己的朋友舍掉性命，再没有比这更大的爱情了。}所以魔鬼也自以为比主本人更高明，因为主在受苦中屈服于他；因为论到主，也明白那在圣咏中所读的：\BibleQuote{祢使他稍微逊于天使。}因此，祂自己虽是无辜的，却被那以仿佛公义的权利攻击我们的不义者处死，祂便能以最公义的权利战胜他，并掳掠那因罪恶而造成的俘虏，藉着涂抹债卷，释放我们脱离那因罪恶而公义的俘虏生活，并用祂自己的义血——不义地倾流的血——救赎我们这些虽为罪人却将要成义的人。\par

18. 因此，魔鬼直到今日仍嘲弄属于他的人，他向他们显现为假的中保，仿佛通过他的仪式他们能被洁净，或更确切地说，被缠住并淹没；因为他很容易说服骄傲的人嘲笑并轻视基督的死亡，而他自己越远离这死亡，就越被他们视为更圣洁、更神圣。但留在他身边的人很少，因为外邦人承认并以虔诚的谦卑汲取为自己付出的代价，依靠它而放弃他们的仇敌，聚集到他们的救赎主那里。因为魔鬼不知道天主的最高智慧如何利用他的陷阱和愤怒，来带来祂忠信者的救恩，从前者——即属灵受造物的开始——直到后者——即身体的死亡，如此\Quote{从一端达到另一端，以威能和甘饴安排了万有。}因为智慧\Quote{因着她的纯洁，穿越并贯通万有，没有不洁之物能落入她内。}由于魔鬼与肉身的死亡毫无关联——这导致了他极度的骄傲——另一种死亡已在永恒的地狱之火中预备，藉此不仅那些有属土身体的灵魂，而且那些有属气身体的灵魂，都可以受折磨。但骄傲的人——因基督死了而轻视祂，祂是以如此大的代价赎买了我们——不仅将前者死亡，也将人们带回到从原罪而来的可悲的自然状态，并将与魔鬼一同被投入后者死亡。他们因此更喜欢魔鬼而非基督，因为前者将他们投入前者死亡，而他自身因本性的差异并未落入其中；基督却因祂的大慈悲为他们而降临到那里。然而他们毫不迟疑地相信自己比魔鬼更好，并用各种咒骂攻击和谴责魔鬼，同时他们知道魔鬼至少与这种死亡——基督因之被轻视的死亡——毫无关系。他们也不考虑一种可能性：天主的圣言，在自身内保持不变，毫无变化，却藉着取了较低的本性，能够承受某种较低级的事物，这是不洁之神无法承受的，因为它没有属土的身体。因此，虽然他们自己比魔鬼更好，但因他们带着血肉之躯，他们能如此死亡，如同魔鬼确实不能死亡一样——因为魔鬼没有这样的身体。他们大大依靠自己祭献的死亡，却没有觉察到他们是向诡诈而骄傲的神祇献祭；或者即使察觉到了，也认为与这些神祇的友谊对自己有益，尽管他们是背信弃义、嫉妒成性的，其唯一的目的就是阻碍我们的回归。\par

% structure: p0036 source=h2 target=subsection reason=relative-heading-rank; base=h2

\subsection{第十四章——基督是为洁净我们的过犯的最完美祭品。在每个祭献中，有四件事应当考虑。}

19. 他们不明白，即使是最骄傲的神祇本身，也不会因祭献的荣誉而欢喜，除非真正的祭献是归于一真天主的，而他们想要取代祂受敬拜；而且这祭献除非由一位圣洁公义的司铎奉献，否则不能正确举行；除非所奉献之物是从那些为其奉献的人所接受的；并且它也必须是无瑕的，以便能为有瑕者作洁净。至少所有希望为自己向天主奉献祭献的人都渴望这些。那么，有谁像天主的独生子那样公义圣洁的司铎呢？祂无需藉祭献来洁净自己的罪，无论是原罪还是人生中添加的罪。有什么比人的血肉更适合被选为献祭之物呢？有什么比会死的血肉更适合作为牺牲呢？有什么比从贞女胎中诞生、未受任何肉欲贪念玷污的血肉更能洁净必死之人的过犯呢？有什么比我们的祭献血肉——成为我们司铎的身体——更悦纳、更可领受的呢？因此，在每个祭献中当考虑的四件事——向谁奉献、由谁奉献、奉献什么、为谁奉献——同一真实的中保，藉和平祭献使我们与天主和好，祂与所奉献的对象保持合一，使祂为之奉献的人在祂内合一，祂自身同时是奉献者和奉献之物。\par

% structure: p0038 source=h2 target=subsection reason=relative-heading-rank; base=h2

\subsection{第十五章——那些以为自己能凭自己的义德被洁净而得以看见天主的人是骄傲的。}

20. 然而，有些人认为自己能凭自己的义德被洁净，从而默观天主并居住于天主内；而正是他们的骄傲本身最严重地玷污了他们。因为没有任何罪是天主法律更反对的，也没有任何罪能使那最骄傲的神祇——他对低处事物是中保，但对高处事物是障碍——获得更大的统治权：除非要么他的秘密陷阱通过走另一条路被避开，要么如果他通过一个有罪的民族（\OriginalTerm{阿玛肋克}{עֲמָלֵק}，按解释，意即如此）公然发怒，并以战斗阻止进入应许之地的通道，就被主的十字架——由梅瑟伸手所预表——所战胜。这些人以自己义德自许能洁净的原因在于，他们中的一些人曾以心灵之眼穿透整个受造物，并接触到——哪怕只触及极小部分——永恒真理之光；而他们嘲笑许多基督徒尚不能做到后者，这些基督徒暂时仅凭信德生活。但骄傲的人——因此耻于登上木船——从远处眺望海外的故乡，又有何用？或者，谦卑的人无法从如此遥远之处看到它，又有何妨呢？因为他正通过那块木头——对方耻于被承载的木头——实际走向它。\par

% structure: p0040 source=h2 target=subsection reason=relative-heading-rank; base=h2

\subsection{第十六章——关于复活和未来的事，不应咨询古代哲人。}

21. 这些人也因我们相信肉身的复活而责备我们，反倒希望我们在这事上相信他们。仿佛因他们能借着被造之物理解那高超而永恒不变的实体，就有权被咨询关乎可变之物的变迁或时代的有序进程。请问，既然他们最真实地辩论，并以最确凿的证据说服我们：一切暂时之物皆是按着永恒的知识而造，难道他们就能清楚看见这知识本身，或由此收集有多少种动物、每种在起初的种子、成长中的尺度、其受孕、出生、年龄、终结中的数字、依其本性追求与回避相反之物的运动吗？难道他们不是借着地方与时代的历史，而非借着那不变之智慧，来探究这一切，或是信赖他人所记载的经验吗？因此，不足为奇，他们在探究更久远时代的连续性上完全失败，也未能找到那条人类如同顺流而下的河流的进程之终点，以及每个人由此过渡到各自适当终点的路径。这些事是历史学家无法描述的，因为它们远在未来，未经任何人经历或叙述。那些在高等永恒知识上比他人更精通的哲学家，也未能以理智把握这些事；否则他们就不会像现在那样去探究过去的历史事实，而会预知未来；那些能这样做的人，在他们称为占卜者，在我们却称为先知。\par

% structure: p0042 source=h2 target=subsection reason=relative-heading-rank; base=h2

\subsection{第十七章——预知未来事物的多种方式。无论是哲学家，还是古代杰出人物，都不应被咨询有关死者复活之事。}

22. ——虽然先知之名在他们的著作中也并非全然陌生。但最大的区别在于：未来之事是由过去之事的经验推测而来（如医生凭经验将许多预见之事写下来；或如农人、水手也预知许多事；若这类预测提前很久做出，便被视为占卜），或是这些事已经开始走向我们，远处可见，并依观察者感官的敏锐度被宣布，空中权势以此被认为能行占卜（好比一个人从山顶远远看见某人前来，预先告诉平原上的居民）；或是借着圣天使，天主借其圣言和智慧（其中包含过去和未来之事）向某些人预先宣告，或由他们听见后再传给他人；或是某些人的心神被圣神如此高举，以致不通过天使，而亲自在那宇宙的最高顶点看见未来之事的不变原因。［我说，看吧］因为空中权势也\Emph{听见}这些事，或是通过天使的传报，或是通过人；而且只听见那统御万有者认为适合的部分。许多事也藉一种不知情者的本能和内在冲动被预言，如盖法不知自己所说，但身为大司祭，他预言了。\par

23. 因此，关于时代的进程，以及死者的复活，我们不应咨询那些尽可能理解了造物主之永恒的哲学家，在祂内\BibleQuote{我们生活、行动、存在}。因为，他们虽然认识天主，却没有以祂为天主而光荣祂，也不感谢祂，反而自称明智，成为愚昧。既然他们未能如此坚定地将心灵之眼注视于属灵不变本性的永恒，以致能在创造与治理宇宙的智慧本身中看见那些时代的变化——那些变化在那智慧中早已存在，且永远存在，但在此世将要发生，故尚未存在——或看见其中不仅灵魂，而且人类身体直至其适当尺度的完善之改善；既不能在其中看见这些事，他们也不配蒙受圣天使的宣告——无论是通过外在的身体感官，还是通过在心神中显露的内在启示，正如这些事实际显示给我们那些具有真虔诚的祖先，他们预告这些事，借现时迹象或随后应验的事件获得信赖，并赢得权威使人们相信他们关于遥远未来的事，直到世界终结。但是，骄傲而欺诈的空中权势，即使通过其占卜者说出一些关于圣徒的共融与公民身份以及真中保的事——这些事是他们从圣先知或天使那里听来的——其目的也是要借此外来的真理引诱天主的信徒（若可能的话）转向他们自己的谬误。然而，天主藉那些不知自己所说的人行事，为让真理从各方面响起，以帮助信徒，并作证反对不敬虔的人。\par

% structure: p0045 source=h2 target=subsection reason=relative-heading-rank; base=h2

\subsection{第十八章——天主子降生成人，为使我们因信德而洁净，得以被提升至不变的真理。}

24. 既然我们不适宜把握永恒的事物，而罪的污秽压倒了我们——这些污秽是因贪恋暂时之物而获得的，仿佛从可朽的根上自然植入我们内——，我们便需要被洁净。但我们无法被洁净到能与永恒之物相调合，除非通过暂时之物，我们已经与之调合并被紧紧束缚。因为健康与疾病处于相反的两端；但治疗中的中间过程并不会将我们引向完全的健康，除非它在一定程度上与疾病有相似之处。无益的暂时之物只会欺骗病人；有益的暂时之物则接纳那些需要治疗的人，并将他们治愈后交付给永恒之物。理性心灵，如同在被洁净时欠永恒之物默观，同样，在需要被洁净时，它欠暂时之物信德。甚至从前在希腊人中被称为智慧的人中有一位说过：\Quote{真理与信德的关系，如同永恒与有始之物的关系。}他这样说无疑是正确的。因为我们所谓暂时的，他描述为有始的。我们也属于这一类，不仅就身体而言，而且就灵魂的可变性而言。因为任何经历变化之物都不能恰当地被称为永恒的。因此，我们既是可变的，便与永恒相隔。但永生是通过真理应许给我们的；而我们对真理的清晰认识与我们的信德之间的距离，正如可朽与永恒之间的距离。我们现在相信那些为我们而成就于时间中的事，并且通过这信德本身，我们被洁净；这样，当我们达到目睹时，就如真理随信德而来，永恒也可随可朽而来。因此，既然我们的信德将变为真理，当我们获得那应许给我们相信者的事物时，而那应许给我们的是永生；并且真理——不是那将要按照我们的信德而成为的真理，而是那常存的真理，因为在它里面是永恒——那真理曾说过：\BibleQuote{永生就是：他们认识你，唯一的真天主，和你所派遣来的耶稣基督。}那时，当我们的信德通过看见而成为真理，永恒将占有我们现在已改变了的可朽性。直到这事发生，并且为了使之发生——因为我们使相信的信德适应于有始的事物，正如我们在永恒的事物中盼望默观的真理，免得可朽生命的信德与永恒生命的真理不相协调——那与圣父同永恒的真理本身从地上取得了开始，当时天主圣子如此来临，以致成为人子，并亲自接受我们的信德，以便由此将我们引向祂自己的真理，祂如此承担了我们的可朽，却没有失去祂自己的永恒。因为真理与信德的关系，如同永恒与有始之物的关系。因此，我们必须被洁净，好使我们能拥有一个永恒存留的开始，使我们的信德与真理不各有一个开始。我们也不能从有始的状态过渡到永恒，除非我们通过永恒者藉由我们自己的开始与我们结合，而被转移至祂自己的永恒。所以，我们的信德如今在某种程度上已跟随祂——我们所相信的那一位已经升天；祂曾诞生、死去、复活、被提升。在这四件事中，前两件我们自身就认识。因为我们知道人既有开始也会死去。但剩下的两件，即复活和被提升，我们正确地希望在于我们，因为我们相信这些已在于祂。因此，既然在祂内那有始的也已过渡到永恒，在我们内也将同样过渡，当信德达到了真理。因为对于那些如此相信的人，为了使他们能持守在信德之言中，并由此被引向真理，通过真理达至永恒，从而得免于死亡，祂这样说道：\BibleQuote{你们如果存留在我的话内，就确是配作我的门徒。}仿佛他们必会问：有什么果实？祂继续说：\BibleQuote{你们必认识真理。}又仿佛他们会说：真理对可朽之人有什么好处？祂说：\BibleQuote{真理必使你们得自由。}从什么得自由？除了从死亡、腐败、可变性之外？因为真理常存不朽、不腐、不变。但真正的不朽、真正的永存不朽、真正的永不改变，就是永恒本身。\OriginalTerm{永恒}{eternity}\OriginalTerm{信德}{faith}\OriginalTerm{真理}{truth}\OriginalTerm{有始}{beginning}\OriginalTerm{可朽}{mortality}\OriginalTerm{永不改变}{unchangeableness}\OriginalTerm{天主圣子}{Son of God}\OriginalTerm{人子}{Son of man}\par

% structure: p0047 source=h2 target=subsection reason=relative-heading-rank; base=h2

\subsection{第十九章——圣子如何被派遣并预先宣告。在其肉身降生的派遣中，祂是如何成为较小的，而不损害祂与圣父的平等。}

25. 看哪，这就是为什么天主圣子被派遣；更确切地说，看哪，什么是天主圣子被派遣。凡是在时间中成就的事，为产生信德，好使我们得以被洁净以默观真理，在那些有始的事物中，它们出自永恒，又归于永恒——这些事或是此派遣的证据，或是天主圣子派遣本身。但这些证据中，有些预告祂将要来临，有些证明祂已经来临。因为那创造万有的本身成为受造物，必然在万有中找到见证。除非通过许多见证者的派遣来宣讲那一位，人就不必为许多的派遣而受约束。除非有这样的证据，在谦卑的人看来似乎伟大，否则人们不会相信，那伟大的祂能使渺小的人变得伟大，而祂以谦卑的身份被派遣给谦卑的人。因为天、地及其中的一切，作为天主圣子的工程，远大于那些在祂见证中爆发的征兆和奇事。然而，为了使谦卑的人相信这些伟大的事是由祂所成就的，他们竟在那些谦卑的事上战抖，仿佛它们是伟大的。\OriginalTerm{受造物}{creature}\OriginalTerm{派遣}{mission}\OriginalTerm{见证}{witness}\OriginalTerm{伟大的}{great}\OriginalTerm{谦卑的}{lowly}\par

\Quote{因此，时期一满，天主就派遣了祂的圣子，生于女人，生于法律之下；}祂如此卑微，以至于祂是\Quote{成了}的；因此，祂是这样被派遣的，即祂是成了的。因此，如果较大的派遣较小的，我们也就承认祂成了较小的；祂既然较小，也就成了；祂既然成了，也就被派遣了。因为\Quote{祂派遣了自己的圣子，生于女人。}然而，因为万物都是藉着祂造的，不仅是在祂成了并受派遣之前，而且在万物存在之前，我们就明认同一位祂与派遣者相等，而祂被我们称为较小的，因为祂被派遣了。那么，在时期一满，祂应当被派遣之前，即祂尚未被派遣之时，祂如何能被视为与圣父相等，而被祖先们看见呢？当时向他们显现的是一些天使的显现。因为祂对\Person{斐理伯}说：\Quote{斐理伯！这么长久的时间，我同你们在一起，而你们还不认识我吗？谁看见了我，就是看见了父；}当时\Person{斐理伯}当然像其余的人一样看见了祂，甚至那些在肉身中把祂钉在十字架上的人也看见了祂；除非是因为祂既被看见，又未被看见？祂被看见了，是如同祂在派遣中成了人那样；祂未被看见，是如同万物藉着祂被造那样。或者，祂又说：\Quote{那接受我的命令并遵守的，就是爱我的；爱我的，必蒙我父爱他，我也要爱他，并将我自己显示给他。}当祂在人的眼前显现的时候，祂如何能说这话呢？除非是因为祂将那在时期一满圣言所成的血肉，奉献给我们信德去接受；但却将那藉着祂万物受造的圣言本身，保留给因信德而洁净的心灵在永恒中默观。\par

% structure: p0050 source=h2 target=subsection reason=relative-heading-rank; base=h2

\subsection{第二十章．派遣者与受派遣者相等。为何说圣子由圣父所派遣。论圣神的使命。祂如何被派遣及由谁派遣。圣父是整个天主性的根源。}

27. 但如果圣子被称为由圣父所派遣，是因为一位是圣父，另一位是圣子，这绝不会妨碍我们相信圣子与圣父是平等、同体和同永恒的，然而仍被圣父作为圣子派遣。并非因为一位较大，另一位较小；而是因为一位是父亲，另一位是儿子；一位是生者，另一位是受生者；一位是派遣者，受派遣者从祂而来；另一位是从派遣者而来的那一位。因为圣子来自圣父，而非圣父来自圣子。按照这种方式，我们现在可以理解：圣子之所以被称为被派遣，不仅因为\BibleQuote{圣言成了血肉}，而且正是因为被派遣，圣言才成了血肉，并透过祂的身体临在完成那些所记载的事；也就是说，不仅祂被理解为作为人而被派遣（圣言成了人），而且圣言本身也曾被派遣，为要成为人；因为祂被派遣，不是基于任何能力、实体或祂内任何与圣父不平等之处，而是基于这一点：圣子来自圣父，而非圣父来自圣子；因为圣子是圣父的圣言，也称为祂的智慧。因此，如果祂被派遣，不是因为祂与圣父不平等，而是因为祂是\BibleQuote{来自全能天主荣耀的纯粹流溢（\Emph{manatio}）}，这有什么可惊奇的呢？因为那里，流溢者与从它流溢出者，具有同一实体。因为它不是像水从土或石的孔隙中流出来那样流溢，而是像光从光中发出。因为经文说：\BibleQuote{因为她是永恒之光的反映，} 这岂不是说：她是永恒之光的本体？因为光的反映若不是光本身，又是什么呢？所以，她与它所出自的光是同永恒的。但说\Quote{光的反映}比说\Quote{来自光的光}更可取，以免流溢者被认为比它所从出的光更暗淡。因为当人听到光的反映就是光本身时，更容易相信前者藉后者照耀，而不认为后者照耀得更弱。但是，因为不需要警告人不要认为生出了另一位的那光较小（因为没有异端者曾敢这样说，也不相信有人敢这样做），圣经以另一种思想回应，即那流溢之光可能显得比它所从出之光更暗淡；圣经排除了这种猜测，说\BibleQuote{它是那光的反映}，即永恒之光的反映，从而表明它是相等的。因为如果它较小，那么它就是暗黑，而非反映；但如果它较大，则不可能从它发出，因为它不能超过它所由出的那一位。因此，由于它从它发出，故它不大于它；又由于它不是它的暗黑，而是它的反映，故它不小于它：因此它是相等的。我们也不应为此困扰：它被称为来自全能天主荣耀的纯粹流溢，似乎它本身不是全能的，而是来自全能者的流溢；因为不久后论及它说：\BibleQuote{她是唯一的，却能成就一切。} 但是，除了能做一切事的那一位以外，谁还能是全能者呢？因此，它是由它所流出的那一位派遣的；因为爱慕和渴望它的人这样寻求它。他说：\BibleQuote{请从祢神圣的诸天，从祢荣耀的宝座派遣她，使她与我同在，与我一同劳作}；也就是说，教导我热心地劳作，使我不至厌倦地劳作。因为她的劳作就是德行。但是，她以某一种方式被派遣，为与人同在；她以另一种方式被派遣，为使她本身成为人。因为，\BibleQuote{进入圣洁的灵魂，她使他们成为天主的朋友和先知}；同样，她也充满圣天使，并藉他们完成适合此类服务的一切事。但是，当时期一满，她被派遣，不是为了充满天使，也不是为了成为天使（除非是宣布圣父的旨意，那也是她自己的旨意）；也不是为了与人同在或住在人内，因为这以前已经发生过，无论是在族长还是先知中；而是为使圣言本身成为血肉，即成为人。这未来的奥迹，当启示时，也要成为那些智慧圣贤（他们在那位由童贞女出生之前，由妇女所生）的救恩；当这奥迹成就并显明时，就是所有信、望、爱之人的救恩。因为这就是\BibleQuote{伟大的虔敬奥迹：祂在肉身内显示，在圣神内成义，被天使看见，被传于外邦人，在世上被相信，被接到光荣中。}\par

28. 因此，天主的圣言是由祂所派遣的，圣言是祂的圣言；圣言是由祂所派遣的，圣言是由祂所生的（\OriginalTerm{生}{genitum}）；派遣者是那生者，被派遣者是那被生的。然后，当祂被各人领会和感知时，祂便被派遣到各人那里，按各人所能够领会和感知的程度，相称于理性灵魂的理解力，或是在走向天主的途中，或已在天主内达于成全。因此，圣子之所以被称为被派遣，严格来说，并非因祂由圣父所生，而是因那成了血肉的圣言显现于世界，故祂说：\BibleQuote{我出自父，来到了世界上}；或是因祂不时被各人的心灵所感知，正如经上所说：\BibleQuote{求你派遣她，使她与我同在，与我一同劳苦}。那从永远所生的（\OriginalTerm{生}{natum}）是永远的，因为\BibleQuote{它是永远光明的辉映}；但不时被派遣的，是各人所领会的。但当天主子在肉身中显现时，祂在时期一满，由女人所生，被派遣到这个世界。\BibleQuote{因为，在天主的智慧中，世界凭着智慧没有认识天主（\InnerQuote{光在黑暗中照耀，黑暗却没有领会它}），天主\InnerQuote{愿意以宣讲的愚妄来拯救那些相信的人}，并使圣言成了血肉，居住在我们中间}。但当时不时地祂被各人的心灵所感知而出来时，祂确实被称为被派遣，但并非进入这世界；因为祂不以感官方式显现，即祂不呈现于肉体感官。因为就连我们自己，就我们以心灵尽其所能把握永恒之事而言，也不在这世界；所有义人的心灵也不在这世界，即使是那些仍在肉身中生活的人，就他们具有对神圣事物的辨识力而言。但圣父却不可被称为被派遣，即使祂不时被某人所领会，因为祂没有由之而存在或由之而出的那一位；因为智慧说：\BibleQuote{我出自至高者的口}，论到圣神也说：\BibleQuote{祂由圣父所发}，但圣父却不来自任何一位。\par

29. 因此，正如圣父生了圣子，圣子被生；同样，圣父派遣了圣子，圣子被派遣。但正如生者和被生者是一体，同样，派遣者和被派遣者也是一体，因为圣父和圣子是一体。圣神也和他们是一体，因为三者是一体。因为对圣子而言，被生意谓来自圣父；同样，被派遣意谓被认识为来自圣父。对圣神而言，作为天主的恩赐意味着由圣父所发；同样，被派遣意味着被认识为由圣父所发。我们也不能说圣神不由圣子所发，因为同一圣神被称为既是圣父的圣神，也是圣子的圣神，并不是没有理由的。我也看不出祂在向门徒们吹气，并说\BibleQuote{你们领受圣神}时，还想表示别的什么意思。因为那从身体发出并带有身体接触感的肉体气息，并不是圣神的实质，而是藉一个合适的标记宣布：圣神不仅由圣父所发，也由圣子所发。因为即使是最疯狂的人也不会说，祂在向他们吹气时所赐的圣神，与祂升天后所派遣的，是两个不同的神。因为天主的圣神是同一的，是圣父的圣神和圣子的圣神，是圣神，祂在一切内运行一切。但祂被赐予两次，确实是一个有意义的安排（\OriginalTerm{安排}{economy}），我们将在适当的时候，按主的恩赐来讨论。因此，主所说的\BibleQuote{我将从父那里派遣到你们那里去的}，显示圣神既属于圣父也属于圣子；因为当祂说了\BibleQuote{父所要派遣的}时，又加上了\BibleQuote{因我的名}。然而祂并没有说\BibleQuote{父将从我派遣}，正如祂所说\BibleQuote{我将从父那里派遣到你们那里去的}，这显示圣父是整个天主性的源头（\OriginalTerm{源头}{principium}）——或者，若更准确地说，是整个神性的源头。因此，那由圣父和圣子所发的，被引回圣子由之而生的（\OriginalTerm{生}{natus}）那一位。至于福音作者所说的\BibleQuote{圣神还没有赐下，因为耶稣还没有受到光荣}，这该如何理解呢？除非因为基督受光荣之后，圣神特殊的赐予或派遣，乃是前所未有的一种形式。因为先前并非完全没有，而是未曾有过这样的。因为如果圣神先前没有赐下，那么说话的先知们是如何被充满的呢？而圣经明说，并在许多地方显示，他们是由圣神说话。同样，论到洗者若翰也说：\BibleQuote{他还在母胎中就要充满圣神}。他的父亲匝加利亚被发现充满了圣神，从而说出有关他的那些话。玛利亚也充满了圣神，从而预言有关她胎中所怀的主的那些事。西默盎和亚纳也充满了圣神，从而认出婴孩基督的伟大。那么，怎么又说\BibleQuote{圣神还没有赐下，因为耶稣还没有受到光荣}呢？除非因为圣神的这个赐予、或恩赐、或使命，在其来临本身具有一种前所未有的特殊性。因为我们从未读到过，人们藉着圣神降临时会说他们不懂的语言；但那时却发生了，因为需要以可见的记号显明祂的来临，以显示普世万民、所有说不同语言的民族，都要藉着圣神的恩赐而信仰基督，以应验圣咏上所唱的：\BibleQuote{没有言语，没有音词，却可听到他们的声音；它们的声音传遍普世，它们的话语达于地极}。\par

30. 因此，当时期一满，人便与天主的圣言结合，并在某种意义上混合为一，成为一位，天主之子由女人所生，被派遣到这世界，为人类的缘故，也成为人子。天使的本性能够预先象征这事，以便预先宣告，但不能据为己有，以致成为这第二位本身。\par

% structure: p0055 source=h2 target=subsection reason=relative-heading-rank; base=h2

\subsection{第二十一章——论圣神的可感显现，及天主圣三的同永性。已述与待述。}

但关于圣神的可感显现，无论是藉着鸽子的形状，还是藉着火焰的舌头，在被征服和服侍的受造物藉着暂时的运动和形式，显现出祂与圣父及圣子同永、且与他们一样不变的实体时，它并没有如此结合以致成为祂的一位，正如圣言成为的肉身那样；我不敢说以前从未有过此类事。但我敢大胆说，圣父、圣子和圣神，同是一个实体，天主创造者，全能的天主圣三，不可分地工作；但这一点不能被远为卑下、尤其是肉体的受造物完全不可分地彰显：正如圣父、圣子和圣神不能被我们的言语（这无疑是肉体的声音）所称呼，除非在它们各自适当的时间间隔中，由每个词语的适当音节所占据的明确分离所划分。因为在他们自己的实体中，三者是一，圣父、圣子和圣神，完全相同，没有任何时间运动，超越整个受造物，没有任何时间和空间的间隔，同时且从永远到永远是一，仿佛永恒本身，这永恒并非没有真理和爱德。但是，在我的言语中，圣父、圣子和圣神被分开，不能同时被称呼，并且在无形文字中占据各自适当的位置。并且，当我称呼我的记忆、理智和意志时，每个名称分别指向每个，但每个都由三者共同说出；因为这三个名称中没有一个是不同时由我的记忆、理智和意志[由整个灵魂]一起说出的；同样，天主圣三共同造就了圣父的声音、圣子的肉身和圣神的鸽子，而这些事物分别被归于每一位。借着这种相似性，可以在一定程度上看出，天主圣三在自身内是不可分的，却借着可见受造物的显现以可分的方式彰显；并且，在那些据称专属圣父、圣子或圣神之显现的每一事物中，天主圣三的运作也是不可分的。\par

31. 如果现在有人问我，在天主的圣言降生成人之前，那些应该预表它将要来临的言语或可感的形体和显现是如何造成的，我回答说，天主借着天使造成了那些事；这一点我想我已经用圣经的见证充分表明了。如果问我降生成人本身是如何实现的，我回答说，天主的圣言本身成了血肉，也就是说，成了人，但并没有转变和改变成为所形成的东西；而是如此形成，以致那里不仅有天主的圣言和人的血肉，而且有人的理性灵魂，并且这整体因天主而称为天主，因人而称为人。如果这难以理解，那么灵魂必须借着信德，借着越来越远离罪恶、行善功，并借着以神圣渴望的呻吟祈祷来炼净；以便在天主助佑下进步，既能理解又能爱慕。如果问我，圣言降生成人之后，圣父的声音或藉以显现圣神的形体是如何产生的：我确实不怀疑这是通过受造物完成的；但是否只是肉体和可感的，或者还有理性的或理智的灵（因为有人选择这样称呼希腊人称为\OriginalTerm{理智的}{\GreekText{νοερόν}}的东西），当然不是如此结合成为一位（因为谁能说，藉以发出圣父声音的受造物可以说是天主圣父，或者藉以显现圣神为鸽子形状或火焰舌头的受造物可以说是圣神，正如天主之子是由童贞女所生的那个人那样？），而只是服务于传达天主认为必要的那些指示；或者是否有其他理解：这很难发现，也不宜贸然肯定。然而我看不出这些事怎能没有理性或理智的受造物而完成。但还不是时候解释我为何如此认为，只要主赐予我力量；因为必须先讨论和驳斥异端者的狡辩，这些狡辩不是从圣经中提出的，而是来自他们自己的理由，并以此他们认为强迫我们按照自己的意愿理解圣经中关于圣父、圣子和圣神的见证。\par

32. 但是，现在，我想已经充分表明，圣子并不因为被圣父派遣而较小，圣神也不因为圣父和圣子两者派遣祂而较小。因为圣经中确立这些事，或是为了可见的受造物，或是为了向我们的思想推荐[天主内的]发出；但并非出于不平等、不同等或实体不相似；因为，即使天主圣父愿意藉着服从的受造物以可见方式显现，那么说祂是被祂所生的圣子或从祂发出的圣神所派遣，也是极其荒谬的。因此，就让这成为本书的结尾。从此以后，在主帮助下，我们将在其余部分看看异端者那种狡猾的论证是什么性质，以及如何驳斥它们。\par

\SourceRecord{Translated by Arthur West Haddan.；Nicene and Post-Nicene Fathers, First Series；Vol. 3.；Edited by Philip Schaff.；Buffalo, NY: Christian Literature Publishing Co.,；1887.；<http://www.newadvent.org/fathers/130104.htm>.}

% structure: p0001 source=h1 target=section reason=page-title

\section{论天主圣三（第五卷）}

\EditorialNote{继续处理异端者所提出的论据，这些论据并非来自圣经，而是来自他们自己的理性。那些认为圣父与圣子的实体不相同的人被反驳，因为他们认为凡论及天主的，都是按其实体而论；因此，生育与被生育，或被生与未生，既不相同，便是不同的实体；然而此处证明，并非凡论及天主的都是按其实体而论，如祂被称为善和大，或任何其他就其自身而论及祂的；但也有一些是相对地论及祂的，即不针对祂自身，而是针对非祂自身的事物，如祂因圣子而被称为父，因服事祂的受造物而被称为主；在此情况下，若如此相对地论及祂，即针对非祂自身的事物，甚至论及发生在时间中的事，例如\Quote{主，祢成了我们的避难所}，但天主并没有发生任何变化以改变祂，祂自身在其本性或本质上绝对不变。}\par

% structure: p0003 source=h2 target=subsection reason=relative-heading-rank; base=h2

\subsection{第一章——作者向天主祈求什么，向读者祈求什么。在天主内不得思想任何有形体或可变化的事物。}

1. 我现在开始谈论一些问题，这些问题无论任何人，或至少我自己，都不能完全如其所想地言说；即使我们思考天主圣三时，我们的思想（如我们所感）也远远不及我们所想的天主，也不能如祂之所是地理解祂；但祂被看见，正如经上所载，即使是像宗徒保禄那样伟大的人，也是\BibleQuote{藉着镜子，模糊不清地}\BibleRef{格前 13:12}。首先，我向主天主自己祈祷，我们应当时时思念祂，却不能相称地思念祂，应当时时赞美祂，却无任何言语足以表达；求祂赐我理解和解释我所计划之事的能力，并宽恕我若在任何事上冒犯。因为我不仅记得我的渴望，也记得我的软弱。我也请求读者宽恕我，当他们觉察到我在言说那些他们自己理解得更清楚、或由于我言语的晦涩而未能理解的事上，有渴望而非能力时；正如我宽恕他们因自身迟钝而不能理解的事一样。\par

2. 我们将更容易彼此宽恕，如果我们知道，或至少坚信并持守，凡论及不变、不可见、绝对拥有生命并自给自足的本性，不可按照可见、可变、有死或非自足之物的习惯来衡量。然而，尽管我们努力却仍未能把握和认识那些属于我们肉体感官范围之内的事物，或我们在内心中自身是什么，但虔诚的信德并不羞耻地渴求那些超性的、不可言说的事物：我说的是信德，并非因自身能力的骄傲而膨胀，而是因创造者和救主自己的恩宠而燃烧。因为人若尚未理解他用来理解天主的那个理解本身，又怎能理解天主呢？如果他已经理解了这个，就让他仔细思考，在他自己的本性中没有什么比这更好；并让他看看，他是否在其中看到任何形式的轮廓、颜色的光辉、空间的大小、部分的距离、体积的延伸、或任何通过位置间隔的运动，或任何此类东西。当然，在我们自身本性中最好的东西，即我们自己的理智（我们按自己的能力借此理解智慧）中，我们找不到这一切。因此，在我们自身最好的东西中所没有的，我们不应在那比我们最好的东西更好得多的祂身上寻找；这样，我们若可能，并尽我们所能，就可以理解天主是善而没有性质，伟大而没有数量，创造者却不缺乏任何东西，统治却不在任何位置，包容万有而不\Quote{拥有}它们，整个地无处不在却没有位置，永恒而没有时间，创造可变之物而自身不变，没有情感。谁这样思想天主，尽管他尚不能完全发现祂是什么，但仍虔诚地留意，尽他所能，对祂思想任何祂所不是的东西。\par

% structure: p0006 source=h2 target=subsection reason=relative-heading-rank; base=h2

\subsection{第二章——天主是唯一不变的实体}

3. 然而，祂毫无疑问是一个实体，或者，若更恰当的话，一个本质，希腊人称之为\OriginalTerm{本质}{\GreekText{οὐσία}}。因为正如智慧从有智慧而得名，知识从知道而得名，同样，从存在而来的是我们所称的本质。那对祂的仆人梅瑟说\BibleQuote{我是自有者}\BibleRef{出 3:14}和\BibleQuote{你要这样对以色列子民说：那\InnerQuote{自有者}派遣我到你们这里来}\BibleRef{出 3:14}的，有谁比祂更\Quote{是}呢？但其他被称为本质或实体的事物，容许有附属之物，由此产生或大或小的变化。但在天主身上，不可能有这种附属之物；因此，那为天主的，是唯一不变的实体或本质，存在本身（本质之名由此而来）最特别、最真实地属于祂。因为变化的事物不保持其自身的存在；而能变化的事物，即使实际上没有变化，也能不是它曾有过的；因此，那不仅不变化，而且根本不能变化的事物，唯一真正地、毫无困难或犹豫地属于\Quote{存在}的范畴。\par

% structure: p0008 source=h2 target=subsection reason=relative-heading-rank; base=h2

\subsection{第三章——亚略异端者从\Quote{生}和\Quote{未生}这两个词提出的论据被驳斥}

4. 因此——现在我们开始回答那些反对我们信德的人，论及那些既不是按所想说的，也不是按实际所想的事——在亚略派人士惯常反对天主教会信仰的许多论点中，他们似乎主要提出这一条，作为他们最狡猾的诡计，即：凡论及或理解天主的，都不是按偶然而言，而是按本体而言，因此，不受生属于圣父的本体，受生属于圣子的本体；但不受生与受生是不同的；所以圣父的本体与圣子的本体是不同的。对此我们回答说：如果凡论及天主的，都是按本体而言，那么\BibleQuote{我与父原是一体}，也是按本体而言。因此圣父与圣子是一个本体。或者，如果这并非按本体而言，那么就有一些关于天主的事不是按本体而言的，因此我们不再被迫按本体来理解不受生与受生。论及圣子也说：\BibleQuote{他虽具有天主的形体，却没有以自己与天主同等为应当把持不舍的}。我们问：相等是在哪方面？如果祂不是按本体而言相等，那么他们就承认有些关于天主的事可以不按本体而言。那么让他们承认，不受生与受生不是按本体而言的。如果他们不承认这一点，理由是凡论及天主的都必须按本体而言，那么圣子就是按本体与圣父相等。\par

% structure: p0010 source=h2 target=subsection reason=relative-heading-rank; base=h2

\subsection{第四章——偶然常暗示事物有所变化。}

5. 所谓偶然通常意味着它可以通过事物本身的变化而失去，事物对它是偶然的。因为虽然有些偶然被称为不可分离的，希腊文称为\OriginalTerm{不能分离的}{\GreekText{ἀχώριστα}}，如黑色对乌鸦的羽毛；但羽毛会失去那种颜色，并非因为它作为羽毛存在时，而是因为羽毛不是永远存在的。因此事物本身就是可变的；每当那个动物或那片羽毛消灭，整个身体改变并化为泥土时，它当然也失去了那种颜色。虽然那种称为可分离的偶然，也可能通过变化失去，而不是通过分离；例如，黑色被称为人的头发的可分离的偶然，因为头发继续存在可以变白；但是，如果仔细考虑，就很明显，并不是有东西通过分离离开头部而变白，好像黑色从那里离开，去了别处，白色取而代之，而是那里的颜色质料被转变和变化了。因此，在天主内没有偶然，因为没有任何可变的或可能失去的东西。但是，如果你愿意也将那虽然不可失去，却可以减少或增加的东西称为偶然——例如，灵魂的生命：只要它是灵魂，它就活着，因为灵魂永远是，它就永远活着；但是因为当它有智慧时活得更多，当它愚昧时活得较少，这里也发生一些变化，并不是生命缺席如同智慧缺席愚昧人一样，而是生命较少——这种变化也不会发生在天主身上，因为祂保持完全不变。\par

% structure: p0012 source=h2 target=subsection reason=relative-heading-rank; base=h2

\subsection{第五章——论天主的事，没有按偶然说的，而是按本体或按关系说的。}

6. 因此，论及祂的事，没有一项是从偶然说的，因为没有任何偶然而有的事属于祂；然而，所有论及祂的事，也不是全按本体说的。因为在受造而变化的事物中，凡不是按本体说的，必然按偶然来说。因为所有可能失去或减少的事，无论是大小还是性质，对它们来说都是偶然的；同样，按关系说的，如友谊、亲属关系、服务、相貌、平等和其他这类；还有位置和状态、地点和时间、行动和受动。但在天主内，没有什么事是按偶然说的，因为祂内没有可变的东西；然而所有论及祂的事，也不是全按本体说的。因为有些事是按关系说的，如圣父对圣子、圣子对圣父的关系，这不是偶然；因为一是永远的父亲，另一是永远的儿子：但\Quote{永远}的意思不是从圣子\OriginalTerm{受生}{natus}的时候起，以至于圣父不再是父亲，因为圣子从不停止为子，而是因为圣子一直是受生的，从未开始为子。但如果祂曾在某个时间开始存在，或在某个时间停止为子，那么祂就是按偶然被称为子。但如果圣父之所以被称为圣父，是相对于祂自己而不是相对于圣子；圣子之所以被称为圣子，是相对于祂自己而不是相对于圣父；那么二者都按本体被称为父和子。但因为圣父之所以被称为圣父，仅仅因为祂有一个圣子；圣子之所以被称为圣子，仅仅因为祂有一个圣父；所以这些不是按本体说的；因为每一方都不是相对于自己，而是相互和彼此相对而言的；也不是按偶然说的，因为被称为圣父和被称为圣子，对它们来说是永远的和不变的。因此，虽然为父和为子是不同，但它们的本体并无不同；因为它们是按关系而不是按本体被称呼的，而这种关系不是偶然，因为它是不变的。\par

% structure: p0014 source=h2 target=subsection reason=relative-heading-rank; base=h2

\subsection{第六章——回答异端者对同一词\Quote{受生}和\Quote{不受生}的诡辩。}

7. 但若他们以为可以这样回答这个推理——即\Quote{圣父之所以被称为圣父，确是相对于圣子，而圣子相对于圣父；但他们说，\OriginalTerm{无始无终者}{unbegotten}和\OriginalTerm{受生者}{begotten}是相对于自身而言，并非彼此相对；因为称祂为无始无终者与称祂为圣父并非同一回事，因为即使祂没有生圣子，也毫不妨碍我们称祂为无始无终者；并且，若有人生了一个儿子，他本人并不因此就是无始无终者，因为由他人所生的人，他们自己也生养他人；因此他们说，圣父是相对于圣子而被称为圣父，圣子是相对于圣父而被称为圣子，但无始无终者是相对于自身而言，受生者也是相对于自身而言；因此，若凡是相对于自身而言的，都是按\OriginalTerm{本体}{substantia}而言，而无始无终者与受生者不同，那么本体也就不同}。——如果这是他们所说的，那么他们就不明白，他们确实在\Quote{无始无终者}这个术语上说了些需要更仔细讨论的话，因为没有人因为是无始无终者而成为父亲，也没有人因为是父亲而是无始无终者，因此他们以为祂被称为无始无终者是相对于自身而非其他；但另一方面，他们以惊人的盲目未能察觉，没有人能被说成是受生者，除非相对于某物。这种诡辩必须这样驳斥：迫使他们自己说明，圣子是在什么方面与圣父相等，是按相对于自身而言的方面，还是按相对于圣父而言的方面？因为圣子之与圣父相等，并不是按相对于圣父而言的方面，因为相对于圣父，祂被称为圣子，而圣父不是圣子，而是圣父。因为圣父和圣子彼此之间的称呼方式，不同于朋友和邻居之间的称呼方式；因为朋友是相对于他的朋友而被称为朋友，若他们彼此相爱相等，那么同一种友谊存在于两者中；邻居是相对于邻居而被称为邻居，因为他们彼此相等地互为邻居（因为每一个人是另一个人的邻居，其程度与另一个人是他的邻居相同），所以同一种邻里关系存在于两者中。但圣子并非相对于圣子而被称呼，而是相对于圣父，因此圣子之与圣父相等，并不是按相对于圣父而言的方面；那么剩下的就是，祂是按相对于自身而言的方面相等。但凡是相对于自身而言的，都是按本体而言；因此，祂是按本体相等；所以两者的本体是同一的。但当圣父被说是无始无终者时，并不是在说祂是什么，而是在说祂不是什么；并且，当一个相对词被否定时，并不是按本体而否定，因为该相对词本身并非按本体而被肯定。\par

% structure: p0016 source=h2 target=subsection reason=relative-heading-rank; base=h2

\subsection{第七章 —— 否定的添加并不改变范畴}

8. 这须用例子来说明。我们首先应注意，\OriginalTerm{受生}{begotten}一词与\OriginalTerm{圣子}{son}一词所指相同。因为所以是圣子，是因为受生；因为是圣子，所以必然是受生。因此，\OriginalTerm{非受生}{unbegotten}一词是宣示祂不是圣子。但受生与非受生都是适当使用的术语；而在拉丁文中，我们可用\Quote{filius}一词，但语言的惯用法却不允许我们说\Quote{infilius}。然而，若称之为\Quote{non filius}，在意义上并无差别；正如称之为\Quote{non genitus}而非\Quote{ingenitus}也完全相同。同样，邻人与朋友作为关系词，但我们不能说\Quote{invicinus}如同我们可以说\Quote{inimicus}那样。因此，在谈论某事物时，我们不应考虑我们本国语言的用法允许或不允许什么，而应考虑事物本身清楚显示的意义。所以，我们不要再称之为非受生，尽管拉丁文中可以这样称呼；但取而代之，我们称之为非受生，这意义相同。难道这不就是等于说祂不是圣子吗？现在，这个否定前缀的添加，并不使原本关系性的陈述变为实体性的陈述；它所否认的，只是原本所肯定的，正如在其他范畴中一样。当我们说祂是人，我们指的是实体。因此，说祂不是人的人，所表达的并不是另一种范畴，而只是否定那个范畴。所以，正如我凭实体肯定说祂是人，同样我也凭实体否定说祂不是人。当问及祂有多大时，我说祂是四足（即四英尺长），我便是凭数量肯定；而说祂不是四足的人，便是凭数量否定。我说祂是白的，便是在性质上肯定；若我说祂不是白的，便是在性质上否定。我说祂是近的，便是在关系上肯定；若我说祂不是近的，便是在关系上否定。我说祂躺着，便是在情状上肯定；若我说祂不躺着，便是在情状上否定。我凭状态说话，当我说祂是武装的；我凭状态否定，当我说祂不是武装的；这等于说祂是未武装的。我凭时间肯定，当我说祂是昨天的；我凭时间否定，当我说祂不是昨天的。当我说祂在罗马，便是凭地点肯定；我说祂不在罗马，便是凭地点否定。我凭行动范畴肯定，当我说祂打击；但若我说祂不打击，便是凭行动否定，以宣告祂不如此行动。当我说祂被打击，便是凭承受范畴肯定；同样，当我说祂不被打击，便是凭同一范畴否定。总之，没有一种范畴，当我们愿意肯定什么时，却不证明若加否定前缀便是否定同一范畴的。既然这样，如果我凭实体说圣子，便是凭实体肯定；说非圣子，便是凭实体否定。但是，因为我说祂是圣子时是关系性的肯定，因为我指向圣父；所以，若我说祂不是圣子，便是关系性的否定，因为我将同一否定指向圣父，意在宣告祂没有父。但如果被称为圣子与被称作受生完全等同（如前所述），那么被称为非受生就完全等同于被称为非圣子。但我们说祂不是圣子时是关系性的否定，因此我们说祂不是受生时也是关系性的否定。此外，什么是非受生呢？难道不是非受生吗？因此，当祂被称为非受生时，我们并未脱离关系范畴。因为正如受生不是就其自身而言，而是说祂\Emph{来自}一位生者；同样，当一个人被称为非受生时，也不是就他自身而言，而是在宣告他不\Emph{来自}一位生者。然而，两个意义都归属于同一范畴，即关系范畴。但关系性的断言并不指称实体，因此，尽管受生与非受生是不同的，它们并不指称不同的实体；因为，如同圣子指向圣父，非圣子指向非圣父，所以必然地，受生指向生者，非受生指向非生者。\par

% structure: p0018 source=h2 target=subsection reason=relative-heading-rank; base=h2

\subsection{第八章——凡凭实体论及天主的，均为分别论及每一位，且共论及圣三本身。天主内一个本质，而在希腊文为三个性体，在拉丁文为三个位格。}

9. 因此，我们首先要持守这一点：凡是就那至高无上、神圣尊威本身所说的，都是就本体而言的；但凡是就关系所说的，都不是就本体，而是相对而言的。并且，同一本体在圣父、圣子和圣神内的效果是：凡就每一位本身所说的，都应理解为单数，而非复数总和。因为圣父是天主，圣子是天主，圣神是天主，无人怀疑这是就本体而言；然而我们不说至高的天主圣三本身是三位神，而是一位天主。同样，圣父是伟大的，圣子是伟大的，圣神是伟大的；但不是三位伟大的，而是一位伟大的。因为经上记载的不仅是关于圣父——正如他们错误地认为的那样——而是关于圣父、圣子和圣神：\Quote{唯独祢是伟大的，祢是唯一天主。}圣父是善的，圣子是善的，圣神是善的；但不是三位善的，而是一位善的，关于祂经上说：\Quote{除一位，即天主外，没有善的。}因为主耶稣，唯恐那称祂为\Quote{善师}的人仅将祂视为人，所以祂不说：除圣父外，没有善的；而是说：\Quote{除一位，即天主外，没有善的。}因为圣父由\Quote{父}这个名号表明自身；而\Quote{天主}这个名号则表明祂自己、圣子和圣神，因为天主圣三是一位天主。至于位置、状态、地点和时间，并非真正用在天主身上，而是借喻和类比。因为经上说祂居于革鲁宾之间，这是就位置而言；又说深渊遮蔽祂如衣服，这是就状态而言；又说：\Quote{祢的岁月没有穷尽}，这是就时间而言；又说：\Quote{我若升到天上，祢也在那里}，这是就地点而言。至于行动（或创造），或许最真实地用于独一天主，因为只有天主创造，祂自己不被创造。就祂是天主的那一本体而言，祂也不受情感的影响。因此，圣父是全能的，圣子是全能的，圣神是全能的；但不是三位全能的，而是一位全能的：\Quote{因为万有出于祂，借着祂，归于祂；愿光荣归于祂。}所以，凡是就天主本身所说的，既单独对每一位说，即对圣父、圣子和圣神，也对天主圣三整体说，不是复数，而是单数。因为对天主而言，存在与伟大并非两回事，存在与伟大是同一回事；因此，正如我们不说三个本质，我们也不说三个伟大，而是一个本质和一个伟大。我说的本质，希腊文称为\OriginalTerm{本质}{\GreekText{οὐσία}}，我们通常称之为本体。\par

10. 他们的确也使用\OriginalTerm{位格}{hypostasis}这个词；但他们意图在\OriginalTerm{本质}{\GreekText{οὐσία}}与位格之间作出某种区别，我不知道是什么区别；因此，我们当中大多数用希腊文讨论这些事的人，惯于说\GreekText{μίαν} \GreekText{οὐσίαν}, \GreekText{τρεῖς} \GreekText{ὑποστάσεις}，即拉丁文\Quote{一个本质，三个本体}。\par

% structure: p0021 source=h2 target=subsection reason=relative-heading-rank; base=h2

\subsection{第九章 三位的称号并非恰当（就人的意义而言）}

但因我们已惯常将本质理解为与本体相同的东西；所以我们不敢说一个本质、三个本体，而说一个本质或本体和三个位格：正如许多以拉丁文论述这些事、并且具有权威的学者所说的，因为他们找不到任何其他更合适的方式，来言说他们用言语所理解的、超越言语的东西。因为事实上，正如圣父不是圣子，圣子不是圣父，而那被称为天主恩赐的圣神既不是圣父也不是圣子，他们确实是三位。所以复数形式说：\Quote{我与父原是一体。}祂并没有像撒伯流派所说的\Quote{\Emph{是}一体}，而是\Quote{\Emph{原是}一体}。然而，当问及\Quote{什么三位}时，人类语言极为贫乏。但答案仍是三位\Quote{位格}，并非为了能够完全说出，而是为了不致完全无言。\par

% structure: p0023 source=h2 target=subsection reason=relative-heading-rank; base=h2

\subsection{第十章 凡绝对属于天主本质的事物，均以单数而非复数论及天主圣三}

11. 因此，正如我们不说三个本质，同样我们也不说三个伟大或三个伟大的。因为，在那些因分有伟大而成为伟大的事物中，对于它们，\Emph{存在}是一回事，成为\Emph{伟大}是另一回事；譬如一所大房子、一座大山和一个伟大的心灵；在这些事物中，我说，伟大是一回事，而因伟大成为伟大的又是另一回事；一所大房子当然不是绝对的伟大本身。但正是那绝对的伟大，不仅使大房子成为伟大，使任何大山成为伟大，而且也使每一个被称作伟大的事物成为伟大；所以伟大本身是一回事，那些因它而被称为伟大的事物是另一回事。而这伟大本身当然是首要的，并且以一种远超过那些因分有它而成为伟大的事物的方式而伟大。但既然天主不是以那不属于祂自己的伟大而成为伟大（否则天主在成为伟大时，仿佛是那伟大的分有者——那样就会有一个比天主更伟大的伟大，而没有任何事物比天主更大），因此，祂是以那使祂自己成为同一伟大的伟大而伟大。所以，正如我们不说三个本质，同样我们也不说三个伟大；因为对天主而言，存在与成为伟大是同一回事。同样地，我们也不说三个伟大的，而只说一个伟大的；因为天主不是通过分有伟大而成为伟大，而是祂自己本身就是伟大，因为祂就是祂自己的伟大。同样地，也适用于天主的良善、永恒、全能，简言之，适用于所有可以述谓于天主的范畴，当按祂自身来说时——不是隐喻地或通过相似性，而是本真地（如果人能够以任何方式本真地谈论祂的话）。\par

% structure: p0025 source=h2 target=subsection reason=relative-heading-rank; base=h2

\subsection{第十一章——论天主圣三中的关系性称谓。}

12. 然而，在同一圣三中，有些事物被个别地特别述谓，这些绝不涉及它们自身的内在，而是要么相互关联，要么相对于受造物；因此，显然这些事物是关系性地被言说，而非以实体的方式。因为圣三被称为唯一的天主、伟大的、良善的、永恒的、全能的；并且同一天主自身可以被称为祂自己的神性、祂自己的伟大、祂自己的良善、祂自己的永恒、祂自己的全能：但圣三不能同样地被称为父，除非或许在比喻意义上，相对于受造物，由于收养为义子。因为经上记载：\BibleQuote{以色列，请听：上主我们的天主，是唯一的上主。}\BibleRef{申 6:4}这绝不应理解为圣子被排除，或圣神被排除；这唯一的主、我们的天主，我们也正当地称祂为我们的父，因祂藉祂的恩宠重生我们。同样，圣三无论如何也不能被称为子，但它可以整体地被称为圣神，按照经上所记载的：\BibleQuote{天主是神。}\BibleRef{若 4:24}因为父是神，子也是神，父是圣的，子也是圣的。因此，既然父、子和圣神是唯一的天主，而天主确实是圣的，且天主是神，圣三也可以被称为圣神。然而，那并非圣三本身、而是被理解为在圣三内的圣神，按其专名圣神被关系性地言说，因为祂既关联于父也关联于子，因为圣神是父的神和子的神。但这关系在名称上并不显明，而当祂被称为天主的恩赐时便显明了；因为祂是父和子的恩赐，正如主所说：\BibleQuote{祂由父而来。}\BibleRef{若 15:26}并因宗徒所说：\BibleQuote{谁若没有基督的圣神，谁就不属于基督。}\BibleRef{罗 8:9}他肯定说的是圣神本身。因此，当我们说恩赐者的恩赐和恩赐的恩赐者时，我们都在相互关联的关系中言说。因此，圣神是父和子某种不可言喻的共融；或许正因如此，祂被称为圣神，因为同一名称适用于父和子。因为祂被特别地称为他们共同被称呼的名字；因为父是神，子也是神，父是圣的，子也是圣的。因此，为了从双方都适合的名称中表明两者的共融，圣神被称为双方的恩赐。而这一圣三是唯一的天主、独一的、良善的、伟大的、永恒的、全能的；它自身就是自己的统一、神性、伟大、良善、永恒、全能。\par

% structure: p0027 source=h2 target=subsection reason=relative-heading-rank; base=h2

\subsection{第十二章——论相互关系中有时缺名。}

13. 既然我们已经说过圣神是这样关系性地被称为圣神（不是圣三本身，而是那在圣三内的那位），那么，祂所关联的那一位的名称，看来并不反过来与祂的名称相对应——这一点不应该影响我们。因为我们不能像我们说仆人的主人和主人的仆人、儿子的父亲和父亲的儿子那样，也在这里这样说——因为这些都是关系性的言说。因为我们论及父的圣神；但另一方面，我们不说圣神的父，以免圣神被理解为祂的儿子。同样，我们也说子的圣神；但我们不说圣神的子，以免圣神被理解为祂的父。因为许多关系中，找不到一个名称使那些相互关联的事物在名称上相互对应。还有什么比\Quote{凭据}一词更明显地是关系性的呢？因为它关联于它所作为凭据的事物，而凭据总是指某事物的凭据。那么，我们能否像说父和子的凭据那样，反过来又说凭据的父或凭据的子呢？但另一方面，当我们说父和子的恩赐时，我们确实不能反过来又说恩赐的父或恩赐的子；但为了使它们相互对应，我们说恩赐者的恩赐和恩赐的恩赐者；因为这里有一个常用的词，那里却没有。\par

% structure: p0029 source=h2 target=subsection reason=relative-heading-rank; base=h2

\subsection{第十三章 论圣三中\Quote{圣言\OriginalTerm{元始}{Principium}}的相对称谓}

14. 因此，圣父是相对地被称为圣父，祂也被相对地称为\OriginalTerm{元始}{Principium}，以及其他类似的名号；但祂是因圣子而被称为圣父，因万物皆来自祂而被称为元始。同样，圣子是相对地被称为圣子；祂也被相对地称为圣言和肖像。在所有这些称谓中，祂都是指向圣父，但圣父没有一个称谓是来自祂们。圣子也被称为元始；因为当有人对祂说：\BibleQuote{你是谁？}祂回答说：\BibleQuote{就是那元始，也向你们说话。}但是，请问，祂是圣父的元始吗？因为当祂说自己是元始时，祂是想表明自己是造物主，正如圣父也是受造物的元始，因为万物都是从祂而来。因为造物主相对于受造物，正如主人相对于仆人。因此，当我们说圣父是元始，圣子也是元始时，我们并不是在说受造物有两个元始；因为圣父和圣子合在一起，相对于受造物是一个元始，如同一个造物主，一个天主。但是，如果任何存有自身保持不变，并生产或制造某物，那么对该物而言它就是元始；那么我们就不能否认圣神也有理由被称为元始，因为我们没有将祂排除在造物主的称号之外。关于祂，经上记载祂工作；确实，祂在工作时自身保持不变，因为祂本身并没有改变或变成祂所工作的任何事物。请看祂所工作的是什么：他说：\BibleQuote{但是圣神的显现，}他说，\BibleQuote{是为了使人获益。这人由圣神获得智慧的言语，那人由同一圣神获得知识的言语，还有的人由同一圣神获得信德，另一人由同一圣神获得治病的奇恩，另一人获得行奇迹，另一人获得说预言，另一人获得辨别神恩，另一人获得各种异语，另一人获得解释异语：可是这一切，都是这同一圣神所行的，随祂的心愿，个别分配给人。}诚然，如同天主——因为除了天主，谁能行这样的大事呢？——但是\BibleQuote{那确实是同一的天主，在一切人身上行一切事。}因为，如果我们逐点被问及圣神，我们会最真实地回答祂是天主；而且祂与圣父和圣子合在一起是一个天主。因此，天主相对于受造物被称为一个元始，而不是两个或三个元始。\par

% structure: p0031 source=h2 target=subsection reason=relative-heading-rank; base=h2

\subsection{第十四章 圣父和圣子是圣神的唯一\OriginalTerm{元始}{Principium}}

15. 但在天主圣三彼此之间的相互关系内，如果生者是相对于他所生者的元始，那么圣父相对于圣子就是元始，因为祂生了圣子；但是圣父是否也是相对于圣神的元始，因为经上说：\BibleQuote{祂由圣父而来，}这是一个不小的问题。因为，如果是这样，那么祂不仅是祂所生或所造之物的元始，也是祂所赐予之物的元始。这里，也出现了一个常使许多人困惑的问题：为什么圣神也不是一个子，因为祂也由圣父而来，正如福音所读到的。因为圣神是作为恩赐而来，而非作为生者；所以祂不被称为子，因为祂既不是像独生子那样受生，也不是受造，好藉天主的恩宠而生于义子地位，如同我们一样。因为凡由圣父所生的，当被称为子时，只指向圣父，所以圣子是圣父的子，而不是我们的子；但凡是赐予的，既指向赐予者，也指向接受者；因此，圣神不仅是父和子（祂们赐予祂）的圣神，也被称为我们的圣神，因为我们接受了祂：正如\Quote{上主的救恩}，赐予救恩的上主，也被称为我们的救恩，因为我们接受了它。所以，圣神既是赐予祂的天主的圣神，也是接受了祂的我们的圣神。当然，这不是我们赖以存在的那精神，因为那是人内在的精神；但这圣神是以另一种方式属于我们，\Emph{即}如我们所说：\BibleQuote{求你今天赏给我们日用的食粮。}虽然我们也确实接受了那被称为人精神的精神。因为他说：\BibleQuote{你有什么不是领受的呢？}但一个是我们为存在而领受的，另一个是我们为成圣而领受的。因此，关于若翰，也记载他\BibleQuote{以厄里亚的精神和能力而来}；这里厄里亚的精神是指厄里亚所领受的圣神。同样也当理解梅瑟，当时主对他说：\BibleQuote{我要把你身上的精神，分给他们}；即我要将我已赐给你的圣神，也赐给他们。因此，如果那被赐予的也有赐予者作为它的元始，因为它从别处没有领受那从赐予者而来的，那么就应当承认圣父和圣子是圣神的元始，但不是两个元始；正如圣父和圣子是一个天主、一个造物主和一个相对于受造物的主一样，相对于圣神，祂们也是一个元始。然而，圣父、圣子和圣神相对于受造物是一个元始，如同一个造物主和一个天主。\par

% structure: p0033 source=h2 target=subsection reason=relative-heading-rank; base=h2

\subsection{第十五章 圣神在赐下前后是否都是恩赐}

但问题更进一步：圣子藉着受生，不仅具有祂是圣子这一点，而且具有祂\OriginalTerm{绝对地}{absolutely}存在；同样，圣神藉着被赐予，不仅具有祂被赐予这一点，而且具有祂\OriginalTerm{绝对地}{absolutely}存在。因此，祂在被赐予之前是否已经存在，但还不是一个\Emph{恩赐}呢？抑或，正因为天主将要赐予祂，所以祂在被赐予之前就已经是一个\Emph{恩赐}了呢？但如果圣神只当被赐予时才发出，而在有一位可领受祂的人之前确实不能发出，那么，祂又怎能\Emph{\OriginalTerm{绝对地}{absolutely}}存在于祂的实体中呢？因为祂若不是由于被赐予就不能存在。正如圣子藉着受生，不仅具有祂是圣子这一点（这是相对而言的），而且具有祂的绝对实体，以致祂存在。圣神是否总是在发出，不是在时间中发出，而是从永恒中发出，但因为祂如此发出以致祂能被赐予，所以即使在有一位可领受祂的人之前，祂就已经是一个\Emph{恩赐}了呢？因为\Quote{恩赐}与\Quote{已被赐予之物}在含义上有区别。因为一个\Quote{恩赐}在它被赐予之前就可能存在；但它只有在被赐予以后才能被称为\Quote{已被赐予之物}。\par

% structure: p0035 source=h2 target=subsection reason=relative-heading-rank; base=h2

\subsection{第十六章。—— 论及天主在时间中所说的话，是相对地，并非偶然地。}

17. 不要因圣神虽与圣父和圣子同永恒，却被称作某个存在于时间中的事物，就如我们所称的祂为那个已赐予的事物而感到困扰。因为圣神永恒地是恩赐，但某个已赐予的事物却存在于时间中。因为如果主也唯有在开始有奴仆时才如此称呼，那么这称谓对天主而言同样是相对的且存在于时间中；因为受造物并非从永恒就存在，而祂是它的主。那么，我们如何能证明相对用语本身并非偶然，既然没有任何事在时间中偶然地发生于天主身上，因为祂是不能改变的，正如我们在本次讨论的开端所论证的？看！为主，对天主而言并非永恒的；否则我们就被迫承认受造物也是永恒的，因为除非受造物也是永恒的奴仆，否则祂就不会是从永恒的主。因为正如没有主的奴仆不能成立，没有奴仆的主也不能成立。如果有人声称天主确实是唯一永恒的，而时间由于其多样性和可变性并非永恒，但时间并非在时间中开始存在（因为在时间开始之前没有时间，因此为主对天主而言并非在时间中发生，因为祂是那些时间本身的主，而时间确实不是在时间中开始的）：那么对于在时间中被造的人，在祂将成为其主的那个人存在之前，祂当然不是他的主，对此人你又该如何回答？为主，对人而言，确实在时间中发生于天主。而且，为了使所有争论似乎被消除，为主（作为你的主或我的主，我们仅在不久前才开始存在）确实在时间中发生于天主。或者，如果这一点由于有关灵魂的模糊问题似乎也不确定，那么对于祂是以色列民族的主该怎么说？因为，虽然灵魂的本性已经存在，那个民族拥有它（我们现在不探究此事），但那个民族尚未存在，而且它开始存在的时间是明显的。最后，祂是这棵或那棵树的主，或是这棵或那棵庄稼的主（它们仅在不久前才开始存在），这是在时间中发生的；因为，虽然质料本身已经存在，但成为\OriginalTerm{质料}{materiæ}的主与成为已经受造的\OriginalTerm{本性}{naturæ}的主是不同的。因为人有时是木材的主，有时又是箱子的主，虽然箱子是用同样的木材制成的；他当然在已经是木材的主时，还不是箱子的主。那么，我们如何能证明有关天主没有任何事物是按偶然性说的，除非因为没有任何事物通过祂可能会改变而发生于祂的本性，以至于那些相对偶然性是随着它们所谈论的事物的某种变化而发生的？正如朋友是相对地称呼的：因为他只有在开始爱时才开始成为朋友；因此意志发生某种变化，他才会被称为朋友。而金钱，当被称为价格时，是相对地说的，然而它在开始成为价格时并没有改变；同样，当它被称为抵押品或任何其他同类事物时也是如此。如果金钱可以在自身没有变化的情况下如此经常地被相对地谈论，以至于既不在它开始如此称呼时，也不在它停止如此称呼时，其作为金钱的本性或形式发生任何变化；那么，对于那不能改变的天主实体，我们更应该容易承认，某些事物可以相对于受造物而被相对地述说，虽然它开始在时间中被如此述说，但应理解为天主实体本身并没有发生任何事，而只是发生在相对于它被述说的受造物上。经上说：\BibleQuote{主，祢成了我们的避难所。}因此，天主被相对地称为我们的避难所，因为祂被关联于我们，而祂在我们投奔祂时成为我们的避难所；请问那时在祂的本性中发生了什么在投奔祂之前没有的事吗？在我们内确实发生了某种变化；因为我们在投奔祂之前更糟，而通过投奔祂我们变得更好；但在祂内没有变化。同样，当我们通过祂的恩宠重生时，祂开始成为我们的父，因为祂赐给我们能力成为天主的子女。因此，当我们成为祂的子女时，我们的实体变得更好；而祂同时开始成为我们的父，但祂自己的实体没有任何变化。因此，那开始在时间中关于天主被述说、从前没有关于祂被述说的事，显然是相对地关于祂被述说的；然而并非根据天主的任何偶然性，以致有什么事发生在祂身上，而显然是依据那个相对于天主开始被称为某物的受造物的某种偶然性。当一个义人开始成为天主的朋友时，他自己改变了；但千万别让我们说，天主在时间中以某种新爱情爱任何人，这在祂从前是没有的——在祂那里，过去的事并未消逝，未来的事已经完成。因此，祂在创造世界之前就爱了祂所有的圣人，正如祂所预定的；但当他们悔改并找到祂时，他们就被说成开始被祂所爱，这是用人类情感可以理解的方式来说的。同样，当祂被称为对不义者发怒、对善者温柔时，是他们改变了，不是祂；正如光线对弱眼是刺眼的，对强眼是舒适的——即因它们的改变，而非它自己的改变。\par

\SourceRecord{Translated by Arthur West Haddan.；Nicene and Post-Nicene Fathers, First Series；Vol. 3.；Edited by Philip Schaff.；Buffalo, NY: Christian Literature Publishing Co.,；1887.；<http://www.newadvent.org/fathers/130105.htm>.}

% structure: p0001 source=h1 target=section reason=page-title

\section{论天主圣三（第六卷）}

\EditorialNote{本卷提出一个问题：宗徒如何称基督是天主的德能和天主的智慧？并引发一个论辩：圣父是否并非智慧本身，而只是智慧之父？或者智慧生了智慧？但此答案稍后暂缓，同时证明圣父、圣子、圣神的合一与平等，以及我们应当相信一个天主圣三，而非三重（triplicem）神。最后，依拉利的话得到解释：永恒在圣父，显现在肖像，应用在恩赐。}\par

% structure: p0003 source=h2 target=subsection reason=relative-heading-rank; base=h2

\subsection{第一章——按宗徒所说，圣子是圣父的德能和智慧。由此，公教徒对早期亚略人的推理。一个困难被提出：圣父是否不是智慧本身，而只是智慧之父？}

1. 有些人认为，由于经上记载：\Quote{基督，天主的德能和天主的智慧}，他们因此被阻碍，不能承认圣父、圣子、圣神的平等；因为基于此，似乎没有平等，因为圣父本身不是德能和智慧，而是德能和智慧的生者。确实，人们通常以不寻常的热忱提出一个问题：天主怎能被称为德能和智慧之父？因为宗徒说：\Quote{基督，天主的德能和天主的智慧}\BibleRef{格前1:24}。因此，我们这一方有些人曾这样推理反对亚略人，至少反对那些最初起来反对公教信仰的人。因为据报道，亚略本人曾说：如果祂是子，那么祂就是受生的；如果祂是受生的，就有一个时候子是不存在的；他不明白，对天主而言，受生也是从永远开始的；所以子与父是永恒的，如同由火产生并散发的光辉与火同时存在，并且，如果火是永恒的，光辉也将是永恒的。因此后来的亚略人中有些人放弃了那种意见，并承认天主子并非在时间中开始存在。但在我们这一方用以反对那些说\Quote{有一个时候子不存在}之人的论证中，有些人惯于这样推理：如果天主子是天主的德能和智慧，而天主从未没有德能和智慧，那么子就与天主父是永恒的；但宗徒说：\Quote{基督，天主的德能和天主的智慧}；而一个无感觉的人不会说天主在任何时候没有德能或智慧；因此，没有一个时候子不存在。\par

2. 这个论证迫使我们说：天主父不是智慧的，除非拥有祂所生的智慧，而不是因父本身是智慧本身。再者，如果是这样，正如子本身也被称为天主的天主、光的光，我们必须考虑：祂是否能被称为智慧的智慧，如果天主父不是智慧本身，而只是智慧的生者。如果我们坚持这一点，那么为什么祂不是祂自己伟大、自己良善、自己永恒、自己全能的生产者？这样，祂就不是祂自己的伟大、祂自己的良善、祂自己的永恒、祂自己的全能；而是借着祂所生的伟大而伟大，借着祂所生的良善而良善，借着祂所生的永恒而永恒，借着祂所生的全能而全能；正如祂本身不是祂自己的智慧，而是借着生于祂的智慧而有智慧。因为我们不必害怕被迫说：除了受造的义子之外，还有许多与父永恒的天主子，如果祂是祂自己伟大、良善、永恒、全能的生者。因为很容易回答这个诡辩：命名许多事物，并不因此意味着祂是许多永恒子的父；正如祂不是两个子的父，因为基督被称为天主的德能和天主的智慧。因为那个德能就是智慧，那个智慧就是德能；同样，其余也是如此，以至于那个伟大就是德能，或者是上述或未来提到的那些事物中的任何一个。\par

% structure: p0006 source=h2 target=subsection reason=relative-heading-rank; base=h2

\subsection{第二章——论哪些关于圣父和圣子一起说，哪些不说。}

3. 但是，如果关于圣父本身，除了论及其与圣子之关系的那一切——即祂是圣子的父、或生者、或根源——之外，并无任何其他言说；并且生者对其所自生者自然也是根源；而无论论及圣父的其他任何事，都是作为与圣子\Emph{一起}或更确切地说在圣子\Emph{内}而言的——无论我们说祂是藉所生的伟大而伟大，或是藉所生的正义而正义，或是藉所生的良善而良善，或是藉所生的能力或德能而全能，或是藉所生的智慧而智慧：但圣父并不被称为伟大本身，而是伟大的生者；而圣子，正如祂被称为圣子本身，并非与圣父\Emph{一起}被称呼，而是在关系上\Emph{对}圣父，所以祂并非凭自身伟大，而是与圣父\Emph{一起}伟大，圣父即祂的伟大；同样，祂被称为与圣父\Emph{一起}有智慧，因为祂自身就是圣父的智慧；正如圣父被称为与圣子\Emph{一起}有智慧，因为祂藉所生的智慧而有智慧；因此，无论就自身而论被称为什么都不可脱离对方——即，凡是显明其本质的任何称呼，二者都是一同被称呼的——如果事情是这样，那么圣父离了圣子就不是天主，圣子离了圣父也不是天主，而二者一起才是天主。并且经上说：\BibleQuote{在起初已有圣言}，意思是指圣言在圣父内。或者若\BibleQuote{在起初}意指在万有之前；那么在接着的\BibleQuote{圣言与天主同在}中，圣言被理解为仅指圣子，而非圣父与圣子一起，仿佛二者同为一位圣言（因为祂是圣言正如祂是肖像，但圣父与圣子并非一同是肖像，只有圣子是圣父的肖像；正如祂也是圣父的圣子，因为二者一同并非圣子）。但在随后补充的\BibleQuote{圣言与天主同在}中，有很大的理由可作如下理解：\BibleQuote{圣言}——仅指圣子——\BibleQuote{与天主同在}，这\BibleQuote{天主}并非仅指圣父，而是圣父与圣子一同的天主。但若在某种截然不同的事物中也能这样表达，这又有什么奇怪呢？因为有什么比灵魂与身体更不同呢？然而我们能说灵魂与一个人同在，即在一个人的身体内；虽然灵魂不是身体，而人是灵魂与身体一同构成的。因此，接下来的圣经语句\BibleQuote{圣言就是天主}可这样理解：圣言——祂不是圣父——与圣父一同是天主。那么，我们是否可以说，圣父是祂自己伟大的生者，即祂自己德能的生者，或祂自己智慧的生者；而圣子是伟大、德能和智慧，但伟大、全能而智慧的天主则是二者一同？那么如何理解\Quote{天主\Emph{出自}天主，出自光明的光明}？因为并非二者一同是出自天主，只有圣子是出自天主，即出自圣父；也非二者一同是出自光明，只有圣子是出自光明，即出自圣父。除非，也许正是为了简要地暗示并强调圣子与圣父是永恒的，才说\Quote{出自天主的天主，出自光明的光明}，或任何同类说法：如同说，这个离了圣父就不是圣子的东西，\Emph{出自}那个离了圣子就不是圣父的东西；即这个离了圣父就不是光明的光明，\Emph{出自}那个光明——\Emph{即}圣父——离了圣子就不是光明；因此，当说\Quote{天主}（离了圣父就不是圣子）\Quote{\Emph{出自}天主}（离了圣子就不是圣父）时，可以完全理解：生者并不先于祂所生者。如果这样，那么关于他们，唯独这一句不能说，即\Quote{这个出自那个}，因为他们二者并非一同是此。正如不能说圣言\Emph{出自}圣言，因为二者一同并不是圣言，只有圣子是圣言；也不能说肖像\Emph{出自}肖像，因为他们一同并不是肖像；也不能说圣子\Emph{出自}圣子，因为二者一同并不是圣子，正如经上所说：\BibleQuote{我与父原是一体}。因为\BibleQuote{我们是一体}意思是：祂是什么，我也是什么；这是按本质，而非按关系。\par

% structure: p0008 source=h2 target=subsection reason=relative-heading-rank; base=h2

\subsection{第三章——从\BibleQuote{我们是一体}这句话领悟圣父与圣子本质的合一。圣子在智慧及其他一切方面与圣父同等。}

4. 而且我不知道，圣经是否曾对性质不同的事物说过\Quote{他们是一体}这样的话。但如果有不止一个事物具有相同的性质，而它们在意见上不同，那么它们就不是一体，并且正是在它们意见不同的程度上不是一体。因为如果门徒们只是由于是人而已经成为一体，那么主在将他们托付给父时就不会说\BibleQuote{使他们合而为一，如同我们合而为一}。但由于保禄和阿颇罗两人都是人，并且也有相同的意见，他说：\BibleQuote{栽种的和浇灌的是一体}。因此，当一个事物被称为一体，而没有附加在什么方面是一体，但不止一个事物被称为一体时，就意味着相同的本质和性质，没有区别或分歧。但当附加了在什么方面是一体时，则可能意味着从两个以上性质不同的事物中合成的一体。就如灵魂和身体肯定不是一体；因为它们是多么不同啊！除非附加或理解为它们在什么方面是一体，即一个人或一个动物个体。因此，宗徒说：\BibleQuote{那与娼妓结合的，成为一体}；他没有说他们是一体或他是一体，而是附加了\Quote{身体}，好像是由两个不同的身体（男性与女性）结合在一起组成的一个身体。而\BibleQuote{那与主结合的}，他说，成为一神。他没有说那与主结合的是\Quote{一}或\Quote{他们是一}，而是附加了\Quote{神}。因为人的精神与圣神性质不同；但通过结合，它们由两个不同的精神成为一神，使得圣神无需人的精神就已是有福而完美的，但人的精神没有天主就不能有福。而且，我认为并非没有原因，当主在\BibleBook{若望}{福音}中多次反复地谈到合一本身——无论是他自己与圣父的合一，还是我们彼此之间的合一时——他从未说过我们与他也是一体，而是说\BibleQuote{使他们合而为一，如同我们合而为一}。因此，圣父与圣子是一体，无疑是根据本质的合一；并且，如我们已论证的，只有一个天主，一个伟大的，一个智慧的。\par

5. 那么，圣父怎么更大呢？因为如果更大，祂是因伟大而更大；但既然圣子是祂的伟大，那么圣子肯定不大于那生了祂的那位，圣父也不大于祂借以伟大的那个伟大；因此他们是平等的。因为如果不是在祂之所是上——对祂而言，存在与伟大不是两回事——那么祂怎么会平等呢？或者，如果圣父在永恒上更大，那么圣子在任何方面都不平等。因为从何平等？如果你说在伟大上，那个不那么永恒的伟大就不平等，其他一切也如此。或者，也许在能力上平等，但在智慧上不平等？但那个智慧较小的能力，怎能平等？或者，在智慧上平等，但在能力上不平等？但那个能力较小的智慧，怎能平等？因此，若在任何一点上不平等，就在所有方面不平等。但圣经宣告：\BibleQuote{祂不以自己与天主同等为僭越}。因此，任何反对真理的人，只要他感到受权威约束，就必须承认圣子在与天主同等的每一点上都是同等的。让他选择他所愿意的一点；从这一点将向他显示，在一切关于祂本质的说法中，祂是平等的。\par

% structure: p0011 source=h2 target=subsection reason=relative-heading-rank; base=h2

\subsection{相同论证继续。}

6. 因为同样，存在于人心中的德性，虽然每一个都有其各自不同的含义，但在任何方面都是不可分离的；所以，例如，凡在勇气上相等的人，在明智、节制和正义上也相等。因为如果你说某些人在勇气上相等，但其中一人在明智上更大，那么另一人的勇气就不那么明智，因此他们在勇气上也不相等，因为前者的勇气更明智。如果你逐一考虑其他德性，你会发现同样的情况。因为这里的问题不在于身体的力量，而在于心灵的勇气。那么，在那种不变而永恒的本质中，它比人心更是无比单纯，这该是何等事实呢？因为，在人的心灵中，存在并不等同于强有力、明智、正义或节制；因为一个心灵可以存在而不具有这些德性中的任何一个。但在天主那里，存在就等于强有力、正义、智慧，或任何用来表示那种单纯的多重性或多重性的单纯性，以便表示祂的本质。因此，无论我们说天主出自天主，只是这名称属于每一位，但两者一起并不是两个天主，而是一个天主；因为它们彼此如此结合，以至于按照宗徒的见证，甚至在相异和不同的本质中也可能发生；因为主独自是一位神，而人的精神无疑也是一个精神；然而，如果它与主结合，\BibleQuote{就成为一神}：那么，在那里存在着绝对不可分离和永恒的结合，使得如果那被称为天主的只是对两者一起说的，那么当祂被称为天主之子时，祂被称作好似两者的儿子，这并非不合理。或者是这样的：凡是对天主说的表示祂本质的话，除非是对两者一起说的，甚至是对天主圣三本身一起说的？因此，无论是这样还是那样（这需要更仔细的探究），从以上所述已经足够看出，如果圣子在任何与表示祂本质有关的事情上被发现不平等，那么祂就绝不在任何方面与圣父平等，正如我们已经表明的。但宗徒说祂是平等的。因此，圣子在一切方面与圣父平等，并具有同一本质。\par

% structure: p0013 source=h2 target=subsection reason=relative-heading-rank; base=h2

\subsection{圣神也在一切方面与圣父和圣子平等。}

7. 因此，圣神也存在于这同一的\OriginalTerm{实体合一}{unity of substance}与同一的平等之中。因为无论祂是那两者的合一，或是圣德，或是爱，或是因此那合一由于爱，因此那爱由于圣德，显然祂不是两者之一；通过祂，两者被结合；通过祂，被生者被生者所爱，并爱那生祂者；并且通过祂，不是由于参与，而是由于他们自身的本质，也不是由于任何高者的恩赐，而是由于他们自己的，他们\BibleQuote{尽力以和平的联系，保持圣神的合一}；我们被命令藉着恩宠，向天主并向我们自己模仿这事。\BibleQuote{这两条诫命，是全部法律和先知所依据的}。所以那三位是天主，唯一、单独、伟大、智慧、圣善、有福的。但我们从祂、藉着祂、并在祂内成为有福的；因为我们自己藉祂的恩赐成为一，并与祂成为一灵，因为我们的灵魂紧贴祂而追随祂。我们紧贴天主是有益的，因为祂必毁灭每一个远离祂的人。因此，圣神，无论祂是什么，是圣父和圣子所共有的某种东西。但那共融本身是\OriginalTerm{同体}{consubstantial}且\OriginalTerm{同永恒}{co-eternal}的；如果它可合宜地被称为友谊，就如此称呼；但它更适宜被称为爱。而这也是一种实体，因为天主是实体，并且正如所记载的：\BibleQuote{天主是爱}。但正如祂与圣父和圣子同为实体，那实体也同他们伟大、同他们善、同他们圣、以及任何其他关于实体所说的话；因为对天主而言，存在是一回事，而伟大或善等等是另一回事，正如我们上面所展示的。因为如果爱在其中（即在天主内）不如智慧伟大，那么智慧被爱的程度就低于它自身所是；因此爱是平等的，以便智慧按其存在而被爱；但智慧与圣父平等，正如我们上面所证明的；因此圣神也平等；如果平等，就在一切上平等，因为那实体中的绝对单纯。所以它们不多于三位：一位爱那来自祂自身者，一位爱那祂所来自者，以及爱本身。如果这最后一位是虚无，那么\BibleQuote{天主是爱}如何成立？如果它不是实体，那么天主如何是实体？\par

% structure: p0015 source=h2 target=subsection reason=relative-heading-rank; base=h2

\subsection{第六章——天主如何既是单纯又是多重之实体。}

8. 但若有人问那实体如何既是单纯又是多重：首先要考虑，为何受造物是多重的，但在任何方面都不是真正单纯的。首先，一切物体肯定由部分构成；因此其中一部分大，另一部分小，而整体比任何部分都大，无论多大。因为天和地是整个世界体积的部分；而单独的地和单独的天由无数部分构成；它的三分之一小于其余部分，它的一半小于整体；而整个物质世界，通常由其两部分（即天和地）称呼，当然比单独的天或单独的地更大。在每个个别物体中，大小是一回事，颜色是另一回事，形状是另一回事；因为相同的颜色和相同的形状可以伴随缩小的大小而保留；相同的形状和相同的大小可以伴随颜色的改变而保留；即使形状不保留，物体仍可能同样大小并有相同颜色。任何其他一起归属于物体的性质，要么全部改变，要么大部分改变而不影响其余的。因此物体的性质被确凿证明是多重的，并且在任何方面都不是单纯的。精神受造物，即灵魂，若与身体相比，确实是两者中较单纯的；但若忽略与身体的比较，它也是多重的，而且本身也不是单纯的。它之所以比身体单纯，正是因为它不因空间的延伸而扩散于体积中，而在每个身体中，它既整个在整体内，又整个在其每个部分内；因此，当身体的任何微小部分发生任何事，如灵魂能感觉到，尽管并不发生在整个身体内，但整个灵魂仍感觉到它，因为整个灵魂并非不知觉它。但尽管如此，在灵魂中，有技巧是一回事，懒惰是另一回事，聪明是另一回事，记忆力强是另一回事；欲望是一回事，恐惧是另一回事，喜乐是另一回事，悲伤是另一回事；而且无数事物以无数方式存在于灵魂的本性中，有些没有其他，有些较多，有些较少；显然其本性不是单纯的，而是多重的。因为没有任何单纯的东西是可变的，而每个受造物都是可变的。\par

% structure: p0017 source=h2 target=subsection reason=relative-heading-rank; base=h2

\subsection{第七章——天主是圣三，但不是\OriginalTerm{三重}{triplex}。}

但天主确实以多种方式被称呼：伟大、善、智慧、有福、真实，以及其他任何似乎可无不相称地论及祂的事；但祂的伟大即是祂的智慧；因为祂不是因体积而伟大，而是因能力；祂的善即是祂的智慧与伟大，祂的真实即是所有那些；在祂内，有福与伟大、智慧、真实或善、或简言之成为祂自身，并不是不同的事。\par

9. 再者，既然祂是圣三，也不应因此被视为\OriginalTerm{三重}{triple (triplex)}，否则圣父单独，或圣子单独，就会小于圣父与圣子一起。尽管，确实，很难看出我们如何能说圣父单独或圣子单独；因为圣父常与圣子同在，圣子常与圣父同在，永不分离：不是因为他们两者都是圣父，或两者都是圣子；而是因为他们总是彼此相关地成为一，且两者都不单独。但因为我们称圣三本身为唯一的天主，尽管祂常与圣洁的灵体和灵魂同在，但我们说只有祂是天主，因为他们并不与祂同是天主；同样，我们称圣父为唯一的圣父，不是因为祂与圣子分离，而是因为他们不一起都是圣父。\par

% structure: p0020 source=h2 target=subsection reason=relative-heading-rank; base=h2

\subsection{第八章——不能在天主的本性上添加任何事物。}

圣父单独、或圣子单独、或圣神单独，都如圣父、圣子、圣神合起来一样伟大，绝不可称祂为三重。因为物体因自身结合而增大。虽然与妻子结合的，成为一体；但那是比丈夫单独或妻子单独更大的身体。但在属神的事上，当较小的依附较大的，如受造物依附造物主，是前者变得更大，而非后者。因为凡不以体积为伟大者，更伟大就是更善。任何受造物的精神，当它依附造物主时，比不依附时变得更好；因此也更伟大，因为更善。所以\Quote{祂}那与主联合的，便是同一心神：但主并不因此变得更大，虽然那与主联合的确实如此。在天主自身内，当同等的圣子，或与圣父圣子同等的圣神，与同等的圣父联合时，天主并不比其中任何一位更大；因为那圆满不能增加。但无论是圣父、或圣子、或圣神，祂是圆满的，而圣父、圣子、圣神的天主是圆满的；因此祂是\OriginalTerm{天主圣三}{Trinity}而非三重。\par

% structure: p0022 source=h2 target=subsection reason=relative-heading-rank; base=h2

\subsection{第九章 —— 是否一位或三位一起被称为唯一的天主。}

10. 既然我们正在展示我们如何能说圣父单独，因为在神性中除祂自己外没有别位圣父，我们也必须考虑那种意见，即认为唯一真天主不是圣父单独，而是圣父、圣子和圣神。因为如果有人问圣父是否单独是天主，怎能回答祂不是呢？除非我们或许要说圣父确实是天主，但祂不是单独的天主，而是圣父、圣子和圣神才是单独的天主？但那样我们如何处理主的那句证言呢？因为祂是对圣父说话，并指明了祂说话的对象是圣父，祂说：\Quote{永生就是：他们认识祢，唯一的真天主。}亚略派的人确实通常以此认为圣子不是真天主。然而，暂且不论他们，我们必须看看，当对圣父说\Quote{他们认识祢，唯一的真天主}时，我们是否被迫理解为祂想要暗示圣父单独是唯一的真天主；免得我们以为只有三者一起——圣父、圣子和圣神——才算是天主。因此，我们是否因主的证言，既要称圣父为唯一真天主，又要称圣子为唯一真天主，又要称圣神为唯一真天主，又要称圣父、圣子、圣神一起，即天主圣三本身一起，为不是三位真天主，而是一位真天主？或者因为祂加上\Quote{并认识祢所派遣的耶稣基督}，我们是否要补充\Quote{唯一的真天主}，使得语序成为：\Quote{他们认识祢，并认识祢所派遣的耶稣基督，唯一的真天主？}为什么祂省略了提到圣神？是否因为紧接着，每当我们称呼一位与另一位因如此大的和谐而结合，以至于通过这和谐两者成为一体时，这和谐本身必须被理解，即使它没有被提及？因为在同一处，宗徒似乎也略过了圣神；然而在那里，圣神也被理解，当他说：\Quote{一切都是你们的，你们是基督的，基督是天主的。}以及\Quote{女人的头是男人，男人的头是基督，基督的头是天主。}但再次，如果天主仅仅是三者一起，那么天主如何是基督的头，即天主圣三是基督的头？因为基督存在于天主圣三中，以便成为天主圣三？那与圣子同在的圣父，是单独圣子的头吗？因为圣父与圣子是天主，但圣子单独是基督：尤其是当那已成血肉的圣言说话时，依照祂的这种屈尊，圣父比祂大，正如祂说：\Quote{因为我的父比我大。}所以，那与圣父共有同一神性的天主本体，本身就是那中保之人的头，那人是祂单独所是的。因为如果我们正确地称心灵为人的首要部分，即好像是人类实体的头，虽然人本身连同心灵是人；那么圣言与圣父，两者同是天主，岂不更合适、更正确地成为基督的头吗？虽然基督作为人，除非连同那成了血肉的圣言，否则无法被理解。但这一点，正如我们已说过的，我们将在以后稍微更仔细地考察。目前，关于天主圣三的平等和同一实体，我们已经尽可能简要地证明，无论那另一个问题（其更严格的讨论我们已推迟）如何决定，没有什么能妨碍我们宣认圣父、圣子和圣神之间的绝对平等。\par

% structure: p0024 source=h2 target=subsection reason=relative-heading-rank; base=h2

\subsection{第十章 —— 论依拉略所归于每一位的属性。天主圣三在受造物中得为代表。}

11. 某位著述家，当他想要简略地指出天主圣三中每一位的特殊属性时，告诉我们说：\Quote{在永恒的圣父内，在肖像中有形式，在恩赐中有使用}。由于他在解释圣经和坚持信仰方面是一位权威不小的人——因为这是\Person{希拉利}在其\WorkTitle{论天主圣三}第二卷中所说的话——，我尽可能探究了这些词的隐微含义，即圣父、肖像和恩赐，永恒、形式和使用。我认为他用\Quote{永恒}一词的意思不过是：圣父没有父从他而生；但圣子是从圣父而来，以致祂存在并与他永恒。因为如果肖像完全充满它所肖似之物的尺度，那么肖像就与它所肖似之物相等，而不是后者与其肖像相等。关于这个肖像，他称之为\Quote{形式}，我相信是由于美的品质，在那里有如此大的适宜、原初的平等和原初的相似，毫无差别，无任何方面不等，无任何部分不相似，而是恰好与其肖像相符：那里有原初而绝对的生命，对祂而言，生活与存在不是两回事，而是同一回事存在与生活；有原初而绝对的理智，对祂而言，生活与理解不是两回事，而是理解即是生活，即是存在，一切都是一：好似一个完美的圣言，\BibleRef{若 1:1}，一无所缺，又是全能智慧之天主的某种技艺，充满一切活生生的、不变的知识，而这一切在祂内是同一的，正如祂本身是从一体而出的一体，与祂同一体。在那里，天主藉此知道祂所造的一切；因此，虽然时间过去又接续，但天主的认知中没有任何事物过去或接续。因为受造之物并非因被造而被天主认识，而是正因为祂不变地认识它们，它们才被造，即使它们是可变的。因此，圣父与其肖像那不可言喻的结合并非没有享用、没有爱、没有喜乐。因此，那爱、喜乐、幸福或福乐——如果真能用任何人的言语恰当表达的话——他简称为\Quote{使用}；那就是天主圣三中的圣神，非受生，乃是生者与受生者的甘饴，以其丰盛的慷慨和充盈充满一切受造物，好使它们保持其秩序，并在其适当的位置上安息。\par

12. 因此，凡由天主技艺所造的一切，都在自身中显示出某种合一、形式和秩序；因为每一物既是一个统一体——如身体的各种本性和灵魂的各种倾向；又具有某种形式——如身体的形状或品质，以及灵魂的各种学问或技能；还寻求或保持某种秩序——如身体的各种重量或组合，以及灵魂的各种爱或喜乐。因此，当我们藉着受造之物领悟造物主时，我们必然要理解天主圣三，因为在受造物中有着适切的痕迹。因为在那至高天主圣三中，存在着万物的至高根源、最完美的美和最幸福的喜乐。因此，那三者似乎彼此规定，而又自无限。但在有形之物中，单独一样不如三者加起来那样多，二者也比一者多；但在那至高天主圣三中，一者与三者在一起一样多，二者也不比一者多。它们自身是无限的。所以，每一个都在每一个内，一切都在每一个内，每一个都在一切内，一切都在一切内，一切都是一。人若看见这一点，无论是部分地，还是\Quote{藉着镜子，在谜语中}，就当因认识天主而喜乐；当尊祂为天主并感恩；若有人看不见，就当藉着虔敬努力去看，不要因盲目而吹毛求疵。因为天主是唯一的，然而却是天主圣三。我们也不可把\Quote{万有归于祂，藉着祂，并归于祂}这话当作笼统的（即表示无区别的一体），也不可当作表明多神，因为\Quote{归于\Emph{祂}，愿光荣至于无穷。阿们。}\par

\SourceRecord{Translated by Arthur West Haddan.；Nicene and Post-Nicene Fathers, First Series；Vol. 3.；Edited by Philip Schaff.；Buffalo, NY: Christian Literature Publishing Co.,；1887.；<http://www.newadvent.org/fathers/130106.htm>.}

% structure: p0001 source=h1 target=section reason=page-title

\section{论天主圣三（第七卷）}

\EditorialNote{此处解释此前书中延迟的问题，即：生圣子的天主父——圣子是天主的能力和智慧——不仅是能力和智慧之父，祂自己也是能力和智慧；圣神也同样如此：然而不是三个能力或三个智慧，而是一个能力和一个智慧，如只有一位天主和一个本质。然后探究为何拉丁人称一个本质、三个位格，而希腊人称一个本质、三个实体或位格；两种表达方式都源于言语的必要性，以便我们在被问及\Quote{什么三个}时能够回答，同时真实宣认有三个：父、子和圣神。}\par

% structure: p0003 source=h2 target=subsection reason=relative-heading-rank; base=h2

\subsection{第一章——奥古斯丁回到问题：天主圣三的每一位本身是否是智慧。所提问题的解决难度或方式。}

1. 现在，蒙天主恩许，我们更仔细地探讨此前稍加延迟的问题：即：天主圣三的每一位，能否单独、不与另外两位一起被称为天主、伟大、智慧、真理、全能、公义，或任何其他可按绝对意义（而非相对意义）论及天主的事；抑或这些事只能在理解圣三时才可说？因为问题被提出——由于经上记载：\Quote{基督是天主的能力和天主的智慧}——祂是否如此是祂自己智慧和能力的父，以至于祂以祂所生的智慧成为智慧，以祂所生的能力成为能力；并且，既然祂永远是能力和智慧，祂是否永远生了能力和智慧。如果是这样，那么正如我们所说，祂为何不是祂自己的伟大（祂因之而伟大）之父，不是祂自己的良善（祂因之而良善）之父，不是祂自己的公义（祂因之而公义）之父，以及其他任何存在的东西之父呢？或者，如果所有这些事——尽管有多个名称——都被理解为存在于同一智慧与能力中，以至于能力即是伟大，智慧即是良善，而能力又是智慧，正如我们已经论证的；那么，让我们记住，当我提及其中任何一个时，应被视为我提及了全部。于是询问：父本身是否是智慧，并且祂自己就是祂自己的智慧本身；还是祂的智慧如同祂的言语一样。因为祂是借着祂所生的圣言发言，不是借着被说出、发声并消逝的言语，而是借着与天主同在的圣言，且圣言就是天主，万物都是借着祂造成的：借着那与祂相等的圣言，祂永远不变地借着祂表达自己。因为祂自己不是圣言，正如祂不是圣子或肖像。但在发言时（暂且不论天主在时间中、在受造物中产生的那些言语，因为它们发声并消逝——那么在发言时），借着那同永恒的圣言，祂不被理解为单独存在，而是与圣言本身同在，没有祂父当然不能发言。那么，祂是否如同祂是一位发言者那样成为智慧，以至于祂就是智慧，正如祂是圣言，且成为圣言就是成为智慧，也就是能力，好使能力、智慧和圣言是同一的，并相对地被称为圣子和肖像：而父不是单独有能力或智慧，而是与祂所生的能力与智慧本身在一起（\OriginalTerm{生}{\Emph{genuit}}）；正如祂不是单独的一位发言者，而是借着祂所生的圣言并与祂所生的圣言在一起；同样，伟大也是借着祂所生的伟大并与祂所生的伟大在一起？如果祂不是借一件事而伟大、借另一件事而成为天主，而是借祂之为天主的那同一事而伟大，因为对祂而言，伟大与为天主不是两回事；那么结论是：祂也不是单独为天主，而是借着祂所生的神性并与祂所生的神性在一起（\OriginalTerm{神性}{\Emph{deitas}}）；好使圣子是父的神性，正如祂是父的智慧和能力，以及是父的圣言和肖像一样。并且，因为对祂而言，存在与为天主不是两回事，圣子也是父的本质，如同祂是父的圣言和肖像。因此，除了祂是父[无生者]之外，父若不是因为有了圣子，就什么也不是；因此，不仅\Quote{父}所指的（显然祂不是相对于自己而被如此称呼，而是相对于圣子，因此祂称为父是因为祂有了圣子），而且祂就其自身实体而言的所是，之所以如此称呼，是因为祂生了祂自己的本质。因为正如祂只与祂所生的伟大而成为伟大，同样祂也只与祂所生的本质而存在；因为对祂而言，存在与伟大不是两回事。因此，祂是否是祂自己本质的父，如同祂是祂自己伟大的父，是祂自己能力和智慧的父？既然祂的伟大等同于祂的能力，祂的本质等同于祂的伟大。\par

2. 这一讨论源于所记载的：\BibleQuote{基督是天主的德能，天主的智慧}。因此我们的言论被压缩在如此狭窄的限度内，而我们却渴望言说那不可言说之事；以致我们要么必须说基督不是天主的德能和天主的智慧，从而无耻且不虔地抗拒宗徒；要么必须承认基督确实是天主的德能和天主的智慧，但圣父并非祂自己德能与智慧的圣父，这同样不虔；因为如此，祂也不会是基督的圣父，因为基督是天主的德能和天主的智慧；要么认为圣父不是凭自己的德能而有能力，或凭自己的智慧而有智慧：谁敢这样说？或者再次，我们必须理解，在圣父内，存在是一回事，智慧地存在是另一回事，以致祂不是凭祂所是的那一位而智慧：这是通常对灵魂的理解，灵魂有时不智，有时智慧；因其本性可变化，并非绝对且完全纯一。或者再次，圣父在自己的实体方面不是什么；不仅祂是圣父，就连祂\Emph{是}（存在），也是相对于圣子而言的。那么，圣子如何能与圣父同一本质？因为圣父相对于自身，既不是自己的实体，也不相对于自身而存在，甚至祂的本质也是相对于圣子的？相反，圣子更是在同一本质内，因为圣父和圣子是同一本质；因为圣父的存有本身不是相对于自身，而是相对于圣子，祂生了那本质，且凭那本质祂\Emph{是}祂所是的一切。因此，[位格]都不是单独相对于自身而\Emph{存在}；两者都相对地彼此存在。或者，是否只有圣父不是相对于自身被称为圣父，而是无论祂被称为什么，都是相对于圣子而被称呼，而同一圣子却可以相对于自身被述说？如果是这样，那相对于自身被述说的，是什么呢？是祂的本质本身吗？但圣子是圣父的本质，正如祂是圣父的德能和智慧，祂是圣父的圣言和肖像。或者，如果圣子相对于自身被称为本质，但圣父不是本质，而是本质的生育者，并且不是相对于自身而\Emph{存在}，而是凭祂所生的那本质而\Emph{存在}；正如祂凭祂所生的那伟大而伟大：因此圣子也同样相对于自身被称为伟大；因此祂也同样被称为德能、智慧、圣言和肖像。但还有什么比祂相对于自身被称为肖像更荒谬的呢？或者，如果肖像和圣言与德能和智慧并非完全相同，而是前者是相对地被述说，后者是相对于自身而非相对于他者被述说；那么我们就得出这样的结论：圣父并不凭祂所生的智慧而智慧，因为祂本身不能相对于智慧被述说，智慧也不能相对于祂被述说。因为凡相对述说的事，都是互相述说的；因此，即使在本质方面，圣子也是相对于圣父而被述说。但由此得出一个最意想不到的意义：本质本身不是本质，或者至少，当它被称为本质时，所指的并非本质，而是某种相对的事物。正如当我们说到主人时，所指的并非本质，而是相对于奴隶的关系；但当我们说到人，或任何如此相对于自身而非他者被述说的事物时，所指的便是本质。因此，当一个人被称为主人时，人本身是本质，但他被称为主人是相对的；因为他相对于自身被称为人，但相对于他的奴隶被称为主人。但就我们开始的问题而言，如果本质本身是相对地被述说，那么本质本身就不是本质。此外，所有相对地被述说的本质，即使除去关系，也仍然是某物；\Emph{例如}，作为主人的人、作为奴隶的人、作为负重的马、作为抵押的钱，那人、马、钱是相对于自身被述说的，它们是实体或本质；但主人、奴隶、负重物、抵押品是相对于某物被述说的。但如果没有一个人，即某个实体，就无所谓相对地被称为主人的人；如果没有具有某一本质的马，就无所谓相对地被称为负重物；同样，如果钱不是某种实体，就不能相对地被称为抵押品。因此，如果圣父也不是相对于自身是什么，那么就根本没有可以相对地述说于某物的东西了。因为这与颜色不同。一物的颜色是相对于该有色之物而言的，颜色根本不被述说为实体，而总是关于某有色之物；但那有色之物，即使就其为有色而论是相对于颜色的，但就其为一身体而论，则是相对于实体而被述说。但无论如何我们不能认为，圣父同样不能相对于祂自己的实体被称为什么，而是无论祂被称为什么，都是相对于圣子而被称呼；而同一圣子，当祂被称为伟大的伟大、德能的德能时，则既相对于祂自己的实体，又相对于圣父，显然是相对于自身，并且是伟大而德能的圣父的伟大和德能，圣父凭此而伟大而德能。事实并非如此；两者都是实体，且两者是同一实体。正如说白不是白是荒谬的，同样说智慧不智也是荒谬的；正如白相对于自身被称为白，同样智慧相对于自身被称为智。但身体的白并非本质，因为身体本身是本质，而白是它的属性；因此，身体也因那属性而被称为白，而身体的存在与白并非同一回事。因为身上的形式是一回事，颜色是另一回事；两者并不在自身内，而是在某一体积内，这体积既非形式也非颜色，而是有形式有颜色的。真正的智慧既是智慧的，又是其自身内智慧的。由于每一个灵魂因分有智慧而成为智慧的，若它再变为愚拙，那智慧本身却保持不变；当那灵魂变为愚拙时，智慧并不随之改变；因此，智慧并不在因之成为智慧的人内，如同白在因之变白的身体内一样。因为当身体变为另一种颜色时，那白色将不再存留，而是完全消失。但如果生育智慧的圣父也因智慧而成为智慧的，并且存在对于祂并不等同于智慧地存在，那么圣子就是祂的属性，而非祂的后裔；天主性中便不再有绝对的纯一性。但绝不可如此，因为事实上在天主性中本质是绝对纯一的，因此存在等同于智慧地存在。但如果存在等同于智慧地存在，那么圣父并不是凭祂所生的智慧而智慧；否则祂并未生育它，而是它生育了祂。因为当我们说存在等同于智慧地存在时，我们还能说什么，除非祂是凭祂所是的那一位而智慧？因此，那为祂成为智慧的原因，也是祂存在的原因；因此，如果祂所生的智慧是祂智慧的原因，那么它也是祂存在的原因；这不可能发生，除非是通过生育或创造祂。但从未有人说过智慧在某种意义上生育了圣父或创造了圣父；因为还有什么比这更无意义的呢？因此，圣父本身是智慧，而圣子之所以被称为圣父的智慧，正如祂被称为圣父的光一样；即，正如光来自光，却仍是同一位光，同样，我们应理解智慧来自智慧，却仍是同一位智慧；因此也是同一本质，因为在天主内，存在等同于智慧存在。对智慧而言，智慧地存在，对德能而言，能够地存在，对永恒而言，永恒地存在，对正义而言，正义地存在，对伟大而言，伟大地存在，那存在本身对本质而言也是如此。由于在天主纯一性中，智慧地存在无非就是存在，因此智慧与本质是同一的。\par

% structure: p0006 source=h2 target=subsection reason=relative-heading-rank; base=h2

\subsection{第二章　圣父与圣子一同是一个智慧，如同是一个本体，虽然不是同一个圣言。}

3. 因此圣父与圣子一同是一个本体、一个伟大、一个真理、一个智慧。但圣父与圣子一同并不是同一个圣言，因为二者一同并不是同一个圣子。因为正如圣子是相对于圣父而言，并非就其本身如此称呼；同样，圣言也是相对于它的那位而言，被称为圣言。既然他是圣子，因为他也是圣言，他是圣言，因为他也是圣子。因此，既然圣父与圣子一同当然不是同一个圣子，那么圣父与圣子一同也就不是二者的同一个圣言。所以，他并非因他是智慧而是圣言；因为他被称为圣言并非就其本身，而是相对地相对于他的那位，正如他被称为圣子是相对于圣父；但他因他是本质而为智慧。因此，因为是一个本体，所以是一个智慧。但由于圣言也是智慧，然而他并不是因他是智慧而是圣言，因为圣言被理解为是相对而言，而智慧则是本质上而言：让我们明白，当他被称为圣言时，是指 \Emph{生出的} 智慧，以至于他也是圣子和肖像；而当这两个词，即 \Emph{智慧}（\Emph{是}）\Emph{生出的} 被使用时，在两个词之一中，即 \Emph{生出的}，圣言、肖像和圣子都被包含在内，而在所有这些名称中，本质并不被表达，因为它们是相对而言的；但在另一个词中，即 \Emph{智慧}，由于它也是按实体而言（因为智慧本身是明智的），本质也被表达出来，以及他那成为明智的存有。因此，圣父与圣子一同是一个智慧，因为是一个本体，且各自是智慧的智慧，正如本质的本质。因此，他们并不是因为圣父不是圣子、圣子不是圣父，或因为圣父是非生的而圣子是生出的，就不是一个本体；因为在这些名称中只表达了它们的 \Emph{相对} 属性。但二者一同是一个智慧和同一个本体；在其中存在与是明智的是同一回事。而二者一同并不是圣言或圣子，因为存在与是圣言或是圣子不是同一回事，正如我们已经充分表明这些术语是相对而言的。\par

% structure: p0008 source=h2 target=subsection reason=relative-heading-rank; base=h2

\subsection{第三章　为什么圣经中主要用智慧的名称来暗示圣子，虽然圣父和圣神都是智慧。圣神与圣父和圣子一同是一个智慧。}

4. 那么，为什么圣经中几乎从未谈及智慧，除非是为了表明它是天主所生或受造的？——就那藉以创造万物的智慧而言，是生的；但就人而言，当他们归向那并非受造而生、而是生的智慧，并如此受光照时，是受造或被造的；因为这些人们自身中也会产生某种可称为他们的智慧的东西：正如圣经预言或叙述：\BibleQuote{圣言成了血肉，居住在我们中间}；因为这样基督成为智慧，因为他成为了人。难道因此智慧在这些书中并不发言，也没有谈及它，除非是为了宣告它是天主所生或由祂所造（虽然圣父本身就是智慧），即因为智慧应当被我们赞美和效法，藉着效法它我们被[正确地]塑造？因为圣父说出它，使它成为祂的圣言：但并非像一个从口中发出声音的言语，或是在说出之前被思想的言语。因为这种言语在特定的时间间隔中完成，但那圣言是永恒的，并通过光照我们向我们说话，讲出关于它自己和关于圣父所应当向人讲的话。因此祂说：\BibleQuote{除了父以外，没有人认识子；除了子和子所愿意启示的人以外，也没有人认识父}：因为圣父藉着子，即藉着祂的圣言启示。因为如果我们所说的言语，是暂时的、易逝的，尚且宣告它自己和它所谈论的对象，那么天主的话，藉以创造万物的圣言，更是如此？因为圣言如此宣告圣父，正如祂是圣父；因为圣言本身也是如此，并且就它是智慧和本质而言，它就是圣父所是的那一位。因为就它是圣言而言，它不是圣父所是的；因为圣言不是圣父，圣言是相对而言的，就像圣子，肯定不是圣父。因此基督是天主的德能和智慧，因为祂自己，也是德能和智慧，是从那本身就是德能和智慧的圣父而来；正如祂是来自那光的光，是来自那生命之源的与天主圣父同等的生命之源。因为祂说：\BibleQuote{在祢那里有生命的源泉，在祢的光明中我们才能看见光明。}因为\BibleQuote{父怎样在自己内有生命，也照样赐给子在自己内有生命}；并且\BibleQuote{祂是真光，照耀每个来到世界的人}；而这光，\BibleQuote{圣言曾与天主同在}；但\BibleQuote{圣言也是天主}；而\BibleQuote{天主是光，在祂内毫无黑暗}：但不是肉体的光，而是属神的光；然而并非如此属神，以致它是因光照而成的，就如对宗徒们说的：\BibleQuote{你们是世界的光}，而是\BibleQuote{那照耀每个人的光}，即那至高的智慧本身，祂就是天主，我们现在所谈论的。因此子来自父是智慧的智慧，正如祂是光的光，天主的天主；所以圣父单独是光，圣子单独是光；圣父单独是天主，圣子单独是天主；因此圣父单独是智慧，圣子单独是智慧。正如二者一同是一个光和一个天主，同样二者是一个智慧。但子成了\BibleQuote{由天主使我们成为智慧、正义和圣化者}，因为我们在时间中归向祂，即从某个特定时间开始，以便永远与祂同在。而祂自己从某个时间开始\BibleQuote{成了血肉的圣言，居住在我们中间}。\par

5. 为此，当圣经中宣告或叙述任何关于智慧的事时，无论是智慧本身说话，或是论及它，主要向我们所指明的就是圣子。并且藉着那位是肖像的榜样，让我们也不离开天主，因为我们也是天主的肖像：不过不是那与祂平等的，因为我们是由父藉着子被造的，而不是从父所生，如同那一位那样。我们是如此，因为我们被光所光照；但那一位是如此，因为祂是光照的光；因此，祂没有模式，却成为我们的模式。因为祂不模仿任何先于祂的，就父而言，祂与父从未分离，因为祂与父是同一实体。但我们却努力模仿那位常存的，追随那位屹立的，并在祂内行走，向祂伸展；因为祂藉着贬抑自己，在时间内为我们成为道路，这道路藉着祂的天主性成为我们永恒的居所。因为为了那些没有因骄傲而堕落的纯洁理智之灵，祂以天主的形式、与天主同等、并作为天主给他们一个榜样；同样，为了给堕落了的人一个回归的榜样，人因罪的污秽和死亡的惩罚不能看见天主，祂\BibleQuote{空虚了自己}，并非改变自己的天主性，而是取了我们的可变化性；并\BibleQuote{取了奴仆的形体}，来到我们这个世界，祂\BibleQuote{本来就在这世界}，因为\BibleQuote{世界是藉着祂造的}；这样，祂成为向上看的人的榜样，向下看的人的榜样，健康者坚持的榜样，病者痊愈的榜样，将死者不惧怕的榜样，死者复活的榜样，\BibleQuote{使祂在一切事上居首位}。因此，因为人不应当追随除天主以外的任何东西以获得幸福，却又不能感知天主；那么通过追随成为人的天主，他就可以同时追随他所能感知的、并应当追随的那一位。因此，让我们爱祂并依附于祂，藉着倾注在我们心中的爱德，藉着赐给我们的圣神。所以，如果因着那位与父平等的肖像所给我们的榜样，为使我们得以重新塑造为天主的肖像，圣经在论及智慧时谈及圣子，我们藉着智慧生活而追随祂，这并不奇怪；虽然父也是智慧，正如祂是光、是天主一样。\par

6. 圣神，无论我们是否应称祂为那绝对的爱，这爱结合圣父与圣子，也从下面结合我们，这样，所写的\BibleQuote{天主是爱}并非不合适；但祂自己怎会不是智慧呢？既然祂是光，因为\BibleQuote{天主是光}？或者，若以任何其他方式单独且适当地称呼圣神的本质，那么，既然祂是天主，祂确实就是光；既然祂是光，祂确实就是智慧。但是圣神是天主，圣经藉着宗徒声明，他说：\BibleQuote{你们不知道你们是天主的宫殿吗？}并立即补充说：\BibleQuote{天主的圣神住在你们内}；因为天主住在自己的宫殿内。因为天主的圣神不是以仆人的身份住在天主的宫殿内，因为他在另一处更清楚地说：\BibleQuote{你们不知道你们的身体是住在你们内的圣神的宫殿吗？这圣神是天主赐给你们的，你们已不属于自己？因为你们是用重价买来的：所以要在你们的身体上光荣天主。}但是什么是智慧，除非是精神且不变的光？因为那边的太阳也是光，但它是物质的；精神受造物也是光，但它不是不变的。因此父是光，子是光，圣神是光；但一起不是三个光，而是一个光。同样，父是智慧，子是智慧，圣神是智慧，但一起不是三个智慧，而是一个智慧；因为在天主圣三内，存在与是有智慧是同一回事，父、子、圣神是一个实体。在天主圣三内，存在与是天主不是两回事；因此父、子、圣神是一个天主。\par

% structure: p0012 source=h2 target=subsection reason=relative-heading-rank; base=h2

\subsection{第四章　希腊人称\OriginalTerm{三个位格}{\GreekText{τρεῖς} \GreekText{ὑποστάσεις}}、拉丁人称\OriginalTerm{三位}{tres personae}的原因何在。圣经从未说在一位天主内有三位。}

7. 那么，为了谈论那不可言说之事，使我们能以某种方式说出我们根本无法完全说出的事，我们的希腊先贤讲论一个本质、三个实体；而拉丁先贤则讲论一个本质或实体、三个位格；因为正如我们已经说过的，在我们的语言，即拉丁语中，本质通常无非就是实体。只要所说的仅是在奥秘中被理解，这样的说法就足够了，为的是当人问起这三个是什么时，能有话可说；真正的信德宣称这三个是三个，因为它既宣告父不是子，又宣告圣神——天主的恩赐——既不是父也不是子。那么，当人问这三个是什么，或这三个是谁时，我们就求助于找出某个专名或通名，用以涵盖这三个；而心中却想不出这样的名称，因为神性的卓越超越了惯常言语的能力。因为天主更真实地被思想而非被言说，并且更真实地存在而非被思想。因为当我们说雅各伯不同于亚巴郎，依撒格既不是亚巴郎也不是雅各伯时，我们确证他们是三个：亚巴郎、依撒格和雅各伯。但若问这三个是什么，我们便答\Quote{三个人}，用一个专名来统称它们；但若我们说\Quote{三个动物}，则是用一个通名；因为人，正如古人所定义的，是理性的、有死的动物；或者，按照我们的\WorkTitle{圣经}通常的说法，\Quote{三个灵魂}，因为以较优的部分即灵魂来称呼整体，即灵魂和身体构成的整个人，是合适的；例如，\WorkTitle{圣经}说随雅各伯下到埃及的有七十五个灵魂，而不是说这么多人。同样，当我们说你的马不是我的，而属于另一个人的第三匹马既不是我的也不是你的时，我们承认有三匹马；若有人问这三个是什么，我们便用专名回答\Quote{三匹马}，或用通名\Quote{三个动物}。再说，当我们说牛不是马，而狗既不是牛也不是马时，我们说到一个三；若有人问我们这三个是什么，我们不再用专名说三匹马、三头牛或三条狗，因为这三者不属于同一个种，而是用通名\Quote{三个动物}；或者，如果取更高的属，则说\Quote{三个实体}、\Quote{三个受造物}或\Quote{三个本性}。但凡是能以专名以复数形式表达的事物，也能以通名表达。但凡是通名称呼的事物，不一定都能以专名称呼。因为三匹马，这是一个专名，我们也可称为三个动物；但一匹马、一头牛和一条狗，我们只能称为三个动物或三个实体（这是通名），或其他可以通称的说法；但我们不能称它们为三匹马、三头牛或三条狗（这些是专名）；因为只有那些在名称所意指的方面有共同性的事物，才能用一个复数名称来表达。因为亚巴郎、依撒格和雅各伯，在\Quote{人}方面有共同性，因此他们被称为三个人；一匹马、一头牛和一条狗，在\Quote{动物}方面有共同性，因此它们被称为三个动物。同样，三棵不同的月桂树，我们称为三棵树；但一棵月桂、一棵桃金娘和一棵橄榄，我们只能称为三棵树，或三个实体，或三个本性；同样，三块石头我们称为三个物体；但石头、木头和铁，我们只能称为三个物体，或用其他更高的通名来称呼。因此，关于父、子、圣神，既然它们是三个，让我们问它们是什么，它们有什么共同性。因为成为父这一点不是它们所共有的，使它们能相互称为父；正如朋友，因彼此相对的关系而被称为三个朋友，因为他们是互相对待的。但三位一体并非如此，因为那里只有父是父，而且不是两个的父，只是子的父。它们也不是三个子，因为父不是子，圣神也不是子。也不是三个圣神，因为圣神在其特有的意义上（也称作天主的恩赐）既不是父也不是子。那么，它们是三个什么呢？如果说三个位格，那么位格所意指的就是它们所共有的；因此这个名称对它们而言，按照说话方式，要么是专名，要么是通名。但在本性没有差异的地方，那些数目上为多的事物，既然能用通名表达，也就能用专名表达。因为本性的差异导致月桂、桃金娘和橄榄，或马、牛和狗，不能分别用前面的专名（三棵月桂或三头牛）来称呼，而只能用通名（前面的用三棵树，后面的用三个动物）。但在这里，既然本质没有差异，就必须有一个专名来称呼这三个，而这个专名却找不到。因为\Quote{位格}是一个通名，甚至连人也可以这样称呼，尽管人与天主之间有如此大的差异。\par

8. 此外，关于那个非常\OriginalTerm{普遍的}{generalis}词语，若因\Quote{位格}一词所指的为三者所共有，我们就称他们为三位格（否则他们绝不可如此称呼，正如他们不被称作三子，因为\Quote{子}所指的非三者所共有）；那么，我们为何不称他们为三天主？因为既然圣父是位格，圣子是位格，圣神是位格，故有三者位格；同样，既然圣父是天主，圣子是天主，圣神是天主，为何不称三天主？或者，既然因他们不可言喻的合一，这三者共为一天主，为何不也称作一位格？这样，我们就不能说三位格，尽管我们单独称每位为位格；正如我们不能说三天主，尽管我们单独称每位为天主，无论是圣父、圣子或圣神。这是否因为圣经并未说三天主？但我们同样找不到圣经任何一处提及三位格。还是因为圣经既未称这三者为三位格，也未称他们为一位格（因为我们在圣经中读到\Quote{主}的位格，却未将主视为一位格），所以，只凭说话与推理的需要，便合法地说三位格——并非因圣经这样说，而是因圣经不反对这样说；而倘若我们说三天主，圣经就会反对，因圣经说：\Quote{以色列，你要听：主，你的天主，是唯一的天主？}那么，为何也不合法地说三个本质？圣经同样未曾说，却也不反对？因为若本质是三者共有的\OriginalTerm{特指的}{specialis}名称，为何不称他们为三个本质，正如\Person{亚巴郎}、\Person{依撒格}和\Person{雅各伯}被称为三人，因为人是所有人共有的特指的名称？但若本质不是一个特指的名称，而是一个\OriginalTerm{普遍的}{generic}名称，因为人、牲畜、树、星宿、天使等都被称为本质，那么为何不称这些为三个本质，正如三匹马被称为三动物，三棵月桂树被称为三树木，三块石头被称为三物体？或者，若因天主圣三的合一，他们不被称为三个本质，而是一个本质，那么为何不同样因天主圣三的同一合一，他们也不被称为三个实体或三个位格，而是一个实体和一个位格？因为正如本质之名是共有的，各自单独称为本质，同样，实体之位或位格之名也是共有的。因为按照我们的用法，对于位格的理解，在希腊用法中即是对实体的理解；因为他们说三个实体、一个本质，正如我们说三个位格、一个本质或实体。\par

9. 那么，除了承认这些说法源于说话之必要——当丰富的论证需要对抗异端者的诡计或谬误时——还剩下什么呢？因为人的软弱，在尽其领悟限度内，借着虔诚的信德或任何辨知力，对于主天主其造主，于心灵隐秘之处所把握的，妄图通过言语向人的感官表达时，它不敢说三个本质，恐怕在那绝对平等中被人理解为存在任何差异。同样，它也不能说没有\OriginalTerm{三个某物}{tria quædam}，因为\Person{撒伯里乌}这样说而陷入异端。因为必须虔诚地相信——从圣经中确知——并以无疑的悟性用心灵之眼把握：既有圣父、圣子、圣神，且圣子与圣父不同，圣神亦与圣父或圣子不同。于是它寻求称他们为哪三个，最后回答是实体或位格；借这些名称，它并非意指差异，而是否认单一性：好叫人们不仅从被称为一个本质而理解其中的合一，也从被称为三个实体或位格而理解其中的天主圣三。因为若在天主内，\OriginalTerm{存在}{esse}与\OriginalTerm{自立}{subsistere}是同一回事，他们就不该被称为三个实体，正如他们不被称为三个本质一样；因为如同在天主内，存在与有智慧是同一事，所以正如我们不称三个本质，也不称三个智慧。因为若在天主内，是天主与存在是同一事，那么称三个本质便不正确，如同称三天主不正确一样。但若在天主内，存在是一回事，自立是另一回事，如同在天主内，存在是一回事，而做父亲或做主人是另一回事（因为祂之所是，是就祂自身而言；但祂被称为父，是因与圣子的关系，被称为主，是因与侍奉祂的受造物的关系）；因此，祂自立是相对地，如同祂生子是相对地，祂统治也是相对地：这样，实体就不再是实体，因为它将是相对的。因为从存在（esse）祂被称为本质（essentia），从自立（subsistere）我们称为实体（substantia）。但将实体称为相对的乃是荒谬的，因为每物都是就其自身而自立；何况天主呢？\par

% structure: p0016 source=h2 target=subsection reason=relative-heading-rank; base=h2

\subsection{第五章——在天主内，实体一词用得不当，本质一词才恰当。}

10. 然而，若说天主应当说是\OriginalTerm{自立}{subsist}——（这个词正确地适用于那些作为\OriginalTerm{主体}{subject}的东西，在其中，那些被说成是在一个主体中的东西，如颜色或形状在身体中。因为身体自立，所以是\OriginalTerm{实体}{substance}；但那些东西在身体中，身体自立并是它们的主体，而它们不是实体，而是在一个实体中；因此，若那颜色或形状消失，并不剥夺身体之为身体，因为身体的存在并不在于它拥有这样或那样的形状或颜色；所以，可变的和\OriginalTerm{单纯}{simple}的事物都不被正确地称为实体。）——如果，我说，天主如此自立，以致祂可以被正确地称为实体，那么祂内就有某物如同在一个主体中，而祂就不是单纯，\Emph{即}：祂的存在等同于任何其它关于祂自身所说的东西；如伟大的、全能的、良善的，以及这类不无合适地用于天主的一切。但是，说天主自立，并相对于祂自己的良善而成为主体，而这良善不是实体或更确切地说\OriginalTerm{本质}{essence}，并且天主自己不是祂自己的良善，而是它在祂内如同在一个主体中，这是不虔敬的。因此，显然天主被不恰当地称为实体，是为了通过更常用的名称\Quote{本质}而理解祂的存在，这本质是祂真正且恰当地被称呼的名称；所以，或许唯独天主应被称为本质，这是正确的。因为祂确实是独一的，因为祂是不变的；祂向祂的仆人梅瑟宣布这是祂自己的名字，说：\BibleQuote{我是自有者}；又说：\BibleQuote{你要对以色列子民这样说：\InnerQuote{那自有者派遣我到你们这里来。}}然而，无论祂被称为本质——这是祂恰当地被称呼的，还是被称为实体——这是祂不恰当地被称呼的，祂都是相对于自身而被称呼，而不是\OriginalTerm{相对}{relative}于任何事物；因此，对于天主，存在与自立是同一回事；所以，天主圣三若是一个本质，也是一个实体。因此，或许称祂们为三个\OriginalTerm{位格}{person}比三个实体更方便。\par

% structure: p0018 source=h2 target=subsection reason=relative-heading-rank; base=h2

\subsection{第六章——为什么我们在天主圣三中不谈一个位格和三个本质。没有接受上述所说的人，应当如何相信天主圣三。人既是按照天主的肖像，也是天主的肖像。}

11. 但为免我显得偏袒我们[拉丁人]，让我们再作探究。虽然他们[希腊人]若愿意，也可如他们称三个实体为三个\OriginalTerm{实体}{\GreekText{ὑποστάσεις}}那样，称三个位格为三个\Quote{\GreekText{πρόσωπα}}，但他们更喜欢那个可能更符合他们语言习惯的词。因为位格一词的情况也是如此；对天主而言，存在与成为一位格并非二事，而是全然同一。因为若存在是就自身而言，而位格是相对而言，如此我们就应说三个位格：圣父、圣子、圣神；正如我们说到三个朋友、三个亲戚或三个邻居，他们彼此相对，而非每一个就自身而言。因此，其中任何一个是另外两个的朋友、亲戚或邻居，因为这些名词具有相对意义。那么如何？我们是否应称圣父为圣子和圣神的位格，或圣子为圣父和圣神的位格，或圣神为圣父和圣子的位格？但\Quote{位格}一词在任何情况下通常并不如此使用；在这圣三中，当我们说圣父的位格时，我们所指的不过是圣父的实体。因此，正如圣父的实体就是圣父自己，不是就其为圣父而言，而是就其存在而言，同样圣父的位格也不外乎圣父自己；因为祂之所以被称为一位格是就自身而言，而非就圣子或圣神而言；正如祂就自身而言被称为天主、伟大、良善、公义及任何此类名称；并且正如对祂而言，存在与为天主、为伟大、为良善是同一回事，同样，存在与成为一位格也是同一回事。那么，为何我们不将这三者合称为一位格，如一个本质和一位天主，却称三位格，而我们不称三个神或三个本质？除非是因为我们希望用一个词来表达那藉以理解圣三的含义，这样当我们被问及\Quote{哪三个}时，我们承认他们是三个，就不至于完全缄口不言。因为，如果本质是属，而实体或位格是种，如某些人所想，那么我必须略去我刚才所说的话，即它们应被称为三个本质，正如它们被称为三个实体或位格；如同三匹马被称为三匹马，同时被称为三个动物，因为马是种，动物是属。因为在这种情况下，种并不以复数形式表达，而属以单数形式表达，仿佛我们可以说三匹马是一个动物；但正如它们因专名而被称为三匹马，它们也因属名而被称为三个动物。但如果他们说，实体或位格之名并非指明种，而是指明某一单独和个体之物；如此，一个人不被称作实体或位格，正如他被称作人，因为人是所有人的共性，而是以同样的方式被称为这个或那个人，如亚巴郎、依撒格、雅各伯，或任何其他若在场可以用手指指出的个体：那么同样的理由也适用于这些。因为正如亚巴郎、依撒格、雅各伯被称为三个个体，他们也因此被称为三个人、三个灵魂。那么为何，如果我们也要根据属、种和个体来推理圣父、圣子、圣神，却不称他们为三个本质，正如他们被称为三个实体或位格？但这一点，如我所说，我略过不提：但我肯定，如果本质是一个属，那么单个本质就没有种；正如，因为动物是一个属，单个动物就没有种。因此，圣父、圣子、圣神并非一个本质的三个种。但如果本质是一个种，如人是种，而我们所称为实体或位格的是三个，那么他们共有同一个种，如同亚巴郎、依撒格、雅各伯共有被称为人的种；但人不是被划分为亚巴郎、依撒格、雅各伯，如同一个人可以被划分为几个单独的人；因为这是全然不可能的，因为一个人已经是一个单独的人。那么，为什么一个本质被划分为三个实体或位格？因为，如果本质是一个种，如人是种，那么一个本质就如同一个人：或者，正如我们所说的，任何三个具有相同性别、相同身体构造、相同心智的人是一个本性——因为他们是三个人，但是一个本性——同样在圣三中，三个实体一个本质，或三个位格一个实体或本质？但这在某种程度上是类似的，因为古代说拉丁语的人，在他们拥有这些术语（即本质或实体，这是最近才使用的）之前，曾用\Quote{本性}一词。因此，我们并非按照属或种来使用这些术语，而是仿佛按照一个共同且相同的质料。就如同如果三座雕像由同一金子制成，我们会说三座雕像一块金子，但我们既不应称金为属、雕像为种，也不应称金为种、雕像为个体。因为没有种超越自己的个体，从而包含任何外在于它们的事物。因为当我定义人是什么（这是一个种名）时，每一个存在的人都被包含在同一单个定义中，没有任何不属于人的东西属于它。但当我定义金时，不仅是雕像（如果它们是金的），而且戒指以及任何由金制成的其他东西，都将属于金；即使没有东西由金制成，它仍被称为金；因为，即使没有金雕像，也并非因此就没有任何雕像。同样，没有种超越它属的定义。因为当我定义动物时，由于马是此属的一个种，每一匹马都是一个动物；但每一座雕像却不是金。因此，尽管在三座金雕像的情况下，我们正确地称三座雕像、一块金子；但我们并不这样说，以致于把金理解为属，把雕像理解为种。因此，我们也不这样称圣三为三个位格或实体、一个本质和一位天主，如同三个某物出自同一质料（留有余地，\Emph{即}）；尽管无论那是什么，它都在三者中展现。因为除了圣三，那本质别无他物。然而，我们说同一本质的三个位格，或三个位格一个本质；但我们不说出自同一本质的三个位格，仿佛在其中本质是一回事，位格是另一回事，正如我们可以说出自同一金子的三座雕像；因为在那里，是金是一回事，是雕像是另一回事。当我们说三个人一个本性，或三个同一本性的人，他们也可以被称为出自同一本性的三个人，因为从同一本性中还可以有另外三个这样的人。但在圣三的本质中，无论如何不可能有任何其他位格出自同一本质。此外，在这些事物中，一个人不等于三个人加起来；两个人多于一个人；在相等的雕像中，三个加起来的金子比每一个单尊的多，一尊的金子比两尊的少。但在天主并非如此；因为圣父、圣子、圣神一起并不比圣父单独或圣子单独有更大的本质；但这三个实体或位格（如果必须这样称呼）加在一起等于每一个单独者：这是属自然的人不能领悟的。因为除非在体积和空间的条件下（或大或小），他无法思考，因为幻象或仿佛物体的影像在他脑海中飞舞。\par

12. 在未从这种不洁中被净化之前，让他相信圣父、圣子、圣神——唯一的天主，独一、伟大、全能、良善、公义、仁慈，一切有形与无形之物的创造者，以及凡按人的能力所能恰当而真实地论及祂的一切。当被告知只有圣父是天主时，让他不要把圣子或圣神与祂分开；因为与祂一起，祂是唯一的天主，同样与祂一起，祂也是唯一的天主；因为当我们被告知圣子也是唯一的天主时，我们必须毫不分离地理解，连圣父或圣神也不分离。让他这样讲论同一性体，即不认为一位比另一位更大、更善或任何方面有所不同。然而，不是指圣父本身既是圣子又是圣神，也不是指每一位相对于其他位所独有的称呼，例如圣言只用于圣子，恩赐只用于圣神。正是因此，他们也承认复数形式，正如福音所载：\BibleQuote{我与父原是一体}。祂既说了\BibleQuote{\Emph{一体}}，又说了\BibleQuote{\Emph{我们是}一体}，这是按性体而言，因为他们是同一天主；又说了\BibleQuote{\Emph{我们是}}，这是按关系而言，因为一位是圣父，另一位是圣子。有时性体的合一未被言明，而只有关系者以复数被提及：\BibleQuote{我父与我将来到他那里，并在那里安置我们的住所}。\Emph{我们将来到}和\Emph{我们将安置住所}是复数形式，因为先前说了\BibleQuote{我与父}，即圣子与圣父，这些术语彼此相对使用。有时意义完全隐含，如《创世记》所载：\BibleQuote{让我们照我们的肖像，按我们的模样造人}。\Emph{让我们造}和\Emph{我们的}都是以复数说出，只应按关系者来理解。因为并非诸神造人，或按诸神的肖像和模样造人；而是圣父、圣子、圣神按圣父、圣子、圣神的肖像造人，使人能作为天主的肖像而存在。天主就是天主圣三。但由于天主的肖像并非完全与祂平等，因为它不是由祂而生，而是由祂所造；为了表明这一点，人这样成为肖像，即他是\BibleQuote{按照肖像}，也就是说，他不是通过等同被造得相等，而是通过某种相似而接近祂。因为接近天主不是藉着空间的间隔，而是藉着相似，远离祂是藉着不相似。因为有些人作此区分，他们认为圣子是肖像，而人不是肖像，只是\BibleQuote{按照肖像}。但宗徒驳斥他们说：\BibleQuote{男人当然不应蒙头，因为他是天主的肖像和光荣}。祂没有说\Emph{按照肖像}，而是\Emph{肖像}。这肖像，既然在别处被称为\BibleQuote{按照肖像}，并不是相对于与圣父平等的圣子而言的；否则祂就不会说\BibleQuote{按照我们的肖像}。因为当圣子只是圣父的肖像时，如何能说\Emph{我们的}？但人被称为\BibleQuote{按照肖像}，如前所述，是由于肖像的不平等；因此是按照\Emph{我们的}肖像，使人成为圣三的肖像；并非像圣子与圣父平等那样与圣三平等，而是如前所述，通过某种相似而接近它；正如遥远的事物之间在某种意义上可以表示亲近，不是就地方而言，而是就某种效仿而言。因为经上也说：\BibleQuote{你们应更新你们的心思意念}。祂同样对这些人说：\BibleQuote{因此，你们应该效法天主，如同蒙爱的儿女}。因为论到新人，经上说：\BibleQuote{这新人按创造他者的肖像，在知识上更新自己}。或者，如果我们为了满足论证的需要而承认复数形式，甚至不顾及关系术语，以便当被问及\Quote{三个什么}时，我们可以用一个术语回答，说三个实体或三个位格；那么，不要让人想到任何体积或间隔，或任何哪怕最小的不相似的距离，以致在天主圣三中，任何一位被认为稍小于另一位，无论事物之间如何可能有小大之别：为了使既无位格混淆，也无任何不平等的区分。如果理智不能把握这一点，就让信德持守它，直到那藉先知发言者照亮人心，祂说：\BibleQuote{假使你们不信，你们必不能存立}。\par

\SourceRecord{Translated by Arthur West Haddan.；Nicene and Post-Nicene Fathers, First Series；Vol. 3.；Edited by Philip Schaff.；Buffalo, NY: Christian Literature Publishing Co.,；1887.；<http://www.newadvent.org/fathers/130107.htm>.}

% structure: p0001 source=h1 target=section reason=page-title

\section{论天主圣三（第八卷）}

\EditorialNote{说明并证明：不仅圣父不大于圣子，而且二者一同不大于圣神，天主圣三中任意二者一同不大于其一，三者一同也不大于各自。进而表明：如何从我们对真理的理解、对至高之善的认识、以及对正义的内在之爱（即使是尚未成义的心灵也爱慕正义的灵魂）来领悟天主本身的本质。但首要的是强调：应当藉着爱寻求对天主的认识，圣经中称天主就是爱；在这爱中也指出了某种天主圣三的痕迹。}\par

% structure: p0003 source=h2 target=subsection reason=relative-heading-rank; base=h2

\subsection{序言——以上论述的总结。在信仰的较为困难的问题中应遵循的规则。}

我们在别处已经说过，在天主圣三中，那些相对地指称各个位格的谓词，如圣父与圣子以及二者的恩赐——圣神，是专属于每一位的；因为圣父不是天主圣三，圣子不是天主圣三，那恩赐也不是天主圣三；但当每一位本身被单独论及时，他们不是以复数形式被说成三位，而是一体——天主圣三本身，如圣父是天主，圣子是天主，圣神是天主；圣父是善，圣子是善，圣神是善；圣父是全能，圣子是全能，圣神是全能：然而不是三位神、三位善、三位全能者，而是一位天主、善、全能者——天主圣三本身；以及任何其他关于他们所说的，不是彼此相对，而是各个本身。因为他们正是按本质被如此论述，因为在它们内，存在与伟大、与善、与智慧，以及任何其他关于每一位或天主圣三本身所说的，都是同一的。因此，他们被称为三位格或三个本体，并非为了表明本质的差异，而是为了若有人问这三个是什么，或哪三样东西，我们能用某个词作答。天主圣三中存在着如此伟大的平等，以致不仅就神性而言圣父不大于圣子，而且圣父与圣子一同也不大于圣神；三位的每一位，无论哪一位，也不小于天主圣三本身。这就是我们所说的；若经常讨论和重复，无疑会变得更加熟悉。然而讨论也必须有某种限度，我们应以最虔诚的热心祈求天主，愿祂开启我们的悟性，除去争辩的倾向，使我们的心智能够辨别那既无体积亦无变化的真理本质。如今，既然造物主自己以祂奇妙的仁慈帮助我们，让我们思考这些主题，我们将比之前更深入地进入它们，尽管它们实际上是相同的；同时持守这一规则：凡尚未被我们的理智清楚阐明的，也不应从我们的信德之坚固中松脱。\par

% structure: p0005 source=h2 target=subsection reason=relative-heading-rank; base=h2

\subsection{第一章 以理性表明：在天主内，三位并不比一位更大。}

2. 因为我们说，在这天主圣三中，两位或三位并不比其中一位大；血肉的知觉不接受这一点，别无其他原因，只因它只是按其所能知觉受造的真事物，却不能辨别创造它们的真理本身；因为若它能，那么连有形之光也丝毫不会比我们所说的更清楚。因为关于真理的本体，既然只有它真正存在，没有什么是更大的，除非因为它更真实地存在。但关于一切理智的、不变的事物，没有一样比另一更为真实，因为全都同样是不变的、永恒的；其中被称为伟大的，并非从别的来源而伟大，而是从它真正存在的来源。因此，在那里伟大本身就是真理，凡具有更多伟大的，必须具有更多真理；所以，凡不具有更多真理的，也不具有更多伟大。此外，凡具有更多真理的，当然更真实，恰如具有更多伟大的更大；因此，关于真理的本体，那更真实的就更伟大。但圣父与圣子一同并不比圣父单独或圣子单独更真实。所以二者一同并不比其中每一位单独更大。既然圣神也同样真实，圣父与圣子一同并不比祂更大，因为他们也不更真实。同样，圣父与圣神一同，既然在真理上并不超过圣子（因他们不更真实），在伟大上也不超过祂。同样，圣子与圣神一同，正如圣父单独一样伟大，因为他们同样真实。同样，天主圣三本身与其内的每一位同样伟大。因为在那里真理本身即是伟大，那不更真实的并不更伟大：就真理的本质而言，真实与存在是同一的，存在与伟大是同一的；所以伟大与真实是同一的。因此，关于它，同样真实的必定也是同样伟大的。\par

% structure: p0007 source=h2 target=subsection reason=relative-heading-rank; base=h2

\subsection{第二章 必须抛弃一切有形体的观念，才能理解天主是如何作为真理的。}

3. 但就物体而言，可能这块金与那块金都是真的[实在的]，但这一块可能比那一块更大，因为在这种情况下，大小与真不是同一回事；成为金是一回事，成为大是另一回事。在灵魂的本性中也是如此：灵魂被称为大，并不在于它被称为真的方面。因为拥有一个真的[实在的]灵魂的人，也可能没有伟大的灵魂；因为身体和灵魂的本质并非真理[实在]本身的本质；而天主圣三，唯一的天主，伟大、真实、信实、真理。如果我们试图思考祂，在祂自己允许和赐予的范围内，不要想象任何触及或拥抱于局部空间，如同三个身体，或任何结合体的紧密，如同神话中传说的三体\Person{格律翁}；反而，任何可能在心灵中出现的这样的事物——即三者中比每个单独更大，且单独中比两个更小——都应当毫不犹豫地被拒绝；因为如此一切有形体的事物都被拒绝了。但在精神事物中，也莫让任何可变的、可能出现在心灵中的事物被当作天主。因为当我们从这个深处向那个高处奋进时，在我们能够知道天主是什么之前，已经能够知道祂不是什么，这已是迈向了不小的知识。因为祂绝不是地，也不是天；也不是，譬如说，地和天；也不是任何我们在天上所见的那样的事物；也不是任何我们未见，但或许在天上的事物。即使你凭借思想的想象把太阳的光放大到你能做到的程度，无论是让它更大，还是更亮，一千倍，或无数倍；这也不是天主。就如我们想象纯洁的天使是神灵，赋予天体生命，并按照他们事奉天主的旨意改变和处理它们；即使所有\Quote{千万的千万}聚在一起，成为一体；也不是任何类似的事物是天主。即使你想象这些神灵没有身体——这对有肉体的思想来说确实是非常困难的事。看，如果你能，啊，被可朽坏的身体所压制、被许多各种属世思想所拖累的灵魂；看，如果你能，天主是\OriginalTerm{真理}{veritas}。因为经上记载说\BibleQuote{天主是光}；并非如同这双眼睛所看见的那样，而是如同当有人说祂是真理[实在]时心灵所看见的那样。不要问真理[实在]是什么，因为属肉体的想象和幻象的云彩会立刻挡在路上，并扰乱那安宁，即当我说真理[实在]时在你心中最初闪烁的安宁。努力留在那最初闪烁之中，如果你能，你就被那闪烁照亮，仿佛被一道闪电，当有人对你说\Quote{真理[实在]}时。但你不能；你会滑回那些平常和属世的事物中。请问，是什么重量使你如此滑回？难道不是你所染上的贪欲污点的粘鸟胶，以及你偏离正道之错误的重量吗？\par

% structure: p0009 source=h2 target=subsection reason=relative-heading-rank; base=h2

\subsection{第三章．如何知道天主是至善。除非转向天主，心灵不能成为善。}

4. 再看，且试试你能否看见。你确实只爱那善的，因为大地是善的，有高山之巍峨、丘陵之适度、平原之平坦；良田是善的，可喜而肥沃；房屋是善的，比例匀称、宽敞明亮；动物和有生命的身体是善的；空气是善的，温和而有益健康；食物是善的，可口且利于健康；健康是善的，无痛苦无疲倦；人的面容是善的，五官端正、神情愉悦、容光焕发；朋友的心灵是善的，有默契的甘甜和爱的信赖；义人是善的；财富是善的，因为它随时有用；苍天是善的，有日月星辰；天使是善的，因他们神圣的顺从；言论是善的，甜美地教导并恰当地劝诫听者；诗歌是善的，音韵和谐、含义深刻。何必再一一列举？这个是善的，那个也是善的；但请撇开这个和那个，若你能够，直视善本身，那么你将看见天主，祂不是因别的善而善，而是万善之善。因为在所有这一切善的事物中，无论是我提到的，还是任何其他可觉察或可思想的，当我们真实地判断时，我们之所以能说一个比另一个更好，除非我们心中已印有善本身的观念，据此我们既赞许某些事物为善，又把一个善置于另一个善之上。所以应当爱的是天主，不是这个或那个善，而是善本身。因为灵魂所当寻求的善，并不是一个它须借判断超乎其上的善，而是它须借爱来依附的善；这除了天主还能是什么？不是善的心灵、或善的天使、或善的苍天，而是善本身的善。因为或许我所说的可以通过这种方式更容易理解。例如，当一个心灵被称为善时，正如有两个词，我从这些词中理解两件事——一是它是心灵，另一是它是善的。它本身并未参与使自己成为心灵，因为那时还没有任何东西能使它成为任何东西；但要使自己成为善的心灵，我看出，必须通过意志来实现：并非因为它作为心灵本身不是任何善——否则它怎能已经被称为，且最真实地被称为，比身体更好呢？——但它之所以尚未被称为善的心灵，是因为还需要意志的行动，使它变得更加卓越；如果它忽略了这个行动，那么它应受责备，且合理地不被称作善的心灵。因为它与履行了此行动的心灵不同；既然后者值得称赞，前者无疑因未履行而应受责备。但当它有目的地这样做，并成为善的心灵时，它仍不能达到如此，除非它转向某物，而这物本身不是它。它要成为善的心灵，能转向什么，除非转向它所爱、所求、所获得的善？如果它从这善再转回，而变得不善，那么正因它从善转离，除非那它曾转离的善仍留在它内，否则它若想悔改，就无法再转回那里去。\par

5. 因此，若没有不变之善，就不会有可变之善。每当有人对你提到这个善那个善时，这些事物在另一方面也可被称为不善；如果你能撇开那些因分有善而成为善的事物，并认出那善本身，因分有它它们才成为善（因为当谈到这个或那个善的事物时，你也同时理解善本身）；那么，我说，如果你能除去这些事物，并看见善本身，你就看见了天主。如果你以爱依附祂，你将立刻有福。然而，其他事物之所以被爱，只是因为它们是善的；所以，你依附它们时，竟不爱心爱那使它们成为善的善本身，该当惭愧。那心灵，正因为它是心灵，当它尚未通过转向不变之善而成为善时，而只是如我所说、仅仅是心灵时；每当它如此使我们喜悦，以至于如果我们正确理解，我们甚至将它置于一切有形之光以上时，它并非在其本身中使我们喜悦，而是在那创造它的技艺中。因为它作为受造物而被认可，是在其中看到它应当被造成之处。这就是真理，是纯善：因为它不是别的，正是善本身，因此也是至善。因为除非那善是从另一个善而来的，否则没有善可以减小或增加。因此，心灵为了成为善，就转向那使它成为心灵的那一位。因此，当意志转向那善而爱它时，意志就与本性和谐，使心灵在善中得以成全，而那善也带来了另一善，即当意志转离它时也不会失去的善。因为从至善转离，心灵便失去了成为善的心灵，但并不失去成为心灵。而这也是一种善，且比身体更好。因此，意志失去了意志所获得的。因为心灵早已存在，它能够渴望转向它所从出的那一位；但尚未存在的是，它能够在存在之前就渴望。而我们的至善就在这里：当我们看见一件事物是否应当存在或曾经存在，我们领悟到它应当存在或曾经存在，并且当我们看见那件事物若非应当存在就不可能存在时，我们也不理解它是以何种方式应当存在的。这善离我们每个人都不远：因为我们生活、行动、存在都在祂内。\par

% structure: p0012 source=h2 target=subsection reason=relative-heading-rank; base=h2

\subsection{第四章．天主必须先由无误的信德认识，才能被爱。}

6. 然而，我们必须藉着爱德站稳于此并紧守于此，好能享受那使我们存在的那一位的临在，而缺了祂我们根本无法存在。因为\BibleQuote{我们现在因信德而行，并非因目睹}，我们的确尚未看见天主，正如同一位宗徒所说，\BibleQuote{面对面}。但除非我们现在就已经爱祂，我们将来绝不会看见祂。然而，谁会爱他所不认识的呢？因为一样东西可能被认识却不被爱，但我要问：不被认识的东西能被爱吗？倘若不能，那么就没有人在认识天主之前爱祂。而认识天主，除了以心神默观并稳固地领悟祂之外，又是什么呢？因为祂不是可以用肉体的眼睛来搜寻的形体。然而，在我们有能力默观和领悟天主（如祂所能被默观和领悟的那样，这是心净者所特许的）之前，因为\BibleQuote{心里洁净的人是有福的，因为他们要看见天主}——除非祂因信德而被爱，否则心就不可能洁净，以便适宜和合宜地看见祂。因为那三样（即信德、望德和爱德）——为使它们建立在心灵中，整个神圣圣经的框架已被建立起来——除非存在于一个相信它尚未看见之物、并盼望和爱它所相信之物的心灵中，又能在哪里呢？因此，即使是那不被认识却被相信的，也能被爱。但无可争辩的是，我们必须小心，免得那相信它未见之物的心灵，为自己虚构了某种不存在的东西，并盼望和爱那虚假的。因为那样的话，它就不是出于纯洁的心、无亏的良心和无伪的信德的爱德了，而爱德正是命令的终向，如同一位宗徒所说的。\par

7. 但是，当我们因阅读或听闻而相信任何我们未曾见过的有形事物时，心灵必然会为自己构造某种带有形体特征和形式的形象，正如它可能偶然进入我们思想的那样；这形象要么不真实，即便真实（这极少发生），我们凭信德相信它也无益处，只是有助于通过它来暗示的某种其他目的。因为，有谁阅读或听闻宗徒保禄所写的，或关于他所写的，而不会在想象中构想出这位宗徒本人以及其中提到名字的所有人的面貌呢？而在众多认识这些书的人当中，每个人以不同的方式想象那些形体的特征和形式，究竟谁想象得更为接近和类似现实，这确实是不确定的。确实，我们的信德并不忙于那些人的形体面貌，而只在于：藉着天主的恩宠，他们如此生活并如此行事，正如圣经所见证的那样；这正是相信有益处、并且不可绝望、必须追求的事。因为连主自己在肉身中的面貌，也被无数想象的不同方式所幻想，而它本来只有一个，无论它是什么。在我们对主耶稣基督的信德中，有益的并非心灵为自己想象的东西（或许与现实大相径庭），而是我们按祂的类别所认作的人：因为我们有一种关于人性的概念，像规则一样根植于我们心中，据此我们立即知道，凡我们看到的这样的东西就是一个人或人的形像。\par

% structure: p0015 source=h2 target=subsection reason=relative-heading-rank; base=h2

\subsection{第五章——如何爱那虽未被认识的天主圣三。}

我们的观念是按照这一概念形成的，当我们相信天主为了我们而成为人，作为谦卑的榜样，并显示天主对我们的大爱时。因为我们相信并坚定地持守在心中的是：天主由女人出生，并被有死之人通过如此巨大的侮辱而引向死亡的谦卑，是医治我们骄傲肿胀的最主要疗法，是解开罪债锁链的深奥奥迹。同样，因为我们知道什么是全能，我们相信全能的天主在祂奇迹和复活的大能中，并根据那些或由自然植根于我们、或由经验收集的事物的种类与类型，来形成关于此类行为的观念，使我们的信德不致虚假。因为我们既不认识童贞玛利亚的面貌——她从无丈夫玷污、且在分娩本身也无玷污地奇妙生下了祂；我们也没有见过拉匝禄身体的轮廓，也没有见过伯达尼，也没有见过坟墓和那块祂命令挪开以便从死者中复活祂的石头；也没有见过那从岩石中凿出的新坟墓，祂自己就是从那里复活的；也没有见过橄榄山，祂是从那里升天的。简而言之，我们中间没有见过这些事的人，不知道它们是否如我们所想象的那样，甚且更可能判断它们并非如此。因为当我们偶然事先想象的一个地方、一个人或任何其他形体的外貌，在它出现在眼前时与先前出现在心中时相同，我们会感到不小的惊奇。这种情况极少且几乎从不发生。然而我们最坚定地相信这些事，因为我们根据我们确定的具体和一般概念来想象它们。因为我们相信我们的主耶稣基督是由一位名叫玛利亚的童贞女所生。但什么是童贞女，什么是出生，什么是专有名词，我们不是相信，而是确实知道。至于在谈论或回忆这些事时出现在心中的是否是玛利亚的面貌，我们既不知道，也不相信。因此，在此情况下，我们可以不违反信德地说：或许她有这样的或那样的面貌，或许她没有；但没有人能不违反基督徒的信德说：或许基督是由一位童贞女所生。\par

8. 因此，既然我们渴望理解天主圣三的永恒、平等与合一，尽我们所被允许的，但应在理解之前先相信；并且我们务必谨慎，使我们的信德不致虚假；既然我们必须享有同一的天主圣三，才能生活幸福；但若我们相信了任何关于它的虚假之事，我们的望德便会毫无价值，我们的爱德也不纯洁：那么，我们怎能藉着相信来爱那我们所不认识的天主圣三呢？是按特殊的或一般的概念，就像我们爱宗徒 \Person{保禄} 那样吗？就他而言，即使他并非我们想起他时在脑海中出现的那副面容（对此我们全然不知），但我们知道人是什么。因为不必远说，\Emph{我们}就是人；显然，他也是人，并且他的灵魂与身体结合，过着必死之人的生活。因此，我们相信他具有我们自己所发现的人性，按所有人类本性同属的种或属。那么，关于那至卓越的天主圣三，我们有什么特殊的或一般的认识呢？仿佛有许多这样的圣三，其中一些我们凭经验得知，以便我们可以相信那个圣三也像它们一样，通过相似性的规则印在我们心中，无论特殊的或一般的概念；并因此而爱那件我们相信但尚未认识的事，从我们所认识之事的相同性？但这当然不是如此。或者，正如我们在主 \Person{耶稣基督} 内爱祂从死者中复活，虽然我们从未见过任何人从那里复活，同样我们可以相信并爱那看不见的、且从未见过其相似体的天主圣三吗？但我们确实知道什么是死，什么是活；因为我们活着，也时常看见并经历过死人和快要死的人。至于复活，除了重新活着，即从死亡返回生命，还能是什么呢？因此，当我们说并相信有一位天主圣三时，我们知道什么是三，因为我们知道三是什么；但这并不是我们所爱的。因为我们只要愿意，随时可以轻易地拥有它，撇开其他事物，只需举起三根手指。或者，我们爱的并非每一个三，而是\Emph{那}一位圣三，即天主？那么我们在圣三中所爱的是：它是天主。但我们从未见过或认识任何别的天主，因为天主是唯一的；祂是那唯一我们尚未看见，却凭信德爱慕的天主。但问题是，从什么已知事物的相似或比较中我们可以相信，以便爱那我们还不知道的天主呢？\par

% structure: p0018 source=h2 target=subsection reason=relative-heading-rank; base=h2

\subsection{第六章——尚未成义的人如何能认识他所爱的义人}

9. 所以，请与我一同回转，让我们思考我们为何爱那位宗徒。难道是因为他的人性——我们清楚地知道，我们相信他曾是一个人吗？当然不是；如果是这样，那么他现在就不是我们所爱的那位，因为他已不再是那个人，他的灵魂已与身体分离。但我们相信我们所爱在他身上的仍然活着，因为我们爱他的正义的心灵。那么，除了我们知道心灵是什么、以及成为正义的是什么之外，我们还有什么普遍的或特殊的标准呢？我们不妨说，我们之所以知道心灵是什么，是因为我们也有心灵。因为我们从未用眼睛看见过它，也未曾从多个所见心灵的相似性中归纳出普遍或特殊的概念；相反，正如我之前所说，是因为我们自己拥有它。因为有什么东西比那借以感知一切其他事物、即心灵自身更亲密地认识自身、感知自身呢？因为我们也通过身体的运动认出他人也活着，这源于我们自身的相似性，因为我们活着时也如此运动我们的身体，正如我们观察到那些身体被运动一样。因为即使当一个活的身体运动时，我们的眼睛也无法看到心灵——这个无法被眼睛看见的东西；但我们察觉到那个体积内包含某种如同我们自身所含有的东西，以便以相似的方式运动我们自己的体积，这就是生命和灵魂。这并非人类远见和理性的特有属性，因为即使禽兽也感知到不仅它们自己活着，而且彼此之间也相互感知对方活着，并且我们也如此活着。它们也没有看见我们的灵魂，除了通过身体的运动，并且是立即且最容易地通过某种自然的一致。因此，我们既从自己认识任何人的心灵，也通过自己去相信我们不了解的人。因为我们不仅能感知到心灵存在，还能通过反省我们自己的心灵知道心灵是什么：因为我们有心灵。但是我们从哪里知道什么是正义的人呢？因为上面我们说过，我们爱那位宗徒，仅仅因为他是一个正义的心灵。所以，我们知道什么是正义的人，就像我们知道什么是心灵一样。但什么是心灵，如已说过的，我们从自己知道，因为我们里面有心灵。然而，如果我们自己不正义，我们从哪里知道什么是正义的人呢？但如果只有正义的人才知道什么是正义的人，那么仅有正义的人才会爱正义的人；因为人不能爱他所相信为正义的人——正因他相信其为正义——除非他知道成为正义的是什么；根据我们上面所展示的，除非凭某种普遍或特殊概念的标准，没有人会爱他所相信却未见之物。如果因此只有正义的人爱正义的人，那么一个尚未成为正义的人怎么会愿意成为正义呢？因为没有人愿意成为他不爱的事物。但无疑，为使不正义的人成为正义，他必须愿意成为正义；而为了愿意成为正义，他爱正义的人。所以，甚至尚未正义的人也爱正义的人。但不知道什么是正义的人，就不能爱正义的人。因此，甚至尚未正义的人也知道什么是正义的人。那么他从哪里知道呢？他用眼睛看见了吗？有什么有形之物是正义的，如同白、黑、方、圆那样？谁能这样说？然而人用眼睛只见过有形之物。但人身上的正义唯有心灵；当人称一个人为正义时，是从心灵而非身体而称的。因为正义在某种程度上是心灵的美，使人美丽；许多身体畸形丑陋的人也是如此。正如心灵不能用眼睛看见，其美也是如此。那么，尚未正义的人从哪里知道什么是正义的人，并爱正义的人以便成为正义呢？难道通过身体的运动有某些迹象显示某人正义吗？但若某人完全不知道什么是正义，他又怎么知道这些是正义心灵的迹象呢？所以他确实知道。那么，即使我们尚未正义，我们从哪里知道什么是正义呢？如果我们从自身之外知道，那就是通过某种有形之物。但这不是有形之物。所以，我们在自身内知道什么是正义。因为当我试图说出正义时，我找不到别处，只有在我自身内；如果我问别人什么是正义，他在自身内寻找答案；凡能真实回答的人，都是在自身内找到了答案。事实上，当我想谈论迦太基时，我就在自身内寻找可说的，并在自身内发现迦太基的概念或影像；但我通过身体、即通过身体的知觉获得了它，因为我曾亲身到过那座城，看见并感知了它，并把它保存在记忆中，以便当我想要谈论它时，能在自身内找到关于它的话语。因为它的语词就是它在记忆中的影像本身，不是当迦太基被念出时的两个音节的声音，甚至不是那个名字本身有时在心中默念时的声音，而是当我说出那个双音节词时，甚至在我说出它之前，我在心中所辨认出的那个东西。同样，当我想谈论从未见过的亚历山大时，它的影像也出现在我面前。因为我曾听许多人说并相信那座城很大，按他们所告诉我的那种方式，我尽力在脑海中形成一个它的影像；这就是当我想要谈论它时，在我心中关于它的语词，在我用声音说出那五个音节——几乎人人都知道这个名字——之前。然而，如果我能把这影像从心中展现给认识亚历山大的那些人的眼前，他们要么都会说：不是那样；要么如果说\Quote{是}，我会大为惊讶；当我凝视心中的影像（如同它的图画）时，我仍不知道它就是它，而是相信那些保留了所见影像的人。但我询问什么是正义时并非如此，我找到它时也非如此，我凝视它并说出它时也非如此；当我被听见时不是我那样被认可，我听见时也非如此认可；好像我曾经用眼睛看见过这类东西，或通过某种身体知觉学过，或从那些这样学过的人那里听说过。因为当我说并深知\Quote{正义的心灵是在生活和行为中有意且明智地将各自应得的给予每个人的心灵}时，我并没有像迦太基那样想到一个不在场的东西，也没有像对亚历山大那样想象它是否如此；而是辨认出某种在场的东西，我在自身内辨认它，尽管我自己并不是我所辨认的那个东西；许多人若听见，会认可它。凡听见我并明智地认可的人，也在他自身内辨认这同一事物，即使他自己也不是他所辨认的。但当一位正义的人说这话时，他辨认并说出他自己之所是。而他从哪里辨认它呢？除了从他自身内？这并不奇怪，因为他除了从自身内还能从哪里辨认自己呢？可奇之处在于，心灵应在自身内看见它从未在别处看见的东西，且真实地看见，并看见那真正的正义心灵，而它自身是心灵，却还不是正义的心灵，而它却在自身内看见了它。难道在一个尚未正义的心灵中存在着另一个正义的心灵吗？或者，如果不存在，那么当它看见并说出什么是正义的心灵、且在自身内而非别处看见它时，而它自己又不是正义的心灵，它在那里看见什么呢？难道那是内在真理，临在于有能力领悟它的心灵中吗？然而并非所有人都有能力；而有能力领悟的人，也并非都是他们所领悟的，即他们自身也并非正义的心灵，正如他们能够看见并说出什么是正义的心灵一样。那么他们除非紧紧依附于他们所领悟的那同一个形式本身，从而由此被塑造并成为正义的心灵，否则怎么能这样呢？不仅是辨认并说出\Quote{正义的心灵是在生活和行为中有意且明智地将每个人应得的给予他的心灵}，而且同样他们自己也能正义地生活并在品格上成为正义，从而将各自应得的给予每人，以致不欠任何人什么，而只是彼此相爱。但除了爱，人怎能依附于那形式呢？那么，为什么我们爱另一个我们所相信为正义的人，却不爱那形式本身——我们在其中看见什么是正义的心灵——以便我们自己也能成为正义呢？难道不是因为我们若不爱那形式，我们根本就不会爱那借着形式而爱的人吗？但当我们不正义时，我们对那形式的爱太少了，以致无法使我们成为正义？因此，那个被相信为正义的人，是通过那形式和真理而被爱的，这形式和真理是爱者在他自身内所辨认和理解的；但那个形式和真理本身不可能从别的来源被爱，只能从它本身被爱。因为除它自身外，我们找不到任何其他这类东西，以致我们可以通过相信而爱它当它未知时，因为我们在此已经知道另一类似的东西。因为凡可能见过这类东西的，就是它本身；没有其他这类东西，因为只有它本身是这样的。所以，凡爱人的，要么因他们是正义而爱，要么为使他们成为正义而爱。因为他也应当如此爱自己，要么因自己是正义，要么为使自己成为正义；这样他爱邻人如同自己就没有风险。因为以别的方式爱自己的人，是错爱自己，因为他爱自己是出于使自己不正义的目的；因此是出于使自己邪恶的目的；由此他就不爱自己了；因为\BibleQuote{爱不义的人，恨自己的灵魂。}\par

% structure: p0020 source=h2 target=subsection reason=relative-heading-rank; base=h2

\subsection{第七章——论真爱，我们借此达到对天主圣三的认识。天主应被寻求，不是向外，藉着寻求借天使行奇事，而是向内，藉着效法善天使的虔诚。}

10. 那么，在我们关于天主圣三和认识天主的探究中，首要关注的除了什么是真爱，倒不如说什么是爱，没有别的。因为那称为爱的应该是真的，否则它就是欲情；因此，那些欲情的人被不恰当地称为爱，正如那些相爱的人被不恰当地称为欲情一样。但这就是真爱：持守真理，使我们能正义地生活，因而能轻看一切可朽之物，相比之下，我们要看重对人之爱，由此我们愿意他们正义地生活。因为这样我们也应准备好为弟兄们而有益地死去，正如我们的主耶稣基督以祂的榜样教导我们的。因为如同有两条诫命，全法律和先知都系于其上，即爱天主和爱近人；圣经并非无因地常用一条代替两条：或者只论及天主，如这段经文：\BibleQuote{我们知道，天主使一切协助那些爱祂的人获得益处}；以及\BibleQuote{谁若爱天主，这人是祂所认识的}；还有\BibleQuote{天主的爱，藉着所赐予我们的圣神，已倾注在我们心中}；以及其他许多章节；因为那爱天主的，必然也做天主所命令的，并且他爱天主的程度与他所做的相称；因此他也必然爱近人，因为天主命令了这一点：或者圣经只提及爱近人，如那段经文：\BibleQuote{你们要彼此承担重担，这样就满了基督的法律}；以及\BibleQuote{全部法律总括在这样一句话里：你要爱你的近人如你自己}；以及在福音中：\BibleQuote{凡你们愿意人给你们做的，你们也要照样给人做：这就是法律和先知。}圣经中还有许多其他段落，其中似乎只命令爱近人以达圆满，而对天主的爱则保持沉默；而法律和先知却系于这两条诫命。但这同样是因为那爱近人的，必然也首先爱爱本身。但\BibleQuote{天主是爱；那存留在爱内的，就存留在天主内。}因此他必然首先爱天主。\par

11. 因此，那些透过统治世界或世界各处的权能来寻求天主的人，远离并被弃绝了天主；不是由于空间的距离，而是由于情感的差异：因为他们努力向外寻找道路，而放弃了自己内在的事物，天主就在其中。因此，即使他们可能曾听到某种神圣的天上权能，或以某种方式想到了它，但他们更贪恋其作为，即人的软弱所惊奇之事，却不效法那获得神圣安息的虔诚。因为他们出于骄傲，宁愿能够像天使那样行事，而不愿通过敬虔成为天使所是。因为没有任何神圣存在以自己的权能为乐，而是以那赐予他相称权能者为乐；他知道，以虔诚的意志与全能者结合，比凭自己的权能和意志去做那些不能做这些事情的人所战栗之事，更是权能的标志。因此，主耶稣基督自己在行这些事情时，为了教导那些惊奇之人更好的事，并引导那些专注于怀疑非常规的暂时之事的人转向永恒和内在之事，说：\BibleQuote{凡劳苦和负重担的，你们都到我这里来，我要使你们安息。你们背起我的轭。}祂没有说：\Quote{跟我学，因为我使死了四天的人复活}；而是说：\BibleQuote{跟我学，因为我是良善心谦的。}因为谦逊，最坚实，比骄傲，最膨胀，更有力、更稳妥。因此祂接着说：\BibleQuote{你们必要找得你们灵魂的安息}，因为\BibleQuote{爱不自大；}并且\BibleQuote{天主是爱；}并且\Quote{那在爱内忠实的人，要在祂内安息}，从外面的喧嚣中被召回，进入宁静的喜乐。请看，\BibleQuote{天主是爱}：我们何必出去，跑到高天和深渊寻求祂呢？祂就在我们内，如果我们愿意与祂同在。\par

% structure: p0023 source=h2 target=subsection reason=relative-heading-rank; base=h2

\subsection{第八章——那爱自己弟兄的，就是爱天主；因为他爱爱本身，而这爱来自天主，且就是天主。}

12. 愿无人说：我不知道我所爱的是什么。愿他爱他的兄弟，他就将爱同一份爱。因为他知道他用以爱的那份爱，胜过他所爱的兄弟。因此，如今他能知道天主，胜过他知道自己的兄弟：显然知道得更多，因为更临在；知道得更多，因为更在他内；知道得更多，因为更确实。拥抱天主的爱，并藉着爱拥抱天主。那就是爱本身，它藉着圣德的纽带，结合所有善天使和所有天主的仆人，将我们与他们彼此结合，也使我们从属地结合于祂自己。因此，我们越是从骄傲的膨胀中痊愈，就越被爱充满；那充满爱的人，除了天主，又能充满什么呢？可是，你会说：我看见爱，并尽我所能以心智凝视它，我相信圣经所说：\BibleQuote{天主是爱；那住在爱内的，就住在天主内。}但当我看见爱时，我并未在其中看见天主圣三。不，你若看见爱，便看见了天主圣三。但若我能，我会提醒你，使你看见你看见了它；只要让爱本身临在，使我们被爱推动向善。既然，当我们爱爱时，我们爱一个爱某物者，而这正是因着这一点，即他确实爱某物；因此，爱爱什么呢，好使爱本身也被爱？因为那不爱任何事物的，就不是爱。但如果它爱自身，它就必须爱某物，好使它能够爱自身作为爱。正如一个词语指示某物，也指示自身，但并不指示自身是词语，除非它指示它确实指示某物；同样，爱也确实爱自身，但除非它爱自身，视为在爱某物，它爱自身就不是作为爱。那么，爱爱什么呢，除了我们用爱所爱的那事物？但这事，从离我们最近的事物开始，就是我们的兄弟。请听宗徒若望如何赞颂兄弟之爱：\BibleQuote{那爱自己兄弟的，就住在光明中，在他身上没有任何绊跌的因由。}显然，他将正义的圆满置于爱我们的兄弟；因为那在他身上\BibleQuote{没有任何绊跌的因由}的，确实是成全的。然而他似乎默然略过了对天主的爱；他绝不会这样做，除非是因为他意在兄弟之爱本身中理解天主。因为在同一书信稍后，他极其清楚地这样说：\BibleQuote{亲爱的，我们该彼此相爱，因为爱是出于天主；凡爱的，都是由天主而生，并且认识天主。那不爱的，也不认识天主，因为天主是爱。}这段经文足够清楚地宣告：这同一兄弟之爱本身（因为我们用以彼此相爱的就是兄弟之爱）由如此大的权威提出，不仅是出自天主，而且就是天主。因此，当我们从爱而爱我们的兄弟时，我们就是从天主爱我们的兄弟；我们也不可能不首先爱那使我们爱兄弟的爱本身：由此可以得出，这两条诫命必须互为前提，不能单独存在。因为既然\BibleQuote{天主是爱}，那爱爱的，必然爱天主；但凡爱兄弟的，必然爱爱。所以稍后他说：\BibleQuote{那不爱自己所看见的兄弟的，怎么能爱自己所看不见的天主呢？}因为他看不见天主的原因，就是他不爱他的兄弟。因为那不爱兄弟的，就不住在爱内；那不住在爱内的，就不住在天主内，因为天主是爱。再者，那不住在天主内的，也不住在光明内；因为\BibleQuote{天主是光，在他内毫无黑暗。}因此，不住在光明内的，他若看不见光，即看不见天主，因为他处于黑暗中，这有什么奇怪呢？但他用肉眼看见兄弟，而天主是不能被肉眼看见的。但若他以属灵的爱爱那用肉眼看见的人，他就会以内在之眼看见那作为爱本身的天主，那藉内在之眼可以被看见的。因此，那不爱他所看见的兄弟的，怎能爱天主呢？他之所以看不见天主，正是因为天主是爱，那不爱兄弟的人是没有这爱的。也不要让那另一个问题困扰我们：我们当用多少爱爱兄弟，当用多少爱爱天主？对天主的爱，无可比拟地超过对我们自己的爱，但对兄弟的爱当如同对自己的爱；我们爱自己越多，就越爱天主。因此，我们从同一份爱爱天主和我们的近人；但我们爱天主是为了天主，爱我们自己和我们近人也是为了天主。\par

% structure: p0025 source=h2 target=subsection reason=relative-heading-rank; base=h2

\subsection{第九章——我们对义人的爱，是由对不变之正义形相的爱本身所燃起的。}

因为为什么我们听到并读到这些时，心中火热呢？\BibleQuote{看，如今正是蒙悦纳的时候；看，如今正是救恩的日子：凡事不给人绊脚石，免得这职务受指责；反倒在各样事上表明自己是天主的仆役，以极大的坚忍、患难、穷乏、困苦、鞭打、监禁、扰乱、勤劳、警醒、禁食；以纯洁、知识、坚忍、恩慈、圣神、无伪的爱、真理之言、天主的德能、左右两手的义德武器；借着荣耀和羞辱、恶名和美名：像是迷惑人的，却是诚实的；似乎不为人知，却是人所共知的；似乎要死，看，我们却活着；似乎受惩罚，却不被处死；似乎忧愁，却常常喜乐；似乎贫穷，却使许多人富足；似乎一无所有，却拥有万有。}为什么我们读到这些，就被对保禄宗徒的爱所激励？除非我们相信他确实这样生活。但我们相信天主的仆役应该这样生活，并不是因为我们从某人那里听到，而是因为我们内在、在我们之上，在真理本身中看见了它。因此，我们因我们所见而爱那位我们相信这样生活的人。除非我们超过一切爱那模式——我们看出它永恒不变——我们就不会因那个原因爱他，因为我们坚信他的生活，当他活在肉身时，是与这个模式适应并和谐的。但不知怎的，借着信仰——我们相信某人这样生活——我们更被激励去爱这个模式本身；并且借着希望——我们不再完全绝望，我们也能这样生活——我们作为人，从这样的事实（有些人这样生活）出发，更热切地渴望它，更自信地祈求它。因此，对模式的爱——他们被认为照着它生活——使得这些人的生活本身被我们爱；而他们这样被相信的生活，又激发起对那个模式更炽热的爱；这样，我们越爱天主，就越确定、越平静地看见祂，因为我们在天主内看见那不变的义德模式，我们据此判断人应当怎样生活。所以，信德有助于认识和爱天主，并不是如同认识一个完全未知或完全不被爱的，而是借此可以更清楚地认识祂，更坚定地爱祂。\par

% structure: p0027 source=h2 target=subsection reason=relative-heading-rank; base=h2

\subsection{十、在爱中有三件事，如同天主圣三的一个痕迹。}

但爱或爱德是什么？圣经如此赞扬和宣扬的，岂不就是爱善吗？但爱是出于爱者，借着爱，某物被爱。看，于是有三件事：爱者、被爱者、爱。那么，爱是什么？岂不就是某种生命，它结合或试图结合两个事物，即爱者和被爱者？这在外在的、肉欲的爱中也是如此。但为了汲取更纯净更清晰的，让我们践踏肉体，上升到心灵。心灵在朋友中爱什么，除了心灵？那么，那里也有三件事：爱者、被爱者、爱。还需要由此上升，去寻求那些上面的事，尽人所能。但在这里让我们的目标稍作停留，不是以为已经找到了所寻求的；而是如同通常必须先找到寻找某物的地方，尽管尚未找到那物本身，但我们已经找到了在哪里寻找它；同样，说了这些就足够了，使我们好像有了某个起点的枢纽，从那里编织余下的论述。\par

\SourceRecord{Translated by Arthur West Haddan.；Nicene and Post-Nicene Fathers, First Series；Vol. 3.；Edited by Philip Schaff.；Buffalo, NY: Christian Literature Publishing Co.,；1887.；<http://www.newadvent.org/fathers/130108.htm>.}

% structure: p0001 source=h1 target=section reason=page-title

\section{论天主圣三（卷九）}

\EditorialNote{在人——天主的肖像——内存在某种三一，即理智，以及理智用以认识自身的知识，以及理智用以爱自身及其知识的爱；并显示这三者彼此相等，同属一个实体。}\par

% structure: p0003 source=h2 target=subsection reason=relative-heading-rank; base=h2

\subsection{第一章——我们应如何探究关于天主圣三。}

1. 我们确实寻求一种三一——并非任何三一，而是那既是天主，而且是真实、至高、唯一的天主的天主圣三。那么，让我的听众们等待吧，因为我们仍在寻求。没有人能公正地指责这样的寻求，只要那寻求者——他所寻求的，无论是认识还是讲述，都是极其困难的——在信德上坚定不移。但无论谁，若他看得更清楚或教导得更好，都会迅速而公正地指责那些对此妄加断言的人。经上说：\BibleQuote{寻求天主，你的心就会生活；}并且，免得有人轻率地因好像已经把握了它而欢喜，经上说：\BibleQuote{要不断寻求祂的面容。}宗徒说：\BibleQuote{若有人自以为知道什么，他还不知道他应该怎样知道。但若有人爱天主，这人是被天主所认识的。}他没有说\BibleQuote{认识天主}，那是危险的僭越，而是说\BibleQuote{被天主所认识。}同样，在另一处，当他曾说：\BibleQuote{但现在你们既然认识了天主}后，他立即纠正自己，说：\BibleQuote{更可以说是被天主所认识。}尤其是在另一处，他说：\BibleQuote{弟兄们，我不以为我已经获得了；我只做一件事：忘记背后，努力向前，向着标竿直跑，要得天主在基督耶稣里从高天召我来得的奖赏。所以我们凡是成全的人，都应存这样的心意。}他告诉我们，今生的成全无非是忘记背后，努力向前，向着标竿直跑。因为寻求者有最安全的目标，直到我们抓住了我们所趋向和奋力追求的东西。但那正确的目标始于信德。因为某种信德在某种意义上就是知识的起点；但某种知识只有在今生之后，当我们面对面看见时，才能成为完全的。因此，我们应当存这样的心意，知道寻求真理的态度比擅自将未知当作已知更为安全。因此，我们应当如此寻求，如同我们将要找到；如此找到，如同我们还要寻求。因为\BibleQuote{人做完的时候，才开始。}对于当信的事，我们应无疑虑地怀疑；对于当理解的事，我们应不轻率地断言：在前者上，必须持守权威；在后者上，须寻求真理。那么，关于这个问题，我们当相信圣父、圣子、圣神是一位天主，是整个受造物的创造者和掌管者；并且圣父不是圣子，圣神也不是圣父或圣子，而是彼此关联的位格的三一，平等实体的合一。让我们祈求祂的帮助，以理解祂，我们愿意理解祂；并在祂所赐的范围内，以如此虔诚的关心和焦虑来解释我们所理解的，以至于即使我们在某些情况下说错，至少不说任何不相宜的话。例如，如果我们说了关于圣父的任何不属于圣父，却属于圣子或圣神，或属于天主圣三本身的话；或说了关于圣子的任何不恰合圣子，却恰合圣父，或圣神，或天主圣三的话；或者再次，说了关于圣神的任何不适合于圣神特有属性，却不异于圣父，或圣子，或天主圣三本身的话。甚至现在，我们希望看圣神是否恰是那至高的爱；如果不是，那么爱是圣父，或是圣子，或是天主圣三本身；因为我们不能抗拒圣经最确定的信德和严肃的权威，圣经说：\BibleQuote{天主是爱。}然而，我们不应背离而陷入亵渎的错误，以致对于天主圣三说出任何不适合于创造者，反而适合于受造物，或者仅凭空虚的幻想所捏造的话。\par

% structure: p0005 source=h2 target=subsection reason=relative-heading-rank; base=h2

\subsection{第二章——在爱中所发现的三种事物必须加以考虑。}

2. 既然如此，让我们把注意力转向我们自以为已经发现的那三样事物。我们此刻并非谈论天上的事，也并非谈论天主圣父、圣子和圣神，而是谈论那虽不完美却仍是肖像的肖像，即人；因为我们软弱的心灵或许能更熟悉、更容易地凝视它。那么，当我——进行此项探究的人——爱任何事物时，便涉及三样事物：我自己、我所爱的对象以及爱本身。因为我并不爱爱，除非我爱一个爱人者；因为没有所爱之物，就没有爱。因此，有三样事物：爱人者、被爱者以及爱。然而，若我只爱我自己，岂不只有两样事物：我所爱的对象和爱？因为当一个人爱自己时，爱人者与被爱者乃是同一的；同样，当一个人爱自己时，爱与 被爱 也是同一回事。因为当说\Quote{他爱自己}和\Quote{他被自己所爱}时，所说的是同一件事。在此情形下，爱与 被爱 并非两件不同的事；正如爱人者与被爱者并非两个不同的人。然而，即便如此，爱与所爱之物仍是两样事物。因为当一个人爱自己时，除非爱本身被爱，否则就没有爱。但爱自己是一回事，爱自己的爱是另一回事。因为爱除非已经爱着某物，否则不被爱；既然没有所爱之物，就没有爱。因此，当一个人爱自己时，有两样事物：爱和所爱之物。因为那时爱人者与被爱者是同一的。由此可见，并非凡有爱之处都必然理解为三样事物。让我们从探究中排除人所构成的其他许多事物；为了尽可能清楚地发现我们目前所寻求的，我们只讨论心灵。那么，当心灵爱自己时，它表露出两样事物：心灵和爱。但爱自己除了愿意帮助自己享用自己的存在以外，还有什么呢？当某人愿意自己恰好如其所是时，意志便与心灵相等，爱便与爱人者相等。如果爱是一个实体，它当然不是身体，而是精神；并且心灵也不是身体，而是精神。然而，爱与心灵并非两个精神，而是一个精神；也并非两个本质，而是一个：然而这里有两样事物成为一体：爱人者和爱；或者，如果你喜欢这样表述，所爱者和爱。确实，这些是彼此相对而言的。因为爱人者相对于爱，爱相对于爱人者。因为爱人者是以某种爱去爱，爱是某个爱人者的爱。但心灵和精神并非相对而言，而是表达本质。因为心灵和精神之所以存在，并非因为某个特定的人的心灵和精神存在。因为如果我们从人——人是连同身体一起称呼的——减去身体，便剩下心灵和精神。但如果减去爱人者，就没有爱；如果减去爱，就没有爱人者。因此，就它们彼此相对而言，它们是两样；但就它们自身而言，各自都是精神，两者一起也是一个精神；各自都是心灵，两者一起也是一个心灵。那么，天主圣三在哪里呢？让我们尽力注意，并呼求永恒的光，愿祂光照我们的黑暗，使我们能在自身内，在允许的范围内，看见天主的肖像。\par

% structure: p0007 source=h2 target=subsection reason=relative-heading-rank; base=h2

\subsection{第三章——认识并爱自己的人灵中天主圣三的肖像。心灵藉自身认识自身。}

3. 因为心灵除非也认识自己，否则不能爱自己；因为怎能爱自己所不认识的呢？若有人说，心灵根据普遍或特殊的知识，相信自己是如它从经验中发现别人那样的，因此爱自己，那么他说得极为愚拙。因为心灵若不认识自己，又怎能认识其他心灵呢？因为心灵并非认识其他心灵而不认识自己，就如肉体的眼睛看见其他眼睛却不看见自己；因为我们通过肉体眼睛看见物体，除非我们照镜子，否则不能将那些眼睛所放射出来的光线折射回自身，并触及我们所辨别的任何东西——这确实是一个极为精细而晦涩的论题，直到它被清楚地证明是否如此。但无论我们通过眼睛辨别的力量的本性是什么，确实，无论是光线还是别的什么，我们无法用眼睛辨别那力量本身；而是用心灵去探究它，如果可能，甚至用心灵去理解它。那么，正如心灵本身通过身体的感官收集对有形事物的知识，同样它通过自身收集对无形事物的知识。因此，它也通过自身认识自己，因为它是无形的；因为如果它不认识自己，它就不爱自己。\par

% structure: p0009 source=h2 target=subsection reason=relative-heading-rank; base=h2

\subsection{第四章——三者是一，且同等，即心灵本身、爱及对自身的认识。此三者同属实体，且相对地述说。此三者不可分离。此三者并非部分般结合与混合，而是同属一个本质，且是相对的。}

4. 然而，正如当心智爱自身时，有\Quote{两个事物}（\OriginalTerm{两个事物}{duo quædam}），即心智及心智对自身的爱；同样，当心智认识自身时，有\Quote{两个事物}，即心智及心智对自身的知识。因此，心智本身、心智对自身的爱及心智对自身的知识，是\Quote{三个事物}（\OriginalTerm{三个事物}{tria quædam}），这三者是一；当它们成全时，它们是相等的。因为若一个人爱自身少于他之所是——例如，假设一个人的心智只爱自身如同人之身体应该被爱一样，然而心智大于身体——那么它就有过犯，其爱不完全。再者，若它爱自身多于他之所是——例如，它爱自身如同天主应该被爱一样，然而心智与天主相比不可同日而语——那么它也有极大的过犯，其自爱不完全。但若它爱身体如同天主应该被爱一样，则它的过犯更为乖谬和错误。同样，若知识小于那可知且能被完全认识之物，则知识不完全；但若知识大于该物，则认识的本性高于被认识之物，如同对身体的知识大于身体本身，身体被那知识所认识。因为知识是认识者理性中的某种生命，而身体不是生命；任何生命都大于任何身体，不是通过体积，而是通过能力。但当心智认识自身时，它自己的知识不高于自身，因为自身认识，自身被认识。因此，当它完全认识自身，且不与任何其他事物一起时，它的知识便等于自身；因为它的知识不是来自另一本性，既然它认识自身。当它完全知觉自身，且知觉不到别的时，它便既不更少也不更大。所以我们正确地说，这三者，即心智、爱和知识，当它们成全时，因而是相等的。\par

5. 类似的推理提示我们——如果我们能以任何方式理解这事——这些事物[即爱与知识]存在于灵魂中，并且，它们仿佛隐含在灵魂内，从灵魂中开展出来，以致可以被知觉并被算作实体性的，或可说是本质性的。它们并非如同在一个主体中，如同颜色、形状或任何其他性质或数量存在于身体中那样。因为这一类[物质]事物并不超出其所在的主体；这个特定身体的颜色或形状也不能是另一身体的。但心智也能用其爱自身的爱去爱自身之外的某物。此外，心智不仅认识自身，也认识许多其他事物。因此，爱与知识并非如同在主体内那样被包含在心智中，而是这些事物也\Emph{实体性地存在，如同心智本身一样}；因为，即使它们是相互关联地述谓的，但它们各自存在于自己的实体中。它们也不是像颜色与有颜色的主体那样相互关联地述谓的；颜色在有颜色的主体内，但它自身没有固有的实体，因为有颜色的身体是实体，而颜色在实体中；但如同两个朋友也是两个男人，他们是实体，而说他们是男人并非相对而言，但说他们是朋友则是相对而言。\par

6. 然而，尽管爱者或知者是一个实体，而且\Emph{知识是一个实体}，\Emph{爱是一个实体}，但爱者与爱，或知者与知识，是彼此相对而言的，如同朋友那样：然而心智或精神却不是相对的，如同人也不是相对的：尽管如此，爱者与爱，或知者与知识，不能彼此分离地存在，如同朋友那样的人可以分离。虽然朋友们看起来也可以在身体上分离，但在心神上，就他们是朋友而言，却不能分离：甚至可能发生这种事，一个朋友甚至开始恨另一个朋友，并因此不再是朋友，而另一方却不知道，仍然爱着他。但是，如果心智用以爱自身的爱不再存在，那么心智也同时不再爱。同样，如果心智用以认识自身的知识不再存在，那么心智也同时不再认识自身。正如任何有头之物的头当然是头，它们彼此相对地述谓，尽管它们也是实体：因为头是一个身体，而有头之物也是一个身体；如果没有头，那么也就没有有头之物。只有这些事物可以通过切割而彼此分离，那些却不能。\par

7. 即使有些物体不能完全分开和分裂，但除非它们由自身适当的部分组成，否则它们就不是物体。部分相对于整体来陈述，因为每一个部分都是某个整体的部分，而整体因其拥有所有部分而成为整体。但由于部分和整体都是物体，这些事物不仅被相对地陈述，而且也是实体性地存在。那么，也许心灵是一个整体，而它用以爱自身的爱，以及它用以认识自身的知识，就像是它的部分，这个整体由这两个部分组成。或者，是否存在着三个相等的部分，共同构成一个整体？但没有任何部分包含它所属的整体；然而，当心灵认识自身为一个整体时，即完美地认识自身，那么它的知识就遍及整体；当它完美地爱自身时，它爱自身为一个整体，它的爱也遍及整体。那么，是否如一种由酒、水、蜂蜜调成的饮料，每一个单独的部分都遍及整体，但它们仍是三样东西（因为饮料中没有不包含这三样东西的部分；它们并非像水和油那样结合，而是完全混合在一起：它们都是实体，由三者构成的整个液体是一个实体），——我们是否要以这种方式来设想这三者在一起：心灵、爱和知识？但水、酒和蜂蜜并非同一个实体，尽管由它们混合而成的饮料成为一个实体。而我看不出另外那三者为什么不是同一个实体，因为心灵自身爱自身，自身认识自身；这三者如此存在，以至于心灵完全不被任何其他事物所爱或所认识。因此，这三者必然具有同一个本质；因此，如果它们仿佛通过混合而混淆在一起，那么它们就完全不能成为三者，也不能彼此相互指涉。正如你用同一块金子制成三个相似的戒指，尽管它们彼此相连，但彼此相互指涉，因为它们是相似的。因为每一相似物都与某物相似，于是便有戒指的\OriginalTerm{三一}{trinity}和同一块金子。但如果它们彼此融合，每一个都通过整个体积与另一个相混合，那么那个\OriginalTerm{三一}{trinity}就会崩溃，根本不存在；不仅不能像那三个戒指那样被称为一块金子，而且现在它根本不能被称作三件金器。\par

% structure: p0014 source=h2 target=subsection reason=relative-heading-rank; base=h2

\subsection{第五章——这三者各自在自身内，又彼此都在一切中。}

8. 但在这三者中，当心灵认识自身并爱自身时，存留着一个\OriginalTerm{三一}{trinity}：心灵、爱、知识；这个\OriginalTerm{三一}{trinity}并不因任何混合而混淆在一起：尽管它们各自在自身内，又彼此都在一切中，或各自在两个中，或两个在各自中。因此一切都在一切中。因为心灵确实在自身内，因为它相对于自身被称为心灵：尽管它被认为是认识、被认识或可认识的，相对于它自身的知识；并且尽管它作为爱的、被爱的或可爱的，被关联到它用以爱自身的爱。而知识，尽管它关联到认识或知道的心灵，但它也相对于自身而被谓述为被认识和认识：因为心灵认识自身的知识对自身并非未知。尽管爱关联到所爱的、拥有其爱的心灵，但它也相对于自身是爱，以便也在自身内存在：因为爱也是被爱的，却只能被爱用爱来爱，即用自身。所以这些事物各自在自身内。但它们也在彼此内；因为爱的心灵\Emph{在}爱内，爱在爱者的知识内，知识\Emph{在}认识的心灵内。同样，每个各自在两个中，因为认识并爱自身的心灵在它自身的爱和知识内；而爱并认识自身的心灵的爱在心灵和它的知识内；认识并爱自身的心灵的知识在心灵和它的爱内，因为它认识自身所爱的，爱自身所认识的。因此，每一个两个也在每一个各自中，因为认识并爱自身的心灵同它的知识一起在爱内，同它的爱一起在知识内；而爱和知识本身也在认识并爱自身的心灵内。但一切如何在一切中，我们已在上面阐明；因为心灵爱自身作为一个整体，认识自身作为一个整体，完全认识它自己的爱，完全爱它自己的知识，当这三者相对于自身完美时。因此，这三者奇妙地彼此不可分离，然而每一个各自是一个实体，而它们一起是一个实体或本质，同时它们相互相对地谓述。\par

% structure: p0016 source=h2 target=subsection reason=relative-heading-rank; base=h2

\subsection{第六章——事物本身中有一类知识，永恒真理本身中另有知识。有形事物也应按永恒真理的规则来判断。}

9. 但当人心认识自己并爱自己时，它认识并爱的并非任何不变之物：每个人以某种说话方式宣告他自己的个别之心，当他考虑自身内所发生的事时；但通过专门的或普遍的知识抽象地定义人心。所以，当他对我谈论他自己个别之心，比如他是否理解这个或那个，或不理解，或他是否愿意这个或那个，或不愿望时，我相信；但当他普遍地或专门地谈论关于人心的真理时，我认出并赞同。由此明显可见，每个人在自身内看见一件事，如此另一人能相信他所说的，却可能看不见；但另一人在真理本身内看见一件事，如此另一人也能够注视它；前者在不同时间经历变化，后者则存在于永恒不变之中。因为我们并非通过用肉身眼睛看见许多心智而凭借相似性收集关于人心的普遍或专门知识：而是注视不朽的真理，从中尽可能完美地定义，不是某一个人的心智是如何的，而是它按照永恒理应是怎样的。\par

10. 由此，即使在那些通过身体感官被引入并以某种方式被注入记忆的物质事物的影像中，由此那些未被见过的事物也通过想象的影像被思考，无论它们与实际不同，还是甚至恰如它们实际——即使在这里，我们也通过始终保持在我们心智之上的其他完全不变的规则，在我们正当地接纳或拒绝时，被证明在我们内接受或拒绝。因为当我回忆我曾见过的\Place{迦太基}城墙，并想象我未曾见过的\Place{亚历山大}城墙，并在虚构的形状之间偏好此胜于彼时，我根据理性理由作出偏好；来自上方的真理的判断仍然是强大而清晰的，并牢固地立足于其自身绝对不可摧毁的规则之上；即使它仿佛被物质影像的迷雾所覆盖，却并不被包裹和混淆在其中。\par

11. 但是，区别在于：我是在那黑暗之下或之中仿佛被隔绝于晴朗的天空之外；还是（如同在高山上通常发生的）享受两者之间的自由空气，既向上仰望最宁静的光明，又向下俯视最浓密的云雾。因为当我听说某人为信仰的卓越和坚定而忍受了痛苦的折磨时，兄弟之爱的热情如何在我内点燃？如果那人被指给我看，我渴望与他结合，结识他，将他结为友谊。因此，如果有机会，我走近他，向他说话，与他交谈，尽我所能用言语表达对他的善意，并希望在他内也能产生并表达对我的善意；我努力通过相信而进行一种精神的拥抱，因为我无法如此快速地探究并完全洞察他的内心深处。因此，我以纯洁而真挚的爱爱那忠诚而勇敢的人。但是，如果他在交谈中向我坦白，或者通过疏忽以任何方式显露出，他相信天主有不当之处，并且也在祂内渴望某种肉性之物，而他是为了这样的错误或为了他所希望的金钱的欲望，或为了虚浮的贪求人的称赞才忍受这些折磨：那么，我先前对他的爱，因被冒犯而反感，仿佛被驱离，从这不配的人身上被移除，便停留在我曾经相信他是那样而爱他的那种形式中；除非我已经爱他到这样的目的，即他可能成为那样，而我发现实际上他并非那样。而那人本身也没有改变：尽管他可以改变，以便成为我已经相信他已是的那样。但在我的心中，确实有某事改变了，\Emph{viz.}，我对他所形成的判断，先前的是一种，现在又是另一种：并且同样的爱，从上方不变的义德的命令，从享用的目的转向了劝告的目的。但那不动摇而稳固的真理形式本身，我本应享用那被我以为良善的人，并且同样我劝告他成为良善，以不动移的永恒，将同样的不朽且最健全理性的光芒，既照射到我的心智之视，也照射到那些影像的云雾上，当我思考我所见过的那同一个人时，我从上方辨别这些影像。再者，当我回忆起某个我曾在\Place{迦太基}见过的、建造得美丽而对称的拱门时，某种通过眼睛被心智认识并转移到记忆中的实在，引起了那想象的景象。但我心中还看到另一事物，根据那事物，那艺术品令我愉悦；并且据此，如果它使我不悦，我就会纠正它。因此，我们根据那永恒真理的形式来判断那些个别事物，并通过理性心智的直觉来辨别那形式。而它们本身，我们或者通过身体感官触摸其在场，或者若不在场则记住它们固定在我们记忆中的影像，或者按照与它们的相似性描绘出我们自己也如果愿意且能够就费力建造的这类事物：以一种方式在心中想象身体的影像，或通过身体看见身体；但以另一种方式，通过简单的理智把握心智之眼之上的东西，\Emph{viz.}，那些形式的理由和难以言表的美丽技巧。\par

% structure: p0020 source=h2 target=subsection reason=relative-heading-rank; base=h2

\subsection{第七章　我们在内在中怀有并生出圣言，来自我们在永恒真理中所看见的事物。圣言，无论是关于受造物还是关于造物主，都由爱所怀有。}

12. 我们于是以心灵的目光，在永恒真理中——一切受造之物皆由此真理而成——看见我们依以存在、并以真确理性在自身或形体内行事的形式；并且由此获得对事物的真知，这真知仿佛在我们内产生为言语，我们藉言谈从内里生出它；且它虽被生出，却未离开我们。当我们向他人讲话时，我们运用声音或某种身体标记的服务，将仍存于我们内的言语应用于他人，以便藉某种可感的记忆，在听者心中也产生某种相似之物——相似，我说，于那说话者心中未曾离开之物。因此，我们在言谈和行为中，藉身体肢体所做的一切——人的行为因此受称赞或责备——无不是先在我们内以言语预言之。因为没有人愿意做任何他未先在心中说出的。\par

13. 而这言语是由爱所孕生，这爱或是对受造物，或是对造物主，即或是对可变的性体，或是对不变的真理。\par

% structure: p0023 source=h2 target=subsection reason=relative-heading-rank; base=h2

\subsection{第八章．欲爱与爱德有何区别}

[所孕生的]因此或是由于欲爱，或是由于爱德：并非受造物不该被爱；但若那[对受造物的]爱被归于造物主，那么它便不再是\OriginalTerm{欲爱}{cupiditas}，而是\OriginalTerm{爱德}{charitas}。因为当受造物为自身而被爱时，便是欲爱。那时它不仅无助于人使用受造物，反而在享用中败坏他。因此，当受造物与我们平等或低于我们时，我们必须使用低等者以归于天主，但必须在主内恰当地享用平等者。因为正如你应在主内、而非在你自身内享用你自己，同样也应在主内享用你所爱如己的人。因此，让我们在主内享用我们自身和我们的弟兄；且由此我们切莫敢向自身松垮、仿佛向下放松自己。当一项思想被深究、且令我们喜悦——或出于犯罪，或出于行义——时，言语便诞生了。因此，爱犹如中介，将我们的言语与孕生它的心灵联结，且毫无混乱地作为第三者与二者结合于无形的拥抱中。\par

% structure: p0025 source=h2 target=subsection reason=relative-heading-rank; base=h2

\subsection{第九章．在爱慕属灵事物上，所生的言语与所孕生的言语相同。在爱慕属肉欲的事物上则不然。}

14. 但在对属肉欲和暂时之物的爱中，如同在动物的后代中，言语的孕生是一回事，诞生是另一回事。因为这里由欲爱所孕生的，靠获得而诞生。对贪财而言，知道和爱金钱是不够的，除非也拥有它；知道和爱吃饭或与某人同寝，除非也去行；知道和爱荣誉和权力，除非它们实际实现。不仅如此，这一切即使获得，也不使人满足。\Quote{凡喝这水的，}祂说，\Quote{还要再渴。}\BibleRef{若 4:13}同样，圣咏作者说：\Quote{他怀了劳苦，生了不义。}\BibleRef{咏 7:15}他所说的劳苦或痛苦是被怀上的，即当那些知道和愿意并不足够的事物被怀上，心灵因匮乏而焚烧和生病，直至达到那些事物，仿佛生出它们。因此在拉丁语中，我们优雅地用\OriginalTerm{获得（使诞生）}{parta}一词同时表示\OriginalTerm{发现}{reperta}和\OriginalTerm{查明}{comperta}，这两个词听起来仿佛源自生产。因为\Quote{私欲怀了胎，就生出罪来。}\BibleRef{雅 1:15}因此主宣告：\Quote{凡劳苦和负重担的，你们都到我这里来；}\BibleRef{玛 11:28}并在另一处说：\Quote{那些日子，怀孕的和哺乳的有祸了！}\BibleRef{玛 24:19}当祂因此将所有正确的行为或罪恶都归于言语的诞生时，祂说：\Quote{凭你的口，}祂说，\Quote{你要被成义，并且凭你的口，你要被定罪。}\BibleRef{玛 12:37}——这里意指的不是有形的口，而是内在无形的思想和心中的口。\par

% structure: p0027 source=h2 target=subsection reason=relative-heading-rank; base=h2

\subsection{第十章．惟独被爱的知识是否为心灵的言语}

15. 于是理应询问：是否一切知识都是言语，抑或唯有被爱的知识才是言语？因为我们也认识我们所憎恨的事物；但所不喜爱的事物，不能说被心所孕育或产生。因为并非一切以任何方式触动心的事物都被它孕育，有的仅止于被认识，却不被说成言语，比如我们现在所谈论的便是如此。因为那些以音节占据时间的事物，无论是说出抑或仅被思想，在一种意义上被称为言语；在另一种意义上，凡是被认识的事物，只要能从记忆中引出并被界定，即便我们厌恶其本身，也被称为印在心上的一言语；在另一种意义上，当我们喜爱那在心上被孕育的事物时，也是如此。宗徒所说的那句话必须按这最后一种言语来理解：\Quote{除非受圣神感动，没有一个人能说\InnerQuote{耶稣是主}}；因为那些人也说这句话，但按\Quote{言语}一词的另一种意义，主自己曾论及他们说：\Quote{不是凡向我说\InnerQuote{主啊，主啊}的人，都能进入天国。}甚至就我们所憎恨的事物而言，当我们正确地厌恶并正确地指责它们时，我们赞同并喜爱所加给它们的指责，而这指责便成了言语。令我们不快的不是对邪恶的认识，而是邪恶本身。因为我喜欢认识并界定何谓不节制，而这便是它的言语。正如在艺术中有已知的缺陷，对缺陷的认识被正确地认可，当行家辨认出形式的种或形式的缺失，以肯定或否定其存在或不存在；然而缺乏形式而陷入缺陷，是应受谴责的。界定不节制并说出它的言语，属于道德艺术；但不节制本身则属于该艺术所谴责的范围。正如认识和界定何谓语法错误属于说话的艺术，而犯语法错误则是该艺术所谴责的缺陷。因此，言语——我们现在所愿认识并指明的——是知识与爱相结合。所以，每当心认识和爱自身时，它的言语便藉着爱与它结合。既然心认识爱且知道爱，那么言语就在爱中，爱也在言语中，两者都在那爱着并说着的心中。\par

% structure: p0029 source=h2 target=subsection reason=relative-heading-rank; base=h2

\subsection{第十一章 那认识自己的心的肖像或受生之言与心本身相等}

16. 但所有据有形式的知识与其所认识的事物相似。因为另有一种据有缺乏的知识，当我们谴责时，只按此说出一言语。而此对缺乏的谴责相当于对形式的称赞，因而被认可。因此，心对所知的形式具有某种相似性，无论是喜爱那形式，还是厌恶其缺乏。由此，就我们认识天主而言，我们肖似祂，但尚未达到相等的地步，因为我们对祂的认识不及祂本身的实有。当我们藉身体感官论及物体时，心中便升起物体的某种相似性，那是记忆中的幻象；因为当思想物体时，物体本身并不在心中，而只有那些物体的相似性。因此，当我们为了物体而认同这些相似性时，便犯了错误，因为认此为彼实属谬误；然而心中的物体影像比物体本身的形式更优越，因为它存在于更优的本性中，\Emph{即}活生生的实体——心。同样，当我们认识天主时，虽然我们比认识祂之前变得更好，尤其当这认识同样被喜爱并配得被爱时便成为言语，因而这认识成为天主的某种相似；然而这认识是较低级的，因为它存在于较低的本性中：因为心是受造物，天主是创造者。由此可推知，当心认识并认同自身时，这认识如此是其言语，以至于它完全与心相等、相同，因为既非较低实体的认识（如对物体的认识），也非较高实体的认识（如对天主的认识）。由于认识与它所认识者（即其认识的对象）相似，在此情况下，当被认识的正是认识的心本身时，这相似便完美而相等。因此，它既是肖像也是言语，因为它被说出关乎那心，并在认识中与那心相等；那受生的与那生者相等。\par

% structure: p0031 source=h2 target=subsection reason=relative-heading-rank; base=h2

\subsection{第十二章 为什么爱不是心的后代，而知识却是。问题的解答。心连同对自身的知识和对自身的爱乃是天主圣三的肖像}

17. 那么，什么是爱？它岂不是图像吗？岂不是话吗？岂不是受生的吗？因为为什么心智认识自身时就生它的知识，而它爱自身时却不生它的爱呢？因为如果它是自己认识的原因，理由在于它是可知的，那么它也是自己爱的原因，因为它是可爱的。那么，很难说为什么它不生两者。因为还有一个更进一步的问题，关乎至高的天主圣三本身，全能的天主造物主，人是按祂的肖像受造的，这问题困扰着人们，而天主真理藉着人的言语邀请人走向信德；\Emph{viz.}为什么圣神也不应当被相信或被理解为由天主圣父所生，以致于祂也可被称为圣子？这个问题我们正试图以某种方式在人的心智中探究，以便从一个较低的肖像——其中我们自己的本性本身，在被追问时，以我们更熟悉的方式回应——我们能够将更练习有素的心灵目光从受光照的受造物转向那不变的光；但要假定，真理自身已说服我们，正如没有基督徒怀疑天主的圣言是圣子一样，圣神是爱。那么，让我们回到对这肖像——即受造物，有理智的心智——的更仔细的追问和思考；其中，一些事物在时间中产生的知识（这些事物以前并不存在），以及对一些以前未被爱过的事物的爱，更清楚地向我们展开该说什么：因为对于言语本身（它必须按时间安排），那在时间顺序中被理解的事物更容易解释。\par

18. 首先，因此，显然一个事物可能是可知的，即能够被认识，然而它可能不被认识；但不可能一个不可知的事物被认识。所以，必须清楚地持守：凡我们所认识的任何事物，同时在我们内产生关于它自身的知识；因为知识是从两者——认识者和被认识的事物——生出来的。因此，当心智认识自身时，它独自是其自身知识的父母；因为它自身既是所知之物，又是它的认识者。但它在认识自身之前，对它自己已经是可知的，只是它自身内的关于自身的知识还不存在，只要它还不认识自身。那么，在认识自身时，它生出与自身相等的关于自身的知识；因为它并不认识自身为比自身更小，它的知识也不是关于他人本质的知识，不仅因为它自己认识，而且因为它认识自身，如同我们上文所说。那么，关于爱该说什么？为什么当心智爱自身时，它不似乎也生出了对自身的爱？因为它在爱自身之前，对自己已经是可爱的，因为它能够爱自身；正如它在认识自身之前，对自己已经是可知的，因为它能够认识自身。如果它对自己不是可知的，它永远不能认识自身；同样，如果它对自己不是可爱的，它永远不能爱自身。那么，为什么不能说它通过爱自身而生出了自己的爱，如同通过认识自身而生出了自己的知识？是不是因为这样确实清楚表明这是爱的本原，爱从它发出？因为爱从心智自身发出，心智在爱自身之前对自身已经是可爱的，所以它是自己借以爱自身的爱的本原；但因此不能说这爱是由心智所生，如同由心智认识自身的知识那样，因为在知识的情况下，事物已经找到，我们称之为产生或发现；而这通常先有一个探求，在达到目的时才安息。因为探求是寻找的欲望，或者同样，是发现的欲望。但那些被发现的事物可以说是被生出来的，因此它们像后代；但在何处？除非在知识本身的情形中。因为在那情形中，它们可以说是被说出和形成的。因为虽然我们寻找而找到的事物已经存在，但关于它们的知识并不存在，这知识我们视为所生的后代。此外，在探求中的欲望\OriginalTerm{欲望}{appetitus}从探求者发出，并以某种方式悬而未决，除非所寻找的事物被找到并与探求者结合，它才在它所指向的终点安息。而这欲望，即探求——虽然它似乎不是爱（被认识的事物借爱而被爱，因为在此我们仍在努力认识）——然而它是同类的东西。因为它可以被称为意志\OriginalTerm{意志}{voluntas}，因为凡寻找的人都愿意找到；如果所寻找的属于知识，凡寻找的人都愿意认识。但如果他热切而诚恳地愿意，他就被称为钻研\OriginalTerm{钻研}{studere}：这个词最常用于追求和获得任何学问的情形。因此，心智的生出先有某种欲望，通过它，经由寻找和找到我们愿意认识的东西，后代，\Emph{viz.}知识本身，诞生了。为此，那使知识受孕并生出的欲望，不能正确地被称为生出和后代；而那引导我们渴望认识事物的同一欲望，当事物被认识时，就变成了对已认识的事物的爱，因为它持守并拥抱所接受的后代，即知识，并将其与它的生者结合。于是，在心智自身、它的知识（即关于自身的后代和言）以及作为第三位的爱之中，有一种天主圣三的肖像；这三者是一，同一实体。后代并不更小，因为心智按自己存在的尺度认识自身；爱也不更小，因为它按自身知识和自身存在的尺度爱自身。\par

\SourceRecord{Translated by Arthur West Haddan.；Nicene and Post-Nicene Fathers, First Series；Vol. 3.；Edited by Philip Schaff.；Buffalo, NY: Christian Literature Publishing Co.,；1887.；<http://www.newadvent.org/fathers/130109.htm>.}

% structure: p0001 source=h1 target=section reason=page-title

\section{\WorkTitle{论天主圣三（卷十）}}

\EditorialNote{本章显示在人灵中有另一种\OriginalTerm{天主圣三}{trinity}，且更为明显，即在于记忆、理智与意志。}\par

% structure: p0003 source=h2 target=subsection reason=relative-heading-rank; base=h2

\subsection{第一章——好学心智之爱，即求知者之爱，并非对所不知之物的爱。}

1. 现在让我们依序以更精确的目的继续前行，更详尽地解释这同一要点。首先，既然无人能爱一个完全无知之物，我们便须仔细考察那些好学之人的爱究竟是怎样的——即那些尚未知晓、却仍渴求学习某门学问之人的爱。当然，在那些不常用\Quote{研习}一词的事物上，爱常因听闻而起，当某物之美的声誉激起心灵去观看并享受它时，心灵便从频繁的视见中普遍知道形体之美的内涵，并且内里存在着一种认可外在所渴望之物的能力。当此事发生时，所产生的爱并非完全未知之物，因其种类已知。但当我们爱一个从未谋面的好人时，我们是从对其德行的认识而爱他，而这些德行我们是在真理本身中（抽象地）认识的。但在学问之事上，通常是他人的权威——他们称赞并推荐这学问——点燃了我们对它的爱；尽管如此，若我们心中不曾对每类学问的知识至少有一丝印象，我们便不可能对这学问燃起丝毫热忱。因为，谁会费心劳力去学习修辞学，除非他事先知道这是一门关于言说的科学？有时，我们又因听闻或经验到的学问成果本身而惊叹，从而因学习而燃起获取这些成果能力的渴望。正如有人不识字，有人对他说：有一种学问，能让人用手在静默中书写话语，传给千里之外的人，而收信的人不是用耳，而是用眼来领会这些话语；若此人亲眼目睹此事发生，难道他不会因渴望知道如何做到这事，而全然被推动去学习，以达到他已知并掌握的那个成果吗？正是如此，学习者的好学热忱被点燃：因为一个人若全然无知，他绝不可能爱它。\par

2. 同样，如果有人听到一个未知的符号，例如某个他不知道含义的词的发音，他渴望知道它是什么；也就是说，他渴望知道通过那个发音约定要想起的是什么事物：就像他听到 \Emph{temetum} 这个词被说出，不知道，就问它是什么。那么他必须已经知道它是一个符号，即这个词不是空洞的声音，而是由它意指某物；因为在其他方面，这个三音节词他已经知道，并且已经通过听觉将其清晰的形态印入他的心灵。那么在他身上还需要什么，使他能够进一步认识那个已经知道其所有字母和各个音节的间隔的事物，除非与此同时他知道它是一个符号，并且也激起了他渴望知道它是何物的符号？因此，事物越是已知，却又未完全知道，心灵就越渴望知道关于它尚未知道的部分。因为如果他只知道它是一个这样的口语词，而不知道它是某物的符号，他就不会再寻求什么，因为可感的事物已尽感官所能被感知。但因为他知道它不仅是一个口语词，而且是一个符号，他渴望完全知道它；而没有符号被完全知道，除非知道它是何物的符号。那么，那个以热忱的谨慎寻求知道这一点，并被研习的热忱点燃而坚持探索的人，能说他没有爱吗？那么他爱什么？因为当然，没有东西能被爱除非它被知道。因为那个人并不爱他已经知道的三个音节。但如果他爱它们中的这一点，即他知道它们意指某物，这不是现在的问题，因为他寻求知道的并非这个。但我们现在要问的是，他爱的是什么，在他渴望知道但确实还不知道的事物中；因此我们奇怪他为什么爱，因为我们确知没有东西能被爱除非它被知道。那么他爱什么，除了他知道并在事物的理性中察觉学习有何等卓越，其中包含所有符号的知识；以及精通这些有何益处，因为通过它们人类团体互相交流自己的感知，免得人类的集会实际上比完全孤独更糟糕，如果他们不通过交谈交流思想？那么，灵魂辨明这种适宜和有用的形式，知道它，并爱它；而寻求任何他不懂的词语含义的人，努力尽可能将那种形式完善于自身：因为在一件事是在真理之光中看到它，另一件事是渴望它作为自己能力之内的事。因为他在真理之光中看到理解和说万民的各种语言，从而不听到任何语言也不被任何人作为外人听到，是何等伟大和美好的事情。那么，这种知识的美已经通过思想被辨明，并且已知的事物被爱；那事物被如此看待，并如此激发学习者的研习热忱，以至于他们被它感动，并热切地渴望它，在他们为获得这种能力而付出的所有辛劳中，以便他们也能在实践中拥抱他们先前通过理性知道的事物。这样，每个人越接近希望中的那种能力，就越热切地以爱渴望它；因为那些学科被更热切地学习，人们不绝望于能够达到的；因为当任何人对达到目的不抱希望时，他要么爱得冷淡，要么根本不爱，无论他多么看到它的卓越。因此，因为所有语言的知识几乎普遍被认为无望，每个人都最努力学习自己民族的语言；但如果他觉得自己甚至不能完全理解它，也没有人在这知识中如此懒惰，以至于不希望在听到一个未知词时，想知道它是什么，并寻求和学习它（如果他能）。而当他寻求时，当然他有一种学习的研习热忱，并且似乎爱一个他不知道的事物；但实际情况并非如此。因为那种形式触动心灵，心灵知道并思考它，其中明显可见从通过已知和说出的词语使心灵彼此联合而来的适宜性。这种形式点燃研习热忱，在他寻求确实不知道的事物时，但他凝视并爱那未知的形式，而该事物属于那形式。那么，例如，如果有人问：什么是 \Emph{temetum}（因为我已经引用了这个词），并对他说：这与你何干？他会回答：恐怕我听到某人说话，却不理解他；或者也许在某处读到了这个词，却不知道作者的意思。请问，有谁会对这样一个询问者说：不要在乎理解你所听到的；不要在乎知道你所读到的？因为几乎每个理性的灵魂都能迅速辨明那种知识的美，通过这种美，人们的思想借助有意义的词语的发音而相互交流；正是由于这种已知并因此被爱的适宜性，这样一个未知词被热切地寻求。那么，当他听到并得知 \Quote{\Emph{temetum}} 是我们祖先用来指称酒，但这个词在我们现在的语言用法中已经过时，他可能会认为由于祖先的这本或那本书，他仍然需要这个词。但如果他也认为这些书是多余的，那么他可能现在认为这个词不值得记住，因为他看到它与那种他用心灵知道、凝视并爱的学问形式无关。\par

3. 因此，在一切情况下，一个勤学头脑的爱——即一个渴望知道其所不知道之事的人的爱——并不是对所不知道之事物的爱，而是对它所知道之事物的爱；因着这已知之物，它渴望知道其所不知道的。或者，如果它有如此强烈的好奇心，以至于被冲走，并不是因为它所知道的任何其他原因，而仅仅出于对未知之事的爱；那么，这样的好奇之人无疑与普通学生有别，但也不比后者更爱他所不知道的事物；恰恰相反，更恰当地说，他憎恨他所不知道的事物，因为他希望不再有未知之事，即渴望知道一切。但唯恐有人提出一个更困难的问题，宣称一个人不可能憎恨他所不知道的事物，正如不可能爱他所不知道的事物一样，我们不会抗拒真理；但必须明白，说\Quote{他爱知道未知之事}与说\Quote{他爱未知之事}并不相同。因为一个人可能爱知道未知之事，但不可能爱未知之事。因为\Quote{知道}一词并非无意义地放在那里；因为爱知道未知之事的人，爱的不是未知之事本身，而是对它们的知道。除非他知道\Quote{知道}意味着什么，否则没有人能肯定地说他知道或不知道。因为不仅说\Quote{我知道}并且说得真实的人必须知道知道是什么；而且说\Quote{我不知道}并且说得自信而真实、且知道他说得真实的人，当然也知道知道是什么；因为当他审视自己并真实地说\Quote{我不知道}时，他区分了不知道的人和知道的人；而且既然他知道他说这话是真实的，如果不知道知道是什么，他又怎能知道呢？\par

% structure: p0007 source=h2 target=subsection reason=relative-heading-rank; base=h2

\subsection{第二章——无人爱未知之事。}

4. 因此，没有勤学者，没有好奇者，爱他所不知道的事物，即使他正以最强烈的渴望急切地想知道他所不知道的东西。因为他要么已经普遍地知道他爱什么，并且也渴望在一些个别事物中知道它，这些事物也许被称赞，但他还不知道；他在脑海中描绘一个想象的形式，藉此他被激发去爱。那么，他从哪里描绘这个，除非从他已经知道的事物？然而，如果他发现被称赞的形式与他心中所想象并最充分知道的另一形式不同，他也许不会爱它。如果他爱了，他将从他知道的那一刻开始爱它；因为稍早之前，所爱的那个形式与那形成它的头脑习惯向自己展示的形式是不同的。但如果他发现它类似于传言所宣称的形式，并且如此相似，以至于他可以说\Quote{我早已爱上了你}；那么，即使那时，他也没有爱一个他不知道的形式，因为他已经在那个相似性中知道了它。或者，我们在永恒理性的原型中看见某些东西，并在此中爱它；当这在一个暂时之物的形象中显现出来，我们相信那些试用过它的人的赞美，并如此爱它，那么根据我们上面已经充分讨论过的，我们并没有爱任何未知的东西。或者，我们爱某个已知之物，并且因着它寻求某个未知之物；所以，绝不是对未知之物的爱支配我们，而是对已知之物的爱，我们知道未知之物属于该已知之物，以便我们也知道那仍然未知的我们所寻求的东西；正如我先前关于一个未知的词语所说的。或者，每个人都爱知道本身，凡渴望知道任何事物的人都不能不知道这一点。由于这些原因，那些想知道他们所不知道的任何事物的人，似乎爱未知之物，并且由于他们强烈的探究渴望，不能说他们没有爱。但事实究竟如何不同，以及没有任何东西能被爱而不被知道，我想我已经说服了每一个仔细审视真理的人。但是，由于我们所给的例子属于那些渴望知道他们自身以外的某物的人，我们必须考虑，当心灵渴望认识自己时，也许会出现某种新的概念。\par

% structure: p0009 source=h2 target=subsection reason=relative-heading-rank; base=h2

\subsection{第三章——当心灵爱自己时，它并非对自己未知。}

5. 那么，理智热切寻求认识自己时，所爱的究竟是什么呢？当时它对自己尚未认识。因为，看啊，理智寻求认识自己，并以求知的热情激发自己。因此，它爱；但爱的是什么呢？是它自己吗？但是，当它还不认识自己时，这怎么可能呢？没有人能爱他所不认识的东西。难道是由于某种传闻向它宣告了自己的形相，如同我们通常听说不在场的人那样吗？也许，它爱的不是自己，而是它所想象出的自己的样子，这也许与自己本身大不相同；或者，如果理智中的幻想与理智本身相似，那么当它爱这个幻想的形象时，它就在认识自己之前爱了自己，因为它凝视着与自己相似的东西；于是，它认识其他理智，以便从中描绘自己，从而在种属上认识了自己。那么，为什么当它认识其他理智时，却不认识自己呢？毕竟没有任何东西比它自己更临在于它。但是，如果正如身体的眼目更了解其他眼目，而不是了解它们自身，那么让它不要再寻求自己，因为它永远找不到自己。因为眼目除非在镜子里，否则永远不能看见自己；并且绝不能以任何方式假定，这类东西也能应用于对无形之物的默观，以致理智仿佛在镜子里认识自己。或者，它是否在永恒真理的理性中看到认识自己是多么美好，因此爱它所见的，并努力在自身中实现它？因为，虽然它自身对自己是未知的，但它知道认识自己是多么好。这确实非常奇妙：它尚未认识自己，却已知道认识自己是多么卓越的事。或者，它是否通过某种隐秘的记忆——这记忆并未抛弃它，尽管它已远行——看到某个最卓越的目的，\Emph{即}它自身的宁静与真福，并相信除非认识自己，否则无法达到该目的？因此，当它爱那目的时，它就寻求这（认识）；它爱那已知的（目的），为了那未知的（认识）而寻求。但是，为什么对自身真福的记忆能够持续，而对自身的记忆却不能同样持续，从而使它既能认识它所渴望达到的（自己），又能认识它所要达到的（真福）？或者，当它爱认识自己时，它所爱的不是它自身（它尚未认识），而是认识的活动本身；并更烦恼于它自身对它想要涵盖万物的知识有所欠缺？它知道什么是认识；而当它爱这已知的东西时，它也渴望认识自己。那么，如果它不认识自己，它又怎能认识自己的认识呢？因为它知道自己认识其他事物，但不知道自己；正是由此它也知道了什么是认识。那么，一个不认识自己的东西，怎能知道自己在认识任何事物呢？因为它不是说某个其他理智在认识，而是它自己在认识。因此，它认识自己。此外，当它寻求认识自己时，它当下就知道自己在寻求。因此，它又认识自己。所以它不可能完全不知道，因为它确实知道自己在不知道。但是，如果它不知道自己在不知道，那么它就不会寻求认识自己。因此，就在它寻求自己的这一事实中，它显然被证明对自己是已知多于未知。因为它知道自己在寻求，并且知道自己尚未认识自己，正因为它寻求认识自己。\par

% structure: p0011 source=h2 target=subsection reason=relative-heading-rank; base=h2

\subsection{第四章 理智如何认识自己，不是部分地，而是整体地。}

6. 那么我们该说什么呢？那部分地认识自己的，是否不认识自己的另一部分？但说它不整体地认识它所认识的，这是荒谬的。我不是说它认识全体；但凡是它所认识的，它就整体地认识。因此，当它认识关于自身的任何事物——这些事物它只能整体地认识——时，它便整体地认识自身。而它确实知道自己认识某物，然而除非整体地，它不能认识任何事物。所以它整体地认识自身。此外，它里面有什么比它活着更被自己所认识呢？它不能同时是心灵而不活着，而它还具有某种超越之物，\Emph{即}理解力；因为野兽的灵魂也活着，却不理解。因此，正如心灵是整个心灵，它也整体地活着。而它知道自己活着。所以它整体地认识自身。最后，当心灵寻求认识自身时，它已经知道自己是一个心灵；否则它不知道是否在寻求自身，也许在寻求一物而意在寻求另一物。因为可能发生这样的事：它本身不是心灵，因此在寻求认识心灵时，它并没有寻求认识自身。因此，既然心灵在寻求认识什么是心灵时，知道自己是在寻求自身，那么它当然知道自己是一个心灵。此外，如果它在自身内知道这一点——即它是一个心灵，并且是完整的心灵——那么它便整体地认识自身。但假设它不知道自己是心灵，只是在寻求自身时只知道自己在寻求自身。因为这样，它也可能以它物代替此物，如果它不知道这一点；但为了不以它物代替此物，它无疑知道自己在寻求什么。但如果它知道自己在寻求什么，并且寻求的是自身，那么它当然认识自身。那么它还寻求什么呢？但如果它部分地认识自身，却仍然部分地寻求自身，那么它所寻求的便不是自身，而是自身的一部分。因为我们说到心灵本身时，是把它当作整体来说的。此外，因为它知道它尚未被自身整体找到，它知道整体有多大。于是它寻求所缺的部分，正如我们惯于回忆那从心灵中溜走、但并未完全离开心灵的事物；因为当它回来时，我们能认出它就是我们所寻求的那同一事物。但心灵如何能进入心灵，就好像心灵可能不在心灵之中一样？此外，如果它找到了部分，而不整体地寻求自身，然而它整体地寻求自身。因此，它作为整体呈现在自身面前，没有什么需要寻求的了；因为所缺的是那被寻求的，而不是寻求的心灵。既然它整体地寻求自身，它便一无所缺。或者，如果它不是整体地寻求自身，而是已被找到的部分寻求尚未找到的部分，那么心灵并不寻求自身，因为没有部分寻求自身。因为已被找到的部分并不寻求自身；同样，尚未找到的部分本身也不寻求自身，因为它是由已被找到的部分所寻求的。因此，既然心灵既不整体地寻求自身，也没有任何部分寻求自身，那么心灵根本就不寻求自身。\par

% structure: p0013 source=h2 target=subsection reason=relative-heading-rank; base=h2

\subsection{第五章 为何命灵魂认识自身。关于其自身本体的心灵错误从何而来。}

7. 那么，为何命令它（灵魂）认识自身呢？我想，是为了让它思考自身，并依照自己的本性生活；也就是说，寻求按照自己的本性被规范，\Emph{即}在祂之下——它应当服从祂，并高于那些事物——它应当被置于它们之上；在祂之下由祂治理，并高于那些它应当治理的事物。因为它借着邪恶的欲望做出许多事，好像忘记了自己。因为它看到某些内在本善的事物，在更卓越的天主性体中；而当它应当坚定不移地享受它们时，却转离了祂，由于想要把这些事物归于自己，并想不靠祂的恩赐而像祂，而是靠自身成为祂所是的，于是开始运动并逐渐滑向越来越低的事物，它认为这些是越来越高的事物；因为它不为自己所满足，也没有任何事物能使它满足，若它离开那唯一充足者的话：因此，由于匮乏和痛苦，它变得过于专注自己的行为以及通过它们获得的不安逸乐；就这样，由于渴望从外在事物中获取知识——它认识并爱慕这些事物的本性，并且感到若不焦虑地紧握就会失去它们——它失去了自己的安稳，并且越觉得自己不会失去自己，就越少想到自己。因此，虽然不认识自己和不想自己是两回事（因为我们不说一个博学多才的人不懂文法，当他只是没有去想它，因为此时他在想医学技艺）——那么，我说，虽然不认识自己和不想自己是两回事，爱的力量如此之大，以至于心灵将那些它长久怀着爱意思索、并通过勤勉研究的紧密依附而长入其中的事物，即使在某种程度上返回思考自身时，也随自身带入。因为这些事物是它通过肉性感官在外部所爱的有形之物；并且由于某种日常的熟悉而与它们纠缠在一起，却又不能将这些有形事物本身如同带入无形本性的领域那样带入内部；因此它组合了它们的某些形象，并将这样由自身形成的形象推入自身。因为它赋予形成这些形象以自身的某些实质，但同时保留下某种东西，借此它可以自由地判断那些形象的形式；而这东西更恰当地说是心灵，即理性理智，它被保留以供判断。因为我们看到，灵魂中那些被有形之物的相似所赋予形式的官能，也是我们与野兽共有的。\par

% structure: p0015 source=h2 target=subsection reason=relative-heading-rank; base=h2

\subsection{第六章 心灵对自身的意见是欺骗性的。}

然而，当心神如此深情而亲密地将自己与这些形象连接起来，以至于甚至认为自己也是同类之物时，它便错了。因为它在一定程度上被塑造成它们，不是由于它就是如此，而是由于它认为自己是如此：并非它认为自己是一个形象，而是直接认为它就是那个怀有形象的事物本身。因为仍有能力分辨它留在外面的有形事物，以及它从其中包含在自身内的那个有形事物的形象：除非当这些形象被投射出来，仿佛是在外部被感觉而非在内部被思考，就像睡着、疯狂或出神的人那样。\par

% structure: p0017 source=h2 target=subsection reason=relative-heading-rank; base=h2

\subsection{第七章。——论哲学家关于灵魂实体的意见。认为灵魂是形体的人之错误并非出于对灵魂认识不足，而是由于他们给它添加了外物。所谓\Quote{发现}的含义。}

因此，当它认为自己是这类事物时，它便认为自己是形体的；并且由于它完全意识到自己的优越性，凭此优越性它管理身体，于是便产生了这样的问题：身体的哪一部分在身体中拥有更大的能力；有人便认为这一部分就是心神，甚至就是整个灵魂。因此有人认为是血液，有人认为是大脑，有人认为是心脏；并非如圣经所说：\Quote{我要全心赞美你，上主；}以及\Quote{你要全心爱你的上主天主；}——因为\Quote{心}这个词是通过误用或比喻从身体转移到灵魂的；而是他们干脆认为它就是那当我们撕裂内里时所见到的身体的一小部分。另一些人则认为灵魂是由极微小而独立的微粒（他们称之为原子）组合凝聚而成。还有一些人说它的实体是气，另一些人说是火。另一些人则认为它根本不是实体，因为他们除非是形体，否则不能设想任何实体，而他们发现灵魂不是形体；但在他们看来，灵魂是我们身体的某种配合，或元素的结合，肉体仿佛由此联系在一起。因此所有这些人都认为灵魂是有死的；因为无论它是形体，还是形体的某种组合，当然都不能常存不死。然而，那些认为它的实体是与形体相反的生命类型的人，由于他们发现它是一种使每个活体具有生命和活力的生命，便继而尽力各自证明它是不死的，因为生命不能没有生命。\par

至于有人在世界四元素之外加了第五种我不知道的什么形体，并说灵魂是由此构成，我认为在此无需花费时间讨论。因为他们要么把形体理解为和我们一样的，\Emph{viz.}其一部分在空间的延展中小于全体，他们应被算作相信心神是形体的人；要么他们称一切实体或一切可变的实体为形体，但他们知道并非所有实体都通过长、宽、高被包含在空间的延展中，我们无需因措辞问题与他们争论。\par

现在，对于所有这些意见，任何人只要看到心神的本质既是实体而又非形体——即它不以自己的较小部分占据较小的空间延展，以较大部分占据较大的空间——就必然同时看到，那些认为它是形体的人之所以犯错，并非由于对心神认识不足，而是因为他们将一些没有它们就无法设想任何实体的性质附加于它。因为如果你命令他们设想没有形体幻象的存在，他们便认为那纯粹是虚无。于是，心神就不会寻找自己，仿佛它缺少自己。因为有什么比那呈现于心神的事物更呈现于认识呢？或者有什么比心神本身更呈现于心神呢？因此，所谓\Quote{发现}，如果我们考察这个词的来源，除了指\Quote{进入}所寻之物外，还有什么意思呢？那些仿佛自行进入心神的事物，通常不被称为\Quote{发现}，尽管可以说被知道；因为我们并未通过寻求而努力\Quote{进入}它们，即发明或发现它们。因此，正如心神本身真正地寻找那些由眼睛或身体任何其他感官所寻求的事物（因为心神甚至引导着肉体的感觉，然后当那感觉到达所寻求的事物时，它便发现或发明）；同样，它也不借助肉体感觉，而是通过自身来发现或发明它应当知道的其他事物，当它\Quote{进入}它们时；无论这发生在天主内的更高实体中，还是发生在灵魂的其他部分；正如它在判断身体形象本身时所做的那样，因为它发现这些形象在灵魂内，通过身体被印入。\par

% structure: p0021 source=h2 target=subsection reason=relative-heading-rank; base=h2

\subsection{第八章。——灵魂如何探究自身。灵魂关于自身的错误源自何处。}

11. 这实在是一个令人惊异的问题：灵魂以什么方式寻找并找到自身？它寻找时以什么为目标，或者它来到何处，以至它能达到或发现？因为有什么比心灵自身更存在于心灵之中呢？但由于它因爱而\Emph{在}于那些它所思念的事物中，又惯常因爱而存在于可感觉的、即有形的事物中，它就不能在没有那些有形之物的影像的情况下存在于自身之内。于是，可耻的错误起来阻挡它的道路，因为它不能将可感觉之物的影像从自身分离，以便单看自己。因为这些影像藉着爱的紧密附着奇妙地与它粘合了。它的不洁就在这里；因为当它努力单独思想自己时，它竟把自己想象成为那个没有它就不能思想自己的东西。因此，当它被命令要认识自己时，不要像脱离了自身那样去寻找自己，而要撤去它加添给自己的东西。因为自身更深地居于内在，不仅比那些显然在外面的可感觉之物更深，也比它们的影像更深；这些影像确实存在于灵魂的某一部分中，\OriginalTerm{即}{viz.}，动物也拥有这一部分，虽然它们缺乏理解，而理解是心灵所特有的。因此，既然心灵是内在的，它就在某种程度上从自身走出，当它对那些如同许多注意行为的足迹般的东西施加爱的情感时。而这些足迹，在外部有形之物被感知的时候，就仿佛印在记忆上，甚至当那些有形之物不在时，它们的影像仍为思想它们的人所获得。所以，让心灵认识自己，不要像缺失了那样去寻找自己；而要固定于自身那[自愿的]注意的行为——它曾藉此漫游于其他事物之中——并让它思想自己。这样它就会看到，它从未在任何时候不爱自己，从未在任何时候不认识自己；但因它与自己一同爱了另一个东西，它便把自己与那东西混淆了，在某种意义上与它合而为一了。于是，当它拥抱不同的事物仿佛它们是一体时，它竟认为那些不同的事物是同一的了。\par

% structure: p0023 source=h2 target=subsection reason=relative-heading-rank; base=h2

\subsection{第九章。——心灵藉着理解认识自己的命令这一行为本身而认识自己。}

12. 因此，不要让它像缺失了那样去分辨自己，而要努力分辨自身为临在。也不要让它好像不知道自己在认识自己，而要它把自己与它所知道是另一个的东西区分开来。因为如果它不知道 \Quote{认识} 是什么意思，也不知道 \Quote{你自己} 是什么意思，它如何能努力遵从那个给予它的命令 \Quote{认识你自己} 呢？但如果它两者都知道，那么它也就知道自己了。因为 \Quote{认识你自己} 对心灵说的话，不像 \Quote{认识革鲁宾和色辣芬}；因为它们是缺席的，我们相信它们，并根据那种信仰被宣告为某种天上的大能者。也不同于说：认识那人的意志；因为这一点我们完全无法通过感官或理解来感知，除非通过实际展示的有形记号；并且这样我们便更相信而非理解。又不同于对一个人说：看你的脸；这他只能通过镜子才能做到。因为我们自己的脸本身也是无法看到的；因为它不在我们目光可及之处。但当对心灵说 \Quote{认识你自己} 时，它就在理解 \Quote{你自己} 这个词的同一行为中认识了自己；这没有别的原因，只因为它是临在于自身的。但若它不理解所说的，那么它当然就不会做所命令的事。因此，它正是被命令去做那件事，而这件事正是它在理解那命令本身时所做的。\par

% structure: p0025 source=h2 target=subsection reason=relative-heading-rank; base=h2

\subsection{第十章。——每一个心灵都确切知道关于自身的三件事：它理解，它存在，它生活。}

13. 那么，当它被命令认识自己时，切不可向它所知道的自我的存在添加任何东西。因为它至少知道，这话是对它自己说的；即对那个存在、活着、且理解的自我说的。但一个死尸也存在，牲畜也活着；但死尸和牲畜都不理解。因此，它如此知道自己如此存在、如此活着，正如一个理解力存在着、活着一样。所以，举例来说，当心灵以为自己是一团空气时，它以为那空气是理解的；然而它知道自己理解，但它并不知道自己是一团空气，只是这样以为罢了。让它将那它所自以为的东西分开；让它辨别出它所知道的东西；让那个东西留给它，关于这一点，即使那些曾以为心灵是这或那有形之物的人也未曾怀疑过。因为当然，并非每一个心灵都把自己当作空气；有些却以为自己是一团火，有些是脑髓，有些是某一种有形之物，有些是另一种，如同我之前提到的那样；然而所有心灵都知道它们自己理解、存在并且活着；但它们把理解归之于它们所理解的东西，而把存在和活着归之于它们自己。没有人会怀疑，要么没有一个不理解的人活着，要么没有一个人活着却不能说他是存在的；因此，那个理解的东西既存在又活着；不像一个死尸存在却不活着，也不像一个灵魂活着却不理解，而是以某种更恰当、更卓越的方式。此外，它们知道自己愿意，它们同样知道没有人能够愿意，除非他存在并且活着；它们也将那意志本身归之于它们用那意志所愿意的某物。它们也知道自己记忆；并且它们同时知道，除非一个人存在且活着，否则他不可能记忆；但我们也把记忆本身归之于某物，因为我们记忆那些事物。因此，在许多事物中的知识和学问包含在这三者中的两者之中：记忆和理解；但意志必须参与，我们才能享受或使用它们。因为我们享受已知的事物，在这些事物本身中，意志因它们的缘故而找到愉悦，并安息其中；但我们使用那些我们用来指向我们要享受的另一事物的东西。人的生命之所以邪恶和应受责备，无他，只是由于错误地使用和错误地享受。但此处并非讨论这个问题的时候。\par

14. 但既然我们讨论心灵的本性，让我们从我们的考虑中除去所有从外部通过身体感官获得的知识；更仔细地注意我们已经提出的论点，即所有心灵都知道并对自己确知。因为人无疑曾怀疑过：活着、记忆、理解、愿意、思考、知道、判断的能力，是属于空气、火、脑髓、血液、原子，还是除了通常的四种元素之外的我不知道的第五种物体；或者我们这肉身的组合或调和本身是否有能力成就这些事。有人试图证实这一点，另一个人试图证实那一点。然而，谁曾怀疑过他自己活着、记忆、理解、愿意、思考、知道、判断呢？因为即使他怀疑，他活着；如果他怀疑，他记起为什么怀疑；如果他怀疑，他理解他怀疑；如果他怀疑，他愿意确定；如果他怀疑，他思考；如果他怀疑，他知道他不知道；如果他怀疑，他判断他不该轻率地同意。因此，凡在其他事情上怀疑的人，不应怀疑这一切；因为如果不是这样，他就不可能怀疑任何事。\par

15. 那些认为心灵或是身体、或是身体的组合或调和的人，会认为所有这些都存在于一个主体中，所以实体是空气、火或其他他们认为是心灵的有形之物；而\OriginalTerm{理解}{intelligentia}作为其性质存在于这个有形之物中，所以这个有形之物是主体，而理解在主体中：\Emph{viz.}心灵是主体，他们认为心灵是有形之物，而理解 [intelligence]，或任何其他那些我们提到作为我们确知的事物，是在那个主体中。那些否认心灵本身是身体、但认为它是身体的组合或调和的人也持有大致相同的意见；因为区别在于：前者说心灵本身是实体，理解如同在主体中那样存在于其中；后者说心灵本身是在一个主体中，即在身体中，而心灵是身体的组合或调和。因此，结果他们还能怎么想，除非认为理解也如同在主体中那样存在于同一个身体中？\par

16. 所有这些人并未察觉到，心灵即使在寻求自身时，也认识自身，正如我们已经指出的。但任何东西，若其本体未被认识，就绝不能说是被正确认识的。因此，当心灵认识自身时，它认识自己的本体；而当它对自己确定时，它对自身本体也是确定的。但它对自己是确定的，正如上面所说那些事物令人信服地证明的那样；尽管它根本不确定自身是空气、是火、是某个物体，还是身体的某种机能。因此，它不是其中任何一个。而那个被命令认识自身的整体，就拥有这一点：它确定自己不是那些它不确定的东西中的任何一个，并且确定自己仅仅是那个它唯一确知它之所是的。因为它是这样思考火、空气以及它所想到的任何其他身体之物。它绝不可能以思考自己所不是之物的方式来思考自己之所是。因为它通过一种想象幻想来思考所有这些东西，无论是火、空气、这个或那个物体，或是身体的某一部分或组合与调和；而且，它当然不能说就是所有这些，而只是其中之一。但如果它是其中任何一个，它就会以不同于其余的方式思考这一个，即：不是通过想象幻想，如同缺席之物被思考那样（那些物体本身或同类之物曾为身体感官所触及），而是通过某种内在的、非虚构的、真实的临在（因为没有什么比它自身更为临在）；正如它思考自身活着、记忆着、理解着、意愿着。因为它在自身内知道这些东西，并不想象它们如同通过外在感官触及它们一样，如同有形物体被触及。如果它从对这些东西的思考中没有附加任何东西给自己，以致于认为自己就是那一类东西，那么，任何从自身留给它的，就是它自身。\par

% structure: p0030 source=h2 target=subsection reason=relative-heading-rank; base=h2

\subsection{第十一章——在记忆、理解（或理智）与意志中，我们需注意能力、学识与使用。记忆、理解与意志在本质上是一，在关系上是三。}

17. 因此，暂且将心灵对自身确定的所有其他事物放在一边，让我们特别考虑并讨论这三者——记忆、理解、意志。因为我们通常也能从这三者中辨别出青年人能力的特点；因为一个男孩记忆越牢固、越容易，理解越敏锐，学习越热切，他就越值得在能力上受到赞扬。但当谈到一个人的学识时，我们所问的不是他记忆得多么牢固和容易，或理解得多么精明；而是他记住了什么，理解了什么。而且，因为心灵不仅被视为学识可嘉，也被视为善良可嘉，所以人不仅要注意他记住了什么、理解了什么，还要注意他意愿什么（\Emph{velit}）；不是他意愿得多么热切，而是首先他意愿什么，然后他意愿得多大程度。因为热切地爱的心灵，当它爱那应当热切地被爱的事物时，才当受赞扬。既然我们谈论这三者——能力、知识、使用——那么第一项要从三个方面考虑：一个人在记忆、理解和意志上能做什么。第二项要考虑的是，任何人通过勤勉的意志在其记忆和理解中所获得的东西。但第三项，即\Emph{viz}使用，在于意志，它处理那些包含在记忆和理解中的东西，无论它们是引向进一步的目的，还是满足于作为目的本身。因为使用是将某物纳入意志的权能；而享用是以喜乐使用，不再是基于希望，而是基于实际事物。因此，每一个享用者都使用；因为他将某物纳入意志的权能，并以此为目的而满足。但并非每一个使用者都享用，如果他寻求纳入意志权能的那物，不是因它本身，而是因其他某物。\par

18. 既然，这三者，记忆、理解、意志，不是三个生命，而是一个生命；也不是三个心智，而是一个心智；那么，它们当然也不是三个实体，而是一个实体。因为记忆，被称为生命、心智和实体，是就自身而言的称呼；但被称为记忆，则是相对于某物而言的。同样，我也要说理解和意志，因为它们被称为理解和意志是相对于某物而言的；但各自就自身而言，是生命、心智和本质。因此这三者是一，因为它们是同一生命、同一心智、同一本质；而它们各自就自身而言所称呼的任何其他名称，也一起称呼，不是复数，而是单数。但它们是三者，在于它们彼此相互关联；而如果它们不平等，不仅彼此之间不平等，而且每个与所有之间不平等，那么它们肯定不能相互包含；因为不仅每一个被每一个包含，而且全部被每一个包含。因为我记得我有记忆、理解和意志；我理解我理解、意志和记忆；我意愿我愿意、记忆和理解；我同时记得我整个的记忆、理解和意志。因为我记忆中我不记得的部分，不在我记忆里；没有什么比记忆本身更在记忆之中。因此我记得整个记忆。同样，凡我所理解的，我知道我理解；凡我所意愿的，我知道我意愿；但凡我所知道的，我记住。因此我记得我全部的理解和全部的意志。同样，当我理解这三者时，我一同理解它们作为整体。因为凡可理解的事物中，除了我不知道的，没有我不理解的；但凡我不知道的，我既不记住也不意愿。因此，凡可理解之物中我不理解的，随之我也既不记住也不意愿。凡可理解之物中我记住并意愿的，随之我也理解。我的意志也包含我全部的理解和全部的记忆，当我使用我理解和记住的整体时。因此，当所有都被彼此相互包含，并作为整体，每一个作为整体等于每一个作为整体，同时每一个作为整体等于所有作为整体；这三者是一，一个生命，一个心智，一个本质。\par

% structure: p0033 source=h2 target=subsection reason=relative-heading-rank; base=h2

\subsection{第十二章——心智在其自己的记忆、理解和意志中是\OriginalTerm{天主圣三}{Trinitas}的肖像。}

19. 那么，我们现在是否应当尽我们所能的决心向上提升，达到那至首要至高的本质，人的心智是它不完全的肖像，却仍是肖像？或者，这些同一的三者是否应当藉由我们从外面通过身体感官所接受的事物，更清晰地显明在灵魂内，在其中，对有形事物的知识在时间中被印在我们身上？因为我们发现心智本身在其自己的记忆、理解和意志中是这样：由于它被认为始终认识自身并始终意愿自身，它也被认为同时始终记住自身、始终理解和爱自身，虽然并不始终思想自己为\Emph{分离的}于那些不是它自身的事物；因此，它对自身的记忆和对自身的理解，在其中难以辨别。因为在这种情况下，这两者非常紧密地结合，一个丝毫不在时间上先于另一个，看起来仿佛不是两样事物，而是一样事物有两个名称；而当需要的感觉不揭示它时，爱本身并不那么明显地被感到存在，因为所爱之物总是近在眼前。因此，这些事物可以更清晰地阐述，甚至对悟性较迟钝的人，如果处理那些在时间中被带入心智触及范围、并在时间中发生在心智上的主题；当它记住它以前没记住的，看见它以前没看见的，爱它以前没爱过的。但鉴于本书的篇幅，这一讨论现在需要另一个开始。\par

\SourceRecord{Translated by Arthur West Haddan.；Nicene and Post-Nicene Fathers, First Series；Vol. 3.；Edited by Philip Schaff.；Buffalo, NY: Christian Literature Publishing Co.,；1887.；<http://www.newadvent.org/fathers/130110.htm>.}

% structure: p0001 source=h1 target=section reason=page-title

\section{论天主圣三（第十一卷）}

\EditorialNote{甚至在外在的人身上，也指出了一种天主圣三的肖像；首先，在那些从外部感知的事物中，即在被看见的形体对象中，以及由它印在观看者视觉上的形式中，以及结合二者的意志的目的中；虽然这三者既不相等，也非同一实体。其次，在心灵本身之中，观察到一种同一实体的三个某种东西构成的三一，仿佛是从外部感知的事物引入那里的；即存在于记忆中的形体对象的肖像，以及当思想者的心灵之眼转向它时由此形成的印象，以及结合二者的意志的目的。这后一种三一也被说成属于外在的人，因为它是从外部感知的形体对象引入心灵的。}\par

% structure: p0003 source=h2 target=subsection reason=relative-heading-rank; base=h2

\subsection{第一章——在外在人身上的天主圣三踪迹}

1. 无人怀疑，正如内在的人赋有理智，外在的人也赋有身体的感觉。那么，让我们尝试，若我们能，也在外在的人身上发现天主圣三的些微踪迹，但并非它本身也同样是天主的肖像。因为宗徒的意见是显而易见的，他声称\Emph{内在}的人在认识天主上，按创造他者的肖像得以更新；而他也在另一处说：\Quote{但我们\Emph{外在}的人虽然渐渐腐朽，但内在的人却日日更新。}那么，让我们在那腐朽的事物上，尽可能寻找天主圣三的肖像，即使不那么明显，却也许更容易辨认。因为那外在的人也不是无缘无故被称为人，而是因为它有几分内在之人的相似。由于我们被造成为必死和肉身的这一状况，我们处理可见事物比处理可理解之物更为容易和熟悉；因为前者是外在的，后者是内在的；前者由身体感觉感知，后者由心智理解；而我们自己，\Emph{即}我们的心智，不是可感觉的事物，即身体，而是可理解的事物，因为我们是生命。然而，如我所说，我们如此熟悉地关注身体，我们的思想以如此惊人的倾向向外投射到身体上，以至于当它从有形事物的不确定性中退出来，以便以更确定和稳定的知识专注于精神时，它却飞回那些身体，并从那里寻求安息，而正是从那里它汲取了软弱。我们必须使我们的论述适应这种软弱；因此，如果我们有时试图更恰当地区分和更便捷地表达内在的属灵事物，我们必须从属于身体的外在事物中取用相似的例子。那么，赋有身体感觉的外在之人，便与身体打交道。这种身体感觉，容易观察到，有五种：视、听、嗅、味、触。但就我们所寻求的而言，探究所有这五种感觉既相当麻烦，也无必要。因为其中一种感觉向我们宣告的，也适用于其余。那么，让我们主要使用眼睛的见证。因为这身体的感觉远超过其他；并且按着它种类上的差异，更接近心灵的视觉。\par

% structure: p0005 source=h2 target=subsection reason=relative-heading-rank; base=h2

\subsection{第二章——视觉中的某种三一。证明在视觉中有三件事，按其本性不同。可见事物如何产生视觉，或被见之物的肖像。此事藉一例子更清楚地显示。这三者如何结合为一。}

2. 当我们看见任何有形物体时，最容易考虑和区别的就是这三件事：首先，我们看见的物体本身；无论是一块石头、火焰，还是任何其他可由眼睛看见的东西；而且这个物体在被看见之前可能早已存在；其次，视觉或看的行为，它在我们感知到呈现于感官的物体本身之前并不存在；第三，那使眼目感官停留在所看见的物体上、只要它被看见就保持着的，即\Emph{viz}精神的注意。在这三者中，不仅有明显的区别，而且有不同性质。因为首先，那可见的物体与眼目感官的性质大不相同，由于感官作用于它而产生视觉。显然，视觉本身不就是被那被感知的事物所成形的知觉吗？虽然如果可见对象撤去，就没有视觉，而且如果没有可见的物体，也根本不可能有任何这种视觉；然而，当那物体被看见时，使眼目感官成形的物体，以及由它印在感官上的称之为视觉的形式本身，绝不是同一实体。因为那被看见的物体就其自身本性而言是可分离的；但感官，它甚至在看见它所能看见的东西之前就已经存在于活着的主体中，当它遇到某种可见之物时——或者当可见物体现在与它联结而被看见时，从可见物体在感官中产生的视觉——那么，我说，感官或视觉，即从外部被成形的感官，属于活着的主体之性质，这性质完全不同于我们凭视觉所感知的物体，并且感官之所以被成形为视觉而非感官，是由于那物体所致。因为在可感对象呈现在我们面前之前，如果我们里面没有感官，那么当我们有时什么也看不见时（无论是在黑暗中还是闭着眼睛时），我们就与瞎子没有区别。但我们与他们的区别在于，即使当我们看不见时，我们里面也有那种使我们能够看见的东西，称为感官；而他们里面没有这个，他们被称为瞎子没有别的原因，就是因为他们没有它。此外，那使感官停留在我们所看见的事物上并将两者结合在一起的精神注意，不仅在性质上与可见物体不同（因为一个是精神，另一个是身体），而且也与感官和视觉本身不同：因为这种注意是精神单独的活动；而眼目感官之所以被称为身体的感官，没有别的原因，只是因为眼睛本身也是身体的肢体；虽然无生命的身体不能感知，但灵魂与身体混合，通过一个物质的工具感知，那个工具被称为感官。而且，当一个人失明时，这感官也因身体的痛苦而被切断和熄灭；而精神保持同一；它的注意由于失去眼睛，固然没有可以借看见与外在身体结合的肉体感官，从而定睛于它并看见它，但通过这努力本身表明，即使身体的感官被移除，它本身既不会消灭也不会减少。因为无论能否实现，仍然保留着一个完好的看的\OriginalTerm{欲望}{appetitus}。那么，这三者：被看见的物体、视觉本身、以及将两者结合在一起的精神注意，不仅因各自的属性，而且因性质的差异，明显可以区分。\par

3. 既然在这种情况下，感觉并非来自所看见的那个身体，而是来自知觉着的活体——灵魂以一种其自身奇妙的方式与之融合——的那个身体；然而，视觉的产生，即官能本身被成形，是由所见的身体造成的；这样，现在不仅有官能的能力（即使眼睛健全在黑暗中也能完好存在），而且有实际被成形的官能，称为视觉。因此，视觉是由可看见的事物产生的；但并非仅由此事物，除非也有观看者在场。所以，视觉是由可看见的事物与观看者共同产生的；其方式为，在观看者方面，有看的官能和观看、凝视对象的意向；然而，那官能的成形（称为视觉）却仅由所见的身体，即可见的事物印刻而成；此事物一旦移除，那在官能中当所见之物在场时存在的形式便不再存留：然而官能本身仍然存在，它甚至在感知任何事物之前就已存在；正如水中的影像，只要印在其上的身体本身还在水中，影像就存在；但若这身体被移除，便不再有任何这样的影像，尽管水仍然存在，它在获得那个身体的形状之前就已存在。因此，我们确实不能说可见事物产生官能；但它产生形式，这形式仿佛是它自身的相似，当我们在观看中感知任何事物时，这相似就在官能中形成。但我们并不通过同一官能区分我们所见身体的形式与它在观看者官能中产生的形式；因为这二者结合得如此紧密，以至于没有区分的余地。但我们理性地推断，除非所看见的某种相似被印刻在我们自己的官能中，否则我们根本不可能有感觉。因为当一枚戒指印在蜡上时，并不能因为没有图像分离出来我们就认为没有产生图像。但既然在蜡被分离后，所造之物仍然留存，以至于可以被看见；因此我们很容易确信，在蜡与戒指分离之前，蜡中就已经有戒指所印的形式了。但如果戒指印在流体上，当戒指移开时就不会出现任何图像；然而尽管如此，理性仍应辨别出，在流体中，在戒指移开之前，已经有一个从戒指产生的戒指形式，它与戒指本身的形式有所区别，后者在戒指移开后仍然存在，而前者则停止存在。所以，不可因为眼睛的（感官）知觉在对象移开后图像不存留，就认为它在看见对象之时不包含所见的身体形象。因此，很难说服头脑较迟钝的人：当我们看见一个可见事物时，它的形象在我们官能中形成，而这个形式就是视觉。\par

4. 但若有人留意我将要提及的事，他们就会发现这项探究并无如此困难。通常，当我们看了一小会儿光然后闭上眼睛时，眼前似乎浮现出某些明亮的色彩，它们变化多端，逐渐暗淡直至完全消失；我们必须理解，这些是当那发光体被看见时在官能中形成的形式的残余，并且这些变化发生在它们缓慢逐步褪去的过程中。因为如果我们碰巧凝视窗格，这些色彩中也常出现窗格的形状；由此可见，我们的官能受到来自被见之物的这些印象的影响。因此，这种形式在我们看见之时就已存在，并且当时更为清晰明确。但它与所感知之物的形式紧密结合，以至于完全无法区分；而这就是视觉本身。甚至，当油灯的小火焰因眼睛发散的光线而某种方式加倍时，就会出现双重视觉，尽管所见之物只是一个。因为同一光线，各自从每只眼睛射出，各自受到影响，因为它们在注视那个物体时被不允许均匀地会合，从而形成两眼的单一联合视觉。所以，如果我们闭上一只眼睛，我们将不会看到两个火焰，而是一个真实的火焰。但为何，如果我们闭上左眼，右边出现的那景象就消失；反之，闭上右眼，左边的那景象就不复存在，这本身既冗长，对当前主题也无必要去探究和讨论。因为对于手头的工作，只需考虑这一点：除非在我们官能中产生一个与我们感知之物完全相似的影像，火焰的外观就不会根据眼睛的数量加倍；因为采用了一种特定的感知方式，能够分离光线的会合。当然，真正单一的事物不可能被一只眼睛看成双倍，无论你如何下拉、按压或扭曲它，如果另一只眼是闭着的话。\par

5. 既然如此，让我们记住这三者（尽管本质不同）如何被调和成一种统一：即所见之\OriginalTerm{形体形式}{forma corporis}，印在感官上的其\OriginalTerm{印象}{imago}（即\OriginalTerm{视觉或成形感官}{visio seu sensus formatus}），以及将感官应用于可感之物并保持视觉于其中的\OriginalTerm{心灵意志}{voluntas mentis}。这三者中的第一个，即可见之物本身，并不属于生物的本性，除非当我们辨识自己的身体时。但第二个属于该本性到如此程度：它是在身体内、并通过身体而在灵魂内完成的；因为它是在感官内完成的，而感官既不能脱离身体也不能脱离灵魂。但第三个唯独属于灵魂，因为是意志。虽然这三者的实体如此不同，但它们融合成一种统一，以至于前两者几乎无法区分，即使有理性的介入作为判断者，即所见之形体形式和感官中形成的其印象，也就是视觉。意志如此有力地结合此二者：既将感官应用于被感知之物以使其成形，又在其成形后保持于该物中。如果这意志如此强烈，可被称为\OriginalTerm{爱、欲望或贪恋}{amor, desiderium, libido}，它也会强烈地影响生物身体的其余部分；并且在较迟钝较坚硬的物质不抵抗之处，将其改变为相似的形状和颜色。可以看到变色龙的小身体随其所见的颜色而迅速变化。至于其他动物，由于肉体的粗笨不易变化，后代大多显露母亲的特殊幻想，即她们特别欣喜地看着的任何东西。因为最初的种子越柔嫩，可以说越易成形，就越有效、越有能力追随母亲灵魂的倾向和通过那身体在其中形成的幻想，即她贪婪注视过的幻想。可以举出许多例子，但一个就够了，取自最可信的书籍：viz。雅各伯所做的，他放置杂色的枝条在水槽里让绵羊和山羊在看的时候喝水，使它们产下各种颜色的后代，当时它们已经怀胎。\par

% structure: p0010 source=h2 target=subsection reason=relative-heading-rank; base=h2

\subsection{第三章——三者合一发生在思想中，即记忆、内在视觉和意志结合两者。}

6. 然而，\OriginalTerm{理性灵魂}{anima rationalis}按照\OriginalTerm{外在人的三一}{trinitas exterioris hominis}生活时，是生活在\OriginalTerm{堕落方式}{degeneratus modus}中；即当它把心灵意志应用于那些从外部形成\OriginalTerm{肉体感官}{sensus corporis}的事物时，不是可赞美的意志（由此将其导向某种有用目的），而是卑劣的欲望（由此依附于它们）。因为即使形体形式（被身体感知到的）被撤去，其\OriginalTerm{相似性}{similitudo}仍留在\OriginalTerm{记忆}{memoria}中，意志可再次将眼目转向它，以便从内部由此成形，正如感官曾由外部被可感身体的呈现所成形。因此，那三一是由记忆、\OriginalTerm{内在视觉}{visio interna}和结合两者的\OriginalTerm{意志}{voluntas}产生的。当这三者结合为一，从该结合本身它们被称为\OriginalTerm{构思}{conceptio}。在此三者中不再有实体的多样性。因为可感身体不在那里（它完全不同于生物的本性），肉体感官也不是成形的以产生视觉，意志本身也没有执行其将感官（尚未成形）应用于可感身体并在成形后保持于其中的职务；但取代从外部被感知的形体形象，记忆出现了，它保留着灵魂通过肉体感官所吸收的那个形象；取代那当感官通过可感身体被成形时的外在视觉，内在出现了一个类似的视觉，同时心灵之眼从记忆所保留的事物成形，而被思想的身体事物不在场；意志本身则和以前一样，曾将尚未成形的感官应用于从外部呈现的身体事物，并在成形后与之结合，现在就转而将回忆心灵之视觉指向记忆，以便心灵视觉可由记忆所保留者成形，从而在构思中出现一个类似的视觉。正如理性将可见外观（肉体感官由此而成形）与其相似性（在感官成形时被造成以产生视觉）区分开（否则它们会如此结合，以至被视为全然同一）；同样，虽然那\OriginalTerm{心象}{phantasia}也来自心灵思想所见过的身体外观，由记忆保留的身体相似性以及从中在回忆心灵之眼中形成的东西组成，但它看起来是单一的，只有通过理性的判断才能发现是二者，正是这判断使我们明白，那留在记忆中的——即使我们从别的来源想到它——不同于当我们回忆（即再次回到记忆并在那里找到同一外观）时所产生的东西。如果这（留在记忆中的事物）现在不存在，我们会说我们完全忘记了而无法回忆。如果回忆者的眼目不是从记忆中的那个事物成形，思想者的视觉就根本无法发生；但两者的结合——即记忆所保留者与由此表达出来以成形回忆者眼目的事物——使它们显得如同一，因为极其相似。但当构思者的眼目从那里转开，停止注视记忆中所感知的东西时，印在其上的形式将什么也不留在那眼目中，它将由重新转向的对象成形，从而产生另一个构思。然而，它在记忆中所留下的东西仍存在，当我们回忆时，它可再次转向，并因转向而被其成形，并与成形之源成为一体。\par

% structure: p0012 source=h2 target=subsection reason=relative-heading-rank; base=h2

\subsection{第四章　—　这种合一如何实现。}

7. 然而，如果那驱动并引导将被成形的眼睛来来回回、并在成形时使之结合的意志，全然转向内在的幻象，并且完全将心智之眼从环绕感官的物体之临在以及从感官本身移开，而完全转向那在内中所觉察的影像，那么，从记忆中表达出来的形体样式就会呈现得如此精确，以致连理性本身也无法分辨所看见的是外在的物体本身，还是仅仅是内心中某种类似的东西。因为人们有时因过度思考可见事物而被诱惑或惊吓，甚至会突然相应地发出言语，仿佛真的置身于这些行为或苦难之中。我记得有人曾告诉我，他在思想中如此清晰且仿佛实在地感知到一个女性身体的形体，以致被感动，好像那是真实的一样。灵魂对自己的身体有如此大的力量，并且对转动和改变它（有形）衣裳的特质有如此大的影响，正如一个人穿着衣服时可以被影响，衣服贴在他身上。我们在睡眠中被想象所欺骗时，也有同样的感受。但有很大的区别：是身体的感官被麻痹而迟钝（如睡眠者的情况），还是被其内在结构扰乱（如疯人的情况），或者以其他方式分散（如占卜者或先知的情况）；因此出于这些原因之一，心智的意向被迫以某种必然性转向那些浮现的影像——它们要么来自记忆，要么通过某种隐秘的力量，经由类似精神实体的某种精神混合而产生；或者，如同健康清醒的人有时会发生的那样，被思想占据的意志转向远离感官，从而以各种可感事物的影像来形成心智之眼，仿佛那些可感事物本身真的被感知到一样。这些影像的印象不仅发生在意志因欲望而指向这些事物时，也发生在心智为了\Emph{逃避}和防范它们而被牵引去看这些事物本身时。因此，不仅是欲望，还有恐惧，都使得身体之眼被可感事物本身所形成，也使\OriginalTerm{心智之眼}{acies}被这些可感事物的影像所形成。相应地，恐惧或欲望越强烈，眼睛就越清晰地成形——无论是对于通过身体感知邻近物体的人，还是对于从记忆中存留的物体影像进行构想的人。因此，物体在位置中所是对身体感官的，就是物体的相似物在记忆中所是对心智之眼的；观看事物之人的视像对所由以形成感官的物体外观的，就是构想者的视像对所建立于记忆中的物体影像的——心智之眼由此形成；意志的意向对所看见的物体和与之结合的视像（以使其中三者虽本性不同却达成某种合一）的，就是同样的意志的意向对记忆中的物体影像和构想者的视像（即心智之眼在返回记忆时所采取的形式）的结合，以使这里也达成三者的某种合一——它们现在不再因本性的不同而被区分，而是属于同一实体，因为这一切都在内，而整个就是同一个心智。\par

% structure: p0014 source=h2 target=subsection reason=relative-heading-rank; base=h2

\subsection{第五章　—　外在之人或外在视觉的三位一体不是天主的肖像。即使在罪中也渴望天主的相似。在外在视觉中，有形之物的形式如同父母，视像如同子女；而结合这两者的意志则暗示圣神。}

8. 但是，正如当[同时]身体的形相和物种已经消失时，意志无法将感知的感觉召回；同样，当记忆所承载的影像被遗忘抹去时，意志将无法通过回忆迫使心灵之眼返回，从而被其形成。但因为心灵有巨大的能力去想象不仅遗忘的事物，而且从未见过或经历过的事物，或通过增加、减少、改变、组合，随心所欲地想象那些尚未从记忆中消失的事物，它常常想象事物如此这般，要么它知道它们不是，要么不知道它们是。在这种情况下，我们必须小心，免得它或说谎以欺骗，或持有意见以致被欺骗。如果它避免了这两个恶，那么想象出的幻象就不会妨碍它：正如通过感官经历或被记忆保留的可感事物，如果既不因愉悦而热烈追求，也不因不悦而卑鄙逃避，就不会妨碍它。但当意志离开更好的事物，贪婪沉迷于这些时，它就变得不洁；它们被有害地思想，当它们在场时，更有害地当它们缺席时。因此，按照\Emph{外在}人的三位格生活的人活得恶劣和堕落；因为使用可感和有形事物的目的产生了那个三位格，尽管它在内部想象，却想象外部的事物。因为没有人能够很好地使用这些事物，除非感官所感知的事物的影像被保留在记忆中。除非意志大部分驻留在更高和内在的事物中，并且除非那意志本身（它适应于外部身体或内部事物的影像）将其中所接受的一切指向一个更好和更真实的生命，并安息于那个目的，通过注视它来判断那些事应该做；否则我们做了什么，除了宗徒禁止我们做的，当他说：\Quote{不要与此世同化}？因此，那个三位格不是天主的肖像，因为它通过身体的感官从最低的、即有形的受造物（心灵高于它）而在心灵本身中产生。然而它并非完全不相似：因为有什么事物没有按它的种类和尺度有天主的相似性呢？既然天主造了万物都很好，且别无原因，只因祂自己是至善？所以，凡存在之物，只要是善的，就显然还具有至善的某种相似性，无论距离多么遥远；如果是自然的相似性，那么肯定是正确而有序的；但若是有缺陷的相似性，那么肯定是堕落而乖张的。因为即使是灵魂在它们的罪恶中，也追求别的什么，只是某种天主的相似性，以一种骄傲和颠倒的、可说是奴隶般的自由。所以我们的原祖父母也不会被说服犯罪，除非说了：\Quote{你们将成为像神一样}。无疑，受造物中任何在某种程度上相似天主的，也不可称为祂的肖像；只有那比祂自身更高的，才是祂的肖像。因为只有那从祂那里完全复制而来，且与祂之间没有其他本性介入的，才是祂的肖像。\par

9. 论到那视觉；即，在观看者的感官中形成的形相；产生它的有形事物的形相，就好比是父母。但它不是真正的父母；因此那也不是真正的子女；因为它并非完全从它而生，因为为了从它形成，还有别的东西应用于有形事物，即观看者的感官。因此，爱这个就是疏远。因此，使两者结合的意志，\Emph{即}准父母和准子女，比两者都更属灵。因为那被辨识的有形事物根本不是属灵的。但在感官中产生的视觉，混合了一些属灵的东西，因为它不能没有灵魂而产生。但它不完全属灵；因为被形成的是身体的感官。所以，使两者结合的意志诚然更属灵，如我所说；因此它开始\OriginalTerm{暗示}{insinuare}，就好比是，在天主圣三中圣神的位格。但它更属于被形成的感官，而不是产生它的有形事物。因为活物的感官和意志属于灵魂，不属于石头或其他被看见的有形事物。因此它并非从有形事物如同从父母而出；也不从那个作为准子女的另一者，即感官中的视觉和形相而出。因为意志在视觉发生之前就已存在，这意志将待形成的感官应用于待辨识的有形事物；但它尚未满足。因为尚未看见的东西如何能满足？满足意味着意志安息满足。因此，我们既不能称意志为视觉的准子女，因为它先于视觉存在；也不能称其为准父母，因为视觉并非从意志形成和表达，而是从被看见的有形事物。\par

% structure: p0017 source=h2 target=subsection reason=relative-heading-rank; base=h2

\subsection{第六章——论在视觉中我们应如何估量意志的 \OriginalTerm{安息}{Requies} 和 \OriginalTerm{终向}{Finis}。}

10. 或许我们可以正确地称之为视觉为意志的终止与安息，但仅针对这一个对象（即可见的物体）而言。因为意志不会仅仅因为看见了它现在想要的东西就不再想要别的东西。因此，这并不是人整个的意志本身——其终止无非是幸福——而是暂时指向这一对象的意志，其终止在于看见本身，无论是否将所见之物指向其他事物。因为如果它不将视觉指向更远之处，而只愿意看见这个，那么毫无疑问，视觉是意志的终止；这是显而易见的。但如果它指向更远之处，那么它肯定还想要别的东西，就不再仅仅是看见的意志；或者如果是看见的意志，也不再是看见这个特定事物的意志。正如有人希望看见伤疤，从而知道曾有过创伤；或者希望看见窗户，从而透过窗户看见行人：所有这些以及类似的其他意志行为，都有它们自己适当的\OriginalTerm{直接}{proximate}终止，而这些终止又被指向意志的\OriginalTerm{最终}{final}终止——我们借此意愿幸福生活，并达到那种不指向任何其他事物、对喜爱它的人而言自足的生命。因此，看见的意志以视觉为终止；看见这个特定事物的意志以看见这个特定事物为终止。所以，看见伤疤的意志渴望它自己的终止，即看见伤疤，而不超过它；因为证明曾有创伤的意志是另一个不同的意志，尽管依赖于前者，其终止也是证明曾有创伤。同样，看见窗户的意志以看见窗户为终止；因为还有一个依赖于它的更进一步意志，即透过窗户看见行人，其终止也是看见行人。但所有彼此联结的各个意志，如果它们所指向的那个意志是善的，那么它们就都是正确的；如果那个是恶的，那么它们就都是恶的。因此，正确意志的连续系列就像一条由某些阶梯组成的路，借此攀登到幸福；而堕落歪曲的意志的纠缠则是一条锁链，如此行动的人将被它捆绑，以致被丢到外面的黑暗中去。所以，那些在行为和品格中歌唱上升阶梯之歌的人是有福的；而那些如拉长绳般拖罪孽的人有祸了。同样，说意志处于安息中——我们称之为它的终止——如果它仍指向更远之处，就好像我们说脚在行走中不动，当它放在那里，却可以从那里再向前迈出一步，沿着人的脚步方向。但如果某物如此令人满足，以至于意志带着某种愉悦安息于其中，那么它仍然不是人最终趋向的目标；而是同样指向更远之处，从而被视为不是公民的祖国，而是旅人休息甚至停留的地方。\par

% structure: p0019 source=h2 target=subsection reason=relative-heading-rank; base=h2

\subsection{第七章　在回忆所见之人的记忆中还有另一个三一}

11. 然而再看另一个三一，它确实比在可感事物和感官中的那个更为内在，但仍从中构思而来；而现在不再是身体的感觉由身体形成，而是心灵的眼由记忆形成，因为我们从外界感知的物体的\OriginalTerm{种相}{species}已经固着在记忆本身之中。我们把记忆中的那个种相称为准父母，它孕育出在构思者之\OriginalTerm{幻象}{phantasy}中形成的那个种相。因为在我们构思之前，它就在记忆中，正如在我们（感官）感知之前，物体也在位置中，以便视觉发生。但当它被构思时，就从记忆所保存的形式中，在构思者的心灵\OriginalTerm{视域}{acies}中通过回忆复制并形成了那个种相，它是记忆所保存之物的准后裔。但既不是真正的父母，也不是真正的后裔。因为当我们通过回忆思考任何事物时，从记忆形成的心灵视觉，并不从我们记得曾经看见的那个种相产生；因为我们不可能记得从未看见的事物；然而，由回忆形成的眼睛，在我们看见所记得的物体之前就已存在；因此，更不用说在我们将其存入记忆之前了。所以，虽然在回忆者心灵之眼中形成的种相是从记忆中的种相形成的，但心灵之眼本身并非由此产生，而是早已存在。由此可知，如果一方不是真正的父母，另一方也不是真正的后裔。但那个准父母和那个准后裔都暗示了某种东西，由此更内在和更真实的事物可以更实际、更确定地显现。\par

12. 此外，更难清晰地辨别：连接视觉与记忆的意志，是否不是其中某一方的父母或子女？而同一本性及实体的相似和平等造成了这种区分的困难。因为在此情况下，无法像处理由外形成的感官（它容易与可感觉的身体区分，而意志又与两者区分）那样，因为这三者之间本性的差别，我们已在上文充分讨论过。因为尽管我们目前谈论的这种三一是从外部进入心灵的，但它发生在内部，没有一部分是在心灵本身的本性之外。那么，如何能证明意志既非记忆中所含形体相似物的准父母，也非准子女，亦非回忆时从中复制的那个相似物的准父母或准子女？当意志在构想行为中将两者如此结合，使它们单独看来如同一体，只能靠理性辨别时，又如何证明？首先需要考虑：除非我们在记忆的深处保留了想要回忆之事的全部或部分，否则不可能有回忆的意志。因为对于一个完全被遗忘的事物，回忆的意志不可能产生；因为我们已记得想要回忆的事物是或曾经在我们的记忆中。例如，如果我想要回忆昨晚吃了什么，要么我已记得我吃了晚饭，要么即使还未记得这个，至少我记起了有关那个时间本身的某些事——如果没有别的，至少我记起了昨天，以及昨天人们通常用晚餐的那段时间，以及什么是晚餐。因为如果我根本没有记起任何这类事，我就不可能想要回忆昨晚吃了什么。由此可见，回忆的意志确实源自记忆中保留的事物，加上通过回忆在辨识行为中从中复制的事物；也就是说，来自我们已记得的某物与思考者在回忆时由此在心灵之眼中形成的视觉的结合。但连接两者的意志本身还需要某种其他事物，它似乎就在记忆者近旁。因此，有多少记忆就有多少这种三一；因为没有哪一个缺少这三样东西：\Emph{即}在思考之前已储存在记忆中的东西、在辨识时在构想中发生的事情，以及连接两者的意志，而意志从两者及自身作为第三者，完成一个单独的事物。或者，我们是否更应这样认识此类中的某个单一的三一：即我们可以一般地说，所有隐藏在记忆中的形体种类作为一个单一的整体，而回忆和构想此类事物的心灵的一般视觉作为另一个单一的整体，连接两者的意志作为第三者加入这两者的结合，使整个事物成为由三者构成的某个统一体？\par

% structure: p0022 source=h2 target=subsection reason=relative-heading-rank; base=h2

\subsection{第八章 不同的构想方式。}

但既然心灵之眼不能一眼看见记忆所保留的全部事物，这些思考的三一便以一系列的退隐与接替交替出现，因此那三一变得极其多数；然而，如果它不超出记忆中所储存事物的数量，便不是无限的。因为，虽然我们从任何人通过任何身体感官对物质事物所获得的最初知觉开始计算，甚至也包括那些已被遗忘的事物，但数量无疑仍是确定而有限的，尽管无法计数。因为我们不仅称无限的事物为不可数的，也称那些虽有限但超出任何人计算能力的事物为不可数的。\par

13. 但由此我们可以更清楚一点地觉察到：记忆所储存和保留的东西，与回忆者的构想中从中复制的东西是不同的，尽管当两者结合在一起时，它们看起来是同一个东西；因为我们只能记住我们实际见过那么多、那么大、那么样的形体；因为心灵从身体感官将它们吸收到记忆中；而构想中所见之物，虽从记忆中的事物提取，却无数倍地增多和变化，完全没有止境。因为我无疑只记得一个太阳，因为事实上我只见过一个；但若我愿意，我可以构想两个、三个、或任意多个；但我构想多个时的心灵视觉，是由我记住一个太阳的同一记忆所形成的。我记住它正如我见过的那样大。因为如果我记得它比我见过的大或小，那么我就不再记得我所见的，所以我不记得它。但由于我记得它，我便记得它如我所见的那样大；然而我可以按我的意志构想它更大或更小。我记住它如我所见；但我可以构想它按我的意愿运行，按我的意愿静止不动，按我的意愿来自某处、去往某处。因为我有能力构想它为方形，尽管我记得它是圆形；而且我可以构想它为我喜欢的颜色，尽管我从未见过绿色的太阳，因此也不记得它；对于太阳是这样，其他一切事物也是如此。但由于这些事物形体的物质和可感觉性，当心灵想象它们以它在内中所构想的同样方式存在于外部时，便会陷入错误——无论它们已经不再存在于外部，却仍保留在记忆中，还是以任何其他方式，我们所记忆的东西不是通过忠实的回忆，而是通过思想的变异性在心灵中形成的。\par

14. 然而，我们常常也会相信别人以真实叙事告诉我们的、他们自己通过感官所知觉到的事物。在这种情况下，当我们听到所叙述的事物时，心灵之眼似乎并没有回转至记忆，以便在我们的思想中引出种种景象；因为我们并非凭自己的回忆，而是凭别人的叙述来构想这些事物；由意志作为第三者将隐藏在记忆中的\OriginalTerm{种相}{species}和回忆者的视象结合而成的那个\OriginalTerm{三一结构}{trinity}，在此似乎并未完成。因为我所构想的并非隐藏在我记忆中的东西，而是我所听到的，当有人向我叙述某事时。我不是在说说话者的话语本身，以免有人以为我已转向那外在的、发生在可感事物或感官中的另一个三一结构；而是我在构想叙述者用言语和声音向我指明的那些物质事物的种相；这种相我当然不是凭回忆，而是凭听闻来构想。但若我们更仔细地思考，即使在这种情形下，记忆的界限也并未被逾越。因为若非我一般地记得他所说的个别事物，即便它们那时是第一次被串联在一个故事中，我也根本无法理解叙述者。举例来说，某人向我描述一座山被伐尽木材，披上橄榄树林，他是向记得山、木材和橄榄树林的种相的我描述；若我忘记了这些，我便完全不知道他在说什么，因此也就无法构想那个描述。于是乎，每个构想有形事物的人，无论他自己想象什么，或听到、读到对过去事物的叙述或对未来事物的预言，都求助于记忆，并在此中找到他思想中所注视的一切形式的界限和尺度。因为没有人能构想他从未见过的颜色或形状、从未听过的声音、从未尝过的滋味、从未嗅过的气味，或从未感受过的任何有形物的触感。但若没有人构想任何有形之物，除非他（通过感官）知觉过它，因为没有人记得任何有形之物，除非他这样知觉过，那么，就如知觉的界限在于物体，思想的界限也在于记忆。因为感官从我们所知觉的物体接收种相，记忆从感官接收种相；而构思者的心灵之眼从记忆接收种相。\par

15. 再者，正如意志将感官应用于身体对象，它也将记忆应用于感官，并将构思者的心灵之眼应用于记忆。但那个协调并联合这些事物的东西，本身也分离并分开它们，即意志。然而，意志通过身体的运动将身体感官与被知觉的物体分开，或者以阻止我们知觉该事物，或者以使我们停止知觉它：如当我们转移目光不看我们不愿看的东西，或闭上眼睛；同样地，耳朵远离声音，或鼻孔远离气味。同样我们也远离滋味，或合上嘴，或将口中的东西吐出。在触觉上，我们或移开身体之物，以免触碰我们所不愿触碰的，或若已经触碰，便将其抛掷或推开。如此，意志通过身体的运动，使身体感官不与可感事物结合。它按自己的能力行此事；因为当它因奴隶性死亡的状态而在此事中忍受艰难时，结果便是痛苦，以致意志除了忍耐别无他存。但意志使记忆远离感官；当意志专注于其他事物时，不允许眼前的事物附着于记忆。这正如任何人可以看到的，当我们因思考别的事而似乎没有听到有人对我们说话时。但这是误解；我们确实听到了，但我们不记得，因为说话者的话语通过意志的命令（意志被转向别处，而通常正是意志将话语固定于记忆）即时从耳朵的知觉中溜走了。因此，在这种情况下，更准确的说法应是：我们不记得，而非没有听到；因为甚至在阅读时，我自己也经常遇到这样的情况：当我读完一页或一封书信，我不知道我读了什么，便重新开始。由于意志的目的被固定于别处，记忆并未如感官本身应用于文字那样应用于身体感官。同样地，任何行走时意志专注于别处的人，不知道他走到了哪里；因为若他没有看见，他就不会在那里行走，或者会更加注意地摸索着行走，尤其若他正在经过一个不熟悉的地方；然而，因为他走得很轻松，他当然看见了；但因为记忆并未如眼睛的感官应用于他所经过的地方那样应用于感官本身，所以他连刚看到的东西也完全不能回忆。现在，意愿将心灵之眼从记忆中的东西转开，无非就是不去思想它。\par

% structure: p0027 source=h2 target=subsection reason=relative-heading-rank; base=h2

\subsection{第九章 种相由种相继生成。}

16. 在这种安排中，当我们从物体的形式出发，最终到达在理解者的\OriginalTerm{直观}{contuitu}中产生的形式时，我们发现四种形式可以说是逐步地一个从另一个出生：第二个从第一个，第三个从第二个，第四个从第三个；因为从物体本身的形式，产生了在感觉者的感觉中出现的形式；由此，产生了在记忆中出现的；再由此，产生了在理解者心灵之眼中出现的。因此，意志三次将父母与后代结合起来：首先，物体的形式与在身体感觉中由它产生的形式；其次，后者又与由之在记忆中产生的形式；第三，这个又与在理解者心灵的直观中由它产生的形式。但中间的组合，即第二种，虽然更接近第一种，却不如第三种那样与第一种相似。因为有两种视觉：一种是\OriginalTerm{感觉}{sentientis}的视觉，另一种是\OriginalTerm{思维}{cogitantis}的视觉。但为了产生思维的视觉，在记忆中，从感觉的视觉中制造出某种相似的东西，以便心灵之眼在思维时可以转向它，就像眼睛的目光（\OriginalTerm{锐利（目光）}{acies}）在感觉时转向物体本身一样。因此，我选择了在这一类中提出两个三位一体：一个是在感觉的视觉从物体中形成时，另一个是在思维的视觉从记忆中形成时。但我没有推荐中间的那个；因为当在感觉者感觉中出现的形象被托付给记忆时，我们通常不称它为视觉。然而，在所有情况中，意志只是作为父母与后代的结合者而出现；所以，无论它从何处出发，都不能被称为父母或后代。\par

% structure: p0029 source=h2 target=subsection reason=relative-heading-rank; base=h2

\subsection{第十章——想象甚至将所见之物添加到未见之物上}

17. 但是如果我们只记得我们所感觉到的，并且只设想我们所记得的；为什么我们经常会设想虚假的东西呢？当然，我们不会错误地记得我们所感觉过的东西，除非是因为那个意志（我早已尽力尽可能表明它是这类事物的结合者和分离者）按照它自己的意愿和乐趣，引导将要形成的设想者的视觉，通过记忆的隐秘库藏；并且为了设想我们不记得的东西，它驱使它从这里取一个东西，从那里取另一个东西，从我们记得的那些东西中；这些东西结合成一个视觉，制造出某种被称为虚假的东西，因为它要么不在外部物体中存在，要么似乎不是从记忆中复制来的，因为我们不记得曾经见过这样的东西。谁曾见过黑天鹅？因此没有人记得黑天鹅；然而谁能不设想它呢？因为很容易将我们在其他身体上见过的黑色应用于我们通过视觉认识的那种形状；因为我们见过两者，所以记得两者。我也不记得一只四条腿的鸟，因为我从未见过；但通过在我见过的某种有翼形状上添加另外两条腿（我也同样见过），我非常容易地凝视这样的幻想。因此，在联合设想我们记得单独见过的东西时，我们似乎不是在设想我们所记得的东西；而实际上我们是在记忆的规律下这样做，我们从记忆中取出一切，按照我们自己的喜好在多种多样的方式中结合起来。因为即使是我们从未见过的物体的巨大体积，我们也不能在没有记忆帮助的情况下设想；因为我们的视线通常通过世界的广袤达到的空间尺度，也是我们在设想物体时尽可能扩大其体积的尺度。然而，理性甚至走得更远，但幻想不跟随她；例如当理性宣告无限的数目时，根据有形之物设想的任何人的视觉都无法把握。同样的理性也教导我们，最小的原子可以无限分割；然而当我们到达我们记得见过的那些微小细小的粒子时，我们便不能再看到更细更小的幻象，尽管理性继续分割它们。因此，我们设想的任何有形之物，要么是我们记得的，要么是从我们记得的那些东西中产生的。\par

% structure: p0031 source=h2 target=subsection reason=relative-heading-rank; base=h2

\subsection{第十一章——数、重量、尺度}

18. 但因那些单独印在记忆中的事物，可以按数目被构想，故尺寸似乎属于记忆，而数目属于视觉；因为，虽然此类视觉的多样性不可胜数，但每个视觉在记忆中都被规定了一个不可逾越的限度。因此，尺寸出现在记忆中，数目出现在事物的视觉中：正如可见物体本身有某种尺寸，观看者的感官极其繁多地适应于该尺寸，且从一个可见对象形成了许多观者的视觉，以致于一个人因两只眼睛的数目，通常以双重外表观看同一事物，正如我们上文所设定的。因此，那些视觉所源自的事物中有某种尺寸，而在视觉本身中有数目。但那结合并调节这些事物，将它们组合成某种统一，并且不将其感知或构想的欲求安静地止息，除非在视觉所源自的事物中，这意志类似于重量。因此，我愿预先注意这三者：\OriginalTerm{尺寸}{measure}、\OriginalTerm{数目}{number}和\OriginalTerm{重量}{weight}，它们也应在其他一切事物中被感知。其间，我已尽可能并尽可能多地向人指明，意志是可见事物与视觉的联结者，如同父与子，无论是在感性知觉还是构想中，而它既不能被称为父，也不能被称为子。因此，时间提醒我们，要在内在的人中寻求这同一天主圣三，并努力从属魂的、属血肉的、如所称的外在人，即我长久以来所谈论的，向内推进。在此，我们希望能找到按天主圣三的天主肖像，祂亲自帮助我们的努力，正如事物本身所显示，也如圣经所见证的：祂以\OriginalTerm{尺寸}{measure}、\OriginalTerm{数目}{number}和\OriginalTerm{重量}{weight}处置一切。\par

\SourceRecord{Translated by Arthur West Haddan.；Nicene and Post-Nicene Fathers, First Series；Vol. 3.；Edited by Philip Schaff.；Buffalo, NY: Christian Literature Publishing Co.,；1887.；<http://www.newadvent.org/fathers/130111.htm>.}

% structure: p0001 source=h1 target=section reason=page-title

\section{论天主圣三（卷十二）}

\EditorialNote{开始区分智慧与知识，指出一种特殊的、属于被称为知识（较低者）的天主圣三；这种天主圣三虽属内在的人，却尚不能称为或视为天主的肖像。}\par

% structure: p0003 source=h2 target=subsection reason=relative-heading-rank; base=h2

\subsection{第一章——论外在的人与内在的人之性质。}

1. 来吧，现在让我们来看一看，外在的人与内在的人之间，仿佛有一条界线在哪里。因为凡是我们心智中与野兽共有的，如此之多都正当地属于外在的人。因为外在的人不应被认为仅仅是身体，还应加上身体的某种特定生命，身体的结构由此获得活力，以及装备来感知外在事物的所有感官；而当这些已感知的外在事物的形象被固定在记忆中，并通过回忆再次被看到时，这仍然属于外在的人。在所有这些事物中，我们与野兽并无不同，除了身体形态上我们不是俯伏的，而是直立的。并且藉此，那创造我们的那一位告诫我们：不要在我们更好的部分——即心智——上像野兽，而身体上的直立却使我们与它们不同。这并不是说我们要把心智投入到那些高耸的肉体事物中；因为为了意志的安息甚至投入这类事物，就是使心智俯伏。但是，正如身体自然地向上直立朝向那些最崇高的肉体事物，即天上的事物；同样，心智，作为一个精神实体，必须向上直立朝向那些精神事物中最崇高的，不是由于骄傲的昂首，而是由于义德的虔敬。\par

% structure: p0005 source=h2 target=subsection reason=relative-heading-rank; base=h2

\subsection{第二章——唯独人能从受造物中感知属于身体之物的永恒理由。}

2. 野兽也能通过身体感官从外部感知肉体事物，并将它们固定在记忆中，记住它们，并在其中寻求适合的、避开不便的。但是，注意这些事物并不仅作为自然摄取而且作为刻意交付给记忆来保留，并通过回忆与概念将它们重新铭刻——当它们正滑入遗忘时——以便正如概念由记忆所包含的内容形成，记忆本身的内容也可通过思想被牢固地固定；通过取用记忆所记住的这样那样的事物，并仿佛将它们彼此拼接，来重新组合想象的可视对象；审视在这一类中，似真之物如何与真物区分，而这不仅在精神事物中，而且在肉体事物本身中——这些行为以及类似的行为，虽然涉及可感事物以及心智通过身体感官所推出的事物，但由于它们结合了理性，因此不是人与野兽所共有的。但是，更高的理性部分是根据无形而永恒的理由来判断这些肉体事物的；这些理由除非高于人心，否则当然不会是永恒的；然而，除非我们自己的某些东西与它们相连，否则我们无法将它们作为判断肉体事物的尺度。但我们是从度量与形状的规则来判断肉体事物的，心智知道这些规则永恒不变。\par

% structure: p0007 source=h2 target=subsection reason=relative-heading-rank; base=h2

\subsection{第三章——属于默观的更高理性与属于行动的较低理性，同在一个心智之内。}

3. 但是我们自身中那与处理肉体与暂时之物有关的部份，确实是理性的，因为它不是人与野兽共有的；但它仿佛是自我们心智的理性实体中被抽出，我们藉此理性实体依赖并紧贴那可理解而不可改变的真理；而它被派来处理和指导较低级的事物。因为正如在所有的野兽中，没有找到与亚当相似的一位助手，除非从他自身取出一个，并形成为他的伴侣；同样，对于那藉以咨询超性与内在真理的心智，在灵魂的、我们与野兽共有的那些部分中，没有相似的一位助手供此类人事所需，以处理肉体之事。因此，我们的理性的某一部分——并非分离以致破坏统一，而是仿佛被分流出来以成为共融的助手——被划分出来执行它自己的适当工作。而且，正如在男女的情况下，二人是一体，同样，在心智中，一个本性包含了我们的理智与行动，或我们的谋略与执行，或我们的理性与理性欲望，或任何其他更重要的术语可以用来表达它们；所以，正如论及前者时所说的，\BibleQuote{他们二人将成为一体}，也可以论及这两个说：他们二人是在一个心智之内。\par

% structure: p0009 source=h2 target=subsection reason=relative-heading-rank; base=h2

\subsection{第四章——天主圣三与天主的肖像仅在于心智中属于默观永恒事物的那一部分。}

4. 因此，当我们讨论人心智的本性时，我们讨论的是一个单一的对象，并不把它加倍成为我提及的那两个，除非就其功能而言。因此，当我们在心智中寻求天主圣三时，我们在整个心智中寻求它，并不将理性在暂时事物中的行动与对永恒事物的默观分离开来，以致进一步寻求某个第三项，由此天主圣三得以完成。但这个天主圣三必然要在心智的整个本性中被发现，以至于即使对暂时事物的行动被撤销——为了那项工作需要那种助手，为此心智的某一部分被转移以处理这些较低级的事物——仍然会在那无处被分割的单一心智中发现一个天主圣三；并且当这种分配已经作出时，不仅可以在只属于默观永恒事物的那一部分中发现一个天主圣三，而且还可以发现天主的肖像；而在那被转移去处理暂时事物的另一部分中，虽然可能有一个天主圣三，却不能找到天主的肖像。\par

% structure: p0011 source=h2 target=subsection reason=relative-heading-rank; base=h2

\subsection{第五章——论在男女婚姻及其子女中构想天主圣三肖像的意见。}

5. 因此，在我看来，那些主张：就人的本性而言，可以在三个位格中发现天主肖像的三一性，并且这种三一性在男女的婚姻及他们的后裔中得以完成——即男人自己似乎指示了圣父的位格，而从男人生出之物指示了圣子的位格，而第三位，他们说，即女人，她从男人而出，却不是儿子或女儿，虽然后裔是因她的怀胎而生——这些人似乎没有提出一个合理的意见。因为主曾论及圣神说，祂由圣父所发，祂却不是儿子。在这个错误的意见中，唯一可能被提出、且确实按圣经的信德充分显示的是这一点——在关于女人被造的原始记载中——凡是从某人产生而成为另一人的，不能都称为儿子；因为女人的位格是从男人的位格产生的，她却未被称作他的女儿。这意见的其余部分实则是如此荒谬，甚至如此虚假，以致于极易驳斥。因为我不提这样的事：认为圣神是天主圣子的母亲，圣父的妻子；因为或许会有人回答说，这些事在肉体的事上冒犯我们，因为我们想到的是肉体的怀孕和生育。虽然这些事本身由纯洁的人以最纯洁的方式思想，对他们而言，一切都是纯洁的；但在污秽和不信的人（他们的心与良心都被玷污）看来，没有什么是纯洁的；以至于就连按肉体由童贞女所生的基督，对他们中的一些人也是绊脚石。然而，对于那些至高的属神之事（低级受造物的种类也按此相似被造，尽管极其遥远），那里没有任何可损伤的、可朽坏的，没有在时间内诞生的，没有从无定形之物形成的，或任何类似的说法；这些都不应扰乱任何人的清醒审慎，以免他在避免空洞的厌恶时落入有害的错误。让他习惯于在物质事物中发现属神事物的痕迹，以致当他开始从那里向上攀登，在理性引导下，为了到达创造这一切的不变真理本身时，他不会把在低级事物中所鄙视的东西也带到高级事物中。因为从没有人因选择智慧为妻而感到羞愧，因为\Quote{妻子}这个名称会让人想到生儿育女的朽坏结合；也因智慧本身在性别上是女人，因为在希腊文和拉丁文中，它都以阴性词来表达。\par

% structure: p0013 source=h2 target=subsection reason=relative-heading-rank; base=h2

\subsection{第六章——为何应拒绝此意见。}

6. 因此，我们不拒绝此意见，并非因为我们害怕将那圣洁、不可侵犯、不变的爱视为天主圣父的配偶——它由祂所出，但不是作为后裔以便生出那创造万有的圣言——而是因为圣经清楚地显示它是假的。因为天主说：\BibleQuote{让我们按我们的肖像，按我们的模样造人；}稍后又说：\BibleQuote{天主于是按天主的肖像造了人。}的确，由于是复数，\Quote{我们的}这个词如果不是人按一个位格（无论是圣父、圣子或圣神）的肖像而造，就不能正确使用；但因他是按天主圣三的肖像而造，故说\Quote{按我们的肖像}。然而，免得我们以为在天主圣三中应相信三个神，而同一天主圣三是一神，故说\BibleQuote{天主于是按天主的肖像造了人}，而不说\Quote{按祂自己的肖像}。\par

7. 这类表达在圣经中是常见的；然而有些人，虽持守天主教信仰，却未加留意，以致认为\Quote{天主按天主的肖像造了人}这句话，是指圣父按圣子的肖像造了人；他们因此想要断言，圣子在圣经中也被称为天主，仿佛没有其他最真实、最清楚的证据证明圣子不仅被称为天主，而且被称为真天主。因为他们旨在解释这经文中的另一个难点时，却陷入困境，无法自拔。因为如果圣父按圣子的肖像造了人，以致人不是圣父的肖像，而是圣子的肖像，那么圣子就不像圣父。但若虔诚的信德教导我们（事实上它确实如此教导），圣子在本质的平等上像圣父，那么按圣子的肖像所造之物，也必然是按圣父的肖像所造。此外，如果圣父不是按自己的肖像，而是按祂圣子的肖像造了人，那么为什么祂不说：\Quote{让我们按祢的肖像和模样造人}，而实际上祂说的是\Quote{我们的}呢？除非是因为天主圣三的肖像被造在人身上，如此人就成了唯一真天主的肖像，因为天主圣三本身就是唯一真天主。圣经中这类表达数不胜数，但引用这些就够了。在圣咏中这样说：\Quote{救恩属于上主；祢的祝福临于祢的百姓}，仿佛这话是对别人说的，而不是对那曾说过\Quote{救恩属于上主}的那一位。此外，祂说：\Quote{因为靠着祢，我将脱离试探，并因寄望于我的天主，我将跳过墙垣}，仿佛是对别人说：\Quote{靠着祢我将脱离试探。}又说：\Quote{在王的心肠中【的仇敌】，万民因此仆倒在祢以下}，仿佛是说：在祢仇敌的心肠中。因为祂曾对那位王，即我们的主耶稣基督说：\Quote{万民仆倒在祢以下}，而祂以\Quote{王}这个词所指的就是祂，当祂说\Quote{在王的心肠中的仇敌}时。这类事在新约中较为罕见。然而宗徒致罗马人说：\Quote{关于祂的圣子，按血肉来说，是从达味的后裔生的，按圣德的神，藉着我们的主耶稣基督从死者中复活，被立为具有大能的天主子}；仿佛他上文是在谈论另一个人。因为\Quote{天主子藉耶稣基督的死者复活被显明}是什么意思呢？难道不正是同一耶稣基督被显明为具有大能的天主子吗？因此，正如在这段经文中，当我们读到\Quote{耶稣基督的大能的天主子}，或\Quote{耶稣基督的按圣德的神的天主子}，或\Quote{耶稣基督的从死者复活的天主子}，而通常的表达方式本该是：藉自己的大能，或按自己圣德的神，或藉自己的死者的复活，或藉他们的死者的复活——如我所说，我们并不被迫理解另一个人，而是同一的、同一位，即天主子我们的主耶稣基督的位格；同样，当我们读到\Quote{天主按天主的肖像造了人}时，虽然更常见的说法是\Quote{按自己的肖像}，但我们并不被迫理解天主圣三中的任何另一位，而是理解同一的、自身同一的天主圣三本身，祂是唯一的天主，人是按祂的肖像造的。\par

8. 既然情况如此，如果我们要接受同样的天主圣三的肖像，不是在一个人的身上，而是在三个人身上，即父亲、母亲和儿子，那么人就不是在为他造了妻子之前，也不是在他们生养了儿子之前，按天主的肖像造的；因为那时还没有一个天主圣三。难道有人会说，那时已有一个天主圣三，因为女人虽然在男人身侧尚未成形，儿子在父亲腰间尚未出世，但已在原始的本性中存在了吗？那么，为什么圣经在说了\Quote{天主按天主的肖像造了人}之后，紧接着说：\Quote{天主创造了他；祂造了他们，有男有女：天主就祝福了他们}？（或者如果这样划分：\Quote{天主创造了人}，然后加上\Quote{按天主的肖像造了他}，第三次再附上\Quote{祂造了他们，有男有女}；因为有些人害怕说\Quote{祂造了他有男有女}，以免理解成某种怪异的东西，比如他们称为阴阳人的那种；虽然即便如此，用单数来理解两者也并非错误，因为经上说\Quote{二人成为一体}。）那么，为什么如我开始所说的，关于按天主的肖像所造的人的本性，圣经只强调男和女呢？当然，为了完成天主圣三的肖像，它本应也提到儿子，即使儿子尚在父亲的腰间，正如女人在男人的身侧一样。或者也许女人已经造好了，而圣经只是用简短概括的陈述，将后来要更仔细说明如何发生的事合并在一起；因此不能提到儿子，因为儿子还没有出生？仿佛圣神不能在扼要的陈述中也包含这一点，同时将在适当的地方叙述儿子的出生；正如圣经后来在其适当的地方叙述了女人是从男人身侧取出的，但在此处并未省略提及她。\par

% structure: p0017 source=h2 target=subsection reason=relative-heading-rank; base=h2

\subsection{第七章 —— 人如何是天主的肖像。女人是否也是天主的肖像。如何理解宗徒所说：\Quote{男人是天主的肖像，但女人是男人的光荣}这一说法，并按其寓意和奥秘的含义来理解。}

9. 因此，我们不应如此理解人被造为至上天主圣三的肖像，即天主的肖像，好像这同一肖像应当被理解为存在于三个人身上；尤其当宗徒说男人是天主的肖像，并因此除去他头上的遮盖，而劝告女人应使用遮盖时，他这样说：\BibleQuote{男人本不应当蒙头，因为他是天主的肖像和光荣；但女人是男人的光荣。} 那么对此我们该说什么呢？如果女人按她自身位格的程度补充了天主圣三的肖像，为什么在女人被从男人肋骨中取出后，男人仍被称为那肖像？或者，如果三个人中的一个人可以被称为天主的肖像，如同在至上天主圣三中每位也是天主一样，为什么女人也不是天主的肖像呢？因为她正是为此受教导要蒙头，而男人被禁止蒙头，因为他是天主的肖像。\par

10. 但我们必须注意，宗徒所说不是女人而是男人是天主的肖像，这与创世记所记载并不矛盾：\BibleQuote{天主创造了人：照天主的肖像创造了祂；造了一男一女；并祝福了他们。} 因为这段经文说人性本身（只在两性中才是完整的）被造为天主的肖像；它并没有将女人从它所表示的天主肖像中分离。因为说了天主按天主的肖像造了人之后，它说\BibleQuote{祂创造了他}，又说了\BibleQuote{一男一女}：或者至少换一种标点方式：\BibleQuote{造了一男一女。} 那么宗徒怎么告诉我们男人是天主的肖像，因此他禁止蒙头；而女人不是，因此命令她蒙头呢？除非，按我已经说过的，当我论及人类心智的本性时，女人与其丈夫一同是天主的肖像，以致那整个实体成为一个肖像；但是当她单独被指为其作为\Emph{助伴}的品质时（这关涉女人本身），那么她不是天主的肖像；但就男人单独而言，他完全地、完整地是天主的肖像，正如女人也与他结合为一。正如我们论人类心智的本性时所说：当它作为一个整体默观真理时，它是天主的肖像；当某一部分从中被分割、转向认知暂时之物时，然而在它注视并咨询真理的那一面，它也是天主的肖像；但在它被导向认知较低事物那一面，它不是天主的肖像。既然它越延伸至永恒之物，就越按天主的肖像被塑造，因此不应限制它，以阻止它并使之远离那永恒之物；所以男人不应蒙头。但是因为理性认知（它从事于有形和暂时之物）过多地向下发展是危险的；这应当在头部有控制力，蒙头即表示这一点，借此表明应当被抑制。因为神圣而虔诚的含义是圣天使所喜悦的。因为天主不按时间的方式观看，当任何事按时间和一时的方式发生时，在祂的视力和知识中并无新事发生，不像这类事物影响感官（无论是动物和人的肉体感官，还是天使的天上感官）的方式那样。\par

11. 因为圣保禄宗徒在外表上谈论男女性别时，预示了某种更隐藏真理的奥迹，这可以从他另一处所说的来理解：他说真正的寡妇是无依无靠的、没有子女和孙辈的，然而她应当信赖天主，并昼夜不停地祈祷，此处他指出，女人因受欺骗而陷入过犯，却因生育而得救；然后他加上了：\BibleQuote{她们若在信德、爱德、圣洁和节制上持守。} 仿佛一个好寡妇可能因她没有儿子，或她所生的儿子不愿持守善工而受损。但因为那些被称为善工的东西，好比是我们生命的儿子——按那种生命的意义，它回答\Quote{人的生命是什么？}的问题，即\Quote{他在这些暂时之事中如何行事？}——希腊人称这种生命不是\OriginalTerm{生命}{\GreekText{ζωή}}而是\OriginalTerm{生计、生活}{\GreekText{βίος}}；又因为这些善工主要在实行慈悲的事工中完成，而慈悲的事工对外教人、不相信基督的犹太人、或任何异端者或分裂者（在他们内找不到信德、爱德和节制圣洁的人）都没有益处：宗徒所要表明的是明显的，并且就形象和奥秘而言，因为他在谈论女人的蒙头，这除非指向某个隐藏的圣事，否则仍将是空洞的话语。\par

12. 因为，不仅最真的理性如此断言，而且宗徒本人的权威也宣告：人不是按身体的外形，而是按理性的心智，按天主的肖像受造的。那种认为天主被身体肢体的轮廓所限制和界定的想法，是卑贱而空洞的。此外，难道同一位蒙福的宗徒不是说过：\BibleQuote{你们应在心思念虑上更新，穿上新人，就是按照天主的肖像受造的}；并在另一处更明确地说：\BibleQuote{脱去旧人及其行为，穿上新人，这新人即是依照创造他者的肖像，在认识天主上更新的}吗？那么，如果我们是在心思念虑上更新，而新人就是按照创造他者的肖像在认识天主上更新的，那么无人可以怀疑：人是按照创造他者的肖像受造的，不是按身体，也不是按心智的任何部分不加区分，而是按理性的心智，其中可以存在对天主的认识。也正是根据这种更新，我们藉基督的圣洗成为天主的子女；穿上新人，当然就是藉信德穿上基督。那么，有谁认为妇女与这种共融无缘呢？她们与我们同为恩宠的承继者；而且同一宗徒在另一处说：\BibleQuote{因为你们众人藉着信德在基督耶稣内，都是天主的子女；你们凡受洗归于基督的，就是穿上了基督：不再分犹太人或希腊人，不再分奴隶或自由人，不再分男或女，因为你们众人在基督耶稣内已合而为一}。请问，有信德的妇女难道失去了她们身体的性别吗？但是，因为她们在无性别之处（即心思念虑中）按天主的肖像得以更新，人在无性别之处（即心思念虑中）是按天主的肖像受造的。那么，为何男人因此不必蒙头，因为他是天主的肖像和光荣，而女人却必须蒙头，因为她是男人的光荣？难道女人不是在心思念虑中更新，而此心思念虑是按创造他者的肖像在认识天主上更新的吗？但因她在身体性别上与男人有别，便可以恰当地用她的身体覆盖来代表那转向管理暂时事物的那部分理性；如此，天主的肖像可以留在人的心智中那依附于默观或咨询事物永恒理性的一面；显然，这不仅男人拥有，女人也拥有。\par

% structure: p0022 source=h2 target=subsection reason=relative-heading-rank; base=h2

\subsection{第八章　偏离天主的肖像。}

13. 因此，在他们（男人和女人）的心智中承认有一个共同的本性，但在他们的身体中却描绘出同一个心智的分裂。于是，当我们通过心智的各个部分，沿着内在的某些步骤上升时，在首先遇见我们与野兽所不共有的东西之处，理性开始了，这样内在的人便可以在此被认出来。如果这个内在的人本身，通过那被委派管理暂时事物的理性，因过度进入外在事物而滑得太远，并且他的头（即男性部分，如其所说，在议事的瞭望塔中主持）也同意，不加以约束或勒住它；那么他就会因他所有的仇敌（即魔鬼及其群魔，他们嫉妒德行）而衰老；并且永恒事物的景象也会从他的头本身被夺走，他与他的配偶吃了那禁物，以致他眼中的光明离他而去；这样，两人都脱离了真理的光照而赤身露体，他们良心的眼睛睁开，看见自己被遗弃而羞耻不雅，如同甜美果实的叶子，却没有果实本身，于是他们便编织美好的言词却没有美好行为的果实，仿佛在邪恶地生活时，用善言掩盖自己的羞耻。\par

% structure: p0024 source=h2 target=subsection reason=relative-heading-rank; base=h2

\subsection{第九章　继续同一论点。}

14. 因为灵魂爱恋自己的能力，便从共同的整体滑向特别属于它自己的部分。这种背信的骄傲被称为\BibleQuote{罪恶的开端}，它本可极好地受天主法律的管理，若它作为统治者跟随天主在普世受造物中，却因寻求超出整体之物，并力图按其自身的法律来管理，而被推入对部分的关心，因为没有任何事物超出整体；于是因贪求更多而变得更少；因此贪财被称为\BibleQuote{万恶的根源}。它通过自己仅占部分的躯体，管理那整体（在其中它力图违反管理整体的法律，做属于自身的事）；于是它因喜爱有形体的形状和运动——因为它本身不拥有这些事物本身，并且因缠绕在它固定在记忆中的这些形象中，因幻想的淫乱而污秽地被玷污——而将自己的一切功能指向那些目的，为此它通过身体的感官好奇地寻求有形体和暂时的事物，或者因傲慢的膨胀而自以为优于其他沉迷于身体感官的灵魂，或者沉入可耻的肉欲漩涡。\par

% structure: p0026 source=h2 target=subsection reason=relative-heading-rank; base=h2

\subsection{第十章　一步一步达到最深的堕落。}

15. 当灵魂为自己或为他人善意地咨询，以感知内在和更高的事物——这些事物在一种纯洁的拥抱中被所有热爱它们的人共同拥有，没有狭隘或嫉妒，并非个别地，而是共同地——那么，即使它在任何事上因对暂时事物的无知而被欺骗（因为它在此时的行为是暂时的），并且它没有坚持所应采取的行动方式，这试探也不过是人所共有的。能够如此度过此生，即我们如同踏上归途的道路一般行走，以致没有试探临到我们，除非是人所共有的，这是一件大事。因为这是一种与身体无关的罪过，不可算为奸淫，因而极易被宽恕。但是，当灵魂为了获得那些通过身体感知的事物，出于证明、超越或触摸它们的私欲，以便将它们的终极之善置于其中时，那么它所做的一切都是邪恶的，并且犯下了奸淫，得罪了自己的身体：它从内在攫取虚妄的形像，并通过虚妄的思想将它们组合，以致除了此类事物之外，似乎再没有什么是神圣的；由于自私的贪婪，它在错误中变得多产，由于自私的挥霍，它被耗尽了力量。然而，它不会从一开始就骤然跳入如此无耻和悲惨的奸淫，而是如经上所记载：\Quote{轻视小事的人，必渐渐跌倒。}\par

% structure: p0028 source=h2 target=subsection reason=relative-heading-rank; base=h2

\subsection{第十一章——人内里的兽像。}

16. 蛇不是迈开大步爬行，而是凭借每一片鳞片最细微的移动而前进；同样，堕落（脱离善）的滑动只有逐渐地占据疏忽者，从对天主肖像的扭曲渴望开始，最终达到兽类的肖像。因此，他们赤身失却了最初的衣服，因可朽性而获得了皮衣。因为人的真正尊荣是天主的肖像和模样，除非通过与那印刻者本身的关系，否则不能保存。因此，一个人越是少爱属于自己的事物，就越能紧贴天主。但为了试探自己能力的欲望，人凭自己的命令下降到自己，仿佛降到一个中间层级。于是，当他希望像天主一样不受任何人的管束时，他甚至从自己的中间层级被惩罚性地推至最低处，即那些野兽所喜好的事物；因此，他的尊荣是天主的肖像，而他的羞辱是野兽的肖像，\Quote{人在尊荣中而不醒悟，就如无知的畜类，被造成与它们相似。} 那么，他怎能从至高之处经过如此遥远的距离而下降到最低之处呢？除非通过自己的中间层级。因为当他忽略了永远不变的智慧之爱，却贪求通过经验认识暂时和可变的万物时，这种知识令人自高自大，并不建立德行；于是，心灵被压垮，仿佛被自己的重量从幸福中抛出来；并通过对自己中间地位的考验，藉着自己的惩罚而学到它所抛弃的善与它所陷入的恶之间的区别；并且，在丢弃并摧毁了自己的力量之后，它不能返回，除非藉其造物主的恩宠呼召它悔改，并赦免它的罪过。因为谁能将不幸的灵魂从这死亡的身体中拯救出来呢？除非是天主藉我们的主耶稣基督的恩宠。这恩宠我们将在适当的地方，按祂恩赐我们的能力加以论述。\par

% structure: p0030 source=h2 target=subsection reason=relative-heading-rank; base=h2

\subsection{第十二章——内在之人有一种隐秘的婚姻。思想中的非法快乐。}

17. 现在，靠着主的助佑，让我们完成我们所承担的那部分涉及知识的理性的讨论，即对可朽可变的世物的认识，这些认识对于管理今生事务是必要的。因为正如在最初被造的两个人的那种可见的结合中，蛇没有吃禁树，而只是劝说他们去吃；女人没有独自吃，而是给了她的丈夫，他们一同吃了；虽然只有女人与蛇说话，只有她被蛇引诱：同样，在那种隐秘而秘密的结合中，即在一个人身上进行并分辨的结合中，属肉体的，或者我可以说，由于它指向身体的感官，那灵魂的感性的运动（这是与我们和野兽相同的），是与智慧的理性隔绝的。因为的确，身体的事物是通过身体的感官来感知的；但属神的事物，即永恒不变的事物，是通过智慧的理性来理解的。然而，知识的理性与它有非常接近的欲望：因为所谓科学或行动的知识，它理性地处理通过身体感官感知的身体事物；如果做得好，是为了将那知识指向至善的最终目的；但如果做得不好，是为了享乐它们，好像它们是那些能使其在虚假的幸福中安息的善物。每当那种属肉体的或动物的感官，借着理性的生命力，向心灵的这个处理时间性和物质性事物（为了人的行动职责）的意向引入某种自我享乐的诱因时，即把自己当作某种私有的善（而非公共的、共同的、不变的善）来享乐；这时，就好像蛇在对女人说话。同意这种诱惑，就是吃禁树。但如果那同意仅以思想的快乐为满足，而肢体被更高智慧的权威所抑制，以致没有被用作不正义的武器去犯罪；我认为，这应被视为好像女人独自吃了禁食。但如果在这同意中，决定要恶用通过身体感官所感知的事物，以至于如果有能力，也由身体去完成；那么，就必须理解为那女人将非法食物给了她的丈夫，与他一同吃了。因为除非心灵的意向——它拥有将肢体应用于外在行动或抑制它们的主要权能——屈服并服务于那恶行，否则心灵不可能决定不仅以喜悦思想罪，而且要有效地实施。\par

18. 然而，当然，当心灵在思想中独自悦乐非法的事物，而确实不决定去行，但却愉快地持有和思想那些本应一触及心灵就应拒绝的事物时，这不能否认是一种罪，但远比如果也决定在外在行为中实现要小得多。因此，也必须为这样的思想寻求宽恕，必须捶胸，并要说：\BibleQuote{宽免我们的罪债}，并且必须做随后的事，并且在我们的祈祷中加上：\BibleQuote{如同我们也宽免得罪我们的人}。因为这不像是那两个最初的人的情况，每个人都有自己的位格；所以，如果女人独自吃了禁食，她肯定会独自被死亡惩罚；我不能说，在现在一个人的情况下，也可以这样说，如果那思想独自留连，被非法的快乐愉快地喂养（它本应立刻转身离开），而同时没有决定要行恶行，只是愉快地记念它们，那么就好像女人可以在没有男人的情况下被定罪。我们切勿相信这一点！因为这里是一个人，一个人类，他整个将被定罪，除非那些被感知为仅是思想之罪（缺乏行动意愿，但有悦乐心灵意愿）的事情，通过中保的恩宠被赦免。\par

19. 那么，这个推理——我们借此在每一个个人的心灵中寻求某种观与行的理性结合，各自有不同的功能，但两者中保持了心灵的统一；同时保留那由神圣见证所传下来的关于最初两个人（即男人和他的妻子，人类从他们繁衍）的\Emph{历史}的真实性——这个推理，我说，必须仅听到这样的程度：即可以理解宗徒有意要指涉某种在单独一人身上所找到的东西，通过将天主的肖像只归于男人，而不归于女人，尽管在仅仅是两个不同性别的人中。\par

% structure: p0034 source=h2 target=subsection reason=relative-heading-rank; base=h2

\subsection{十三、——那些认为男人代表心灵、女人代表身体感官之人的意见。}

20. 我也并非不知道，在我们之前一些杰出的天主教信仰维护者和天主圣言诠释者，他们在一个人身上寻找这两样东西，并将整个灵魂视为某种卓越的乐园，断言男人是心智，女人是身体感官。根据这种划分——即假定男人是心智，女人是身体感官——如果加以适当注意，似乎一切都能恰当地协调一致；只是经上记载说，在一切走兽和飞鸟中，没有给男人找到一个相配的助手；于是女人是从他的肋骨中造出来的。为此，我个人并不认为身体感官应理解为女人——我们看到这一点是我们与野兽所共有的；而是希望找到某种野兽没有的东西；我更倾向于认为身体感官应理解为蛇，经上记载说蛇比田野里的一切走兽都更狡猾。因为在我们看到与无理性动物共有的那些天然美好事物中，感官因某种生命活力而卓越；但这并非致希伯来人书信中所写的那种感官，那里说：\BibleQuote{坚固的食物是属于成年人的，就是那些因习惯而运用其官能、以辨别善恶的人}；因为这里的\Quote{官能}属于理性本性，关乎理解；而是那种在身体中分为五部分的感官，通过它，不仅我们，而且野兽也能感知到物质的形状和运动。\par

21. 然而，宗徒称男人是天主的肖像和光荣，女人是男人的光荣，这一说法究竟应以这种、那种或任何其他方式来理解，尚可讨论；但显然，当我们按照天主生活时，我们专注于祂不可见之事的理智，应当从祂的永恒、真理和爱德中得到熟练的塑造；而我们理性意向的一部分——也就是同一种理智——必须用于使用可变和物质的事物，没有这些事物，此生便无法延续；这并不是要我们同化于这个世界，将我们的目的置于这类善事中，并将幸福的欲望强加于它们，而是我们在使用暂时事物时，所做的一切都要以得享永恒事物的默观为依归，经过前者，却紧贴后者。\par

% structure: p0037 source=h2 target=subsection reason=relative-heading-rank; base=h2

\subsection{第十四章——智慧与知识的区别是什么。对天主的敬拜就是爱祂。如何通过智慧获得对永恒事物的理智认识。}

因为知识也有其良好的尺度，如果其中那令人自高自大、或通常令人自高自大的东西，被对永恒事物的爱所征服，而爱并不自高自大，反而如我们所知，能建造人。诚然，没有知识，就无法拥有那些使人正直生活的德行，而这些德行可以管理这悲惨的人生，使我们得以达到那真正有福的永生。\par

22. 然而，我们善用暂时事物的行动，与对永恒事物的默观是有区别的；后者归于智慧，前者归于知识。虽然那称为智慧的也可以称为知识，如宗徒所说的：\BibleQuote{我现在所知道的，只是局部的，那时我就要全知道了，如同我全被知道一样}；显然他是指对天主的默观知识，那将是圣人们最高的赏报。但他说：\BibleQuote{这一人从圣神蒙受了智慧的言语，另一人却由同一圣神蒙受了知识的言语}时，无疑明确区分了这两者，尽管他未在那里解释其区别，也未说明如何辨别。然而，在查阅了《圣经》的许多章节后，我在《约伯传》中读到，那位圣人说：\BibleQuote{看，虔敬就是智慧；远离邪恶就是知识}。在这个区分中，必须理解为智慧属于默观，知识属于行动。因为在这里，他所说的\Quote{虔敬}是指对天主的敬拜，希腊文称为\OriginalTerm{虔敬}{\GreekText{θεοσέβεια}}。因为在希腊文手稿中，该句子用了这个词。而在永恒事物中，有什么比本性不变的天主更卓越呢？敬拜祂不就是爱祂吗？我们藉此爱现在渴望见到祂，并相信和希望将来必见到祂；随着我们进步，现在对着镜子观看，模糊不清，那时则要面对面？这就是保禄宗徒所说的\Quote{面对面}。这也是若望所说的：\BibleQuote{亲爱的，现在我们是天主的子女，但我们将来如何，还没有显明；可是我们知道：一显明了，我们必要相似祂，因为我们要看见祂实在怎样}。在我看来，谈论这些以及类似的主题，就是智慧本身的论述。而远离邪恶——约伯称之为知识——无疑属于暂时的事物。因为我们是在时间中处于邪恶之中，应当戒避邪恶，以求达到那永恒的善。因此，凡是我们明智地、勇敢地、节制地、公正地所行的，都属于那种知识或学问，我们的行动藉之在避免邪恶、渴望善事中运转；同样，凡是我们通过探究所获得的、应提防或效法的榜样，以及与我们用途相称的任何必要证明，也属于知识。\par

23. 当论述涉及这些事物时，我认为它属于知识的论述，并与属于智慧的论述相区别；智慧所涉及的是那些既未曾有、也将不会有、而是永恒存在的事物，并且由于它们所处的永恒性，它们被说成曾经有、现在有、将来有，没有任何时间的变化。因为它们的存在并非如此，以至于它们会停止存在；它们的存在也非如此，仿佛它们现在不存在；但它们始终拥有并将永远拥有那绝对的存在。它们常存，但不是像物体那样固定在某个地方；而是作为无形体本性中的可理解之物，如此贴近心灵的瞥视，就像可见或可触之物在地方上贴近身体感官一样。不仅对于置于地方的可感事物，还有它们可理解的、无形体的理由，它们超越地方空间而常存；而且对于在连续时间中经过的运动，超越任何时间中的过渡，也有类似的理由屹立，这些理由本身诚然是可理解的，而非可感的。以心灵之眼达到这些是少数人的福分；当它们尽可能达到时，达到者本人并不常驻于其中，而是仿佛被心灵之眼本身的反弹所斥退，于是对一个非暂时之物产生了一个暂时的思想。然而，这个暂时的思想通过心智所受教导的指令被委托给记忆；这样，被迫离开那里心智能够再次返回那里；尽管，如果思想没有返回记忆并在那里找到它所委托之物，它就会像一个未受教导的人一样被引导到那里，就像之前被引导那样，并在它最初发现它的地方找到它，也就是说，在那无形体的真理中，从那里它可能再次被写下并固定在心中。因为人的思想，例如，并不像那个正方形物体的无形体的、不变的理由本身那样常驻于其中：当然，如果它能在没有地方空间幻想的情况下达到它的话。或者，如果有人领会了某种人工或音乐声响的节奏，它经过一定的时间间隔，而在无声的秘密和深沉的寂静中无时间地静止，那么至少只要那歌声能被听见，它就能被思想；然而心灵之眼从那里捕捉到的（尽管是暂时的）并像吸收进肚子一样储存在记忆中，它就能通过回忆在一定程度上反刍，并将所学得的东西转化为系统的知识。但如果这被彻底遗忘抹去，那么在教导的引导下，他会再次来到那完全失去的东西面前，并且会发现它如同从前一样。\par

% structure: p0041 source=h2 target=subsection reason=relative-heading-rank; base=h2

\subsection{第15章——反驳\Person{柏拉图}与\Person{毕达哥拉斯}的回忆说。萨摩斯人\Person{毕达哥拉斯}。论智慧与知识的区别，以及在暂时事物的知识中寻求天主圣三。}

24. 因此，那位高贵的哲学家\Person{柏拉图}竭力使我们相信，人的灵魂在进入这些身体之前就已经存在；因此，那些所学的东西与其说是作为新知识被接受，不如说是被回忆起来，好像事先已知。因为他告诉我们，一个男孩——我不知道他问了什么关于几何学的问题——回答时仿佛精通那门学问。因为通过逐步而巧妙的提问，他看到了应当看到的东西，并说出了他所看到的。但如果这是对先前已知事物的回忆，那么无疑每个人或几乎每个人被问到时都不会这样回答。因为并非每个人在前世都是几何学家，因为几何学家在人类中如此稀少，以至于几乎在任何地方都找不到一个。但我们更应当相信，理智心灵按其本性如此形成，以至于它能通过一种独特而非物质的光，看到那些事物；这些事物按照创造者的安排，以自然顺序隶属于可理解的事物；就像肉体之眼在这身体的光中看到邻近自身的东西，而这光正是它被造来接受并适合它的。因为肉体之眼在辨别黑白时也并非无需教师，因为它在这肉体中被创造之前就已经知道它们了。最后，为什么只有在可理解的事物中，即使无知的人也能在适当提问下按照任何学问回答？为什么没有人在可感事物中能做到这一点，除非那些人曾在他现在的身体中看到过它们，或者相信了知道它们的人的传闻——无论是某人的著作还是言语？因为我们不能苟同那些人的故事，他们说萨摩斯人\Person{毕达哥拉斯}回忆了这类事情，是他先前在另一个身体里时经历过的；还有其他关于其他人的故事，说他们在心灵中经历过类似的事情：但可以推测，这些回忆是虚假的，如同我们通常在梦中经历的那样，我们以为记得我们从未做过或见过的事情，好像我们做过或见过似的；而这些人的心灵即使在清醒时，也在邪恶欺骗之灵的暗示下受到这种影响，这些邪灵关心的是要证实或散布关于灵魂变化的某些虚假信仰，以欺骗人类。我说，可以从这一点推测：如果他们真的记得他们在其他身体里时在这里看到过的事情，那么同样的事情会发生在许多人、甚至几乎所有人身上；因为他们认为死者从生者而来，反之亦然，无休止地不断有生者从死者而来，就像醒者从睡者而来，睡者从醒者而来一样。\par

25. 因此，如果智慧与知识的正确区分在于：对永恒事物的理智认知属于智慧，而对暂时事物的理性认知属于知识，那么不难判断两者孰先孰后。但若我们必须采用其他区分来辨识这两者——无疑宗徒教导我们它们是不同的，他说：\BibleQuote{这人蒙圣神赐予智慧之言，那人蒙同一圣神赐予知识之言}——然而我们提出的这两者之间的区别仍是最明显的，即对永恒事物的理智认知是一回事，对暂时事物的理性认知是另一回事；无人怀疑前者优于后者。因此，当我们舍弃那些属于外在人的事物，并渴望从我们与野兽共有的那些事物中上升至内在，在达到对可理解且至高、永恒的事物认知之前，对暂时事物的理性认知便呈现出来。那么，让我们也在此寻找一个三一，如果可能，如同我们在身体感官中，以及通过感官以图像方式进入我们灵魂或精神的事物中找到一个三一那样；这样，代替我们通过身体感觉所触摸的、置于我们之外的有形事物，我们在记忆内有了身体的相似物印记，由此可以形成思想，而意志作为第三者将它们结合；就像眼睛的视觉从外部形成，意志将其应用于可见物体以产生视觉，并将两者结合，同时意志本身也作为第三者加入其中。但此主题不可压缩进本书；因此，在接下来的论述中，若蒙天主助佑，它会被妥善探讨，我们得出的结论也得以展开。\par

\SourceRecord{Translated by Arthur West Haddan.；Nicene and Post-Nicene Fathers, First Series；Vol. 3.；Edited by Philip Schaff.；Buffalo, NY: Christian Literature Publishing Co.,；1887.；<http://www.newadvent.org/fathers/130112.htm>.}

% structure: p0001 source=h1 target=section reason=page-title

\section{\WorkTitle{论天主圣三（第十三卷）}}

\EditorialNote{本卷继续探讨关于知识的论题，奥古斯丁在前一卷中已开始从知识（有别于智慧）中寻找某种三一。此间也借机赞扬基督信仰，并解释信徒的信仰如何是唯一而共同的。其次，所有人都渴求幸福，但并非所有人都拥有那使我们达至幸福的信仰；而这信仰是在基督内界定的，祂曾肉身复活；除非藉着祂，没有人能藉罪之赦免脱离魔鬼的辖制。本书也详细说明，基督必须藉正义而非权势战胜魔鬼。最后，当此信仰的话语被记在记忆里时，心中便有某种三一，因为首先是话语的声音在记忆里，即使人没有想到它们时亦然；其次，当他想到它们时，回忆的心灵之眼便由此而形成；最后，意志当他如此思想并记忆时，将二者结合。}\par

% structure: p0003 source=h2 target=subsection reason=relative-heading-rank; base=h2

\subsection{第一章——尝试从圣经中区分智慧与知识的职务。在若望福音序言中，有些话属于智慧，有些属于知识。有些事只能凭信仰认识。我们如何看见我们内的信仰。在同一若望的叙事中，有些事由身体感官认识，有些只凭心灵理性认识。}

1. 在前一卷，\Emph{viz}即本著作的第十二卷中，我们已经足够区分理性心灵在暂时事物中的职务——其中不仅涉及我们的认识，也涉及我们的行动——与同一心灵的更卓越的职务，后者专注于默观永恒之事，且仅限于认识。但我认为更合适的是从圣经中插入一些内容，使两者更易区分。\par

2. 若望福音作者这样开始他的福音：\BibleQuote{在起初已有圣言，圣言与天主同在，圣言就是天主。圣言在起初就与天主同在。万有是藉着祂而造成的；凡受造的，没有一样不是由祂造成的。在祂内有生命，这生命是人的光。这光在黑暗中照耀，黑暗不能胜过祂。曾有一人，由天主派遣，名叫若翰。这人来，是为作证，为给光作证，为使众人藉着他而信。他不是那光，而是来给光作证的。那才是真光，照亮每一个来到世上的人。祂已在世界上，世界原是藉祂造成的，但世界却不认识祂。祂来到自己的领域，自己的人却没有接受祂。但是，凡接受祂的，祂给他们权能，成为天主的子女，就是那些信祂名字的人，他们不是由血，也不是由肉欲，也不是由人的意愿，而是由天主所生的。于是，圣言成了血肉，寄居在我们中间；我们见了祂的光荣，正如父独生者的光荣，满溢恩宠和真理。}这一整段经文，是我从福音中引用的，其较早部分包含不变而永恒的事，默观这些事使我们有福；但随后部分则提到永恒之事与暂时之事相连。因此，根据我们在第十二卷中先前的区分，其中有些事属于知识，有些属于智慧。因为这些话：\BibleQuote{在起初已有圣言，圣言与天主同在，圣言就是天主。圣言在起初就与天主同在。万有是藉着祂而造成的；凡受造的，没有一样不是由祂造成的。在祂内有生命，这生命是人的光。这光在黑暗中照耀，黑暗不能胜过祂。}——需要默观的生活，须藉理性心智来辨别；任何人越是在这方面进步，无疑就越是明智。但关于那句：\BibleQuote{光照在黑暗中，黑暗不能胜过祂，}信仰当然是必要的，以便相信那未看见的事。因为藉着\Quote{黑暗}，他意指远离这种光且几乎无法看见它的必死之人的心灵；为此他接着写道：\BibleQuote{曾有一人，由天主派遣，名叫若翰。这人来，是为作证，为给光作证，为使众人藉着他而信。}但在这里，我们遇到了一件在时间内发生的事，属于知识，知识包含在事实的认知中。我们按照从人类自然概念印在记忆中的那种幻影来思考若翰这个人。无论人信不信，他们都以同样方式思想。因为两者都知道什么是人——其外在部分，即身体，他们通过身体之眼认识；内在部分，即灵魂，他们自身内拥有知识，因为他们自己也是人，并通过与人的交往；所以他们能够思考\Quote{曾有一人，名叫若翰}这句话，因为他们通过言语交流也知道名字。至于那里还说的\BibleQuote{由天主派遣，}凡持守它的人都是以信仰持守；凡不以信仰持守的人，或因怀疑而犹豫，或因不信而讥笑。但两者，如果他们不是属于那些过于愚蠢、心里说\BibleQuote{没有天主}的人，当他们听到这些话时，都会思考两件事，即：天主是什么，以及由天主派遣是什么；如果他们不如事物本身那样去想，那至少也尽其所想。\par

3. 此外，我们从其他来源也知道那为人所见在自己心中的信德本身——若人相信，这信德便在他心中；若他不相信，便不在他心中。但我们认识信德，不像认识物体那样，物体我们用肉眼看见，即使不在眼前，我们也通过它们留在记忆中的影像来思想；也不像认识那些我们未曾见过的事物那样，我们凭所见之物随意在心中构想它们，并将它们存入记忆，以便我们愿意时能回想，藉以同样通过回忆来辨别它们，或者更确切地说，辨别它们的影像，无论这些影像在我们心中是怎样的；也不像认识一个活人那样，我们固然看不见他的灵魂，却从自身来推测，并藉着身体的动态也在思想中凝视活人，如同我们通过视觉所认识的那样。信德在它所居住的心中，对拥有它的人而言，并非如此被看见；然而最确实的知识紧紧把握着它，良心宣告着它。因此，虽然我们奉命相信这叙述，因为我们看不见我们所奉命相信的事物；然而我们自己看见自己心中的信德本身，当这信德在我们内时；因为信德即使对于不在场的事物也是在场，信德对于在外的事物是在内，信德对于看不见的事物本身是被看见的，它同样在时间中进入人心；如果有人不再忠信而成为不信者，它便从他们那里消逝。有时信德甚至被用于虚假的事；因为我们有时这样说：我信了他，他却欺骗了我。这种信德——如果它确实也可被称为信德的话——当真理被发现并将它驱逐时，便从心中消逝而无咎。但对于真实之事的信德，应按人的愿望消逝，转化为事物本身。因为我们不能说，当所信的事物被看见时，信德就消逝了。因为当信德——按\BibleBook{}{希伯来书}的定义是\Quote{未见之事的证据}\BibleRef{希 11:1}——时，它还能被称为信德吗？\par

4. 在接下来的话中：\BibleQuote{这人来，是为作证，为给光作证，为使众人藉着他而信}\BibleRef{若 1:7} 行动，如我们所说，是在时间中完成的。因为即使为那永恒的事物——如那可理解的之光——作证，也是在时间中完成的。正是为此，若翰前来作证，他\BibleQuote{不是那光，而是被派遣来为光作证}\BibleRef{若 1:8} 因为接着他说：\BibleQuote{那普照每人的真光，正在进入这世界。祂已在世界上，世界原是藉着祂造的，世界却不认识祂。祂来到自己的领域，自己的人却没有接受祂}\BibleRef{若 1:9-11} 现在，懂拉丁语的人，从他们所知道的事物中理解所有这些话；其中一些是通过身体感官得知的，如人，如世界本身——其伟大如此明显于我们的视觉；又如言语的响声本身，因为听觉也是身体的一种感官；还有一些是通过心灵的理性得知的，如所说的\Quote{自己的人却没有接受祂}；\Quote{接受}意味着他们不信祂；而什么是信德，我们不是通过身体感官知道的，而是通过心灵的理性。我们也学习了词汇本身——不是声音，而是意义——部分通过身体感官，部分通过心灵的理性。现在我们也不是第一次听到这些话，它们是我们以前听过的。我们将它们作为已知之物保留在记忆中，并在此认出了，不仅是词汇本身，还有它们所意味的。因为当双音节词\OriginalTerm{世界}{mundus}被说出时，那么某种确实是形体的事物——因为那是声音——便通过身体，即通过耳朵被认识。但它所意味的东西，也是通过身体，即通过肉体的眼睛被认识的。因为就世界为我们所知而言，它是通过视觉认识的。但四音节词\OriginalTerm{他们信了}{crediderunt}，就其声音——因为那是形体之物——通过肉体的耳朵到达我们；但它的意义不是通过身体感官发现的，而是通过心灵的理性。除非我们通过心灵知道\OriginalTerm{他们信了}{crediderunt}这个词是什么意思，否则我们不会理解那些被说\Quote{自己的人却没有接受祂}\BibleRef{若 1:11}的人没有做什么。因此，语词的声音从外面震动肉体的耳朵，并到达被称为听觉的感官。人的形像也既在我们自身中被认识，又从外面通过身体感官呈现在其他人中：对眼睛，当被看见时；对耳朵，当被听见时；对触觉，当被握住和触摸时；而且它在我们的记忆中也有一个影像，非形体却类似形体。最后，世界本身的奇妙美丽从外面呈现在我们的注视下，也呈现在被称为触觉的感官下——如果我们接触它的任何部分；这也在我们内部有一个影像，在记忆中，当我们回想它——无论是在房间的围蔽中还是在黑暗中——时，我们便回到那影像。但我们已在\WorkTitle{第十一卷}中充分讨论了这些形体之物的影像；它们虽非形体，却具有形体的相似性，并属于外在之人的生活。但现在我们正在讨论内在之人及其知识，即关于暂时和可变化之物的知识；为此目的和范围，当任何事物——即使属于外在之物——被采纳时，必须为此目的而被采纳，以便可以从中教导某些有助于理性知识的东西。因此，对那些我们与无理性的动物共有的东西的理性使用，属于内在之人；也不能正确地说，这为我们与无理性的动物所共有。\par

% structure: p0008 source=h2 target=subsection reason=relative-heading-rank; base=h2

\subsection{第二章——信德是心的事，不是身体的事；如何在所有信徒中它是共同且同一的。信徒的信德是同一的，无异于那些愿意者的意志是同一的。}

5. 但是信德——出于我们主题的安排，我们不得不在这卷书中更详细地讨论它——我说的是信德，拥有它的人称为信友，没有它的人称为不信者，就像那些没有接受来到自己地方的天主子的人；虽然信德是通过听而在我内形成，但它并不属于那称为听觉的身体感官，因为它不是声音；也不属于我们这血肉的眼睛，因为它既不是颜色也不是形体；也不属于那称为触觉的，因为它没有体积；它根本不属于任何身体感官，因为它是心灵的事，而非身体的事；它也不是在我们外，而是深藏于我们内；没有人在他人身上看见它，每人只在自己内看见它。最后，它可以藉伪装假装，也可以被认为存在于没有它的人身上。因此，每人在自己内看见自己的信德；但在他人身上，他看不见，而是相信那里有信德；他越了解信德藉爱常常产生的果实，就相信得越坚定。因此，这种信德为所有人共有，福音作者接着说到：\BibleQuote{凡接受他的，他赐他们权能成为天主的子女，就是那些信他名字的人：他们不是由血，也不是由肉欲，也不是由男欲，而是由天主生的。} 我说共有，不是像有形物体的形式就视觉而言为所有看见的人的眼睛共有，因为所有观看者的注视在某种程度上都被同一个形式所形成；而是像人的面容可以说为所有人共有；这样说，但每人当然有自己的面容。我们确实完全真实地说，信友的信德是从同一道理铭刻在每个相信同一事物的人的心上。但所信的对象不同于用以相信的信德。因为前者在于那些被说成是、曾经是或将要存在的事物；而后者在信友的心中，并且只有拥有它的人才能看见；虽然它本身并非独一无二，但相似的信德也在其他信友身上。因为它不是数量上的一，而是种类上的一；但由于相似且毫无差异，我们更称它为一而非多。因为当我们看见两个极其相似的人，我们会惊讶地说两人有同一面容。因此，更容易说灵魂是多个——每人当然有一个灵魂——正如我们在宗徒大事录中读到，他们\BibleQuote{是一个灵魂}。而宗徒说到\BibleQuote{一个信德}时，任何人敢说有多少信德者就有多少信德。然而他说：\BibleQuote{妇人，你的信德真大}，又对另一个人说：\BibleQuote{小信德的人哪！你为什么怀疑？}，这表明每人有自己的信德。但信友相似的信德被称为一个，正如意愿者相似的意愿被称为一个；因为在那些有相同意愿的人身上，每人的意愿对自己可见，但对他人不可见，尽管他意愿同一件事；如果它藉任何标记表明自己，它是被相信而非被看见。但每人意识到自己的心灵，当然不是相信，而是明明白白地看见这就是自己的意愿。\par

% structure: p0010 source=h2 target=subsection reason=relative-heading-rank; base=h2

\subsection{第三章——某些欲望在众人中相同，为每人所知。诗人\Person{恩尼乌斯}。}

6. 确实，同一兼具理性和生命的天性如此紧密地协同一致，以致虽然一人不知另一人所愿，但所有人的某些意愿却也彼此皆知；并且虽然每人不知任何其他一人所愿，但在有些事情上，他却能知道众人所愿。由此便产生了那位喜剧演员的诙谐笑话：他答应在另一出戏剧中，在剧场里说出众人心中所想、众人所愿之事；到了约定的日子，更多的人怀着极大的期待聚集而来，众人屏息静默，据说他说出了这样的话：\Quote{你们想贱买贵卖。}尽管他是个平庸的演员，但所有人在他的话中都认出了自己所意识到的真实情况，并以惊人的善意为他喝彩，因为他在众人眼前说出了人所共知、却无人预料到的实话。为何他承诺说出众人的意愿会引起如此巨大的期待？难道不是因为无人知道他人的意愿吗？但难道他不知道那个意愿吗？又有谁会不知道呢？然而，为何如此？除非是因为有些事情，每个人都可以通过同情或认同（无论是恶习还是美德）相当合理地从自身推测到他人身上。但看见自己的意愿是一回事；无论多么确切地猜测他人的意愿，则是另一回事。因为，在人间事务中，我确信\Place{罗马}城建立了，正如我确信\Place{君士坦丁堡}城建立了一样，尽管我亲眼见过\Place{罗马}，而对另一座城，除了相信他人的见证外，一无所知。事实上，那位喜剧演员相信，贱买贵卖是人人共有的意愿，这可能是通过观察自己或也通过试验他人而得出的。但由于这样的意愿实际上是一种缺陷，每个人都能获得相反的美德，或陷入另一种与它相反的缺陷的祸害，以抵抗并战胜它。因为我自己就知道一个例子：有人向一个人出售手稿，那人看出卖主不知其价值，因而要价甚低，于是他出乎卖主意料地付了公平的价钱，远比要价高。再设想一个人，其邪恶如此之大，以致贱卖父母留给他的财产，而贵买东西，以浪费在自己的情欲上。我想，这种放纵挥霍并非难以置信；如果寻找这样的人，可能会找到，甚至不找也会遇见。他们以比剧场更邪恶的邪恶，嘲笑了那个剧场的假设或宣告：他们高价购买耻辱，而低价出卖土地。我们也听说有人为了分发，以高价买进粮食，再以低价卖给同城居民。还要注意古代诗人恩尼乌斯所说的话：\Quote{所有凡人都希望自己被称赞}；其中，他无疑是根据自己和那些他凭经验认识的人，推测了他人的想法；因而似乎宣告了所有人所愿之事。最后，如果那位喜剧演员自己也说：\Quote{你们所有人都希望被称赞，没有一个人希望被辱骂}，他似乎在类似情况下表达了众人所愿。然而，有些人憎恨自己的过错，并不希望别人因为那些连自己也不喜欢的事而称赞他们；他们感谢那些责备他们的人的善意，如果责备的目的是要他们改正。但假如他说：\Quote{你们所有人都希望有福，不希望受苦}，他就说出了每个人都会在自己的意愿中认出来的事。因为无论一个人还可能私下里有什么别的意愿，他都不会脱离那个众人皆知、且众所周知存在于众人心中的意愿。\par

% structure: p0012 source=h2 target=subsection reason=relative-heading-rank; base=h2

\subsection{第四章——获得幸福的意愿在众人中是同一的，但关于幸福本身，各种意愿却极其多样。}

7. 然而奇妙的是，既然获享并保持幸福的意愿在众人中是同一的，那么另一方面，关于幸福本身，何以会产生如此多样和分歧的意愿呢？并非有人不愿拥有幸福，而是并非人人都认识它。因为如果人人都认识它，就不会有人认为幸福在于心灵的良善，有人认为在于身体的快乐，有人认为在于两者兼备，有人认为是这一样，有人认为是那一样。正如人们在某事或某事上感到特别的喜悦，便将幸福生活的观念寄托于此。那么，既然并非人人都认识幸福，又何以人人都如此热切地爱慕它呢？谁能爱慕他所不认识的东西呢？——这一问题我在前几卷中已经讨论过。因此，何以幸福为众人所爱，却不为众人所认识？或许是因为：众人知道幸福本身是什么，但不知道在哪里可以找到它，而争论正是由此产生？——就好像问题涉及世界上某个地方，每个愿意幸福生活的人都应当愿意在那里生活；也好像幸福在哪里这个问题并不隐含在幸福是什么之中。因为，确实，如果幸福在于身体的快乐，那么享受身体快乐的人就是幸福的；如果在于心灵的良善，那么享有这种良善的人就是幸福的；如果在于两者兼备，那么两者都享有的人就是幸福的。因此，当一个人说：\Quote{幸福生活就是享受身体的快乐}，而另一个人说：\Quote{幸福生活就是享受心灵的良善}时，难道两人都认识或都不认识什么是幸福生活吗？那么，倘若没有人能爱慕他所不认识的东西，两人又怎能爱慕它呢？或者，我们先前所假定为最真实、最确定的——即人人都愿意幸福生活——难道其实是假的？因为，假若幸福生活（为了论证起见）就是按照心灵的良善生活，那么那个不愿意这样生活的人又怎能愿意幸福生活呢？我们岂不更该说：那个人不愿意幸福生活，因为他不想按照良善生活，而只有按照良善生活才是幸福生活？因此，并非人人都愿意幸福生活；相反，只有少数人愿意，如果幸福生活无非是按照心灵的良善生活，而许多人不愿意这样做的话。那么，我们是否该认为连学园派的\Person{西塞罗}本人也不曾怀疑的（尽管学园派对一切事物都怀疑）那句名言是假的呢？他在\WorkTitle{荷尔滕西乌斯}对话录中，想要找某一确定无疑的东西作为论证的起点，他说：\Quote{我们确实都愿意幸福。}我绝不能说这是假的。但那么怎样呢？我们是否该说：尽管没有别的生活方式比按照心灵的良善生活更幸福，但连那个不愿意这样生活的人也愿意幸福生活？这似乎太荒谬了。这几乎等于说：\Quote{连那个不愿意幸福生活的人也愿意幸福生活。}谁能听得了、谁能忍受这样的矛盾呢？然而，如果以下两点都成立——人人都愿意幸福生活，然而并非人人都愿意以唯一能使生活幸福的方式生活——那么必然性就将我们逼入了这种困境。\par

% structure: p0014 source=h2 target=subsection reason=relative-heading-rank; base=h2

\subsection{第五章　论同一主题}

8. 也许，摆脱困境的办法在于：既然我们说过，每个人都把幸福生活的观念寄托在使他最感快乐的事物上——例如\Person{伊壁鸠鲁}寄于快乐，\Person{芝诺}寄于良善，其他人寄于别的什么——那么我们就说，幸福生活无非是按照自己的喜好生活。这样，\Quote{人人都愿意幸福生活}这句话就不假，因为人人都愿意自己所喜欢的东西。因为如果连这一点也在剧场中向众人宣告，众人都能在自己的意愿中找到它。但\Person{西塞罗}本人也曾将此作为反对自己的论点提出，并如此驳斥它，使那些这样想的人感到羞愧。他说：\Quote{但是，看啊！有些人虽然不是哲学家，却好争辩，他们说，凡按自己意愿生活的人都是幸福的。}——这也就是我们所说的\Quote{按各人所喜好}。但紧接着他又补充道：\Quote{然而，这确实是假的。因为意愿不当之事本身就是最不幸的；得不到所意愿之事固然不幸，但意愿得到不该意愿之事更为不幸。}他说得极为精彩，也完全正确。因为谁能在思想上如此盲目，如此远离一切正直之光，包裹于不正直的黑暗之中，以至于称那邪恶无耻地生活、无人约束、无人惩罚、甚至无人敢于责备、反而受到多数人称赞（正如圣经所言：\BibleQuote{恶人因心中的欲望受赞美，作恶的人被祝福}）的人为幸福呢？因为这样的人满足了所有最罪恶、最可耻的欲望；虽然即使如此，他仍然是可怜的，但如果他那些错误意愿所指向的事物一样也没有得到，他可能会更少可怜一些。因为每个人都会因邪恶的意愿而成为可怜的，即使意愿止于意愿，但若邪恶意愿的渴望因得到能力而满足，则更为可怜。因此，既然以下事实是真实的：人人都愿意幸福，并以最热烈的爱寻求这一件事，且为此寻求他们所寻求的一切；也没有人能爱慕一个他完全不知道是什么或什么样的事物，也不可能不知道他所知道所愿意的是什么；那么由此可推出，人人都认识幸福生活。但所有幸福的人都拥有他们所愿意的，尽管并非所有拥有所愿意的人立刻就是幸福的。那些要么没有他们所愿意的，要么拥有他们所不愿正当地意愿的人，立刻就是可怜的。因此，只有那既拥有一切他所愿意的，又无所恶意意愿的人，才是幸福的人。\par

% structure: p0016 source=h2 target=subsection reason=relative-heading-rank; base=h2

\subsection{第六章　为何当人人都愿意幸福时，反而选择了使人远离幸福的事物}

9. 既然真福生活包含这两件事，且为众人所知、为众人所爱；那么，我们当如何思量其原因？当人不能两者兼得时，从这两者中，他们宁愿选择拥有一切自己所愿意的，也不愿愿意一切事都善，即使他们并不拥有它们？这正是人类自身的堕落所致，以致他们并非不知道：那没有获得其所愿意的人不是有福的，那以错误的方式获得其所愿意的人也不是有福的，惟有那既拥有一切所愿意的善事、又不愿任何恶事的人才是真福的。然而，当构成真福生活的这两者不能同时赐予时，人选择了那更远离真福生活的一项（因为他得到他所邪恶地渴望的东西，比那未能得到他所渴望的更为远离真福）；然而良好的意愿更应被选择、更应被优先，即使它未能获得它所寻求的事物？因为那善愿一切他所愿的人，并愿那些一旦获得就能使他幸福的事物的人，已经接近成为有福的人。当然，使人有福的并非恶事，而是善事，当他们使人成为有福时。而在善事中，他已经拥有某种东西，而且并非微不足道的——即那良好的意愿本身；他渴望以人的本性所能达到的善事为乐，而不是以任何恶事的行为或成就为乐；他以明智、节制、勇敢和正直的心，殷勤追求并尽可能获得当前悲惨生活中可能得到的善事；以致在恶中为善，当一切恶都被终结、一切善都得以圆满时，他就有福了。\par

% structure: p0018 source=h2 target=subsection reason=relative-heading-rank; base=h2

\subsection{第七章 信德是必要的，好使人能在某个时候得到真福，这仅仅在来生才能达到。骄傲哲学家的真福是可笑的且可悲的。}

10. 为此，在这充满错误和艰难的必死生命中，信德，人藉以相信天主，是超越一切必需的。因为无论什么美善，尤其是那些使人成为善或使人将来成为有福的美善，除了从天主而来，没有任何其他来源能让人获得并增加。但是，当那在这苦难中既善且信的人从这生命进入福乐的生命时，那时将真正实现现在绝对不可能的事——即人可以按自己的意愿生活。因为在那种福乐中，他不会再愿意过坏的生活，也不会愿意有任何缺乏，也不会有任何他愿意的东西缺乏。凡是被爱的，都将临在；凡是不临在的，也不会被渴望。那里的一切都将美好，而至高的天主将是至高的美善，并为那些爱祂的人所享见；而最福乐的是，将确知如此直到永远。然而现在，哲学家们确实按照各人的喜好为自己制造了各自理想中的福乐生活，以便他们能藉着自己的能力，做那因凡人共同处境而无法做到的事——即按自己的意愿生活。因为他们感到，除非人拥有他所愿意的，并且不遭受他所不愿意的，否则无人能成为有福。谁不愿意那使他喜悦并因此称为有福的生命（无论它是什么）掌握在他自己的权下，以致他能持续拥有它呢？然而谁处于这种状况呢？谁愿意遭受苦难以便勇敢地忍受它们，尽管他愿意并且能够忍受它们（若他确实遭受了）？谁愿意生活在痛苦中，即使他能通过忍耐在痛苦中坚守正义而可赞美地生活？那些忍受了这些邪恶的人，或是出于希望拥有或害怕失去他们所爱的（无论是邪恶地还是可赞美地爱着），都认为这些邪恶是暂时的。因为许多人勇敢地穿过暂时的邪恶，走向那永恒的美善。这些人在希望中确实是有福的，即使在遭受这些暂时的邪恶时，他们通过这些邪恶到达那不会逝去的美善。但藉希望而有福的人尚未有福：因为他藉着忍耐等待那尚未把握的福乐。反之，那没有任何这样的希望、没有任何这样的奖赏而受苦的人，无论他运用多少忍耐，他不是真正有福的，而是勇敢地不幸的。因为他并不因此不是不幸的，因为如果他不能忍耐地承受不幸，他会更加不幸。此外，即使他不遭受在自己身体上不愿意遭受的那些事，他也不应被认为有福，因为他没有按自己的意愿生活。且不说其他事情——在身体未受伤害时，属于心灵烦恼的那些事情（没有这些我们仍愿意生活），并且数不胜数——他至少会愿意，如果他能，使自己的身体安全健康，不受任何不便，以致它完全在他自己的控制下，甚至拥有身体本身的不可朽坏性；因他没有获得这一点，且对此悬疑不定，他当然没有按自己的意愿生活。因为虽然他可能出于刚毅而预备接受并以平静的心承受可能发生的任何逆境，但他仍宁愿它们不发生，并尽量防止它们；他这样预备应对两种可能，即尽其所能，他愿望其一而避免另一；倘若他陷入他所避免的，他因而甘愿承受，因为所愿之事未能实现。因此他承受，为使自己不被压垮；但他不愿连受重压。那么，他如何按自己的意愿生活呢？是因为他甘愿坚强地承受他不愿加于己身的事吗？那么他只愿他能做到的，因为他不能拥有他所愿的。这就是骄傲凡人的福乐总纲：不知是该嘲笑，还是该怜悯，他们夸耀自己按意愿生活，因为他们甘愿忍耐地承受他们不愿发生的事。因为他们说，这就像\Person{泰伦提乌斯}的睿智之言——\par

\Quote{既然你所愿之事不能实现，就愿你所能之事。}\par

谁否认这话说得恰当呢？但这是给不幸之人的忠告，使他不再更加不幸。而对于福乐之人（所有的人都愿自己如此），则不能正确或真实地说：你所愿之事不能实现。因为如果他有福，他所愿的都能实现；因为他不会愿那不能实现之事。但这种生命不在今世必死的状态中，除非不朽也来到，否则它不会实现。如果这根本不能赐予人，那么福乐也是徒然寻求，因为没有不朽就不可能拥有福乐。\par

% structure: p0022 source=h2 target=subsection reason=relative-heading-rank; base=h2

\subsection{第八章 没有不朽就不能有福乐}

11. 因此，既然所有的人都愿意有福，那么，如果他们真正愿意，他们肯定也愿意不朽；否则他们就不能有福。但是，无论何种质量的真福，若不是真正的，而只是如此称呼的真福，在此生中被寻求，甚至被假装，同时不朽被绝望，而没有不朽，真正的真福就不能存在。因为，如我们先前已说过并充分证明和结论的，一个人生活有福，是当他如他所愿地生活，并且一无所愿于不当。但没有人不当愿不朽，如果人的本性因天主的恩赐而能拥有它；而如果它不能拥有它，它就不能拥有真福。因为，为了使人能有福地生活，他必须活着。如果生命因他的死亡而离他而去，有福的生命怎能与他同在？而当它离他而去时，无疑它要么是违背他的意愿离开他，要么是顺从他的意愿，要么是两者都不。如果违背他的意愿，那有福的生命如何能如此在他意愿之内却不在他能力之内？而既然没有人是有福的，如果他愿意某物却没有得到，那么当他不情愿地被他所爱的一切中最重要的——并非荣誉、财产或任何其他东西，而是有福的生命本身——所离弃时，他怎能是有福的呢？因为他将根本没有生命。因此，尽管没有感觉留下使他的生命因此悲惨（因为有福的生命离他而去，因为生命完全离他而去），但他只要还在感觉，就是悲惨的，因为他知道违背他的意愿，那为了它他爱一切其他东西、并且他爱它超过一切的东西正在被毁灭。因此，一个生命不能同时是有福的，却又违背人的意愿离他而去，因为没有人在违背自己意愿的情况下成为有福的；因此，违背他的意愿离他而去，使一个人变得悲惨，这比如果他违背意愿拥有它时更甚！但如果它顺从他的意愿离他而去，那么那有福的生命又怎能是那样：拥有它的人愿意它消亡？那么剩下要说的是，有福之人的意识中两者都没有；也就是说，他既不情愿也不情愿被有福的生命离弃，当死亡使生命完全离他而去时；因为他以平静的心坚定不移，对两种可能性都准备好。但那样的生命不是有福的生命，它所值得它所使之有福的人的爱。因为，那使有福之人不爱的，怎能是有福的生命？或者，那被漠不关心地接受（无论它是兴盛还是消亡）的，怎能被爱？或许，那些我们仅因真福而以这种方式爱的德行，竟敢说服我们不爱真福本身？但如果它们这样做，我们当然会停止爱德行本身，当我们不爱那因它而唯独爱它们的东西时。再者，那经过如此检验、筛选、彻底过滤并且如此确定的观点，\Emph{viz.}即所有的人都愿意有福，如果那些已经是真福的人既不意愿也不不意愿有福，这观点又怎能真实？或者，如果他们愿意有福（如同真理所宣告、本性所强制——那至善、不变而有福的造物主确实将这一意愿铭刻其中），那么，我说，如果他们愿意有福，他们当然不愿意不有福。但如果他们不愿意不有福，那么无疑他们不愿意在真福方面被消灭和消亡。但他们不能有福除非他们活着；因此他们不愿意在生命方面如此消亡。因此，无论谁，要么真正有福，要么渴望如此，都愿意不朽。但如果没有他所意愿的，他就不可能有福地生活。因此，结论是：生命除非是永恒的，否则决不能真正有福。\par

% structure: p0024 source=h2 target=subsection reason=relative-heading-rank; base=h2

\subsection{第九章——我们说未来的真福是真正永恒的，并非通过人的推理，而是藉着信德的帮助。真福的不朽因天主圣子的降生成人而成为可信的。}

12. 人性是否能够领受这被公认为可欲求的事物，这是一个不小的问题。但如果有信德存在——即在那些耶稣赐予他们权力成为天主儿女的人身上——那么便不成问题。确实，那些试图从人类理性中探究此事的人，仅有极少数，且是天赋异禀、闲暇充裕、学识极为精深者，才能独自达到对灵魂不朽的探究。即使对于灵魂，他们也未能找到稳固的、即真正的幸福生命；因为他们声称，即使在获得幸福之后，灵魂仍会回到今生的苦难中。他们中间那些对此见解感到羞愧，并认为被净化的灵魂应当无身体地置于永恒福乐中的人，其对世界过去永恒性的看法，足以驳斥他们关于灵魂的这种见解；这一点在此处详述过于冗长，但我想，我们在\WorkTitle{天主之城}第十二卷中已充分阐明了。然而，信德应许的——并非基于人类理性，而是基于神圣权威——是整个人，即由灵魂和身体组成的人，将是不朽的，并因此真正有福。因此，当福音中记载耶稣赐予那些接待祂的人\Quote{成为天主的儿女的权力}；并且通过说\Quote{就是那些信祂名字的人}简要解释了接待祂的含义；接着又补充了他们如何成为天主的儿女，即\Quote{他们不是由血统，也不是由肉情，也不是由人意，而是由天主而生}——以免我们大家都能看到并承受的人性软弱，对达到如此崇高的卓越感到绝望，同处又加上一句\Quote{圣言成了血肉，寄居在我们中间}；相反地，使人们确信那看似不可信的事。因为如果那本性是天主之子的，为了人的缘故，甘愿借怜悯降生为人子——这就是\Quote{圣言成了血肉，寄居在我们中间}之于人的含义——那么，本性为人子者，因天主的恩宠成为天主的儿女，并居住在天主内——唯有在祂内并出自祂，有福者才能分享那不朽——这有多么可信呢？为了使我们确信这一点，天主之子分担了我们的有死性？\par

% structure: p0026 source=h2 target=subsection reason=relative-heading-rank; base=h2

\subsection{第十章——没有比圣言降生成人更合适的方式来使人脱离有死性的苦难。所谓的我们的功勋，其实是天主的恩赐。}

13. 那些说：难道天主没有别的方式使人脱离这有死性的苦难，以至于祂愿意自己独生的、与祂自己同永恒的圣子，藉取人性灵魂与肉体而成为人，成为有死的并经历死亡？——我再说，仅仅这样反驳他们是不够的：即断言天主通过天主与人之间唯一的中保，基督耶稣这个人，所乐意拯救我们的方式，是好的，适合天主的尊严；我们还必须表明，并非没有其他方式为天主所可能（祂的权力对万物一律平等），而是没有任何其他方式比这更适合医治我们的苦难。因为对于建立我们的望德，并使因有死性本身而沮丧的必死之心从对不朽的绝望中获得释放，还有什么比向我们证明天主视我们为多么昂贵的代价，以及祂多么爱我们更为必要呢？而在这伟大的明证中，还有什么比这更显明昭著呢：天主子，不变地善，保持祂自身所是的，却从我们并为我们将祂所不是的领受过来，毫无损伤祂自己的本性，并屈尊进入与我们的共融；首先，毫无自身任何恶的功德，承受我们的恶；然后，以不受约束的慷慨，将祂自己的恩赐赠予我们——我们如今相信天主多么爱我们，并如今希望我们曾经绝望的事物，毫无我们自身的功德，甚至我们恶的功德先我们而行。\par

14. 既然那些被称为我们功绩的，也是祂的恩赐。因为，为使信德藉着爱德而工作，\BibleQuote{天主的爱，藉着所赐与我们的圣神，已倾注在我们心中了}。那时，当耶稣藉复活受光荣后，圣神便被赐下。因为那时祂应许要亲自派遣祂，并派遣了祂；正如经上所记载和预言祂的：\BibleQuote{祂上升高天，掳掠了俘虏，并将恩赐赐予人}。这些恩赐构成了我们的功绩，藉此我们到达不朽福乐的至善。可是，宗徒说：\BibleQuote{但天主却向我们证明了自己的爱，因为当我们还是罪人的时候，基督就为我们死了。那么，我们如今既因祂的血而成义，更要藉着祂脱免天主的义怒}。接着他又补充说：\BibleQuote{因为如果我们还在为仇敌的时候，因着祂圣子的死，与天主和好了；那么，在和好之后，我们更要因祂的生命得救了}。他起初称为罪人的，后来称为天主的仇敌；他起初说因祂的血而成义的，后来称为因天主圣子的死而和好的；他起初说藉祂脱免义怒的，后来要说因祂的生命得救。因此，我们在此之前不仅仅是某种罪人，而是如此之罪以致成了天主的仇敌。但同一宗徒在前面多次用两个称呼称呼我们，\Emph{即}罪人和天主的仇敌——一个似乎最温和，另一个显然最严厉——他说：\BibleQuote{因为当我们还是软弱的时候，祂就在预定时期为不虔敬的人死了}。他所称为软弱的，也正是他所称为不虔敬的。软弱似乎轻微；但有时它大到被称为不虔敬。然而除非它是软弱，它就不需要医生，这位医生在希伯来语是耶稣，在希腊语是\OriginalTerm{救主}{\GreekText{Σωτήρ}}，而在我们的语言中是救主。而拉丁语先前并没有这个词，但可以有，因为当需要时便能造出。宗徒上面的这句话，即\BibleQuote{因为当我们还是软弱的时候，祂就在预定时期为不虔敬的人死了}，与随后两句相连；其中一句论罪人，另一句论天主的仇敌，仿佛他将每句分别对应，\Emph{即}罪人对软弱者，天主的仇敌对不虔敬者。\par

% structure: p0029 source=h2 target=subsection reason=relative-heading-rank; base=h2

\subsection{第十一章——一个难题：我们如何因天主圣子的血而成义？}

15. 但所谓\BibleQuote{因祂的血而成义}是什么意思？请问，这血有什么力量，使相信的人因它而成义？所谓\BibleQuote{因祂圣子的死而与天主和好}又是什么意思？难道当真如此：天主圣父向我们发怒，看见祂圣子为我们而死，便对我们息怒了？难道那时祂圣子已经如此垂顾我们，甚至屈尊为我们而死，而圣父却仍然如此愤怒，除非祂圣子为我们而死，祂就不肯息怒？那么，外邦人的导师自己在另一处所说的话又如何理解：\BibleQuote{面对这些事，我们可说什么呢？若是天主偕同我们，谁能反对我们呢？祂既然没有怜惜自己的儿子，反而为我们众人把祂交出来，岂不也把一切与祂一同白白赐给我们吗？}请问，除非圣父早已和解，祂会交出自己儿子而不怜惜祂吗？这观点难道似乎与那个相反吗？在那处，圣子为我们而死，圣父因祂的死而与我们和好；在这处，仿佛圣父先爱了我们，祂为了我们不惜自己的儿子，祂为我们将祂交出至死。但我知道，圣父也早已爱了我们，不仅是在圣子为我们死之前，而且在创造世界之前；宗徒亲自作证说：\BibleQuote{就如祂在创世以前，在祂内拣选了我们}。圣子为我们被交出也不是不情愿的，圣父自己并不怜惜祂；因为论到祂也说：\BibleQuote{祂爱我，并为我舍弃了自己}。所以，圣父、圣子和两者的圣神共同平等和谐地行事；然而我们是在基督的血中成义，并因祂圣子的死与天主和好。现在，我将尽我所能，在此也解释这事是如何成就的，依照可能足够的程度。\par

% structure: p0031 source=h2 target=subsection reason=relative-heading-rank; base=h2

\subsection{第十二章——众人因亚当的罪都被交在魔鬼的权下。}

16. 因着天主的公义，从某种意义上说，人类被交付于魔鬼的权势；原祖的罪通过婚姻结合，在诞生时原样地传给了所有两性的人，我们原祖的债务束缚了他们所有的后代。这种交付首先在\BibleBook{}{创世纪}中预示，当时对蛇说：\BibleQuote{你必吃土}，对人说：\BibleQuote{你原是土，还要归于土}。在\BibleQuote{你必归于土}这句话中，身体之死被预告，因为如果他始终保持受造时的正直，他也本不会经历死亡；但对他活着时所说的\BibleQuote{你原是土}，则表明整个人变坏了。因为\BibleQuote{你原是土}与\BibleQuote{我的神不永远住在这些人里面，因为他们也是肉体}几乎相同。因此，当时表明他被交付给他，是因为对他说了\BibleQuote{你必吃土}。但宗徒更清楚地宣告了这一点，他说：\BibleQuote{你们从前因过犯和罪过是死的，那时你们在过犯中生活，随从这世界的潮流，顺服空中权势的首领，就是现今在悖逆之子身上运行的灵；我们从前也都在他们中间，放纵肉体的私欲，随着肉体和心中所喜好的去行，本为可怒之子，和别人一样。}\Quote{悖逆之子}就是不信者；在成为信徒之前，谁不是这样呢？因此，所有人原本都在空中权势的首领之下，\BibleQuote{他现今在悖逆之子身上运行}。而我用\Quote{原本}所表达的，正是宗徒说他们自己\Quote{在本性上}和别人一样时所表达的；\Emph{即}本性已被罪恶所败坏，而非如起初被造为正直的那样。但人这样被交付于魔鬼的权势，不应理解为天主如此行或命令如此行，而是祂只准许了这事，且是公义地准许。因为当祂弃绝罪人时，罪的作者立即进入。然而天主并未如此弃绝祂自己的受造物，以致不向祂显示为创造和给予生命的天主，并在惩罚的恶中也将许多善物赐给恶人。因为祂未曾发怒封闭祂的仁慈。祂也未将人从祂自己能力的法则中释放，当祂准许人处于魔鬼的权势之下时；因为连魔鬼自己也没有与全能者的能力分离，也没有与祂的良善分离。因为甚至恶天使无论拥有怎样的生命方式，不也都是通过那赋予一切生命者而存在的吗？因此，如果罪的犯下因天主公义的愤怒使人隶属于魔鬼，那么罪的赦免因天主仁慈的和解无疑将人从魔鬼手中救出。\par

% structure: p0033 source=h2 target=subsection reason=relative-heading-rank; base=h2

\subsection{人应当从魔鬼的权势中被救出，不是凭大能，而是凭公义。}

但是，魔鬼是要被战胜的，不是靠着天主的权能，而是靠着祂的正义。因为有什么比全能者更有权能呢？或者有什么受造物，其能力可与造物主的能力相比呢？但由于魔鬼因自己的乖戾之罪，成了权能的爱好者，正义的抛弃者和攻击者——因为人也如此效法他，愈是倾心于权能，忽视甚至仇恨正义，愈是因获得权能而欢喜或因贪恋权能而焦躁，其效法程度就愈甚——天主喜欢这样：为了将人从魔鬼的掌握中解救出来，魔鬼应被战胜，不是靠权能，而是靠正义；如此，人效法基督，也当以正义而非权能来战胜魔鬼。并不是说权能应被当作邪恶而回避，而是秩序必须保持，正义优先于权能。因为凡人的权能能有多大呢？因此，让凡人持守正义吧；权能将赐予不朽者。与此相比，那些在地上被称为有势者的人，无论其权能多么巨大，都显得可笑而软弱，在那里为罪人挖了陷阱，正是邪恶者似乎最有权势的地方。义人在诗歌中唱道：\BibleQuote{上主，祢所惩戒并用祢的法律教训的人，是有福的：使他在祸患的日子得享安宁，直到为恶人挖了陷阱。因为上主不会抛弃祂的人民，也不会遗弃祂的产业，直到正义归于审判，凡心里正直的人都跟随它。}\BibleRef{咏 94:12-15} 因此，在目前这个时期，当天主子民的强力被延迟时，\BibleQuote{上主不会抛弃祂的人民，也不会遗弃祂的产业}，无论它在谦卑和软弱中遭受多么痛苦和不堪的事情；\BibleQuote{直到正义}——就是虔诚者的软弱现在所拥有的——\BibleQuote{归于审判}，也就是说，要接受审判的权能；这权能在最后为义人保留，那时权能将在正义先行之后，按适当的次序跟随。因为权能与正义结合，或正义加于权能，就构成审判的权威。但正义属于善意的意志；因此当基督诞生时，天使们说：\BibleQuote{天主在天受光荣，主爱的人在世享平安。}\BibleRef{路 2:14} 但权能应当跟随正义，而不是先行于它；因此它被置于\Quote{第二等}，即顺遂的福分；这被称为\Quote{第二等}，源自\Quote{跟随}。因为正如我们之前所论证的，使人有福的是两件事：愿意善和能够行所愿，人不应当如此乖戾（正如同一讨论中指出的），从使人有福的两件事中选择能够行其所愿，而忽略愿其所当行；其实他应当先有善意，然后有大能。此外，善意必须从恶习中净化，人若被恶习战胜，他就这样被战胜以至于他愿意邪恶；那么他的意志怎么还是善的呢？因此，现在但愿权能赐予我们，但权能是用来战胜恶习的，人并不希望有战胜恶习的权能，反而希望有战胜人的权能；这是为什么？岂不是因为他们实际上被战胜，却虚伪地战胜，成为名不副实的胜利者吗？但愿人愿意明智，愿意坚强，愿意节制，愿意公正；并且为了能够真实地拥有这些，他当然要渴望权能，并寻求对自己有权能，而且（虽然奇怪）为了自己而与自己抗争。但是其他一切他正义地渴望却未能拥有的东西，例如不朽和真实圆满的幸福，愿他不要停止渴望，并耐心期待。\par

% structure: p0035 source=h2 target=subsection reason=relative-heading-rank; base=h2

\subsection{第十四章 基督无偿的死亡释放了那些应死者}

18. 那么，那使魔鬼被征服的义是什么？除了耶稣基督的义，还有什么？祂是如何被征服的？因为，当魔鬼在祂身上找不到任何该死之处，却仍杀了祂。这无疑是公义的：我们这些被魔鬼视为债务人的人，因相信这被祂无辜杀害的那一位，而得释放。如此，我们被称为在基督的血中成义。因为那无辜的血为赦免我们的罪而流。因此祂在圣咏中自称\BibleQuote{在死人中自由。}因为只有死了的人才免于死亡的债务。因此在另一篇圣咏中祂说：\BibleQuote{那时我偿还了我未曾抢夺的；}祂以\OriginalTerm{抢夺的}{seized}指罪恶，因为罪恶是非法夺取的。因此，祂也藉祂肉身的口说出，正如福音所载：\BibleQuote{因为这世界的首领来了，在我身上一无所有，}即没有罪；但祂说：\BibleQuote{为使世界知道，我遵行父的命令；起来，我们走吧。}于是祂走向苦难，为偿还我们这些债务人祂自己所不欠的。那么，如果基督愿意用权能而非公义与魔鬼打交道，魔鬼会被这最公正的权利征服吗？但祂保留了祂所能做的，为先做应做的事。因此祂必须同时是人和天主。因为如果祂不是人，祂就不能被杀；如果祂不是天主，人就不会相信祂不是不能做祂愿做的，而是不愿做，也不会认为祂将公义置于权能之上，而是缺乏权能。但现在祂为我们承受了属于人的事，因为祂是人；但如果祂不愿意，祂本有能力不受这样的苦，因为祂也是天主。因此，公义在谦卑中更加可亲，因为祂的天主性中有如此大的权能，若祂不愿意，本可不承受谦卑；这样，藉这如此有能力的祂的死亡，公义被赞扬，权能被应许给我们软弱的凡人。因为祂藉死亡做了这两件事中的一件，藉复活做了另一件。因为有什么比为了公义甚至死在十字架上更公义的呢？有什么比从死者中复活，并带着被杀的肉身升天更有权能的呢？因此祂首先以公义征服了魔鬼，然后以权能：即以公义，因为祂没有罪，却被魔鬼最不义地杀害；以权能，因为祂死后又活了，从此不再死。但即使祂没有被魔鬼杀害，祂也会以权能征服魔鬼；虽然透过复活征服死亡本身比通过活着避免死亡属于更大的权能。但我们因基督的血而成义，当我们藉罪赦从魔鬼的权下被解救时，原因其实不同：这属于魔鬼被基督以公义而非权能征服。因为基督被钉十字架，不是由于不朽的权能，而是由于祂在必死肉身上所取的软弱；然而宗徒论到这软弱说：\BibleQuote{天主的软弱总比人强壮。}\par

% structure: p0037 source=h2 target=subsection reason=relative-heading-rank; base=h2

\subsection{论同一主题。}

19. 看见魔鬼在被他所杀的那位复活时被征服，这并不困难。而看见魔鬼在他自以为得胜时，即当基督被杀时，就被征服了，则是一件更深刻、更难理解的事。因为那时，那属于毫无罪恶之祂的血，为赦免我们的罪而倾流；因魔鬼理应束缚那些作为罪人而被死亡条件所捆绑的人，他也理应藉那毫无罪恶却无辜承受死亡惩罚的那位而释放他们。这义征服了那壮士，并用这锁链束缚他，使他掠夺的器皿——这些器皿连同他自己和他的天使，曾是忿怒的器皿——得以被夺回，变为怜悯的器皿。因为宗徒保禄告诉我们，这些话是主耶稣基督自己在最初蒙召时从天上传给他的。在他所听到的其他事中，他也说到有这话指示他：\BibleQuote{我显现给你，正是为了这个缘故：要拣选你作仆人和见证人，对于你由我所见的事，以及我将来要显示给你的事；我要救你脱离这百姓和外邦人的手，我现在打发你到他们那里去，开他们的眼，叫他们从黑暗中转向光明，并从撒殚的权力下归向天主，好使他们因信我而获得罪赦，并在圣化的众人中蒙受产业。}因此，同一宗徒也劝勉信徒感谢天主父说：\BibleQuote{祂拯救了我们脱离黑暗的权势，把我们迁到祂爱子的国度里；我们藉祂的爱子获得救赎，就是罪过的赦免。}在这救赎中，基督的血被赐下，如同为我们付上的赎价；魔鬼藉此并非被致富，而是被束缚：好使我们脱离他的锁链，他不能将那些基督——无所负债者——藉无偿倾流自己的血所赎买的人，与他一同卷入罪恶的网罗，并交付给第二次的、永久的死亡；而那些属于基督恩宠的人，在创世以前就被预知、预定、蒙拣选，只当如基督为他们而死一样死去，即只经历肉身的死亡，而非灵魂的死亡。\par

% structure: p0039 source=h2 target=subsection reason=relative-heading-rank; base=h2

\subsection{死亡的残余与世界之恶转为选民之善。基督之死如何被合宜地拣选，使我们因祂的血而成义。天主的义怒是什么。}

20. 因为，虽然肉身本身的死亡最初来自第一个人的罪，但对此死亡的善用却造就了最光荣的殉道者。因此，不仅死亡本身，而且这世界的一切灾祸、人的忧伤与劳苦，虽然都来自罪的惩罚，尤其是原罪——因这原罪，生命本身也被死亡的锁链束缚——但在罪获得赦免之后，仍适当地保留着，好让人有可为之奋斗的真理，并使信友的善行得以锻炼；以致新人藉着新盟约，在这世界的灾祸中，为新的世界做好准备，明智地承受这被定罪的生命所应得的苦难，清醒地因它将结束而喜乐，并忠诚而耐心地期盼那在未来的生命中，获得释放后将永远享有的福乐。因为魔鬼被逐出他的领域，也离开信友的心——在他被定罪和不信的人心中，他自己虽已被定罪，却仍在那里统治——只按此有死之身的状况被允许作为仇敌，一如天主知道为他们有益的那样。关于这一点，圣经藉宗徒的口说：\BibleQuote{天主是信实的，必不叫你们受试探超过你们所能受的；但在受试探的同时，也必开一条出路，使你们能忍受得住。}信友虔诚忍受的那些灾祸，或有助于纠正罪过，或有助于锻炼和证明义德，或彰显此生的苦难，使那真正永恒福乐的生命被更热切地渴望、更努力地寻求。但这些灾祸之所以仍在持续，是为了那些人的缘故；关于这些人，宗徒说：\BibleQuote{而且我们也知道：天主使一切协助那些爱他的人，就是那些按他的旨意蒙召的人，获得益处。因为他所预知的人，也预定他们与他儿子的肖像相同，好使他在众多弟兄中作长子。他所预定的，也召叫了他们；他所召叫的，也使他们成义；他所成义的，也使他们受光荣。}在这些预定的人中，没有一人会与魔鬼一同丧亡；没有一人会在魔鬼的权势下甚至至死。接着便是如上引述的话：\BibleQuote{面对这一切，我们可说什么呢？若是天主偕同我们，谁能反对我们呢？他没有怜惜自己的儿子，反而为我们众人把他交出来，岂不也把一切与他一同赐给我们吗？}\par

21. 那么，为何基督的死亡不应发生？不，更确切地说，为何不应当选择这死亡本身，超越其他无数的、全能者可以用来解放我们的方式，而使之成就？我说的是，这死亡既未减少或改变祂的天主性，又如此大地有益于人——从祂所取的人性而来——即永恒的天主子，也是人子，献上了一个暂时的、本不该受的死亡，借此祂可以使他们脱离那本应受的永远死亡。魔鬼曾紧紧抓住我们的罪，并藉此将我们应得地固定在死亡中。祂（基督）没有自己的罪，却由魔鬼带领，无辜地受死。那血的代价是如此之大，以至于那曾暂时杀死基督（以一种不该受的死亡）的魔鬼，不能将任何一个穿上基督的人，应得地扣留在那本应受的永远死亡中。因此，\BibleQuote{天主就在我们还作罪人的时候，为我们死了，藉此证实了祂对我们的爱。何况现在，我们既已因祂的血而成义，更要藉着祂得救，免去义怒。} 他（保禄）说，因祂的血而成义——显然，成义就是我们从一切罪中得到释放；而从一切罪中得到释放，是因为天主子，那位不知罪的，为我们被杀。因此，\BibleQuote{我们要藉着祂得救，免去义怒。} 这义怒当然是指天主的义怒，无非是公义的报应。因为天主的义怒不像人的义怒，不是心灵的激动；而是祂的义怒，圣经在另一处向祂说：\BibleQuote{但是，主，祢掌握着祢的能力，以宁静施行审判。} 因此，如果天主的公义报应有了这样的名称，那么对天主的和好，正确的理解难道不是这种义怒的终结吗？我们成为天主的仇敌，除非因为罪与义为敌？这些罪过被赦免，这样的敌对就终结了，而那些祂亲自称义的人，就与那位公义者和好了。然而，祂当然仍爱他们，即使他们还是仇敌的时候，因为\BibleQuote{祂没有怜惜自己的儿子，反而为我们众人把祂交出，} 那时我们还是仇敌。因此，宗徒正确补充道：\BibleQuote{如果我们还是仇敌的时候，藉着祂儿子的死，得与天主和好，} 藉此，罪过得赦免，\BibleQuote{那么，和好之后，我们更要藉着祂的生命得救。} 藉着死亡和好，却要在生命中得救。因为谁能否认，祂会为朋友赐下祂的生命？祂曾为仇敌赐下祂的死亡！\BibleQuote{不仅如此，} 他说，\BibleQuote{我们还藉着我们的主耶稣基督，因着祂，我们现在已领受了和好，而以天主为乐。} \BibleQuote{不仅，} 他说，\BibleQuote{我们要得救，} 而且\BibleQuote{我们还以……为乐；} 不是在我们自己内，而是\BibleQuote{在天主内；} 也不是藉着我们自己，\BibleQuote{而是藉着我们的主耶稣基督，因着祂，我们现在已领受了和好。} 正如我们上文所论述的。然后宗徒补充道：\BibleQuote{因此，就如罪恶藉着一人进入了世界，死亡藉着罪恶也进入了；这样死亡就殃及了众人，因为众人都犯了罪；} 等等；他在此相当长地论述了两个人：第一个亚当，藉着他的罪与死亡，我们他的后裔被遗传的邪恶所束缚；第二个亚当，祂不仅是人，也是天主，藉着祂为我们付清祂所不欠的，我们从先祖和我们自己的债务中获得释放。此外，既然因着那一个人，魔鬼持有所有通过他败坏的肉体贪欲而生的人，那么因着这一个人，魔鬼理当释放所有通过祂无玷的属灵恩宠而重生的人。\par

% structure: p0042 source=h2 target=subsection reason=relative-heading-rank; base=h2

\subsection{第十七章——降生成人的其他益处。}

22. 在基督的降生成人中，还有许多其他事，虽然骄傲者不悦，却应当被观察并有益地思考。其中之一是，它向人展示了人在天主所造之物中的地位，因为人的本性竟能与天主如此结合，以致两性可以成为一位格，实际上更是三性——天主、灵魂与肉身——成为一位格：因此那些骄傲的恶灵，它们假装帮助人，实则欺骗，却不敢因为人没有肉身而自居于人之上；尤其因为天主子也屈尊在这同一肉身中死去，以免它们因似乎不死而得以被当作神崇拜。此外，天主的恩宠在基督的人性中，毫无先行功劳地传达给我们；因为甚至祂自己也没有任何先行功劳而获得与真实天主如此伟大的合一，并成为天主子，与祂同一位格；而是从祂开始成为人之时起，祂也是天主；因此经上说：\BibleQuote{圣言成了血肉。} 还有，人的骄傲（这是人依附天主的主要障碍）能够通过天主如此伟大的谦卑而被驳斥和治愈。人也懂得自己远离了天主多远；以及当他通过这位中保归来时，疼痛对他有何价值；这位中保以祂的天主性助人，又以祂的人性与人相同。因为有什么更大的服从榜样能赐给我们，这些因不顺从而丧亡的人，比天主子服从天主圣父，以至于死在十字架上更大呢？再者，如何能更好地显示服从本身的赏报，乃是在这位伟大中保的肉身中，它复活得到永生？这也属于造物主的公义和良善：魔鬼应当被同一理性的受造物征服，这受造物曾因征服人而欢喜，并且这受造物出自那同一族类，这族类因一人的起源腐败而被魔鬼完全控制。\par

% structure: p0044 source=h2 target=subsection reason=relative-heading-rank; base=h2

\subsection{第十八章——为何天主子从亚当的后裔并自童贞玛利亚取了人身。}

23. 因為的確，天主本可以承擔為人，在那人性中成為天主與人之間的\OriginalTerm{中保}{Mediator}，從另一個源頭而來，而不是從那以自己罪惡束縛人類的亞當的種族而來；正如祂最初創造的那個人，並非從別人的種族中創造的一樣。所以，祂能夠，要麼如此，要麼以祂願意任何其他方式，創造另一個（人），使第一個征服者被其征服。但天主判定，從已被征服的種族本身中，承擔一個人來征服人類的仇敵，這樣做更好；並且要由一位童貞女來完成，她的懷孕並非由肉體，而是由精神，並非由情慾，而是由信德所先導。那因之使其餘的人——他們承襲原罪——得以繁衍和受孕的\OriginalTerm{肉慾}{carnal concupiscence}，並未介入；而是聖潔的童貞懷孕了，不是通過婚姻交合，而是通過信德——情慾完全缺席——使那從第一人根源所生者，僅承襲種族的起源，而非罪責。因為所出生的，不是被過犯的傳染所敗壞的本性，而是所有這些敗壞的唯一補救。所出生的，我說，是一個人，完全沒有罪，並且永遠沒有罪；通過祂，那些生來不能無罪的人，將得以重生，從罪中獲得自由。因為，雖然婚姻的貞潔正當地運用我們肢體中的肉慾，但它仍會受非自願的衝動的影響，由此表明，要麼在罪之前在樂園中它根本不可能存在，要麼即使存在，也不曾有時會抵抗意志。但現在我們感覺到它如此這般，違反心神的法律，甚至在不涉及生育時，也在我們內激起性交的衝動；如果人在此屈服於它，則它通過一個罪惡行為得到滿足；如果不屈服，則通過一個拒絕行為加以約束：這兩者，誰能懷疑在罪前的樂園中是不可能的呢？因為那時的貞潔不會做任何有失體面的事，那時的愉悅也不會遭受任何騷動。因此，當童貞女的子嗣受孕時，這肉慾必須完全缺席；在其中，死亡之君找不到任何該當死亡的事，卻仍要殺害祂，為能藉生命之君的死亡而被征服；第一亞當——他緊握人類——的征服者，被第二亞當征服，失去基督徒的種族，從人類中釋放了那些因人的罪孽而獲自由的人，藉著那雖在種族內卻不在罪孽中的祂；為使那欺騙者被那曾被他以罪孽征服的種族所征服。這樣做是為了使人不得自高自大，而\Quote{誇耀的，應因主而誇耀。}因為被征服的只是人；他之所以被征服，是因為他驕傲地貪求成為神。但征服者是神而人；因此祂如此征服了，生於童貞女，因為天主在謙卑中，並非如祂管理其他聖人那樣管理那個人，而是親自懷有祂（作為兒子）。天主的這些如此偉大的恩賜，以及其他所有在此我們無法詳細探究和討論的，除非聖言成了血肉，否則都不能存在。\par

% structure: p0046 source=h2 target=subsection reason=relative-heading-rank; base=h2

\subsection{第十九章——在降生成人的聖言中，什麼屬於知識，什麼屬於智慧。}

\OriginalTerm{降生成人的圣言}{Word made flesh}所做、为我们承受的一切，发生在时间和地点之中，按照我们所致力区分的区别，属于知识，不属于智慧。正如圣言无时间无地点，与圣父\OriginalTerm{同永恒}{co-eternal}，在其整体中无所不在；如果有人能够，并尽可能，真实地谈论这圣言，那么他的论述将属于智慧。因此，降生成人的圣言，就是基督耶稣，拥有智慧和知识的宝藏。因为宗徒写信给哥罗森人说：\BibleQuote{我愿意你们知道我为你们和为在劳狄刻雅的人，以及为一切没有亲眼见过我面的人，有多么大的奋斗，为使他们的心受到安慰，在爱德中互相结合，以获得理智的一切富足，得以认识天主的奥秘，就是基督耶稣，在他内蕴藏着智慧和知识的一切宝藏。}宗徒知道所有这些宝藏达到什么程度？他探入其中多少？在其中达到何等伟大之事？谁能知道？然而，按照所记载的：\BibleQuote{圣神的表现赐给各人，是为人的益处；藉圣神，有人得了智慧的言语，有人得了知识的言语，都出于同一圣神；}如果这两者如此区别：智慧归于神圣之事，知识归于人性之事，那么我承认两者都在基督内，所有祂的信徒也与我一样。当我读到：\BibleQuote{圣言成了血肉，居住在我们中间，}我藉圣言理解天主的真子，在血肉中承认人的真子，两者因恩宠的不可言喻的丰盈而结合为一神人位格。因此，宗徒接着说：\BibleQuote{我们见了祂的光荣，如同父独生者的光荣，满溢恩宠和真理。}如果我们把恩宠归于知识，把真理归于智慧，我认为我们就不会偏离我们所推荐的这两者之间的区别。因为在那起源于时间的事物中，这是最高的恩宠：人与天主结合于位格的合一；而在永恒事物中，最高的真理正确地归于天主的圣言。然而，这位本身就是\OriginalTerm{圣父的独生子}{Only-begotten of the Father}，满溢恩宠和真理——这事发生了，为的是使祂自己在为我们在时间中所行的事中，成为我们藉同一信仰而被洁净的对象，好使我们能在永恒事物中坚定地默观祂。那些外邦著名的哲学家，能够藉所造之物理解并辨别天主不可见的事，但如经上论他们所说的：\BibleQuote{以不义抑制了真理}因为他们没有通过中保而哲思，即没有通过那人基督；他们既不信仰祂会按先知的话而来，也不相信祂已按宗徒的话而来。因为他们处在这最低下的事物中，不得不寻求某种媒介，藉以通达他们所理解的那些崇高事物；于是他们落入了欺骗的灵，藉此发生了：\BibleQuote{他们将不朽天主的荣耀变为像可朽的人、飞鸟、四足兽和爬虫的形像}因为他们也以这些形像设立或敬拜偶像。因此，基督是我们的知识，同一基督也是我们的智慧。祂亲自在我们内灌输有关暂时之事的信仰，祂亲自显示有关永恒之事的真理。藉着祂，我们达到祂自己：我们通过知识伸展到智慧；但我们没有离开同一基督，\BibleQuote{在祂内蕴藏着智慧和知识的一切宝藏}但现在我们谈论知识，以后将谈论智慧，按祂自己所赐的。但我们不要如此理解这两者，好像谈论人性事物中的智慧或神圣事物中的知识是不允许的。因为按照较宽松的说话习惯，两者都可称为智慧，两者都可称为知识。然而，宗徒绝不可能写下\BibleQuote{有人得了智慧的言语，有人得了知识的言语}除非这些不同的事物各自被恰当地以不同的名称称呼，我们正在讨论其区别。\par

% structure: p0048 source=h2 target=subsection reason=relative-heading-rank; base=h2

\subsection{第20章——本书所论内容。我们如何逐步达到在实践知识和真信仰中所发现的某种\OriginalTerm{天主圣三}{Trinity}。}

25. 现在，因此，让我们看看这篇冗长的论述成就了什么，收集了什么，达到了什么。所有的人都愿意得真福；然而并非所有的人都有信德，藉以心灵得以洁净，从而获得真福。因此，藉着并非所有的人都愿意的信德，我们必须达到每一个人都愿意的真福。所有人在自己心中都看到他们愿意得真福；在这一点上人性的共识如此之大，以至于一个人参照自己的心去揣测别人的心，是不会受骗的：简而言之，我们知道我们所有人都愿意这一点。但许多人对不朽感到绝望，尽管不如此，任何人都不能成为所有的人都愿意的，即真福。然而他们如果能够的话，也愿意不朽；但由于不相信他们能够，他们并没有那样生活以便能够。因此信德是必要的，好使我们能在人性的一切美善中，即灵魂和身体两方面，获得真福。但同样的信德要求这信德局限于基督内，祂曾带着肉身从死者中复活，不再死亡；并且除非藉着祂，没有人能通过罪过的赦免从魔鬼的统治下获得自由；并且，在魔鬼的居所里，生命必然同时是痛苦和永无止境的，这与其称为生命，不如称为死亡。所有这一切我也在这本书中尽可能论述了，同时我也在这部作品的第四卷中对此说了很多；但在那里是为了一个目的，在这里是为了另一个——即在那里，是为了说明为何以及如何基督在时期一满时由圣父派遣，针对那些说派遣者与被派遣者在本质上不能平等的人；但在这里，是为了区分实践的认知与默观的智慧。\par

26. 因为，我们愿意仿佛逐步上升，在内在的人身上，无论是在认知还是在智慧方面，寻找一种它自己特有的三位一体，正如我们以前在外在的人身上寻找的那样；好使我们能以在这些低级事物上更熟练的心智，按照我们微小的程度，来默观那作为天主的圣三，如果我们甚至能够做到这一点的话，至少是在谜语中，如同透过镜子。如果，那么，有人仅仅凭声音记住了这信德的话语，而不明白它们的意思，就像那些不懂希腊文的人记住希腊文，或者类似地拉丁文，或其他任何他们不懂的语言，他心中是否有一种三位一体呢？因为，首先，那些话语的声音是在他的记忆中，即使他不去想它们；其次，当他想象这些时，他回忆行为的精神视域（\OriginalTerm{眼力}{acies}）由此形成；第三，那位回忆和思考的人的意志将两者结合。然而，我们决不能因此说这个人是在从事内在之人的三位一体，而是外在之人的三位一体；因为他所记忆的，并且当他想时，他尽可能默观，仅仅属于身体的感觉，即听觉。在这样的思维活动中，他除了处理有形事物，即声音的影像之外，什么也不做。但是，如果他记住并回忆那些话语所表示的意义，那么现在确实内在之人的某些东西被引入了行动；然而，如果他并不爱那些在那里被宣告、命令、应许的事物，他还不应该被称为或被认为按照内在之人的三位一体生活。因为他也有可能记住和想象这些事情，假设它们是假的，以便努力反驳它们。因此，那种意志，在这种情况下，将记忆中持有的东西与从那里印在心灵之眼上的东西在概念中结合起来，确实完成了一种三位一体，因为它是加给另外两者的第三个；但是当所设想的被认为假而不被接受时，人并不按照这个三位一体生活。但是，当那些事物被相信为真，并且其中应被爱的被爱时，那时人才最终按照内在之人的三位一体生活；因为每个人都是按照他所爱的而生活。但是，那些不为人所知而仅为人所信的事物如何能被爱呢？这个问题在前几卷中已经讨论过了；我们发现，没有人爱他完全不知道的东西，但是当说到未知的事物被爱时，它们是从已知的事物中被爱的。现在我们这样结束本书：我们劝告义人因信德而生活，这信德以爱德行事，这样，使人明智、勇敢、节制、公义地生活的那些美德本身，全都归于同一信德；因为否则它们不能成为真美德。然而，在今生，这些美德的价值并不如此之大，以至于某些罪过的赦免在这里有时不是必要的；而这种赦免，除非通过那位以自己的血战胜了罪人之王的祂，是绝不会发生的。无论什么样出自这信德和如此生活的观念在信友心中，当它们被保存在记忆中，被回忆所审视，并令意志喜悦时，便呈现一种它自己特有的三位一体。但是天主的肖像，我们将靠祂的助佑之后谈论，尚未在那三位一体中；这一点在它被指明在哪里时将会更加明显，读者可在下一卷中期待。\par

\SourceRecord{Translated by Arthur West Haddan.；Nicene and Post-Nicene Fathers, First Series；Vol. 3.；Edited by Philip Schaff.；Buffalo, NY: Christian Literature Publishing Co.,；1887.；<http://www.newadvent.org/fathers/130113.htm>.}

% structure: p0001 source=h1 target=section reason=page-title

\section{论天主圣三（第十四卷）}

\EditorialNote{此处论及人的真智慧；并指出人按其心灵乃是天主的肖像，但这肖像并非恰当地置于短暂的事物中，如记忆、理解和爱——无论是存在于时间中的信仰本身，还是心灵忙于自身——而是置于永恒的事物中；并且当心灵按照那按自己的肖像创造人的那一位的形象，在认识天主的知识中得以更新时，这智慧得以成全，从而达至智慧，其中所默观的是永恒的。}\par

% structure: p0003 source=h2 target=subsection reason=relative-heading-rank; base=h2

\subsection{第一章。——论我们将要讨论的智慧是什么。哲学家名称的起源。关于知识与智慧的区别已经说过的内容。}

1. 我们现在必须论及智慧；并非论及天主的智慧——那无疑就是天主，因为祂的独生子被称为天主的智慧——而是论及人的智慧，却是真智慧，即按照天主的智慧，是祂真实而首要的敬拜，这在希腊文中以单数术语\OriginalTerm{敬拜天主}{\GreekText{θεοσέβεια}}表示。正如我们已观察到的，当我们的同胞自己也希望用一个术语来翻译它时，他们将它译成了\Quote{虔诚}，而\Quote{虔诚}更通常是指希腊人所说的\OriginalTerm{虔诚}{\GreekText{εὐσέβεια}}。但由于\OriginalTerm{敬拜天主}{\GreekText{θεοσέβεια}}无法用一个词完美翻译，最好用两个词翻译，即译为\Quote{敬拜天主}。这就是人的智慧，正如我们在本作品的第十二卷中已阐明的，圣经的权威也证明了这一点：在天主的仆人约伯的\BibleBook{约伯}{传}上，我们读到天主的智慧对人说：\BibleQuote{看，虔诚就是智慧；远离邪恶就是知识}；或者，如有些人翻译希腊词\OriginalTerm{知识}{\GreekText{ἐπιστήμην}}，译为\Quote{学问}，它当然来自\Quote{学习}一词，因此也可称为知识。因为一切事物都是为了被认识而学习的。尽管同一个词的确也有不同的用法，如当一个人因自己的罪遭受苦难以便被纠正时。这可见于\BibleBook{希伯来人}{书}中：\BibleQuote{因为哪一个儿子，父亲不惩戒他呢？}；同书更为明显：\BibleQuote{凡惩戒的事，当时似乎不是喜乐，而是忧愁；后来却为那些借此操练的人，结出平安的义果。}因此，天主本身就是至高的智慧；但敬拜天主是人的智慧，我们现在论及的正是这智慧。因为\BibleQuote{这世界的智慧，在天主面前是愚妄的}。因此，关于这智慧——即敬拜天主——圣经说：\BibleQuote{智者的众多是世界的福祉。}\par

2. 但如果论及智慧属于智者，我们将做什么？我们敢自称智慧吗？免得我们论及智慧竟是厚颜无耻？我们岂不被毕达哥拉斯的榜样所警示？——他不敢自称智者，而是回答说他是哲学家，即智慧的热爱者；由此产生了这个名称，后来变得如此普遍，以致无论一个人在智慧方面有多大学问——无论在他自己或他人看来——他仍只被称作哲学家。或者，即使是这样的人也没有人敢自称智者，是因为他们认为智者是无罪的？但我们的圣经并非这样说，圣经说：\BibleQuote{你责备智者，他必爱你}。因为无疑，认为人应受责备的人，判断他有罪。然而，就我而言，即使在这个意义上，我也不敢自称智者；对我来说，只要假定他们自己也无法否认的一点就够了：论及智慧也属于哲学家，即智慧的热爱者。因为他们并没有停止这样的讨论，他们自称是智慧的热爱者，而非智者。\par

3. 因此，在论及智慧时，他们如此定义：智慧是有关人类和神圣事物的知识。因此，在上一卷书中，我没有否认对这两类主题——无论是神圣的还是人类的——的认知都可称为知识和智慧。但按照宗徒的话的区分：\BibleQuote{这人得了智慧之言，那人得了知识之言}，这一定义应当如此划分：使关于神圣事物的知识称为智慧，而关于人类事物的知识则归自己以知识之名；我在第十三卷中已论述了后者，但并非将此知识归属于人类所能知道的关于人类的一切事物——其中有极多的虚空和有害的好奇心——而只归属于那些产生、滋养、保卫、坚固那导向真福的最有益的信仰的事物；在这一点知识上，大多数信徒并不坚强，尽管在信仰本身非常坚强。因为只知道人应当相信什么才能获得有福的生命——这生命必须是永恒的——是一回事；而知道这信仰本身如何既能帮助虔诚的人，又能抵御不敬虔的人，则是另一回事；后者宗徒似乎用知识的专名来称呼。当我以前论及这知识时，我的特别任务是首先简要区分永恒事物与暂时事物，并在那里论述暂时事物；同时将永恒事物推迟到本书，我也指出，关于永恒事物的信仰本身是暂时的，在时间中存于信徒心中，但为了获得永恒事物本身却是必要的。我还论证，关于那永恒者作为人为我们所做和所受的暂时事物的信仰——祂所取的人性在时间中并通向永恒——对于获得永恒事物也是有益的；并且，在此暂时的、必死的生命中，人谨慎、勇敢、节制、公正地生活所凭的德行本身，并非真德行，除非它们都指向那同一信仰——虽是暂时的——却仍导向永恒事物。\par

% structure: p0007 source=h2 target=subsection reason=relative-heading-rank; base=h2

\subsection{第二章。——在持有、默观和爱慕暂时的信仰中有一种三一，但尚未恰当地达到天主的肖像。}

4. 因此，就如经上所载：\Quote{我们住在肉身内，就是离主远行，因为我们是在信德中行走，而不是在目睹中}；无疑地，义人既因信德而生活，无论他如何按内在的人而生活，虽然他以这暂时的信德为目标，并达于永恒之事，但在持有、默观和爱慕这暂时的信德时，我们尚未达到那堪称为天主肖像的天主圣三；免得那本应属于永恒之事的，似乎被建立在暂时之事中。因为当人心看见自己的信德——藉此相信那未见之事时，它并非看见永恒之事。因为那不会永远存在，当这使我们远离主、必须凭信德而行的旅居结束时，当那使我们得以面对面目睹的看见接替它时，这信德便不复存在；正如现在，因为我们相信虽未看见，我们将配得看见，并因经由信德而被带至目睹而喜乐。那时将不再是信德——借此相信未见之事，而是目睹——借此看见所信之事。因此，虽然我们记起这过往的尘世生命，并因回忆而想起我们曾相信那未见之事，那信德将被算作已过时与结束之事，而非现时且恒常持续之事。因此，那现在由记起、默观和爱慕这现时且持续的信德而构成的天主圣三，那时将被发现是已完结和过去的，而非仍然持久的。由此可推断，若那天主圣三真是天主的肖像，则这肖像本身须被算作非恒常存在之物，而是属于短暂之物。\par

% structure: p0009 source=h2 target=subsection reason=relative-heading-rank; base=h2

\subsection{第三章——移除一个阻碍前述论点之困难。}

然而，我们切不可认为：灵魂的本性是不朽的，且从受造之初起便永不停止存在，而其本性中最佳之物却不能随着自身的不朽而永存。受造的本性中，有什么比按照造物主的肖像受造更好呢？因此，我们必须寻找那堪称为天主肖像之物，不在于持有、默观和爱慕那不会永远存在的信德，而在于那永远存在之物。\par

5. 我们要不要更仔细、更深入地审视一下，情况是否果真如此？因为可能有人说，即使信德本身已经过去，这个天主圣三也不会消亡；因为正如现在我们既通过记忆持有它，又通过思想辨别它，又通过意志爱它；同样，到那时，当我们既在记忆中持有、又将回忆我们曾拥有它，并以第三者——即意志——结合这二者时，同一个天主圣三仍将继续存在。因为，如果它在过去时仿佛没有在我们身上留下痕迹，那么毫无疑问，我们在记忆中也不会拥有它，以便在回忆它作为过去之事时能够返回，并通过第三者，\Emph{即}意图，结合这两者——即我们记忆中虽未思考却存在的那个，以及通过概念从中形成的那个。但这样说的人，没有察觉到，当我们持有、观看并爱我们现在的信德时，我们所涉及的是一个现在存在的不同天主圣三，与将来当我们通过回忆沉思（而不是信德本身）犹如记忆中存储的它的想象痕迹时，并通过意志作为第三者，\Emph{即}结合这两者——即保留者的记忆中所存在的，和由此印在回忆者理智视觉上的那个——所存在的那个天主圣三是不同的。为了理解这一点，让我们从有形事物中取一个例子，我们在第十一卷中已充分讨论过。因为当我们从低处上升到高处，或从外部进入内部时，我们首先在所见的有形对象、观看者的视觉（当他看到时，视觉由此形成）、以及结合二者的意志意图中发现一个天主圣三。让我们假设一个类似的天主圣三：当现在在我们内的信德像我们所说的有形对象位于空间中那样被牢固地记忆时，从这信德形成回忆中的概念，就像从那个有形对象形成观看者的视觉一样；而为了完成这个天主圣三，意志要算作第三者，它将记忆中牢固的信德与印在回忆视觉上的信德肖像联结和结合起来，正如在那个有形视觉的天主圣三中，所见有形对象的形状和在观看者视觉中形成的相应形状由意志意图结合一样。那么，假定这个被观看的有形对象已经解体并消亡，它哪里都不再存留，使目光能够求助于其视觉；我们能说，因为那个已过去且不复存在的有形对象的形象仍留在记忆中，由此形成回忆时的概念，并由意志作为第三者结合二者，所以它和先前那个天主圣三相同吗？那个天主圣三是在原位摆放的有形对象的外观被看到时所形成的吗？当然不是，而是完全不同的一个：因为，且不说前者来自外部，后者来自内部；前者显然由现有有形对象的外观产生，后者则由那个现已过去对象的形象产生。同样，在我们现在所讨论的情况下（为说明这一点，我们以为提出这个例子是合适的），甚至现在在我们理智中的信德，就像那个有形对象在空间中一样，当被持有、观看、爱慕时，产生一种天主圣三；但是，当这信德在理智中不再存在时（正如那个有形对象在空间中不再存在一样），这个天主圣三也将不复存在。而那时当我们记得它曾经存在，但现在不在我们内时，所存在的那个天主圣三，无疑将是另一个。因为现在的这个是由实际存在的事本身产生的，它当前在场并附着于信仰者的理智；而将来的那个将由回忆者记忆中留下的过去之事的想象产生。\par

% structure: p0012 source=h2 target=subsection reason=relative-heading-rank; base=h2

\subsection{第四章——天主的肖像应寻求于理性灵魂的不朽。如何在理智中显示天主圣三。}

6. 因此，那现在不存在的三一形态不是天主的肖像，那将来也不存在的三一形态也不是；但我们必须在那不朽地植入其不朽性中的人灵魂——\Emph{即}理性或理智的灵魂——内，找到造物主的肖像。因为灵魂本身的不朽是有条件的；灵魂也有其适当的死亡，当它缺少幸福的生命时——这幸福的生命可称为灵魂的真生命；但灵魂之所以被称为不朽，是因为它永不停止以某种生命活着，即使是在最悲惨的状态下也是如此。同样，尽管理性或理智有时在灵魂中沉睡、有时显得微小、有时显得伟大，但人的灵魂从来都仅仅是理性的或理智的；因此，如果灵魂因能运用理性和理智来理解和瞻仰天主而按天主的肖像受造，那么从这如此奇妙而伟大的本性开始存在的那一刻起，无论这肖像是磨损得几乎不存在，还是模糊且被玷污，或是明亮而美丽，它确实始终存在。此外，神圣圣经怜悯这尊严被玷污的状态，告诉我们：\Quote{人纵然行走在肖像中，却徒然自扰；他积蓄财宝，却不知谁来收取。}因此，它不会将虚妄归给天主的肖像，除非它看出那肖像已被玷污。然而，它充分表明，这种玷污并未扩展到取消其肖像是肖像，因为经上说：\Quote{人纵然行走在肖像中。}因此，那句话可以两方面真实地宣告；正如经上所说：\Quote{人纵然行走在肖像中，却徒然自扰}，同样可以说：\Quote{人纵然徒然自扰，却仍行走在肖像中。}因为灵魂的本性虽是伟大的，却可以败坏，因为它不是至高的；虽然它可以败坏，因为它不是至高的，但因为它能承受并能分享那至高的本性，所以它是一种伟大的本性。那么，让我们在这天主的肖像中，藉着那按自己的肖像造了我们的祂的帮助，来寻找一种特殊的三一形态。否则，我们无法健康地探究这一主题，也无法按照来自祂的智慧得出任何结果。但若读者或能记起并回想我们在前几卷书中，尤其是在第十卷中，关于人的灵魂或心智所说的话，或愿意仔细重读这些段落，他就不必在此对如此伟大之事的探究要求更冗长的论述。\par

7. 因此，我们在第十卷中，与其他事物一起说过，人的心智认识自身。因为心智认识任何事物都不如认识那接近它自己的事物；而没有什么比心智本身更接近心智。我们也引用了其他证据，足够充分地证明了这一点。\par

% structure: p0015 source=h2 target=subsection reason=relative-heading-rank; base=h2

\subsection{第五章——婴儿的心智是否认识自身。}

那么，对于婴儿的心智该说什么呢？它如此幼小，对事物如此深邃无知，以至于任何知道一点东西的人的心智看到这黑暗都退缩。难道也要相信它认识自身吗？但因为它过于专注于那些通过身体感官开始感知的事物，且因其新奇而更加喜悦，它确实不能对自己无知，但也无法思维自己。此外，它对外部可感觉事物是多么专注，可以从这一事实推测：它对可感觉的光是如此贪婪，以至于若有人在夜间因疏忽或不知可能后果，将一盏灯放在婴儿躺卧处，使得婴儿的眼睛能朝那一侧注视，但脖子不能转动，那孩子的目光就会如此固定在那方向，以至于我们曾知道有些人因此成了斜视，因为眼睛保留了习惯在它们柔软时以某种方式印下的形状。同样，对于其他身体感官，婴儿的灵魂——就其年龄所允许的程度——仿佛收缩自身并专注于它们，以至于它们只强烈厌恶或强烈渴望那通过肉身冒犯或引诱之物，却不思考自己的内在自我，也不能因劝告而被引导这样做；因为它们还不知道表达劝告的符号——其中言语是主要的——而它们对言语和其他事物都完全无知。而我们在同一卷中已经说明，不认识自己与不思维自己是两回事。\par

8. 但让我们略过婴儿时期，因为我们无法询问它里面发生了什么，而自己又已几乎完全忘记。让我们只满足于此，即可以肯定：当人能够思维自己心智的本性，并发现何为真理时，他除了在自己内，无处可寻。而且他将发现的，不是他所不知道的，而是他所未思维的。因为如果我们不知道我们自己心智中的事物，我们又知道什么呢？我们所能知道的一切，除非通过心智，否则一无所知？\par

% structure: p0018 source=h2 target=subsection reason=relative-heading-rank; base=h2

\subsection{第六章——在思维自身的心智中如何存在一种三一。思维在这三一中的角色是什么。}

但思维的功能如此伟大，以至于连心灵本身也不能，可以这么说，把自己呈现在自己的视线中，除非它思考自己；因此，情况既然如此，即没有什么东西是在心灵的视线中，除非那正在被思考的东西，以至于连心灵本身（我们藉以思考我们所思考的一切）也不能以别的方式呈现在自己的视线中，除非通过思考自己。然而，当它不思考自己时，它\Emph{如何}不在自己的视线中，尽管它永远不能脱离自己，好像自己是一回事，而对自己的视线是另一回事，这是我无法发现的。因为对身体的眼睛这样说并非不合理；因为身体的眼睛本身固定在身体中自己适当的位置上，但其视线延伸到自身以外的事物，甚至达到星辰。而眼睛不在自己的视线中，因为它不看自己，除非借助镜子，如上所述；当心灵通过思考自己将自己置于自己的视线中时，肯定不会发生这种事。那么，当它在思想中注视自己时，它是否藉自己的一部分来看另一部分自己，就像我们藉我们的其他成员（\Emph{即}眼睛）来看我们的一些成员（这些成员可以在我们的视线中）？有什么比这更荒谬可说的或可想的？因为除了藉自己，心灵被挪开了吗？或者，为了进入自己的视线，它被放在哪里，除非在自己的面前？因此，当它不在自己视线中时，它将在那里，而不在它先前所在之处；因为它从一处被移开，然后被放置在另一处。但如果它为了被看见而迁移，那么它将在哪里停留以便观看？它是否仿佛加倍了，以便同时在此处和彼处，\Emph{即}既在它能观看之处，又在它能被看见之处；这样在自身内它观看，而在自身前它被看见？如果我们询问真理，它会告诉我们绝无此类事情，因为当我们如此构想时，我们构想的只是物质对象虚构的形象；而心灵并非如此，对于少数能够探究这类真理的心灵来说，这是非常确定的。因此，看来心灵的观看属于其本性中的某种东西，并且当它构想自身时，它被召回到那本性中，不是仿佛通过空间的移动，而是通过非物质的回转；但当它不构想自身时，它似乎确实不在自己的视线中，也不是以其形成的知觉，但尽管如此，它认识自己，仿佛对自己有一种记忆。就像一个精通多种学问的人：他所知道的东西都包含在他的记忆中，但其中没有任何东西在他的心灵视线中，除了他正在构想的东西；而其余一切都被储存为一种隐秘的知识，称为记忆。因此，我们所阐述的三位一体的构成方式如下：首先，我们在记忆中放置了对象，知觉者的知觉由之而形成；其次，是形构或由此印下的图像；最后，是爱或意志，作为结合两者的东西。那么，当心灵在构想中注视自身时，它理解和认知自己；因此，它产生这自身的理解和认知。因为一件非物质的东西在被注视时被理解，在被理解时被认知。然而，心灵确实不是这样产生对自己的这种知识，即当它通过构想注视自己而被理解时，就好像它先前对自己是未知的；而是它认识自己，就像被包含在记忆中但不被想到的东西那样被认识；因为我们说一个人认识文字，即使他在想别的事情而不是文字。而这两者，产生者和被产生者，由爱作为第三者结合，爱无非是意愿，寻求或持守对某物的享受。因此，我们认为，也应藉记忆、理智、意志这三个术语来提示心灵的三位一体。\par

9. 但既然心灵，如我们在同一第十卷接近结尾处所说，总是记住自己，总是理解和爱自己，尽管它并非总是思考自己，以区别于那些不是它自身的东西；那么，我们必须探究，\OriginalTerm{理解}{intellectus}在何种意义上属于构想，而每个东西的\OriginalTerm{概念}{notitia}存在于心灵中时，即使一个人没有在想它，也仅被称为属于记忆。因为如果是这样，那么心灵就没有这三样东西：即对自己的记忆、理解和爱；而它只记住自己，然后当它开始思考自己时，它才理解和爱自己。\par

% structure: p0021 source=h2 target=subsection reason=relative-heading-rank; base=h2

\subsection{第七章——藉一个实例使事情变得清楚。为帮助读者而处理问题的方式。}

因此，让我们更仔细地思考我们所举的例子，其中表明不知道一件事与不去想（构想）它不同；并且可能发生一个人知道某事但并未在想它，当他正在想别的事而不是那件事时。那么，当精通两种或多种学问的人正在想其中一种时，尽管他没有想其他一种或多种，但他仍然知道它们。但我们能正确地说：\Quote{这位音乐家当然知道音乐，但他现在不懂它，因为他没有在想它；但他现在懂几何，因为他正在想几何}吗？这样的断言，就表面看来，是荒谬的。再者，如果我们说：\Quote{这位音乐家当然知道音乐，但他现在不爱它，因为他现在没有在想它；但他现在爱几何，因为他正在想几何}——这岂不也同样荒谬吗？但我们相当正确地说：\Quote{你所看见正在讨论几何的这个人也是一位完美的音乐家，因为他既记得音乐，又理解并爱它；但尽管他知道并爱它，他现在并没有在想它，因为他正在想他所讨论的几何。}因此我们得到警告：我们对某些事物有一种知识，贮藏在心灵的深处；当它被想到时，它就好像公开走出来，并仿佛置于心灵的眼目之前；因为那时心灵本身发现它既记得、又理解、并爱它自己，即使它没有在想自己，当它在想别的事时。但对于那些我们长久没有想、并且除非被提醒就不能想起的事物，如果这个说法可以允许的话，以某种我不知道的奇妙方式，我们并不\Emph{知道}我们知道。简言之，提醒者对他所提醒的人说得正确：\Quote{你知道这事，但你不知道你知道它；我会提醒你，你会发现你知道你以为你不知道的事。}书籍也能导致同样的结果，\Emph{即}，那些关于读者在理性引导下发现为真实之主题的书籍，并非那些读者因叙述者的信任而相信为真实之主题（如同历史的情况），而是那些他本人也发现为真实的主题，无论是出于他自己，还是出于那作为心灵之光的真理本身。但即使被提醒也不能思考这些事的人，因心灵极大的盲目，深深埋藏于无知之黑暗中，并且极需天主的帮助，才能获得真正的智慧。\par

10. 为此，我愿提出某种证据，无论它可能是什么，关于构想的行为，以显示出自记忆中所包含的事物，心灵的眼目在回忆时如何被形成，并且当一个人构想时，某种东西如同被生出，正如在他构想之前他已经记在心里的那样；因为区别那些在不同时间发生的事情，以及父母在时间上先于子女，是更容易的。因为如果我们把自己引向心灵的内在记忆（借此它记得它自己），以及内在理智（借此它理解它自己），以及内在意志（借此它爱它自己），在这三者总是共存，并且自从它们开始存在以来就一直共存，无论它们是否被想到；那么，这一圣三的肖像确实似乎只属于记忆本身；但由于在这种情况下，一个言语不能没有思想（因为我们\Emph{想}我们所说的一切，即使是由那不属于任何特定语言的内在言语说出的），这一肖像更应在记忆、理智、意志这三者中被辨认。我现在所称的理智，是指我们借以在思想中理解的能力，即当我们的思想由那些已在记忆中但未被想到的东西的发现而形成时；我所说的意志，即爱或偏好，它结合了这个子女与父母，并且在某种程度上是两者所共有的。因此，我也曾尝试，\Emph{即}在第十一卷中，通过外在可感觉的事物（由肉眼所见）引导迟钝的读者；然后我继续与他们一同进入那内在人的能力，借此他推理暂时的事物，而推迟考虑那作为更高权能的主宰，借此他默观永恒的事物。我在两卷书中讨论了这一点，在第十二卷中区分了二者，一个较高而另一个较低，并且较低者应当服从较高者；在第十三卷中，我尽其所能地以真理和简洁讨论了较低者的职责，其中包含了关于人类事物的有益知识，以便我们在此暂时生命中如此行动，以求得那永恒的事物；因为，事实上，我已经在一个单卷中简要地涵盖了如此多样而丰富的主题，一个为许多伟大人物及其伟大论证所熟知的主题，同时表明其中也存在一个圣三，但还不是一个可以被称为天主的肖像的圣三。\par

% structure: p0024 source=h2 target=subsection reason=relative-heading-rank; base=h2

\subsection{第八章——现在应在心灵最高贵部分寻求作为天主肖像的天主圣三。}

11. 但我们如今已进入我们所承担的论题，即考察人心最高贵的部分，人借此认识或能够认识天主，以便在其中找到天主的肖像。因为人心虽不与天主有同一性体，但那比任何性体都善的性体的肖像，应当在我们的本性中比我们本性中任何更好的东西中寻求和找到。然而，必须先就人心自身来考察它，在它成为天主的分享者之前；并在其中找到祂的肖像。因为，正如我们所说，虽然因失去对天主的参与而磨损和玷污，但天主的肖像仍然存留。因为正是在这一点上它是祂的肖像，即它能够容纳祂，并能分享祂；如此大的善之所以可能，仅因它是祂的肖像。那么，人心记忆、理解、爱自身；如果我们看出这一点，我们就看出一个三一，虽然不是天主，但终究是天主的肖像。记忆不是从外部接收它所要持守的；理解也不是在外部发现它所要注视的，如同身体的眼睛那样；意志也不是从外部将两者结合，如同它结合身体对象的形相与由此在看者视觉中形成的形相那样；也不是概念在转向它时，发现了外部所见到的事物的影像，此影像被某种方式捕获并存入记忆，从而回忆者的直观得以形成，意志作为第三者结合两者：正如我们在第十一卷中所示，这在那些于有形事物中发现的、或通过身体感官从有形对象中引入的三一中发生的情形。也不是如我们在进一步讨论那种现在属于内在人的运作、并与智慧相区别的知识时发生或似乎发生的那样；这种知识的对象仿佛是外来的，或是通过历史信息带入的——如时间和地点中发生并过去的事迹和言语，或在事物本性中于其自身时空中确立的事物——或是从人自身中产生以前不存在的，无论通过他人的信息还是自己的思考——如信德，我们在第十三卷中详尽论述过；或如德性，如果它们是真正的德性，人便在此可死之生中如此善生，以致能在天主所许诺的不朽中幸福地生活。这些和诸如此类的事物在时间中有其自身的秩序，并且在该秩序中我们更容易看出一个记忆、视觉和爱的三一。因为这类事物中的一些先于学习者的知识。它们在被人认识之前就是可知的，并在学习者中产生对它们自身的知识。它们或存在于各自适当的地方，或在过去的时间中发生；尽管过去的事物本身并不存在，只有某些作为过去的记号，看到或听到它们便知道它们曾存在并已过去。这些记号或位于地方本身，如\Emph{例如}死者的纪念碑或类似物；或存在于可信的书面著作中，如一切有分量且经权威认可的历史；或存在于已经知道它们的人的心中；因为已经为他们所知的事物当然对其他人也是可知的，他们的知识先于后者，后者能够通过知道者提供的信息来认识它们。所有这些事物，当被学习时，产生某种三一，\Emph{即}由它们自身的种相（在未认识前也是可知的），并由学习者的知识对其的应用（当学习者学习时才开始存在），以及作为第三者的意志将两者结合；而当它们被认识后，在回忆它们时又产生另一个三一，现在内在地在人心自身中，来自那些在学习时被铭刻于记忆中的影像，以及当记忆中的目光因回忆而转向这些影像时思想的赋形，以及作为第三者将两者结合的意志。但那些在人心自身中以前不存而出现的事物，如信德和诸如此类的，虽然似乎是外来的，因为是通过教导植入的，但并非外在或在外发生，如所信的事物；而是完全在人心自身之内开始存在。因为信德不是所信的，而是用以信的；前者被信，后者被看。然而，由于它在心中开始存在，而心在它之前已经存在，它似乎有些外来，并将被算作过去的事物，当看见取代它而它本身止息时。它现在通过其当下驻留、被持守、被看到、被爱，产生一个不同的三一，与它将来通过自身的一些痕迹——在它过去时留在记忆中——所产生的三一不同：正如上文已经说过的。\par

% structure: p0026 source=h2 target=subsection reason=relative-heading-rank; base=h2

\subsection{第九章——义德及其他德性在来世是否止息。}

12. 然而有人提出一个疑问：在此现世中赖以善度生活的诸德，既然它们本身也开始存在于那此前没有它们时也仍然是心灵的心灵中，那么当它们把我们带到永恒事物时，是否也会终止存在呢？因为有些人认为它们将终止，而就其中三种——智德、勇德、节德——而言，这种说法似乎有些道理；但义德是不朽的，它更会在那时在我们内得以成全，而非终止存在。然而，伟大的雄辩家图里乌斯在对话录\WorkTitle{Hortensius}中论述这四种德性时说：\Quote{如果我们被允许在离开此世后，如传说所说，永远生活在幸福岛上，那么既无审判，何需雄辩？或者，甚至诸德本身又有何用？因为既无劳苦也无危险临到我们，我们就不需要勇德；既无他人之物可贪，就不需要义德；既无将不存在的情欲需加管制，就不需要节德；既无善恶之间的选择，也就不需要智德。因此，我们将仅仅通过学习和认识本性而蒙福，诸神的生活也唯独由此而值得赞美。由此我们可以看出，其他一切都是出于必要，而这一件事则是自由选择。}这位伟大的演说家，在颂扬哲学之卓越时，复述了他从哲学家们那里学到的一切，并以出色而愉快的方式加以解释，断言所有这四种德性仅在此生——我们目睹它充满烦恼和谬误——是必需的；而当我们离开此生后，若蒙允许生活在有福生命之处，则其中任何一种都不再需要；但蒙福的灵魂只是在学习和认识中蒙福，\Emph{i.e.}在对本性的默观中，没有比这更好、更可爱的了。这本性是创造并安排了所有其他本性的那一本性。如果义德在于服从这一本性的治理，那么义德当然是不朽的；它在那福乐中不会终止存在，反而会如此完美和伟大，以至于不可能更成全或更伟大。或许其他三种德性——智德，不再有任何错误的危险；勇德，没有忍受邪恶的烦恼；节德，没有情欲的阻挠——也会存在于那福乐中：这样，智德的作用可能是将任何善事物都不置于天主之上或与之同等；勇德的作用可能是最坚定地依附于祂；节德的作用可能是不会因任何有害的缺陷而喜悦。然而，义德如今所从事的帮助困苦者，智德所从事的防备诡诈，勇德所从事的耐心忍受患难，节德所从事的克制邪恶情欲，这些行为将不会存在那里，因为那里根本没有邪恶。因此，这些德性在此现世生活中所必需的行为，如同它们所指向的信德一样，将被算为过去之事；它们如今在我们持守、观看并爱慕它们为现在时，构成一种不同的三一，不同于它们那时将要构成的三一，那时我们将发现它们已不再存在，而是已经过去，并通过它们在经过记忆时留下的某些痕迹来认识；因为那时也将有一个三一，当那痕迹——无论它是何种——被保存在记忆中并被真实地认出时，然后这两者由意志作为第三者结合。\par

% structure: p0028 source=h2 target=subsection reason=relative-heading-rank; base=h2

\subsection{第十章 心灵如何借记忆、理解和爱慕自身而产生三一。}

13. 在我们所提到的所有可朽事物的知识中，有些可知事物在认识它们之前有一个时间间隔，如那些在认识之前已经存在的可感对象，或所有通过历史学习的事物；但有些事物与对它们的认识同时开始存在——例如，若有个先前根本不存在的可见物体出现在我们眼前，它当然不先于我们的认识；或者若有个声音发出而有人听到，毫无疑问，声音和听到那声音同时开始和结束。尽管如此，无论是时间上在先还是同时开始存在，可知事物产生知识，而非由知识产生。然而，一旦知识形成，当已知并储存在记忆中的事物通过回想被回顾时，谁看不见记忆中的保留在时间上先于回想中的观看，以及由意志作为第三者结合这两者？但在心灵中，却并非如此。因为心灵并非外在于自身，仿佛有一个已经存在的心灵，而那个相同的心灵先前并不存在，从别处来到自身；或者并非从别处来，而是在已经存在的心灵本身中，生出了那个相同的心灵（先前并不存在）；正如信德，先前没有的，在已经存在的心灵中产生那样。心灵也不是在认识自身之后，通过回想看见自己仿佛被安置在自身的记忆中，好像它在认识自己之前并不存在；而实际上，从它开始存在之时起，它就从未停止记忆、理解和爱护自身，正如我们已经表明的。因此，当它通过思想转向自身时，就产生了一个三一，此时我们终于也能辨别出一个言语，因为它是由思想本身形成的，意志结合两者。这样，我们就能比迄今为止更清楚地认出我们所寻求的肖像。\par

% structure: p0030 source=h2 target=subsection reason=relative-heading-rank; base=h2

\subsection{第十一章 记忆是否也关乎现在的事物。}

14. 有人会说，心智永远临在自己，说它记起自己，这不是记忆；因为记忆是关于过去的，不是关于现在的。因为有些人，其中包括\Person{西塞罗}，他在论及德性时，将智德分为这三者——记忆、理解、预见：即把记忆归于过去，理解归于现在，预见归于未来；后者只有那些预知未来的人才是确定的；而这不是人的恩赐，除非从上赐予，如对先知们。因此，\WorkTitle{智慧篇}论及人说：\Quote{凡人的思念是惧怕的，我们的预见是不确定的。}但记忆过去和理解现在，是确定的：我指关于非形体之物，它们是现在的；因为形体之物对肉眼是现在的。但若有人否认有对现在的记忆，且注意世俗文献中的用语，那里更讲究词语的精确而非事物的真理。\Quote{尤利西斯没有遭受这样的事，伊塔刻人在如此巨大的危险中也没有忘记自己。}因为当\Person{维吉尔}说\Person{尤利西斯}没有忘记自己时，除了说他记起自己外，还指什么？而由于他临在自己，他不可能记起自己，除非记忆涉及现在的事物。因此，正如在过去的物中，那个使人能够回想并记住它们的东西被称为记忆；同样，在现在的物中，如心智对自己，那个使心智随时临在自己、以便通过自己的思想被理解，然后两者通过对自己的爱结合的东西，也不无理地被称为记忆。\par

% structure: p0032 source=h2 target=subsection reason=relative-heading-rank; base=h2

\subsection{第12章 心智中的天主圣三是天主的肖像，在于它记忆、理解和爱天主，而行此即为智慧。}

15. 因此，心智的这种天主圣三并非是天主的肖像，因为心智记起自己、理解并爱自己；而是因为它也能记起、理解和爱那造它的那位。这样做，它自己就成了智慧的。但如果它不这样做，即使它记起、理解并爱自己，那也是愚蠢的。让它记起它的天主，它是按祂的肖像造的，并让它理解和爱祂。或者更简单地说，让它崇拜天主，祂不是被造的，因为心智本身是被造的，所以它能承受并分有祂；因此经上记载：\BibleQuote{看，敬畏天主，这就是智慧。}然后它将变得智慧，不是靠自己的光，而是分有那至高之光；并且凡永恒之处，它将在福乐中为王。因为人的智慧之所以如此称呼，在于它也是出于天主。因为那时它是真智慧；如果它是人的，便是虚妄的。然而天主的智慧并非如此，不像天主藉以智慧的那种。因为祂不是因分有自己而智慧，如同心智因分有天主而智慧。但正如我们称之为天主的义，不仅指祂自己为义的那义，也指祂称不虔敬者为义时赐给人的义，宗徒称赞后一种义，论及某些人说：\BibleQuote{他们不知道天主的义，并企图建立自己的义，就不服从天主的义；}同样，也可以论及某些人说，他们不知道天主的智慧，并企图建立自己的智慧，就不服从天主的智慧。\par

16. 存在一种非受造的本性，它造了所有其他本性，无论大小，并且无疑比它所造的更高贵，因此也比我们谈论的本性更高贵；\Emph{即}比理性理智的本性，也就是人的心智，它是按造它的那一位的肖像造的。那比其余所有更高贵的本性就是天主。确实，\BibleQuote{祂离我们每人都不远，}如宗徒所说，他接着说：\BibleQuote{因为我们生活、行动、存在都在祂内。}如果这是就身体而言，它甚至可以被理解为有形世界；因为我们在其中也按身体生活、行动、存在。因此，应当以更卓越的方式，即属灵而非可见的方式，来理解关于心智，它是按祂的肖像造的。因为有什么不在祂内呢？关于祂神圣地写着：\BibleQuote{万有都出于祂、藉着祂、并在祂内。}如果所有都在祂内，那么那些活着的和运动的能在谁内生活、运动，除非在祂内？然而，并非所有人都以那种方式与祂同在，如对祂说的：\BibleQuote{我常与祢同在。}祂也并非以那种方式与所有人同在，如我们说：愿主与你们同在。因此，人的特殊不幸是与祂不同在，而没有祂人就不能存在。因为毫无疑问，人并非没有祂，因为人是在祂内；然而，如果人不记起、理解和爱祂，人就不与祂同在。当一个人完全忘记一件事时，当然不可能甚至提醒他。\par

% structure: p0035 source=h2 target=subsection reason=relative-heading-rank; base=h2

\subsection{第13章 人如何能忘记并记起天主。}

17. 让我们从可见之物中取一实例。一个你不认识的人对你说：\Quote{你认识我}；并为了提醒你，告诉你他在何处、何时、以何种方式与你相识；若在提及所有可能唤起记忆的记号之后，你仍不认识他，那么你已如此健忘，以致那知识全然从你心中抹去；此外别无他法，只能相信那告诉你你曾认识他的人的话；若你甚至不相信那对你说话的人堪可信赖，则连此亦不可行。但若你确实记得他，那么你无疑会回到自己的记忆，并在其中找到那尚未被遗忘完全抹去的东西。让我们回到引出此从人际交往中取例的那一点上。在其他经文之外，\BibleBook{}{圣咏} 第9篇说：\BibleQuote{恶人应堕入地狱，以及所有忘记天主的万民；}又\BibleBook{}{圣咏} 第22篇：\BibleQuote{普世地极将记起，并归向主。}这些万民，不会如此忘记天主，以致在受提醒时不能记起祂；然而，他们因忘记天主，仿佛忘记了自己的生命，而堕入死亡，\Emph{即}地狱。但当受提醒时，他们归向主，如同因记起自己所忘记的正当生命而复活。\BibleBook{}{圣咏} 第94篇也这样写道：\BibleQuote{你们这些在百姓中愚昧的，应当领会；你们这些愚人，何时才能明智？那造耳朵的，难道自己听不见吗？}等等。因为这话是对那些因不认识天主而说虚妄之事的人说的。\par

% structure: p0037 source=h2 target=subsection reason=relative-heading-rank; base=h2

\subsection{第十四章—— 灵魂在正当地爱自己时爱天主；若它不爱天主，就必得说它恨自己。即使软弱和迷失的灵魂，在记忆、理解与爱自己方面也始终坚强。愿它归向天主，好能因记忆、理解并爱慕祂而得享真福。}

18. 然而，圣经中关于爱天主还有更多证据。因为在这爱中，其余二者（即记忆与理解）自然被包含在内，因为没有人爱他不记得或完全不知道的事物。因此有一条著名且首要的诫命：\BibleQuote{你当爱主，你的天主。}人的灵魂如此被造，以致它无时不在记忆、理解并爱自己。但既然恨人者欲加害于人，那么当灵魂伤害自己时，说它恨自己也不为过。因为它因无知而对自己怀有恶意，并不认为自己所愿之事对其有害；但若它愿有害之事，它照样对自己怀有恶意。因此经上记载：\BibleQuote{喜爱不义的人，恨自己的灵魂。}所以，知道如何爱自己的人，爱天主；而不爱天主的人，即使他爱自己——这本是自然赋予的——仍可适切地说他恨自己，因为他做有害于己的事，且如同攻击自己的仇敌。这无疑是一种可怕的迷妄：尽管人人都愿有益于己，但许多人所做的却正是最有害于己的事。当诗人描述哑兽的类似疾病时，他说：\Quote{愿诸神将更好的事赐予虔诚者，并将那迷妄归给仇敌。它们用自己的牙齿撕扯自己破碎的肢体。}既然他谈论的是身体的疾病，为何称之为迷妄？除非因为，虽然本性驱使每种动物尽其所能照料自己，但那种疾病却使那些动物撕扯它们本欲保全完好的肢体。当灵魂爱天主，并因此如已述也记忆和理解祂时，那么它也被正当地命令要爱人如己；因为它此时已正当地、而非邪僻地爱自己，当它爱天主时，藉分享祂，那肖像不仅存在，而且得以更新，不再衰老；得以恢复，不再毁损；得以福乐，不再不幸。因为尽管它如此爱自己，以致若面临两难选择，它宁愿失去一切它所爱少于自己的事物，也不愿自己消亡；然而，由于它离弃了那高于它的天主——唯有依赖祂才能保守自己的力量，并享受祂作为光明——（对那在\BibleBook{}{圣咏}中歌颂的：\BibleQuote{我必依靠祢保守我的力量，}以及\BibleQuote{亲近祂，并得蒙光照，}）它已变得如此软弱和黑暗，以致可悲地也从自身堕落，陷入那些非己所是、且低于自身的事物，经由它无法征服的情感和它看不出归途的迷妄。因此，当它如今因天主的怜悯而痛悔时，在\BibleBook{}{圣咏}中呼喊：\BibleQuote{我的力量衰竭；我眼睛的光彩，也离我而去。}\par

19. 然而，在这些软弱和错谬的邪恶中，尽管它们极大，灵魂仍不能失去其自然的记忆、理解和对自己本身的爱情。因此，我上面所引述的话可以正当地说：\BibleQuote{人虽然行走在肖像中，但他实在只是徒然骚扰：他积蓄财宝，却不知谁将收取它们。}因为他为何积蓄财宝呢？岂不是因为他的力量已离弃了他，使他不能拥有天主，因而一无所缺？为何他不能说出他为准聚敛这些东西呢？岂不是因为他眼中的光明已被夺去？因而他看不见真理所说的话：\BibleQuote{糊涂人哪，今夜就要索回你的灵魂；你所预备的，将归谁呢？}然而，即使这样一个人行走在肖像中，人的心灵仍有对自身的记忆、理解和爱情；如果向他表明他不能同时拥有两者，而允许他选择其一、失去另一，即要么他所积蓄的财宝，要么他的心灵，那么谁竟如此完全失去心灵，以致宁愿拥有财宝而不愿拥有心灵呢？因为财宝通常能颠覆心灵，但未被财宝颠覆的心灵，即使没有任何财宝，也能更轻松、更无挂虑地生活。但除非凭借心灵，谁能拥有财宝呢？因为一个婴儿，无论出生时多么富有，纵然是他一切合法之物的主宰，若心灵无知，便一无所有；那么，一个完全失去心灵的人，又怎么可能拥有任何东西呢？但何必谈论财宝呢？如果说任何人，若给他选择，宁愿舍弃财宝也不愿失去心灵；然而，没有一个人宁愿舍弃，甚至没有人将之与身体的眼睛相比——藉着眼睛，不仅是某个人（像对待黄金那样），而是每个人都拥有这穹苍本身。因为每个人都藉着身体的眼睛拥有他所乐于看到的一切。那么，若有谁不能同时保持两者，而必须失去其一，谁不宁愿失去财宝而不愿失去眼睛呢？然而，如果在同样的条件下问他，是宁愿失去眼睛还是失去心灵，有心灵的人谁看不出他宁愿失去前者而不是后者呢？因为一个没有肉体眼睛的心灵仍然是人的；但眼睛没有心灵则是野兽的。谁不愿做一个人，即使在肉体的视觉上瞎眼，也不愿做一头能看见的野兽呢？\par

20. 我说了这么多，是为了使那些较为迟钝的人——这本小书可能到达他们的耳目——即使只被简短地提醒，一个软弱而错误的心灵是多么爱它自身，甚至在错误地爱慕和追求低于它的事物时。现在，它若完全无知于自身，即它若不记得自身、不理解自身，它就不能爱自身。它藉着自身内的天主的肖像，拥有如此强大的能力，以至于能依附于它既是其肖像的那位。因为它在次序中被纳入——不是位置的次序，而是本性的次序——除非在天主之上，没有在其上者。最后，当它完全依附于祂时，它将与主成为一神，正如宗徒所证，说：\BibleQuote{但那与主联合的，便是与祂成为一神。}这是由于它接近分享祂的本性、真理和幸福，而不是由于祂在自己的本性、真理和幸福中增长。那么，当它幸福地依附于那本性时，它将不动地生活，并将不动地看到它所看到的一切。那时，正如圣经所应许的：\BibleQuote{他的愿望必以美物满足}——不动之美物，即天主圣三本身，它自己的天主，它是其肖像。并且，为使它从此不再受苦，它将在祂的临在的隐密处，被祂的伟大丰满所充满，以至罪恶从此不再诱惑它。但现在，当它看见自身时，它看见某种可变之物。\par

% structure: p0041 source=h2 target=subsection reason=relative-heading-rank; base=h2

\subsection{第十五章 虽然灵魂盼望幸福，但它不记得已失的幸福，却记得天主和正义的规条。不变的生活规条连恶人也知道。}

21. 对此，灵魂无疑毫不怀疑它是悲惨的，并渴望成为有福的；它之所以能希望实现这一点，别无他由，唯因其自身的可变性。因为如果它是不可变的，那么，正如它不可能在成为有福后变成悲惨一样，它也不可能在成为悲惨后变成有福。在一个全能而良善的天主之下，除了它自身的罪和它的主的正义之外，还有什么能使它变得悲惨呢？又有什么能使它成为有福的呢？除非是它自己的功德和它的主的赏报？但是，它的功德也是祂的恩宠，而祂的赏报将是它的福乐；因为它无法自己给予它所失去的正义，所以它没有正义。因为这份正义是在人受造时领受的，而确实因犯罪而失去了。因此，它领受正义，为能因这正义而配得领受福乐；因此当它仿佛因自己的善而开始骄傲时，宗徒真实地对它说：\BibleQuote{你有什么不是领受的呢？既然是领受的，为什么你夸耀，好像不是领受的呢？} 但当它正确地记起它的主，并领受了祂的圣神时，由于内在的教导，它完全感到自己唯有藉着祂白白赐予的爱才能上升，也唯有藉着自己自由选择的背弃才能堕落。当然，它并不记得它自己的福乐；因为那福乐曾经存在，但现在不在了，且它已完全忘记，因此甚至不能被提醒。但它相信它天主的可信赖圣经所讲述的关于那福乐的事，这些事由祂的先知记载，讲述了乐园的福乐，并将那关于人最初善与恶的历史信息传授给我们。它记起它的主、它的天主；因为祂常存，既非曾经存在而现已不在，亦非现在存在而未曾存在；而是正如祂永远不会不存在，祂也从未不曾存在。祂圆满地临在于一切。因此，它在祂内生活、行动、存在；所以它能记起祂。这不是因为它回忆起在亚当内或在现在身体的生命之前的任何地方认识过祂，或是在它初受造以便植入这身体时认识过祂；因为关于这一切它什么也不记得。凡有此等事，都被遗忘抹去了。但它得蒙提醒，为能转向天主，如同转向那光，即使当它背离祂时，也在某种程度上被这光所触动。因为正是因此，连不敬神的人也会思想永恒，并在人的品行中正确指责和正确称赞许多事。他们依据什么规则如此判断呢？除非依据那些规则，在其中他们看见人应当如何生活，即使他们自己并不如此生活？他们在哪里看见这些规则？因为他们并不在自己的本性中看见它们；无疑，这些事是由心灵看见的，而他们的心灵显然是可变的，但这些规则，凡能看见它们的人，都视之为不变的；他们也不在自己的心灵特性中看见它们，因为这些规则是正义的规则，而他们的心灵显然是不义的。这些规则究竟写在哪里，连不义的人也能在其中认出什么是正义，在其中他辨别出他应当拥有他自己所没有的东西？那么，它们写在哪里呢？除非在那称为真理的光明之书中？一切正义的法则都是从那里被抄写和转移（不是迁移到它那里，而是仿佛被印在它上面）到行正义的人心中的；正如戒指的印记转到蜡上，却不离开戒指。但那不行正义、却看见他应当如何行的人，就是那背离那光、却仍被那光所触动的人。但那人甚至看不见他应当如何生活，他犯罪时更有可原谅之处，因为他不违犯他所知道的律法；但他也偶尔被无处不在的真理的光辉所触动，当他听到劝告而承认时。\par

% structure: p0043 source=h2 target=subsection reason=relative-heading-rank; base=h2

\subsection{第十六章 天主的肖像如何在人内重新形成。}

22. 然而，那些因受提醒而归向主的人，从他们藉世俗的情欲而依同于此世的丑陋中，被更新而脱离此世，当他们听从宗徒说：\BibleQuote{不可与此世同化，却要在心意更新而变化，}好使那肖像能由最初塑造它的那一位重新开始被塑造。因为那肖像不能自己重塑自己，正如它能自己扭曲自己。他别处又说：\BibleQuote{你们应在心思的心神里更新，并穿上新人，就是按照天主在正义和真实的圣洁中创造的。}所谓\BibleQuote{按照天主创造}，另一处表达为\BibleQuote{按照天主的肖像}。但它因犯罪而失去了正义和真实的圣洁，那肖像因而被玷污和黯淡；当它被重塑和更新时，便恢复这些。但当他讲\BibleQuote{在你们心思的心神里}时，他并非意指两件事，好像心思是一事，心思的心神是另一事；他这样说，是因为一切心思都是\Quote{心神}，但并非一切心神都是\Quote{心思}。因为还有一个心神就是天主，它不能更新，因为它不会变老。我们也说在人与心思不同的\Quote{心神}，属于这心神的是那些按形体相似形成的图像；宗徒在\BibleBook{格林多}{前书}中论及此，说：\BibleQuote{若我以舌头祈祷，我的心神祈祷，但我的理智却没有结果。}他这样说，是因为所说的话不被理解；因为除非借着心神的思维，使有形有声言语的图像先于口头的声响，否则连说出也不能。人的灵魂也称为\Quote{心神}，正如在福音中：\BibleQuote{祂便低下头，交付了心神。}这表示身体因心神离去而死亡。我们也说兽的心神，正如在\Person{所罗门}称为\BibleBook{训道}{篇}的书中所写：\BibleQuote{谁知道人的心神是向上升，兽的心神是向下入地？}在\BibleBook{创世}{纪}中也写道，洪水消灭了所有\BibleQuote{有生命的心神}在内的血肉。我们也说\Quote{心神}指风，这是最明显的有形之物，如\BibleBook{圣咏}{集}：\BibleQuote{火与冰雹，雪与冰，暴风的心神。}既然心神有这么多含义，宗徒用\Quote{心思的心神}表示那称为心思的心神。正如同一宗徒说：\Quote{脱去肉身的身体}时，当然不是指两件事，好像肉身是一事，肉身的身体是另一事；而是因为\Quote{身体}是许多没有肉身的事物的名称（因为除了肉身，还有许多天上的身体和地上的身体），所以他用\Quote{肉身的身体}来表示那身体就是肉身。同样，用\Quote{心思的心神}表示那心神就是心思。别处他也更明白地称之为\Quote{肖像}，同时用其他词语强调同一事。他说：\BibleQuote{你们脱去旧人及其行为，穿上新人，这新人是在天主的认识上，按造他的肖像而更新的。}一处经文说：\BibleQuote{穿上新人，就是按天主创造的}，另一处说：\BibleQuote{穿上新人，就是按造他的肖像更新的。}\par

一处说：\BibleQuote{按天主}；另一处说：\BibleQuote{按造他的肖像。}但在前者说\BibleQuote{在正义和真实的圣洁中}，后者则说\BibleQuote{在天主的认识上}。因此，这心神的更新和重塑或是按天主，或是按天主的肖像。但说\BibleQuote{按天主}，是为了不以为是按另一个受造物；说\BibleQuote{按天主的肖像}，是为了明白这更新发生在有天主肖像的地方，即心思中。正如我们说，一个信德者义人离开了身体，是按身体死亡，不是按心神死亡。因为我们说\Quote{按身体死亡}是什么意思？岂不是按身体或在身体中死亡，而不是按灵魂或在灵魂中死亡？或者我们想说\Quote{按身体俊美}或\Quote{按身体强壮}，不是按心思？这岂不是说他身体俊美或强壮，而不是心思？其他无数例子也是如此。因此，我们不要理解\BibleQuote{按造他的肖像}这句话，好像更新所依据的肖像是不同的，而不是那自身被更新的同一肖像。\par

% structure: p0046 source=h2 target=subsection reason=relative-heading-rank; base=h2

\subsection{第十七章——心思中的天主肖像如何被更新，直到天主的相似在真福中得以圆满。}

23. 这种更新当然不会在皈依的瞬间发生，如同圣洗中的更新在一瞬间藉着赦免所有罪过而发生；因为没有一件罪过，即使是最小的，没有被赦免。但是，正如摆脱发烧是一回事，而从发烧引起的虚弱中重新强壮又是另一回事；又如从身体中拔出刺入的武器是一回事，而通过成功的治疗愈合由此造成的伤口又是另一回事；同样，第一次治疗是除去虚弱的原因，这是藉着赦免所有罪过完成的；但第二次治疗是医治虚弱本身，这是通过在那肖像的更新中逐步进展而发生的：这两件事在圣咏中清楚地显示出来，我们读到：\BibleQuote{祂赦免你所有的罪愆}，这在圣洗中发生；随后是：\BibleQuote{祂治愈你的一切病弱}；这通过每日的添加而发生，同时这肖像正在被更新。宗徒极其明确地谈到这一点，说：\BibleQuote{我们外在的人虽然渐渐损坏，但我们内在的人却日日更新。}并且\BibleQuote{它是在认识天主的方面被更新，即在正义和真实的圣洁中}，依据稍前引用的宗徒的证言。因此，那藉着在认识天主、在正义和真实的圣洁中进步而日日更新的人，将他的爱从暂时的事物转向永恒的事物，从可见的事物转向可理解的事物，从属肉体的事物转向属神的事物；并勤奋地坚持抑制和减少对前者的欲望，而藉着爱将自己与后者结合在一起。他这样做，是依他蒙天主帮助的程度。因为天主亲自说：\BibleQuote{离了我，你们什么也不能做。}当生命的最后一天发现任何人在这样的进步和成长中坚守对中保的信德时，他将被圣天使迎接，被带领到他所崇拜的天主面前，并由祂使之完善；如此，在世界穷尽时将得到一个不朽的身体，不是为了惩罚，而是为了光荣。因为那时天主的肖像将在这形象中达到完美，当天主的观看达到完美时。关于这一点，\Person{圣保禄}宗徒说：\BibleQuote{我们现在是藉着镜子观看，模糊不清，但那时就要面对面了。}又说：\BibleQuote{我们众人以揭开的脸面，反映主的光荣，渐渐地光荣上加光荣，都变成了与主同样的肖像，正如由主，即圣神在我们内所完成的。}这就是在那些进步良好的人中天天发生的事。\par

% structure: p0048 source=h2 target=subsection reason=relative-heading-rank; base=h2

\subsection{第十八章 \Person{若望}的话是否应被理解为指我们将来在身体的不朽中与天主子的相似}

24. 但\Person{若望}宗徒说：\BibleQuote{可爱的诸位，现在我们是天主的子女，但我们将来如何，还没有显明；可是我们知道：一显明了，我们必要相似祂，因为我们必要看见祂实在怎样。}由此可见，天主的完全相似将发生在那天主的肖像中，即当它获得对天主的圆满观看之时。然而，这也可能被\Person{若望}宗徒理解为指身体的不朽。因为我们在这一点上也要相似天主，但只是相似圣子，因为只有圣子在\OriginalTerm{天主圣三}{Trinity}中取了身体，在其中死而复活，并带着这身体升天。因为这也被称为天主子的肖像，在这肖像中我们将拥有如同祂一样的不朽身体，在这一点上被塑造，不是相似圣父或圣神的肖像，而只相似圣子，因为只有关于祂，我们以纯正的信德读到并接受：\BibleQuote{圣言成了血肉。}为此，宗徒说：\BibleQuote{祂所预知的人，也预定他们与祂儿子的肖像相同，使祂在众多弟兄中作长子。}当然，\BibleQuote{从死者中首生者}，按照同一宗徒的话；藉着这死亡，祂的肉体被撒在羞辱中，却复活在光荣中。根据圣子的这肖像，我们在身体上藉着不朽被塑造成相似祂，我们也做同一宗徒所说的事：\BibleQuote{我们怎样带了那属于土的肖像，也要怎样带那属于天上的肖像}；也就是说，我们这些在亚当之后是属土的人，可以藉着真实的信德和确定无疑的望德持守：我们要在基督之后成为不朽的。因为我们现在就能带着同样的肖像，不是在实际看见中，而是在信德中；不是在实际拥有中，而是在望德中。因为宗徒说这话时，是在谈论身体的复活。\par

% structure: p0050 source=h2 target=subsection reason=relative-heading-rank; base=h2

\subsection{第十九章 \Person{若望}更应被理解为指我们在永生中与天主圣三的完美相似。智慧在幸福中达致完美}

25. 但在那肖像方面，关于它曾说过：\BibleQuote{让我们按我们的肖像和模样造人}，我们相信——并且经过我们能做的最大努力，我们也理解——人是按天主圣三的肖像造的，因为不是说的\Quote{按我的肖像}或\Quote{按你们的肖像}。因此，\Person{若望}宗徒的那个地方也必须按照这肖像来理解，当时他说：\BibleQuote{我们将相似祂，因为我们将看见祂实在怎样}；因为他也说到了那一位，关于祂他曾说过：\BibleQuote{我们是天主的子女}。肉身的复活将要在复活的那一刻完成，关于这复活，\Person{保禄}宗徒说：\BibleQuote{在一剎那，一眨眼之间，在最后的号筒时；死人将复活成为不朽的，我们将被改变}。因为在那眨眼之间，在审判之前，属神的身体将要在德能、不朽、光荣中复活，而这身体现在是被播种为属生灵的身体，在软弱、可朽、羞辱中。但那在内心精神上于认识天主时日日更新的肖像，不是外在的，而是内在的，将要在那种看见本身中被成全；这看见在审判之后将面对面，但现在却如同在镜子里，在谜语中前进。我们必须理解，由于这个成全，才说了：\BibleQuote{我们将相似祂，因为我们将看见祂实在怎样}。因为这个恩赐将在那时赐给我们，当这话被说时：\BibleQuote{来吧，我父所祝福的，承受为你们预备的国度}。因为那时不敬神的人将被除去，使他不得见主的荣耀，当左边的人进入永罚，而右边的人进入永生。但\BibleQuote{这就是永生}，正如真理告诉我们；\BibleQuote{认识你}，祂说，\BibleQuote{唯一的真天主，和你所派遣的耶稣基督}。\par

26. 这种默观的智慧，我相信它在圣经中恰当地被称为智慧，以区别于知识；但这只是人的智慧，而人除非从那一位获得，是不能拥有它的，有理性和理智的心灵藉着分有祂，才能成为真正有智慧的——这种默观的智慧，我说，正是\Person{西塞罗}在对话录\WorkTitle{霍尔滕修斯}的末尾所赞扬的，当时他说：\Quote{当我们夜以继日地思考这些事情，磨砺我们的理解力——它是心灵的眼睛——注意不要让它变得迟钝，也就是说，当我们生活在哲学中时；我们，我说，这样做，就抱有很大的希望：如果一方面，我们的这种情感和智慧是必朽的和短逝的，那么在我们尽完人性职责之后，仍将有一个愉快的落幕，一个无痛的熄灭，以及一种生命的安息；或者如果另一方面，如同古代哲学家——而且是最伟大、最著名的——所认为的那样，我们拥有永恒神圣的灵魂，那么我们必须认为，这些灵魂愈是始终保持在它们本身的适正道路上，\Emph{即}理性和探究的渴望中，愈是少与人的邪恶和错误混杂纠缠，它们上升回返天堂就愈容易。}然后他说，加上这句简短的话，并重复它来结束他的论述：\Quote{因此，为了最终结束我的论述，如果我们希望要么在追求这些学问之后获得一个宁静的熄灭，要么毫不迟延地从现世寄居所迁移到另一个好得多的居所，我们必须把全部劳力和关心都用在追求这些学问上。}在这里我感到惊奇，如此伟大的人竟会向那些生活在哲学中的人——哲学通过默观真理使人有福——应许\Quote{在尽完人性职责之后有一个愉快的落幕，如果我们的这种情感和智慧是必朽和短逝的}，好像那件我们不爱，或者更确切地说，我们深深恨恶的东西，那时就会死去归于虚无，以致它的落幕对我们而言会是愉快的！但事实上，他并非从那些他大加赞扬的哲学家那里学到这一点；这种情感带有新学园派的味道，在那学派中，他乐于怀疑甚至最明显的事物。但是，从他称为最伟大、最著名的哲学家那里，他学到灵魂是永恒的。因为永恒的灵魂，当现世生命结束时，被劝勉要保持在\Quote{它们本身的适正道路}上，即\Emph{即}理性和探究的渴望中，并少与人的邪恶和错误混杂纠缠，以便更容易回归天主，这样的劝勉并非不合适。但是，那种由热爱和探求真理构成的道路，并不足以应对痛苦的人——即所有仅有这类理性、却对中保没有信德的必朽者；正如我尽己所能，在这部著作的前几卷，尤其是第四卷和第十三卷中，已经努力证明的那样。\par

\SourceRecord{Translated by Arthur West Haddan.；Nicene and Post-Nicene Fathers, First Series；Vol. 3.；Edited by Philip Schaff.；Buffalo, NY: Christian Literature Publishing Co.,；1887.；<http://www.newadvent.org/fathers/130114.htm>.}

% structure: p0001 source=h1 target=section reason=page-title

\section{《论天主圣三》（第十五卷）}

\EditorialNote{本章首先简短而概括地陈述了前十四卷的内容。随后表明，论述已进展到这样的地步，即我们现在可以在那些永恒、无形和不变的事物本身中探究天主圣三，对之完满的默观为我们许诺了幸福的生命。但作者指出，这天主圣三在此世是由我们如同在镜中、在谜语里所看见的，即它是借着我们是其肖像的天主而看见的，是一种晦暗且难以辨别的相似。同样，他指出，关于圣言的受生，可以从我们理智的话语中得出某种推测和解释，但只能困难地达成，因为在这两种话语之间存在着极大的差异；关于圣神的发出，则可以从与之相连的意志之爱中得出。}\par

% structure: p0003 source=h2 target=subsection reason=relative-heading-rank; base=h2

\subsection{第一章——天主超越理智}

1. 为激励读者在受造物中得以认识创造它们的那一位，我们如今已进展到祂的肖像，就是人，论及人超越其他动物的方面，即理性或理智，以及凡有关理性或理智灵魂而可归于所谓心神的一切。因为有些拉丁作家，按他们独特的表达方式，以此名称区分在人里面卓越的、不在牲畜里面的东西与也存在于牲畜里的灵魂。那么，如果我们寻求任何超越这本性的事物，并且真诚地寻求，那就是天主——即非受造而是创造的本性。这天主圣三是否就是天主，如今我们的任务是不但为信徒凭借天主圣经的权威，而且也为有理解力的人，若可能的话，通过某种理性来证明。我为什么说\Quote{若可能}，当我们开始在我们探讨中论及它时，事情本身会更好地显示出来。\par

% structure: p0005 source=h2 target=subsection reason=relative-heading-rank; base=h2

\subsection{第二章——天主虽不可理解，但应常被寻求。在受造物中寻求天主圣三的踪迹并非徒然。}

2. 因为我们将寻求的天主本身，我希望，会帮助我们免于徒劳，并使我们理解\BibleBook{}{圣咏}中所说的话：\BibleQuote{愿寻求上主的人心喜乐。寻求上主和祂的能力，时常寻求祂的慈容。} 因为似乎总在寻求的东西好像从未找到；那么寻求祂的人的心如何能喜乐，而不是悲哀，如果他们找不到所寻求的呢？但经上不是说\BibleQuote{找到的人的心}要喜乐，而是\BibleQuote{寻求上主的人}。然而\Person{依撒意亚}先知证明，当主被寻求时，祂是可以被找到的，因为他说：\BibleQuote{你们要寻求上主，当你们找到祂时，就呼求祂；当祂临近你们时，让恶人离弃自己的道路，不义之人离弃自己的思想。} 那么，如果寻求祂时，祂可以被找到，为什么又说\BibleQuote{时常寻求祂的慈容}呢？也许连找到的祂也要被寻求？因为不可理解的事物必须如此探究，以致无人会认为，当他能发现他所寻求的是多么不可理解时，他什么都没有找到。那么，如果他理解他所寻求的是不可理解的，他为何还要如此寻求呢？除非因为只要他在对不可理解之物的探究本身中进步，他就不可停止寻求，并在寻求如此伟大的善时变得越来越好，这善既是为更甜美地找到而被寻求，也是为更热切地寻求而被找到。\BibleBook{}{德训篇}中智慧的话可作此解：\BibleQuote{那吃我的，还要饥饿；那喝我的，还要口渴。} 因为他们因找到而吃喝；他们仍继续寻求，因为他们饥渴。信德寻求，明悟找到；先知因此说：\BibleQuote{除非你们相信，你们将不会明白。} 然而，明悟仍继续寻求他所找到的；因为正如\BibleBook{}{圣咏}所唱的：\BibleQuote{天主垂顾人子，看看有没有明白的、寻求天主的。} 人因此应当有这样的理解，以便寻求天主。\par

3. 那么，我们将足够长久地停留在天主所造的万物之中，好使我们借以认识创造它们的天主自己。\BibleQuote{因为自从创造世界以来，祂那看不见的美善，即祂永远的大能和神性，都可凭祂所造的万物，被人辨认出来。} 因此，\BibleBook{}{智慧篇}中斥责那些\BibleQuote{不能由看得见的美善认识那存在者，也不曾因注意工程而认识工匠，却以为火、风、流动的空气、星辰的轨道、汹涌的水流或天上的光体，是统治世界的神祇。如果他们因喜爱这些物的美丽而奉之为神，就让他们知道，它们的主宰更为美好，因为那美丽的创始者创造了它们。如果他们因它们的权能和作为而惊讶，就让他们由此明白，创造它们的那一位更为有力。因为从受造物的伟大和美丽，人可以推断出它们的创造者。} 我引述\BibleBook{}{智慧篇}的这些话，是为使任何信友不会认为我徒然而空洞地先按步伐在受造物中寻找各种相应类型的某些三一，直到最终到达人的心神——那我们在寻求天主的最高天主圣三的踪迹。\par

% structure: p0008 source=h2 target=subsection reason=relative-heading-rank; base=h2

\subsection{第三章——前各卷的简要重述}

4. 但由于我们讨论和论证的需要，迫使我们在十四卷中说了许多事情，但我们无法在一瞥中同时看到它们，以便能迅速在思想中将它们引向我们想要把握的东西，我将尽力在天主的帮助下，尽我所能，不加以论证，简短地汇集我在各卷中通过论证作为已知而确立的一切，并好像置于一种精神视野之下，不是我们如何被说服接受每一点的方式，而是我们已被说服的各点本身；为使后面所述不致与前面所述如此分离，以致阅读前者会使人忘记后者；或者，如果它确实造成了这种遗忘，那么所遗忘的可以通过重读而迅速回忆起来。\par

5. 在第一卷中，从圣经显示了那至高的天主圣三的合一与平等。在第二、第三和第四卷中同样如此；但关于圣子和圣神的派遣问题，经过仔细探讨后写了三卷书；我们已证明，被派遣者并不因此小于派遣者，因为一位派遣，另一位被派遣；既然天主圣三在一切事物中平等，且同样在其本性上是不可改变、不可见、无所不在的，便不可分割地运作。在第五卷中——针对那些认为圣父与圣子的实体因此不同的人，因为他们假定凡论及天主的谓词都是按实体谓词，因而争辩说，生与被生，或受生与不受生，既然不同，便是不同的实体——我们证明：凡论及天主的，并非都是按实体谓词，例如祂被称为善与大是按实体，或是其他论及祂自身的谓词；但有些也是相对地谓词，\Emph{i.e.}不是论及祂自身，而是论及非祂自身的东西；例如祂因圣子而被称为父，或因服事祂的受造物而被称为主；并且在这里，如果任何这样相对谓词的东西，\Emph{i.e.}论及非祂自身的东西，也被作为时间性的谓词，例如\BibleQuote{主，祢成了我们的避难所，}那么祂并没有发生任何变化，而是祂自己在自己的本性或本质中完全不变。在第六卷中，关于基督如何被宗徒称为\BibleQuote{天主的德能和天主的智慧，}的问题，进行了讨论，但对该问题的更仔细处理被推迟，\Emph{viz.}：基督由祂所生的那位，不是智慧本身，而只是祂自己智慧的父亲，还是智慧生出了智慧。但无论怎样，在这卷书中也清楚显示了圣三的平等，并且天主不是三倍，而是一个圣三；而且圣父和圣子并非好像是相对于单独圣神的一个二重：因为其中三并不比一多什么。我们也思考了如何理解希拉里主教的这句话：\Quote{在圣父内的永恒，在肖像中的形式，在恩赐中的使用。}在第七卷中，解释了被推迟的问题：那位生了圣子的天主，如何不仅是他自己能力和智慧的父亲，而且他自己也是能力和智慧；同样，圣神也是如此；然而他们不是三种能力或三种智慧，而是一个能力和一个智慧，如同一个天主和一个本质。接着又探讨了，他们如何被称为一个本质、三位，或者被某些希腊人称为一个本质、三个实体；我们发现这些词语是因言说的需要而如此使用，以便当有人问我们所真实宣认为三的那三者是什么时，可以有一个术语来回答，\Emph{viz.}：圣父、圣子和圣神。在第八卷中，通过理性也向那些理解的人表明：不仅圣父在真理的实体上不大于圣子，而且二者一起也不大于单独圣神，在同一圣三中任何两个也不大于一个，三者一起也不大于每一个单独。接着，我指出：通过被理智所见的真理，以及通过从中产生一切善的至善，以及通过义德——一个正义的心灵甚至被尚未正义的心灵所爱——我们能够理解，按理解的可能限度，那不仅是无形体而且是不变的本性，即天主；并且也通过爱，在圣经中被称为天主，通过它，那些有理解的人首先开始，尽管非常微弱地，辨认出圣三，即：一个爱者、被爱者和爱。在第九卷中，论证推进到天主的肖像，\Emph{viz.}：人按其心灵而言；在其中我们找到了一种三一，\Emph{i.e.}：心灵，以及心灵借以认识自己的知识，以及心灵借以爱自己及其自我知识的爱；并且这三个被显示为彼此平等，且属同一本质。在第十卷中，同样的问题得到了更仔细和精妙地处理，并且达到了这一点：我们在心灵中发现了一种更显著的心灵三一，\Emph{viz.}：记忆、理解和意志。但既然也发现，心灵永远不可能处于不记忆、不理解、不爱自身的情况，尽管它并不总是思考自身；但每当它思考自身时，它并没有在同一思想行为中将自己与有形事物区分开来；关于圣三（此为其肖像）的论证被推迟，以便也在由身体所见的事物本身中找到一种三一，并更清晰地训练读者的注意力。因此，在第十一卷中，我们选择了视觉感官，其中本应在那里发现为真的东西，也可以在其他四种身体感官中被识别，尽管没有明确提及；于是，一个外在人的三一首先显示在那些从外部辨别的事物中，\Emph{viz.}：被看见的身体对象，由此印在观看者眼中的形式，以及结合二者的意志目的。但这三者，如显然的，并非彼此平等且属同一实体。接着，我们在心灵本身中发现了另一种三一，仿佛由外部感知的事物引入心灵；其中同样的三者，如所见，是同一实体：记忆中的身体对象肖像，以及当思想者心灵之眼转向它时由此印上的形式，以及结合二者的意志目的。但我们发现这个三一属于外在人，理由是从外部感知的身体对象引入心灵。在第十二卷中，我们觉得最好将智慧与知识区分开来，并首先在特别被称为知识的东西中寻找一种适当而特殊的三一，因为它是两者中较低者；但尽管我们现在因此达到了与内在人相关的某种东西，却还不能称之为或视为天主的肖像。这一点在第十三卷中通过赞扬基督徒的信德进行了讨论。在第十四卷中，我们讨论了人的真智慧，\Emph{viz.}：在分享天主本身时，由天主的恩赐所赐予人的智慧，它不同于知识；讨论达到了这一点：在天主的肖像中发现了三一，即人按其心灵而言，这心灵\BibleQuote{在认识上更新}于天主，\BibleQuote{照着创造人的那一位的肖像}；\BibleQuote{按祂自己的肖像}；如此获得智慧，其中有对永恒事物的默观。\par

% structure: p0011 source=h2 target=subsection reason=relative-heading-rank; base=h2

\subsection{第四章——普遍本性关于天主教导我们什么}

6. 那么，现在让我们在永恒、无形体、不变的事物本身中，寻求天主圣三，即天主本身；对此的完美默观应许给我们真福生活，这生活只能是永恒的。因为不仅神圣典籍的权威宣告天主存在，而且围绕我们的整个宇宙本性——我们也属于其中——也宣告有一位最卓越的造物主，祂赐给我们心智和自然理性，使我们能看出：有生命的事物优于无生命的事物；有感官的优于无感官的；有理智的优于无理智的；不朽的优于可朽的；有力的优于无力的；正义的优于不义的；美丽的优于丑陋的；善的优于恶的；不可朽坏的优于可朽坏的；不变的优于可变的；无形体的优于有形的；无体的优于有体的；真福的优于悲惨的。因此，既然我们毫无疑问地将造物主置于受造物之上，我们必须承认：造物主以最高意义活着，感知并理解一切，祂不能死、不能朽坏、不能改变；祂不是形体，而是神体，是万有中最有力、最正义、最美丽、最善良、最真福的。\par

% structure: p0013 source=h2 target=subsection reason=relative-heading-rank; base=h2

\subsection{第五章——以自然理性证明天主圣三何其困难}

7. 但我所说的一切，以及任何其他以类似人类言语似乎合宜地论及天主的话，都适用于整个天主圣三，即独一的天主，也适用于天主圣三内的几个位格。因为谁敢说：那独一的天主——即天主圣三本身——或圣父、或圣子、或圣神，要么不是生活的，要么没有感觉或理智；或者，在他们被宣告彼此相等的本性中，其中任何一位是可朽的、或可朽坏的、或可变的、或有形体的？又有谁否认天主圣三中的任何一位是最有力、最正义、最美丽、最善良、最真福的呢？那么，如果这些以及一切此类属性既能用于天主圣三本身，也能用于天主圣三中的每一位，那么天主圣三在哪里或怎样彰显自身呢？因此，让我们先将这些众多的谓词缩减为有限的数量。因为在天主内所谓的生命，就是祂的本质和本性。所以，天主不是借着一种与祂自身不同的生命而生活；除非祂自己就是生命。这生命不像树中的生命，既无理解也无感觉；也不像野兽中的生命，虽具有五感，但无理智。而作为天主的生命，感知一切、理解一切，并且以心智感知，而非以身体感知，因为\Quote{天主是神}。天主不是像动物那样通过身体感知，因为动物有身体，而祂并非由灵魂和身体组成。因此，那唯一的本性如何理解，也就如何感知；如何感知，也就如何理解；祂的感知和理解是同一回事。然而，并非祂在某一时刻会停止或开始存在；因为祂是不朽的。祂不徒然地被称为：\InnerQuote{只有祂是不死的。}因为真正的永生属于祂，祂的本性不容改变。那也是真实的永恒，天主藉此是不变的，无始无终，因而也是不可朽坏的。因此，称天主为永恒、或不朽、或不可朽坏、或不变，是同一回事；同样，说祂是生活的，有理智的，亦即真正智慧的，也是同一回事。因为祂不是领受了智慧而成为智慧的，祂本身就是智慧。而这就是生命，又是能力或权能，又是美丽，因而被称为有力而美丽。因为有什么比智慧更有力、更美丽呢？智慧\InnerQuote{强有力地伸展到地极，并柔和地安排一切}。再者，良善与正义在天主的本性中彼此有别吗？如同在祂的工程中那样，似乎良善与正义是天主两个不同的属性——良善是一个，正义是另一个？当然不是。因为正义本身就是良善，良善本身就是真福。因此，天主被称为无形体的，为使人相信并理解祂是神体，而非形体。\par

8. 此外，如果我们说：永存、不死、不朽、不变、生活、智慧、大能、美丽、正义、良善、有福的神体；这些词中似乎只有最后一个表示实体，其余的表示该实体的属性；但在那不可言说和单纯的本性中并非如此。因为凡是在其中似乎按属性被述说的，都应按实体或本质来理解。因为我们绝不会把神体按实体归于天主，而把良善按属性归于天主；两者都是按实体。同样，我们提到过的所有那些，我们在前几卷中已经详细讨论过的，也都是如此。那么，让我们从我们列举和排列的前四个中选取一个，即永存、不死、不朽、不变；因为这四个，正如我已经论证过的，具有同一含义；这样我们的目标就不会因为对象的多样而分散。最好选取放在第一个的，即永存。让我们对接下来的四个遵循同样的方式，即生活、智慧、大能、美丽。既然某种生命也属于野兽，它没有智慧；而接下来的两个，即智慧和大力，在人的情形中如此相互比较，以至于圣经说：\BibleQuote{有智慧的人胜过强壮的人}；而美丽，通常也归于物质对象。在这四个我们选出的之中，让我们取智慧。虽然这四个在谈论天主时不应被称为不平等；因为它们是四个名称，却是一件事。至于第三组也是最后一组的四个——虽然在天主内，正义与良善或有福是同一回事；而作为神体与正义、良善、有福也是同一回事；然而，因为在人内，可以有一个神体不是有福的，也可以有一个神体既正义又良善，但尚未有福；但有福的神体无疑既是正义、良善，也是神体——那么，让我们更愿意选取那一个即使在人类中也不能没有其他三个而存在的，即有福。\par

% structure: p0016 source=h2 target=subsection reason=relative-heading-rank; base=h2

\subsection{第六章 如何在天主极其单纯的本性中存在天主圣三。是否以及如何从已在人内显示的三位一体中显明天主圣三。}

9. 那么，当我们说：永存、智慧、有福，这三个就是被称为天主的天主圣三吗？我们确实将那些十二个简化为这三个少数；但也许我们可以更进一步，将这三个也简化为其中之一。因为如果在天主的本性中，智慧和大力，或生命和智慧，可以是同一回事，为什么永存和智慧，或有福和智慧，在天主的本性中就不能是同一回事呢？因此，正如我们当初把众多简化为少数时，谈论这十二个或那三个没有区别一样；现在，谈论那三个或那一个——我们已经表明那三个中的另外两个可以简化为那一个——也没有区别。那么，什么论证方式，什么理解的可能力量和活力，什么理性的生动，什么思想的敏锐，能够阐明（现在跳过其他）这一件事：天主被称为智慧，如何是一个三位一体？因为天主不像我们从祂接受智慧那样从任何人接受智慧，而是祂自己就是祂的智慧；因为祂的智慧不是一回事，祂的本质是另一回事，因为对祂来说，存在就是有智慧。的确，基督在圣经中被称作\BibleQuote{天主的德能，天主的智慧}。但我们在第七卷中讨论过如何理解这一点，免得圣子似乎使圣父有智慧；我们的解释归结为：圣子是智慧的智慧，正如祂是光之光，天主之天主。我们也无法发现圣神以别的方式存在，除非祂自己也是智慧，并且完全是一个智慧，如同一个天主，一个本质。那么，我们如何理解这智慧——这智慧即是天主——是一个三位一体？我不是说：我们如何相信这个？因为在信徒中，这应当不容置疑。但假设有任何方式，我们能用理智看见我们所相信的，那是什么方式？\par

10. 因为我们若回忆在这些书中，圣三最初是在何处开始向我们的理解显现的，那便是第八卷；在那里，我们曾竭尽全力提升心智的意向，以理解那至卓越、至不变的本性——我们的心智并非如此。我们默观这一本性时，视之为离我们并不遥远，且在我们之上，不是按位置，而是按其本身可畏而奇妙的卓越，并且如此地，它以其本身的光明与我们同在。然而，在此尚未有圣三向我们明显呈现，因为在那光耀中，我们未能恒心保持心灵之眼专注寻求它；我们只是以某种方式辨认出它，因为那里没有体积，以致我们必须认为二或三的大小大于一。但当我们论及爱——在圣经中被称为天主的爱——时，一个圣三便开始微微向我们显露，\Emph{即}：爱者、被爱者与爱。然而，因为那不可言喻的光明击退了我们的注视，并且在一定程度上显明，我们心智的软弱尚不能与之相称，我们便在已开始的途中折返，并规划根据（可以说）对我们自己心智的更熟悉的思考——按此，人是按天主的肖像被造的——以缓解我们过度紧张的心神；从第九卷到第十四卷，我们便专注于考虑我们所是的受造物，为能藉着所造之物理解并看见天主那不可见的事。现在，我们已在较低的事物中操练了理解，尽其所需，或也许超过所需，看哪！我们渴望，却没有力量，提升自己去瞻仰那最高圣三——即是天主。因为正如我们看见最确定无疑的各种三一一样——无论是那些由外在有形事物所造的，或是当这些从外在感知到的事物本身被思想时；或是当那些源自心智、不涉及身体感官的事物（如信德，或构成生活艺术的各种德行）被清晰的理性所辨明，并被知识所把握时；或是当那心智本身——我们藉以认识我们真实知道的一切——被自身认知或思考自身时；或是当那心智瞻仰任何永恒不变、非其自身的存有时——如此，我说，正如我们在所有这些实例中看见最确定无疑的各种三一一样，因为它们在我们内被造，或在我们内，当我们记起、注视或渴求这些事物时——我们如此也看见那即是天主的圣三吗？因为在那里，我们也藉着理解看见祂如同说话，以及祂的圣言，\Emph{即}圣父与圣子；然后，由此而出，二者共有的爱，\Emph{即}圣神？这些属于我们感官或心智的各种三一，我们是看见它们多于相信，却相信天主是圣三多于看见？但若是如此，那么无疑，我们根本不能藉着所造之物理解并看见天主那不可见的事；或者，若我们果真看见，我们并没有在其中看见圣三；并且其中有些事可以看见，有些事即使看不见，我们也应当相信。正如第八卷表明我们看见那不变之善——我们并非那善，同样，第十四卷在我们论及人从天主所得的智慧时，提醒了我们这一点。那么，为什么我们不在其中认出圣三呢？那被称为天主的智慧，岂不自我认识、自我爱慕吗？谁会说这样的话？或者，哪里没有知识，哪里就绝无智慧，这有谁看不见呢？或者，我们真以为那即是天主的智慧，知道其他事物，却不自知；或爱其他事物，却不自爱？但若这是愚蠢且不虔敬的说法或想法，那么看哪，我们便有了一个三一——\Emph{即}智慧、智慧对自身的知识，以及智慧对自身的爱。因为同样，在人身上我们也发现一个三一——\Emph{即}心智、心智用以自知的知识，以及心智用以自爱的爱。\par

% structure: p0019 source=h2 target=subsection reason=relative-heading-rank; base=h2

\subsection{第七章 —— 从我们谈到的各种三一中并不容易发现天主圣三}

11. 但这三者在人身上是如此存在，以致它们本身并非人。因为人，如古人所定义的，是有理性的、有死的动物。因此，这些事物是人身的主要部分，但并非人本身。任何一个人，\Emph{即}每一个个体的人，都在其心智中拥有这三者。但若我们再如此定义人，说：人是由心智和身体构成的理性实体，那么无疑，人拥有一个非物质灵魂和一个非灵魂身体。因此，这三者并非人，而是属于人，或存在于人内。若我们再撇开身体，只思考灵魂本身，心智是灵魂中某种东西，如同其头、或眼、或面容；但这些不能被视为物体。因此，心智不是灵魂，而是灵魂中为首的部分。但能说圣三在天主内是如此存在，如同某种属于天主而非天主本身的东西吗？因此，每一个个体的人——他被称为天主的肖像，不是按属于其本性的所有事物，而只是按他的心智——是一个人，并且在他的心智中是圣三的肖像。但他是其肖像的那个圣三，在其整体中无非就是天主，在其整体中无非就是圣三。也没有任何事物属于天主的本性而不属于那圣三；并且三位同一本体，不像每个个体的人是一个位格。\par

12. 此外，在这一点上也有重大区别：无论我们谈论人心中的心灵及其知识与爱，还是谈论记忆、理解、意志，我们除非藉记忆，否则不能记忆心灵的任何事物；除非藉理解，否则不能理解任何事物；除非藉意志，否则不能爱任何事物。但在那一位天主圣三中，谁敢说：圣父除非藉圣子，便不理解祂自己、圣子或圣神，或者除非藉圣神，便不爱他们；而且祂只藉自己记忆祂自己、圣子或圣神；同样，圣子除非藉圣父，便不记忆祂自己或圣父，也不藉圣神爱他们，而是只藉自己理解圣父、圣子和圣神；同样，圣神藉圣父记忆圣父、圣子和祂自己，藉圣子理解圣父、圣子和祂自己，但只藉自己爱祂自己、圣父和圣子？——仿佛圣父既是祂自己的记忆，也是圣子和圣神的记忆；圣子既是祂自己、圣父和圣神的理解；而圣神既是祂自己、圣父和圣子的爱？谁敢对那一位天主圣三如此想象或断言？因为如果在那一位内，只有圣子为自己、为圣父并为了圣神而理解，我们就回到了旧日的荒谬，即圣父不是从自己而是从圣子成为智慧的，且智慧并非生智慧，而是圣父藉祂所生的那智慧被称为智慧。因为没有理解之处，也就没有智慧；因此，如果圣父不为自己理解自己，而是圣子为圣父理解，那么圣子肯定使圣父成为智慧的。但若对天主而言，存在即是智慧，本质等同于智慧，那么就不是圣子从圣父获得本质（这是真理），而是圣父从圣子获得本质（这是最为荒谬的虚假）。这一荒谬论调，我们在第七卷中毫无疑问地讨论过、驳斥过并拒绝了。因此，天主圣父是藉着那智慧而成为智慧的，这智慧就是祂自己的智慧；而圣子是来自圣父的智慧，圣父就是智慧，圣子由祂所生；由此推出，圣父也是藉着那理解而理解，这理解就是祂自己的理解（因为不理解的人不能成为智慧）；而圣子是圣父的理解，由这理解所生，这理解就是圣父。同样，关于记忆，也可以恰当地这样说。因为一个人若什么也不记忆，或不记得自己，怎能是智慧的呢？因此，既然圣父是智慧，圣子也是智慧，那么正如圣父记忆自己，圣子也记忆自己；正如圣父记忆自己和圣子，不是藉圣子的记忆，而是藉祂自己的记忆，同样，圣子记忆自己和圣父，不是藉圣父的记忆，而是藉祂自己的记忆。再者，没有爱的地方，谁能说有什么智慧呢？因此，我们必须推断，圣父是祂自己的爱，正如祂是自己的理解和记忆一样。所以，在这至高而不变的本体即天主中，这三者——即记忆、理解、爱与意志——我们看见，并非圣父、圣子和圣神，而只是圣父。因为圣子也是由智慧所生的智慧，正如圣父和圣神不为祂理解，而是祂为自己理解；同样，圣父不为祂记忆，圣神也不为祂爱，而是祂为自己记忆和爱：因为祂自己也是祂自己的记忆、祂自己的理解和祂自己的爱。但祂之所以如此，来自祂所由生的圣父。又因为圣神也是由智慧所发出的智慧，祂也没有圣父作为记忆，圣子作为理解，自己作为爱：因为若另有一位为祂记忆，又有一位为祂理解，而祂只为自己爱，祂就不是智慧了；但祂自己拥有这三者，且拥有它们的方式是它们就是祂自己。但祂之所以如此，来自祂所由出者。\par

13. 那么，有谁能够理解那种智慧——天主籍此认识万有，以至于既非我们所谓的过去在那里是过去，也非我们所谓的未来在那里被当作将要来临的、仿佛缺席的事物等待，而是过去和未来与现在的事物全都呈现；也并非逐个思想事物，以致思想从一物转到另一物，而是所有事物同时在一瞥之中呈现？有谁，我请问，能理解那种智慧，以及类似的明达和类似的知识？因为实际上连我们自己的智慧也超出我们的理解。因为我们以某种方式能够察觉那些呈现在感官或理解力面前的事物；但那些缺席的、曾经呈现过的事物，我们若未遗忘，便通过记忆知道。我们也推测——不是从未来推测过去，而是从过去推测未来——但凭着所有不稳定的知识。因为有些思想，虽然是未来的，我们却仿佛以更清晰和更确定的方式向前展望，因为它们非常接近；我们通过记忆做到这一点，每当可能时，尽我们所能，尽管记忆似乎不属于未来，而属于过去。这一点可以在任何词句或歌曲中检验，我们凭记忆依序说出它们；因为我们肯定不能顺序说出每一个，除非我们在思想中预见到下一个。然而，使我们预见的不是预见，而是记忆；因为直到词句或歌曲的结尾，所说的每一件事都是被预见和期待的。但在这样做时，我们不是被说成凭预见说话或歌唱，而是凭记忆；如果有人特别擅长以这种方式说出许多作品，他通常不是因为他的预见，而是因为他的记忆而受到称赞。我们知道并绝对确信这一切发生在我们心中或借由我们的心；但它如何发生，我们越想仔细审视，就越发感到措辞失效，意图本身也失败，当我们想通过我们的理解（如果不能通过舌头）达到某种清晰时。像我们这样的人，在如此大的心灵软弱中，是否认为可以理解天主的预见与祂的记忆和祂的理解是同一的？祂并不逐个思想每一事物，而是在永恒、不变、不可言喻的一瞥中拥抱祂所知的一切。因此，在这种困难和困境中，我们不妨向永生天主呼求：\BibleQuote{这样的知识，为我太奇妙，太高深，我无法达到。}因为我通过自己理解祢的知识是多么奇妙和不可理解，祢籍此创造了我，而我甚至不能理解祢所创造的我！然而，\BibleQuote{我默思时，烈火熊熊，}以至于\BibleQuote{我不断寻求祢的面容。}\par

% structure: p0023 source=h2 target=subsection reason=relative-heading-rank; base=h2

\subsection{第八章——宗徒如何说如今我们是通过一面镜子看见天主。}

14. 我知道智慧是一种非物质的实体，它是那光，籍此看见那些肉眼看不见的事物；然而一位如此伟大而属神的人（如保禄）说：\BibleQuote{我们现在是借着镜子观看，模糊不清，到那时，就要面对面了。}若问这\Quote{镜子}是什么、是怎样的，我们心中自然想到：在镜子里，除了影像什么也看不见。我们便努力这样去做，以便能通过这个我们自己的影像，如同通过一面镜子，看见创造我们的那一位。这也被同一位宗徒的其他话所暗示：\BibleQuote{我们众人以揭去面纱的面容，好像镜子一样反映主的光荣，被变成同一肖像，光荣上加光荣，都是出于主的神。}他说\BibleQuote{被变成同一肖像}，即通过镜子观看，而不是从瞭望台眺望：这在希腊文里没有歧义，宗徒书信就是从希腊文译成拉丁文的。因为在希腊文中，镜子（看得见事物影像的）与瞭望台（从高处远眺）在字音上完全不同。显然，宗徒用\OriginalTerm{观看}{speculantes}一词来指主的光荣，意思是来自\OriginalTerm{镜子}{speculum}，而非来自\OriginalTerm{瞭望台}{specula}。但他说\BibleQuote{被变成同一肖像}，无疑是指天主的肖像；他称之为\Quote{同一}，意思是我们在镜子中看见的那个肖像本身，因为同一肖像也是主的光荣；正如他在别处所说：\BibleQuote{男人当然不应蒙头，因为他是天主的肖像和光荣}——这段经文已在第十二卷中讨论过。他说的\Quote{被变成}，意思是我们从一种形式变成另一种形式，从暗淡的形式变成光明的形式：因为暗淡的形式也是天主的肖像；如果是肖像，那当然也是\Quote{光荣}，我们作为人受造于其中，优于其他动物。因为关于人性自身，圣经说：\BibleQuote{男人不应蒙头，因为他是天主的肖像和光荣。}这受造物在受造物中是最卓越的，当它被自己的创造者从不虔诚中成义时，就从损毁的形式变成美丽的形式。因为即使在不信本身中，缺陷越应受谴责，本性就越应受赞美。因此他加了\BibleQuote{光荣上加光荣}：从受造的光荣到成义的光荣。虽然这些词\Quote{光荣上加光荣}也可理解为其他方式——从信德的光荣到享见的光荣，从我们作为天主子女的光荣到我们将像祂一样的光荣，因为\BibleQuote{我们将看见祂实在怎样。}但他加了\BibleQuote{都是出于主的神}，表明如此可欲的改变之恩是由天主的恩宠赐予我们的。\par

% structure: p0025 source=h2 target=subsection reason=relative-heading-rank; base=h2

\subsection{第九章——论\Quote{谜语}一词及比喻性言说方式。}

15. 以上所述涉及宗徒的话，即\BibleQuote{我们现今藉着镜子观看}；但他又加了\BibleQuote{藉着谜语}，这加添的意思，凡不熟悉那些包含言辞方式之学的书籍（希腊人称这些方式为\OriginalTerm{转义}{Tropes}，我们拉丁语也使用这个希腊词）的人，都无法知晓。因为我们更常称\Emph{\OriginalTerm{形相}{schemata}}而非形象，同样更常称\Emph{\OriginalTerm{转义}{tropes}}而非方式。而且，用拉丁文表达各个方式或转义的名称，并为每个指定其适当名称，是一件非常困难和罕见的事。因此，有些拉丁译者，因不愿使用希腊词，在宗徒说\BibleQuote{这些都是寓意}之处，便用绕圈子的说法译出——即这些事以另一物表示一物。但这种称为寓意的转义有好几种，其中之一便是所谓的谜语。属类名称的定义必然涵盖其所有种类；因此，正如每匹马都是动物，但并非每个动物都是马，同样，每个谜语都是寓意，但并非每个寓意都是谜语。那么，寓意不就是一种转义，在其中一物从另一物被理解吗？例如在\WorkTitle{得撒洛尼人前书}中有：\BibleQuote{所以我们不要像别人一样睡觉，但要警醒谨守：因为睡觉的人是在夜间睡觉，醉酒的人是在夜间醉酒；但我们属于白昼，应当谨守。}但这个寓意并非谜语，因为这里的意思对所有人（除了极其愚钝者）都是明白的。而谜语，简言之，是隐晦的寓意，例如\Quote{水蛭有三个女儿}，以及其他类似例子。但当宗徒提到寓意时，他并非在言辞中，而是在事实中找到了它；因为他藉着亚巴郎的两个儿子——一个由婢女所生，另一个由自由妇人所生——表明了两个盟约应如何理解，这不仅是言语，也是行动。而在解释之前，这是隐晦的；因此，这样的寓意（即属名）可特别称为谜语。\par

16. 但不仅那些不熟悉包含转义之学的书籍的人会询问宗徒所说\BibleQuote{我们现今在谜语中看见}的意思，那些熟悉此学但想知道我们现在所见的谜语是什么的人也会询问；我们必须为这两个短语找到一个共同含义，\Emph{即}\BibleQuote{我们现今藉着镜子观看}和\BibleQuote{藉着谜语}。因为当整句话如此说出时，它只构成一个句子：\BibleQuote{我们现今藉着镜子在谜语中观看}。因此，依我判断，正如他用\Quote{镜子}一词意指一个图像，同样他用\Quote{谜语}意指任何你想要的相似性，然而却是隐晦且难以看透的。因此，当宗徒谈到镜子和谜语时，尽管任何相似物都可以被理解为所指，以便适应于对天主的理解（以祂可被理解的方式）；然而，没有比那不徒然地被称为祂肖像的东西更适合此目的。所以，不要有人惊奇，我们努力以任何方式看见，甚至是在此生所赐予的看见方式中，\Emph{即}藉着镜子，在谜语中。因为如果看见是容易的，我们就不会在此处听到谜语了。而这是一个更大的谜语：我们看不见我们不能不看见的东西。因为谁看不见自己的思想？然而谁能看见自己的思想？我所说的不是用肉身之眼，而是用内在的视觉本身？谁看不见它，谁又能看见它？因为思想是心灵的一种视觉；无论是那些也被肉眼所见或由其他感官所感知的事物当前呈现；或它们不在场，但它们的相似性被思想所辨识；或者两者都不是，所思之物既非形体之物，也非形体之物的相似性，如德行和恶习；或如思想本身被思想；或是那些属于教学和自由学科的内容；或者是在不变的本性中所有这些事物的更高原因和理由本身被思想；或者甚至是邪恶、虚妄和虚假的事物被我们思考，而感官要么不同意，要么在同意中犯错误。\par

% structure: p0028 source=h2 target=subsection reason=relative-heading-rank; base=h2

\subsection{第十章——论心灵的言语，我们在其中如同在镜子和谜语中看见天主的圣言。}

17. 然而，现在让我们来谈谈那些我们认为已知的、即使没有思考也存在于我们知识中的事物；无论它们属于默观的知识（如我所论证的，这本应被称为智慧），还是属于行动的知识（这本应被称为知识）。因为两者同属一个心智，且同是天主的肖像。但当我们分别且单独地处理这两者中较低者时，它便不应被称为天主的肖像，尽管即使在那个时候，其中也能找到天主圣三的某种相似，如我们在第十三卷中所指出的。因此，我们现在谈论的是人的整体知识，其中我们所知道的一切都是已知的；至少是真实的知识；否则它就不会被知道。因为没有人知道虚假，除非他知道那是虚假的；而如果他知道这一点，那么他就知道了真实：因为知道那是虚假的乃是真实的。因此，我们现在讨论那些我们认为已知的事物，即使它们没有被思考，也是为我们所知的。但确实，如果我们想用言语表达它们，我们只能通过思考来做到。因为即使没有说出口，思想者也在他心里说话。因此，\WorkTitle{智慧篇}中有这样的话：\BibleQuote{他们心里说，思想不正。}因为那些话，\BibleQuote{他们心里说，}通过加上了\BibleQuote{思想}而得到解释。类似的一处是在福音中——某些经师听到主对瘫子说：\BibleQuote{孩子，放心！你的罪赦了。}便心里说：\BibleQuote{这人亵渎了。}他们除了通过思考，又如何能\BibleQuote{心里说}呢？接着是：\BibleQuote{耶稣看见了他们的思念，就说：你们为什么心里思念恶事？}以上是\Person{玛窦}所记。但\Person{路加}这样叙述同一件事：\BibleQuote{经师和法利塞人开始心里想，说：这人是谁？竟说亵渎话！除了天主一个外，谁能赦罪？但耶稣看透了他们的心思，就回答他们说：你们心里忖度什么呢？}在\WorkTitle{智慧篇}中，\BibleQuote{他们心里说，思想}与这里的\BibleQuote{他们心里想，说}是相同的。因为两处都宣称，他们在自己心里说话，\Emph{即}通过思考说话。因为他们\BibleQuote{心里说}，就有人对他们说：\BibleQuote{你们想什么呢？}主自己也论到那个地土出产丰富的财主说：\BibleQuote{他就在心里想，说。}\par

18. 因此，有些思想是内心的话语，主也指示存在一个口，当祂说：\BibleQuote{不是入口的污秽人，而是出口的才污秽人。}在一句话里，祂包含了人的两个不同的口：一个是身体的口，一个是内心的口。因为确实，那些人以为使人污秽的东西进入身体的口；而主所说的使人污秽的东西是从内心的口出来的。所以，主亲自解释了自己所说的话。因为稍后，祂又对门徒论及同一件事说：\BibleQuote{你们还不明白吗？难道你们不晓得，凡入口的，是进到肚腹，然后排泄到厕所里去吗？}这里祂明确地指向身体的口。但在随后的话中，祂清楚地谈到内心的口，祂说：\BibleQuote{但从口里出来的，是由心里发出来的；这些才污秽人。因为从心里发出来的是恶念，}等等。还有什么比这解释更清楚的呢？然而，当我们称思想为内心的话语时，这并不意味着当思想是真实的时候，它们不是同样出于知识的视觉的视觉行动。因为当这些事通过身体外在进行时，言语和视觉是不同的东西；但当我们内心思考时，两者是一——正如视觉和听觉在身体感官中是相互区别的两件事，但在心智中看见和听见却是同一件事；因此，虽然言语在外不是被看见而是被听见，然而内在的言语，\Emph{即}思想，却被\WorkTitle{福音}记载说是被主看见，而非听见的。\WorkTitle{福音}说：\BibleQuote{他们心里说：这人亵渎了。}随后又加上：\BibleQuote{耶稣看见了他们的思念。}因此，祂看见了他们所说的话。因为祂以自己的思想看见了他们的思想，而这些思想原是他们认为除了自己以外无人看得见的。\par

19. 凡能在思想中理解话语的人，不仅是在它发出声音之前，甚至是在考虑其声音的意象之前——因为这不属于任何语言，即那些所谓的民族语言（其中我们的拉丁语就是其中之一）——凡能理解这一点的人，现在就能藉着这镜子，在这谜中看到那圣言之相似的某方面，关于那圣言，经上说：\BibleQuote{在起初已有圣言，圣言与天主同在，圣言就是天主。} 因为必然地，当我们说真话时，即说出我们所知道的东西，从记忆所保持的知识本身，就生出一个与它产生的知识完全同类的语言。因为由我们所认知的事物而形成的思想，就是我们心中所说的语言；这种语言既不是希腊语，也不是拉丁语，也不是任何其他语言。但是，当需要将这个知识传达给我们所说话的对象时，就会采用某个标记来指示它。通常是一种声音，有时是一个点头，前者给耳朵，后者给眼睛，使我们心中所怀的语言也能藉着身体的标记，为身体的感官所知。因为点头或示意不就是以某种方式向视觉说话吗？圣经也为此作证；因为我们读到\BibleBook{若望}{福音}记载：\BibleQuote{我实实在在告诉你们：你们中有一个要出卖我。}门徒便互相观望，猜疑祂说的是谁。耶稣所爱的那个门徒那时斜倚在耶稣的怀里。西满伯多禄便向他示意说：\BibleQuote{祂说的是谁？} 在这里，他藉着示意说出了他不愿用声音说出的话。但是，当我们向在场的人展示这些及类似的身体标记（要么给耳朵，要么给眼睛）时，为了也能与不在场的人交谈，就发明了文字；但它们是话语的标记，正如话语本身是我们谈话中那些我们所想的事物的标记。\par

% structure: p0032 source=h2 target=subsection reason=relative-heading-rank; base=h2

\subsection{第十一章——圣言的相似，正如它那样，不应在我们外在和可感觉的话语中寻求，而应在内在和心智的话语中寻求。我们的话语和知识与神圣的圣言和知识之间有最大的可能不相像。}

20. 因此，那外在发声的言语，是那内在光照之言语的标记；后者更有权利被称为言语。因为那用血肉之口发出的，是言语的清晰声音；它本身也被称为言语，是因为它被用以将内在言语显现于外。原来我们的言语，藉着采用那清晰声音，以向人的感官显现，在某种意义上被制成一个身体的声音，正如天主的圣言藉着採用那肉体，以向人的感官显现，而成了血肉。正如我们的言语变成了清晰声音，却没有变成声音；同样，天主的圣言成了血肉，但我们绝不能说祂被变成了血肉。因为我们的言语变成了清晰声音，而那个圣言成了血肉，都是藉着采用，而不是藉着消耗自身以致变成那东西。因此，凡愿意以某种相似——无论何等不完全——来理解天主圣言的人，尽管在许多方面不同，都不应注视那在耳中发声的我们的言语，不论它是清晰发出还是默默思想。因为所有语言的话语，凡在声音中发出的，也都被默默思想；当肉体的口沉默时，心灵会默念诗句。不仅如此，音节的数量和歌曲的曲调，由于它们是属物质的，并属于那称为听觉的身体感官，对于那些思想它们、并默默转动这一切的人来说，也借着某些与之相对应的无形影像而呈现。但我们必须越过这些，以便达到人的言语，借着它的相似——无论何等不完全——能好像在奥秘中\OriginalTerm{奥秘}{\GreekText{αἴνιγμα}}看见天主圣言。不是那向这位或那位先知所说的言语，关于它经上记载说：\Quote{天主的言语越传越广，日见增多}；又说：\Quote{信德是出于听闻，听闻是出于基督的话}；又说：\Quote{你们从我这里听受了天主的言语，你们接纳了，并不以为是人的言语，而实在是天主的言语}（还有无数类似的经文提到天主的言语，它藉着许多不同语言的声音，通过人的心和口传播；它之所以被称为天主的言语，是因为所传授的道理不是人性，而是神性）；——但我们现在寻求的是，藉着这个相似，以某种方式看见那被称为\Quote{圣言就是天主}的天主圣言，关于它说\Quote{万有是藉着祂造成的}，关于它说\Quote{圣言成了血肉}，关于它说\Quote{至高者的圣言是智慧的泉源}。因此，我们必须进到人的言语，即理性动物的言语，即那按天主的肖像被造——不是由天主所生，而是由天主所造——的言语；这言语既不能以声音说出，也不能在声音的相似之下被思想（如同任何语言之言语必然如此），而是先于所有用以指明它的记号，并从那存留在心灵中的知识所生，当这同一知识按照它真实的样子被内在说出时。因为思想的视觉与知识的视觉极其相似。当它藉着声音或任何身体记号被说出时，它并非按照它真实的样子被说出，而是按照身体所能看见或听见的样子。因此，当言语中有知识中所拥有的，那么就是真实的言语和真理——这正是人身上所寻求的；如此，知识中所拥有的也在言语中，知识中所没有的也不在言语中。这里可以认出\Quote{是就说是，非就说非}。于是，这被造的肖像的相似之最大可能接近于那被生的肖像的相似——藉此，圣子天主被宣告在实质上完全相似于圣父。我们还必须在此奥秘中注意天主圣言的另一个相似，\Emph{\OriginalTerm{即}{videlicet}}，正如关于那圣言所说\Quote{万有是藉着祂造成的}，那里宣告天主藉着祂的独生子造成了宇宙，同样，人的任何工作都首先在内心被说出；因此经上记载：\Quote{言语是一切工作之始}。但在这里，只有当言语是真实的，才是良好工作之始。真实的言语是从行善的知识所生，以便那里也保持\Quote{是就说是，非就说非}；以便那使我们生活的知识中所有的，也都在我们藉以工作的言语中，而一者中没有的，另一者中也没有。否则，这样的言语是谎言而非真理；由此产生的便是罪，而非善工。在这个我们言语的相似中，还有另一个与天主圣言的相似：我们的言语可以没有后续的工作，但任何工作都不能没有先行的言语；正如天主圣言即使没有受造物也能存在，但任何受造物都不能存在，除非藉着那造成万有的圣言。因此，不是天主圣父，不是圣神，也不是圣三本身，而只有那作为天主圣言、藉着万有被造的圣子，成了血肉；虽然圣三是创造者。这是为了使我们能通过我们的言语，跟随并效法祂的榜样，即在我们言语的思想和工作上都没有谎言，而正当地生活。但这一肖像的完美要在将来某个时刻才能实现。为达到此，那善师藉着基督信仰和敬虔的道理教导我们：\Quote{以没有蒙蔽的脸面}，脱离法律的面帕（法律是未来之事的阴影），\Quote{如同镜子观看主的荣耀}，\Emph{\OriginalTerm{即}{id est}}通过镜子凝视，\Quote{我们成为同一肖像，从光荣到光荣，正如由主的神}；正如我们上面所解释的。\par

21. 因此，当这肖像借此转化得以完美更新时，我们就要相似天主，因为我们看见祂，不是透过镜子，而是\BibleQuote{祂怎样；}宗徒保禄称之为\BibleQuote{面对面。}但现在，在这镜子中，在这谜语中，在这如此的肖像中，又有何其大的不相似，谁能解释呢？然而，我将尽己所能，指出某些要点来表明它。\par

% structure: p0035 source=h2 target=subsection reason=relative-heading-rank; base=h2

\subsection{第十二章——学院派哲学}

首先，人所能达到的知识本身，无论他何等精通博学，其性质如何、有多大程度？当我们讲述我们所知道的事物时，我们的思想便是由此以真理形成的。暂且不谈那些来自身体感官、进入心灵的事物——其中有许多与它们所显现的迥异，以至于过分被它们与真理的相似性所压迫的人，在自己看来似乎是清醒的，实际上却并不清醒；正因如此，学园派哲学才如此盛行，以至于因怀疑一切而更加可怜地疯狂——暂且抛开这些通过身体感官进入心灵的事物，那么，还有多少事物是我们以知道自己活着的那种方式所知道的？关于这一点，我们绝对毫不担心，唯恐我们或许被某种真理的相似性所欺骗；因为可以肯定，被欺骗的人仍然活着。而这又不属于那些从外部呈现的视觉对象之列，以致眼睛可能在其中受骗；就像桨在水中看似弯曲，航船驶过时塔楼似乎在移动，以及成千上万其他与表象不同的事物那样：因为这不是肉体之眼所能辨别的。我们知道自己活着的那种知识，是一切知识中最内在的，就连学园派也无法暗示：\Quote{也许你睡着了，自己却不知道，你在梦中见到这些事物。}因为谁不知道，人们在梦中见到的东西恰恰与醒时所见相同？但一个确知自己生命知识的人，不会因此说：\Quote{我知道我醒着，}而是说：\Quote{我知道我活着；}所以，无论他是睡着还是醒着，他都活着。他也绝不会因做梦而在那种知识上受骗；因为睡觉和做梦都属于活着的人。学园派也不能为了反驳这种知识而说：\Quote{也许你疯了，自己却不知道：}因为疯人所见恰恰与神志清醒者所见相同；但疯人仍然活着。他也不回答学园派说：\Quote{我知道我没有疯，}而是说：\Quote{我知道我活着。}因此，说自己知道他活着的人，既不会被欺骗，也不会说谎。那么，就让一千种欺骗性的视觉对象呈现在那说\Quote{我知道我活着}的人面前吧；他并不会惧怕其中任何一个，因为被欺骗的人仍然活着。但如果只有这类事物属于人类知识，那它们确实很少；除非它们能在每一类中如此倍增，以至于不但不少，而且最终达到无限。因为说\Quote{我知道我活着}的人，说他只知道一件事。此外，如果他说\Quote{我知道我知道我活着}，那么就有两件事；而他知道这两件事，这是第三件要知道的。如此，他可以加上第四件、第五件，乃至无数件，只要他继续下去。但是，由于他既不能通过逐个相加来把握无限的数量，也不能说出无限次数的事物，他至少把握了这一点，并且完全确实地把握了，\Emph{即}：这既是真实的，又是如此无限，以至于他无法真实地把握并说出其无限数目。同样的事情也可以在一个确定的意志上注意到。因为如果有人说\Quote{我愿幸福，}那么回答\Quote{也许你被欺骗了}将是厚颜无耻的。如果他说\Quote{我知道我愿如此，并且我知道我知道它，}那么他还可以给这两者加上第三件，\Emph{即}：他知道这两件事；以及第四件：他知道他知道这两件事；如此\OriginalTerm{至无限}{ad infinitum}。同样，如果有人要说：\Quote{我不愿被欺骗，}那么无论他是否被欺骗，他确实不愿被欺骗，这难道不是真实的吗？对他说\Quote{也许你被欺骗了}，岂不是极其厚颜无耻？因为毫无疑问，无论他在何处被欺骗，他在认为自己不愿被欺骗这一点上却绝没有被欺骗。如果他说他知道这一点，他就添加了他所选择的任意数量已知事物，并察觉到那个数量是无限的。因为说\Quote{我不愿被欺骗，我知道我不愿如此，并且我知道我知道它}的人，现在已经能够在此处也陈述一个无限的数量，尽管表达起来可能有些笨拙。还有其他事物也可以用来反驳那些主张人一无所知的学园派。但我们必须限制自己，尤其因为这并非我们在当前著作中所担负的主题。关于那个主题，我们有早年归化时所写的三卷书；任何能够并且愿意阅读并理解它们的人，无疑都不会被他们为了反对真理的发现而提出的许多论点所动摇。因为可知事物有两种——一种是心灵通过身体感官所觉察的，另一种是心灵通过自身所觉察的——这些哲学家对肉体感官大加喋喋不休，却从未能动摇那些心灵自身知道的、关于真实事物的最确实的知觉，例如我所提到的\Quote{我知道我活着}。但切莫让我们怀疑我们通过身体感官所学到的真理；因为通过它们，我们得以认识天和地，以及其中为我们所知的那些事物，正如创造我们和它们的那一位愿意它们被我们所认识那样。也切莫让我们否认，我们知道那些通过他人见证所学到的事物；否则我们就不知道有海洋；我们就不知道那由大量报道所传颂的土地和城市存在；我们就不知道那些人和他们的作品，这些是我们通过阅读历史所学到的；我们就不知道每天从各处带来的消息，并由一致和相合的证据所确认；最后，我们就不知道我们在何处出生、从何人所生：因为在这一切事情上，我们都相信了他人的见证。如果这样说最为荒谬，那么我们就必须承认，不仅我们自己的感官，而且他人的感官，确实极大地增加了我们的知识。\par

22. 这一切，无论是人心藉自身所知的，或是藉身体感官所知的，或是从他人的见证所接受而知的，都存留并保存在记忆的宝库中；由此产生一个真实的话语，当我们说出我们所知时，这话语先于一切声音，先于一切对声音的思想。因为当思想之视源自知识之视时，话语最像那被认知的事物，从中也产生了它的肖象；当它是不属于任何语言的话语，而是关于真实事物的真实话语时，它不拥有任何属于自己的东西，而全然源自产生它的知识。一个人说出他所知时，何时学会的并不重要；因为他有时在学会后立即说出；只要话语是真实的，\Emph{即}源自已知之物。\par

% structure: p0038 source=h2 target=subsection reason=relative-heading-rank; base=h2

\subsection{第十三章——再论我们心灵的知识与话语不同于天主的知识与圣言。}

但情况是这样吗？天主圣父，从祂生出那圣言，\Emph{即}出自天主的圣言——那么，天主圣父，就祂所是的智慧而言，是藉身体感官学到一些事物，还是藉祂自身呢？谁能这样说，谁将天主视为并非理性的受造物，而是超越理性灵魂的一位？至少那些将祂置于一切受造物和一切灵魂之上的人，虽仍在镜中模糊地猜测，尚未面对面看见祂，也只能如此思想。难道天主圣父所知道的那些事物，不是藉身体（因祂没有身体）或藉自身从别处得知，而是从某位那里得知的吗？或者祂需要使者或见证才能知道吗？当然不是；因为祂自身的圆满使祂能知道一切祂所知的事。无疑，祂有使者，\Emph{即}天使；但并非向祂报告祂不知道的事，因为没有任何事是祂不知道的。天使的好处在于就自己的作为咨询真理。所谓他们向祂报告某些事，意即不是祂从他们学习，而是他们藉祂的圣言（无声音的）从祂学习。他们也向祂报告祂所愿意的事，被祂派遣到祂所愿意的人那里，并藉祂的圣言听到一切，即在祂的真理中发现他们自己该做什么：向谁、在何时报告什么事。因为我们也向祂祈祷，却并非告知祂我们的需求。\Quote{因为你们的父，}祂的圣言说，\Quote{在你们求祂以前，已知道你们需要什么。}祂并非在某个特定时间才认识这些事而知道；而是毫无开端地预知了时间中将要发生的一切，包括我们将求祂什么、何时求，以及祂将应允或不允何人何事。关于祂所有的受造物，无论是精神的还是物质的，祂认识它们并非因为它们存在，而是它们存在因为祂认识它们。因为祂并非不知道祂将创造什么；因此祂创造是因为祂知道；祂不是因为创造才知道。祂在创造后认识它们的方式，与创造前认识它们的方式毫无不同，因为祂的智慧并未从它们那里增加什么；智慧保持不变，而它们按合适的时间出现。同样，在\BibleBook{德训}{篇}中记载：\BibleQuote{万物在受造之前，祂都认识；在完成之后，也是如此。}如此，他说，并非不同；在受造前和完成后，都是这样被祂认识。因此，这种认识远非我们的认识可比。天主的认识本身就是祂的智慧，祂的智慧就是祂的本质或实体。因为在那个本性的奇妙单纯中，有智慧与存在不是两回事，而是有智慧就是存在；正如我们之前几卷中多次说过的。但我们的认识在大多数情况下可能失去也可能重新获得，因为对我们来说，存在不同于认识或有智慧；即使我们不知道、不智慧于从别处学来的东西，我们仍可能存在。因此，正如我们的认识不同于天主的认识，同样，从我们认识所生的话语也不同于从圣父本质所生的天主的圣言。这就如同我说：从圣父的知识、从圣父的智慧所生；或者更准确地说，从是知识的圣父、从是智慧的圣父所生。\par

% structure: p0040 source=h2 target=subsection reason=relative-heading-rank; base=h2

\subsection{第十四章——天主的圣言在一切方面与生出祂的圣父同等。}

23. 于是，天主的圣言、圣父的独生子、在一切事上与父相似且相等，出自天主的天主，出自光明的光明，出自智慧的智慧，出自本体的本体，全然是父所是的一切，然而不是父，因为一位是子，另一位是父。因此，凡父所知道的，子都知道；但子知道如同子存在一样，都是来自父，因为在那里知道与存在是同一的。所以，正如存在不是子给予父的，知道也不是。因此，父如同发出自己，生了与自己完全相等的圣言；因为如果在他的圣言中有任何多于或少于他自己，他就不会完全而圆满地发出自己。此处以最高意义认出了：\Quote{是，是；非，非。}因此，这圣言是真正的真理，因为在那个生出它的知识中所有的，也都在它内；不在那知识中的，也不在圣言内。这圣言绝不可能有虚假，因为它是不可改变的，如同生它的那一位。因为\Quote{子不能由自己作什么，惟看见父所作的，子才能作。} 这不是由于能力不足，而是由于力量，真理不能是虚假。因此，天主圣父在自己内知道一切，在子内知道一切；但在自己内如同知道自己，在子内如同知道自己的圣言，这圣言是有关他自己内的一切而说的。同样，子也知道一切，\Emph{即}在子自己内，作为由父在自己内所知道者而生的那一切；在父内，作为由父那里生出、且子在自己内知道的那一切。因此，父与子互相知道；但一位是由于生，另一位是由于被生。他们每一位都同时看见他们知识、智慧、本体中的一切：不是分部分地或逐一地，仿佛时而看这边、时而看那边，又从这边或那边转向这边或那边，以致无法同时看见某些而不忽略另一些；而是，如同我说过的，同时看见一切，其中没有一样不是他始终看见的。\par

24. 那么，我们那既无声音也无声音之思的言语，乃是出于我们在内心看见并说话时的那件事，因此不属于任何语言；故此，在这谜语里，某种程度上相似于那亦是天主的圣言；因为我们的言语也生于我们的知识，正如那圣言生于父的知识：这样的言语（我说我们的言语），我们发现它某种程度上相似于那圣言，让我们也不迟缓地思量它多么不相似，正如我们可能尽力说出的那样。\par

% structure: p0043 source=h2 target=subsection reason=relative-heading-rank; base=h2

\subsection{第十五章 我们的言语与天主圣言之间的不同何其大。我们的言语不能是或被称为永恒的。}

那么，我们的言语仅从我们的知识而生吗？我们不也说许多我们不知道的事吗？而且说时并非怀疑，而是以为是真的；若偶然就事物本身而言是真的，但对于我们的言语却不是真的，因为言语除非生于已知的事物，否则不是真的。就此意义而言，我们的言语是假的，不是当我们说谎时，而是当我们受骗时。当我们怀疑时，我们的言语还不是关于我们所怀疑的事物的，而是关于怀疑本身的言语。因为虽然我们不知道所怀疑的是否为真，但我们确实知道我们在怀疑；因此，当说\Quote{我们怀疑}时，我们说的是真话，因为我们说的是我们所知道的。还有我们可能说谎的事呢？当我们说谎时，我们确实既有意又明知地有一个虚假的言语，其中却有一个真实的言语，\Emph{即}我们说了谎，因为我们知道这个。当我们承认我们说了谎时，我们说的是真话；因为我们说的是我们所知道的，因为我们知道我们说了谎。但那圣言是天主，能比我们做更多，它却不能做这事。因为它\Quote{不能作什么，惟看见父所作的}，并且它\Quote{不凭自己说话}，而是从父那里得到它所说的一切，因为父以特别的方式说它；那圣言的大能在于它不能撒谎，因为那里不能有\Quote{是而非}，只有\Quote{是，是；非，非。} 那么，那不真的，甚至不应该被称为言语。我欣然同意，若如此。那么，若我们的言语是真的，因此能正确地被称为言语，难道我们也能像说\Quote{看见之看见}、\Quote{知识之知识}那样说\Quote{本体之本体}，如同那圣言特别被说且特别应当被说的吗？为何如此？因为对我们而言，存在不等于知道；因为我们知道许多事情，它们以某种方式靠记忆活着，并以某种方式因遗忘而死去；因此，当那些事物不再在我们的知识中时，我们仍然存在；而当我们的知识从我们的心灵中滑落消逝时，我们仍然活着。\par

25. 至于那些以如此方式被知晓，以致永不能脱离记忆的事物——因为它们当前存在，且属于心灵自身的本性——\Emph{例如}，我们知道自己活着（这事持续多久，只要心灵持续；并且，因为心灵恒常持续，这事也恒常持续）——我说，对于此事以及任何其他类似的实例，我们更应在其中默观天主的肖像，却难以辨明：尽管它们始终被知晓，但正因为它们并不始终被思想，如何能就它们说出一个永恒的言语，当我们的言语在我们的思想中被说出时。\newline{}因为灵魂活著是永恒的，知道它活著是永恒的。\newline{}然而，思想自己的生命，或思想对自己生命的认知，对灵魂而言却不是永恒的；因为当它进入这样或那样的活动时，它会停止思想此事，尽管它不会停止知道它。\newline{}由此可知，如果心灵中能有任何永恒的知识，而对此知识的思虑却不能永恒，并且我们内在真实的言语唯有通过我们的思想说出，那么只有天主能理解有一个与祂自身同永恒的圣言。\newline{}除非我们或许要说：那思想的可能性——因为所知之物即使在不被思想时也能被真实地思想——构成了一个与知识本身同样永恒的言语。\newline{}但那尚未成形于思想观照中的东西如何是言语呢？如果它没有那知识的形态，并且仅仅因为它能拥有它而被称为言语，它如何肖似生出它的知识呢？这好比有人说，一个东西之所以应称为言语，是因为它能成为言语。但那能成为言语，因而已被认为配得上言语之名的东西又是什么呢？我说，那可成形却尚未成形的东西，除了我们心灵中的某个东西，当我们忽而思想这事，忽而思想那事时，我们通过来回翻转它而投掷它，直到它们被找到或向我们呈现，那东西又是什么呢？真实的言语只有在这样的情况下才产生：如我所说，我们来回翻转的东西到达我们所知的上面，并被那所知形成，以此取得它的完全相似；以至于每一事物如何被知晓，就如何被思想，\Emph{即}，以这种方式在心中被说出，没有清晰的声音，没有清晰声音的思想——这种思想无疑属于某种特定的语言。因此，即使我们为了不在名称上争论不休而承认：我们心灵中的这个从我们的知识中能被形成的东西，甚至在它被成形之前就已经可以称为言语，因为它是，可以说，已经是可成形的，那么谁看不出它与天主的圣言之间的极大不同？天主的圣言是在天主的形态中，以致它在被成形之前不是可成形的，也从未在任何时候可能无定形，而是单纯的形式，与祂所来自的那位单纯相等，并与祂奇妙地同永恒？\par

% structure: p0046 source=h2 target=subsection reason=relative-heading-rank; base=h2

\subsection{第16章——我们的言语永远不能与天主圣言相等，即使当我们肖似天主时也不能。}

因此，天主的圣言之所以如此称呼，是为了不被称为天主的思虑，免得我们相信在天主内有任何可以旋转的东西，以致它有时接受、有时恢复一种形式，以便成为一个言语，并且又能失去那形式并在某种意义上无定形地旋转。确实，那位卓越的言辞大师深知语词的效力，并洞察了思想的本质，他在诗中说道：\Quote{他与自己思虑着战争的各种结果}，\Emph{即}思想它们。因此，天主圣子不被称为天主的思虑，而被称为天主的圣言。因为我们自己的思虑，达到我们所知的，并由那所知成形，便是我们真实的言语。所以，天主的圣言应被理解为没有任何思虑在天主方面，以致它被理解为单纯形式本身，却不包含任何既可成形又可无定形的东西。确实，圣经中有一些段落提到天主的思虑；但这是按照同样的说话方式来谈的，如同圣经也提到天主的遗忘，然而在严格的语言用法中，天主内当然没有遗忘。\par

26. 因此，既然我们如今在这谜中发现了与天主和天主圣言如此大的不相似，而其中先前却找到了一些相似，这也必须承认：即使当我们肖似祂时，当我们\Quote{看见祂，正如祂所是的}（说这话的人无疑清楚我们现在的不相似），那时我们也不会在本质上与祂相等。因为受造的本性始终小于那创造的本性。那时，我们的言语确实不会是虚假的，因为我们既不撒谎，也不受骗。或许，我们的思虑也不再通过从一个事物到另一个事物来回穿行而旋转，而是一下子、一眼就看到我们所有的知识。然而，即使这事发生了，如果真的发生了，那曾是可成形的受造物确实会被成形，以至于它应当达到的那种形式一无所缺；然而，它仍然不能与那单纯性相等，在那单纯性中没有任何可成形的东西曾被成形或重新成形，而只有形式；它既不是无定形也不是已成形，本身是永恒不变的本体。\par

% structure: p0049 source=h2 target=subsection reason=relative-heading-rank; base=h2

\subsection{第17章——圣神如何被称为爱，以及是否唯独祂如此被称呼。圣神在圣经中恰当地被称为爱。}

27. 关于圣父和圣子，我们已经说了足够多，正如在这面镜子和这个谜中所能见到的。现在我们必须在天主所许可的范围内，论述圣神。根据圣经，圣神既非单独出自圣父，也非单独出自圣子，而是出自二者；这向我们启示了圣父与圣子彼此相爱的共融之爱。然而，天主的圣言为操练我们，使那些并非表面可见、而需在隐藏深处探寻并从中汲取的事物，得以更热切地寻求。因此，圣经并没有说：\Quote{圣神是爱。}若它如此说，就会消除这场探讨的不少部分。但圣经说：\Quote{天主是爱；}因此不确定且有待探究的是：究竟是圣父是爱，圣子是爱，圣神是爱，还是整个天主圣三本身是爱。因为我们不会说天主被称为爱，是因为爱本身是一个堪当天主之名的实体，而是因为爱是天主的恩赐，正如对天主所说：\Quote{祢是我的忍耐。}这并非因为我们的忍耐是天主的本质，而是因为他亲自将这忍耐赐予我们；正如另一处所读：\Quote{因为我的忍耐从祂而来。}因为圣经本身的用语就足以驳斥这种解释；因为\Quote{祢是我的忍耐}与\Quote{主，祢是我的希望}以及\Quote{上主我的天主是我的仁慈}和许多类似经文属于同一类。并且没有说\Quote{主，我的爱}或\Quote{祢是我的爱}或\Quote{天主我的爱}；而是这样说：\Quote{天主是爱}，正如说\Quote{天主是神}。凡不明白此事的人，必须向主祈求理解，而非向我们要求解释；因为我们无法说得更清楚了。\par

28. 那么，\Quote{天主是爱}，但问题是：是圣父、圣子、圣神，还是圣三本身？因为圣三不是三位神，而是一位天主。但我在本书上文已经论证过，那作为天主的圣三，不可从我们心中所呈现的那三样来理解，仿佛圣父是三位的记忆，圣子是三位都有的理智，圣神是三位都有的爱；好像圣父自己不领悟也不爱，而是圣子为祂领悟，圣神为祂爱，而祂自己只为祂自己和它们记得；圣子也不为自己记得或爱，而是圣父为祂记得，圣神为祂爱，而祂自己只为祂自己和它们领悟；同样，圣神也不为自己记得或领悟，而是圣父为祂记得，圣子为祂领悟，而祂自己只为祂自己和它们爱。相反，应当如此理解：所有三位以及每一位，都各自按其本性拥有全部三样。而且这些在祂们身上并不像在我们身上那样，记忆是一回事，理智是另一回事，爱或爱德又是另一回事，而是某种等同一切的同一事物，即智慧本身；并且如此包含在每一位的本性中，以致拥有它的那位就是它所是的，作为不变而单纯的实体。如果这一切已被理解，并且就我们能在如此伟大的事物中所见或推测的，已变得显然真实，那么我不知道为何圣父、圣子和圣神不应都称为爱，三者一起称为一个爱，正如圣父、圣子和圣神都被称为智慧，三者一起不是三个智慧，而是一个智慧。因为同样，圣父是天主，圣子是天主，圣神是天主，三者一起是一位天主。\par

29. 然而，在这圣三中，圣子——而非其他——被称为天主的圣言，圣神——而非其他——被称为天主的恩赐，而唯有天主圣父是圣言所由生、并且圣神主要发自的那一位，这并非没有道理。因此我加上了\Emph{主要}一词，因为我们发现圣神也发自圣子。但圣父也赐予圣子这一点，并非赐予一个已经存在而尚未拥有它的那一位；而是凡祂赐予独生圣言的，祂都是藉着生育而赐予的。因此祂如此生育祂，使得共同的恩赐也从祂发出，并且圣神是二者的圣神。因此，不可分离的圣三的这个区分，不仅要粗略地接受，而且要仔细地思考；因为正是由此，天主的圣言也被特别称为天主的智慧，虽然圣父和圣神也是智慧。那么，如果三位中的任何一位要被特别称为爱，还有什么比圣神更合适的呢？——也就是说，在那单纯而至高的本性中，实体不是一回事，爱是另一回事，而是实体本身就是爱，爱本身就是实体，无论是在圣父、圣子还是在圣神中；然而，圣神应被特别称为爱。\par

30. 正如有时全部旧约的言语在圣经中一同以法律之名来表示。因为宗徒在引用\Person{依撒意亚}先知的一段经文说：\Quote{我要以各种舌头和各种嘴唇向这百姓说话}，却先以\Quote{法律上记载了}作为前言。主自己也说：\Quote{在他们的法律上记载了：他们无缘无故地恨我}，而这却是读于圣咏中。有时由\Person{梅瑟}所传授的特别称为法律：正如经上说：\Quote{法律和先知书直到\Person{若翰}为止}；又说：\Quote{这两条诫命是全部法律和先知书所依据的}。在这里，那由\Place{西乃山}所出的特别称为法律。圣咏也以先知之名来表示；然而在另一处，救主亲自说：\Quote{凡关于我的事，都必须应验，即法律、先知书和圣咏上所记载的}。此处，祂却意指先知之名不包括圣咏。因此，法律连同先知书和圣咏总称为法律，而特别所说的法律则是由\Person{梅瑟}所传授的。同样地，先知书连同圣咏通称为先知，而特别说的先知则不包括圣咏。还有许多别的例子可以教导我们，许多事物名称既可普遍使用，也可特别用于特定事物，若不是为了避免在浅显的案例上作冗长的讨论，本可举出更多。我说了这么多，是怕有人认为我们称圣神为爱是不合适的，因为天主父和天主圣子也可以被称为爱。\par

31. 因此，虽然普遍而言圣父本身和圣神也都是智慧，我们却特别以智慧之名称呼天主的独一圣言；同样，虽然普遍而言圣父和圣子也都是爱，我们却特别以爱之名称呼圣神。但是天主的圣言，\Emph{即}天主的独生子，却由宗徒的口中明确地称为天主的智慧和天主的德能，如他说：\Quote{基督是天主的德能，天主的智慧}。至于圣神被称为爱，这可由仔细查考\Person{若望}宗徒的话而发现。他在说过：\Quote{可爱的诸位，我们应当彼此相爱，因为爱出于天主}之后，接着又说：\Quote{凡爱人的，都是由天主而生，也认识天主；那不爱人的，不认识天主，因为天主是爱}。这里他显然把那出于天主的爱称为天主，因此，爱是出自天主的天主。但因为圣子是由天主圣父所生，圣神是由天主圣父所发，故应追问，我们该宁愿认为二者中哪一位是被称为爱的那位天主。因为唯独圣父是这样的一位天主，并非出自天主；所以，那为爱且身为天主而又出自天主的，或是圣子，或是圣神。但当宗徒在下面提及天主对我们的爱时，并不是指我们爱祂的爱，而是指祂\Quote{爱了我们，且打发祂的圣子为我们的罪作赎罪祭}，于是随即也劝勉我们要彼此相爱，好使天主住在我们内——因为他已称天主为爱——他立刻愿意更明确地说这件事：\Quote{我们所以知道我们住在祂内，祂也住在我们内，是因为祂将祂的圣神赐给了我们}。因此，那由祂赐予我们的圣神，使我们住在天主内，天主也住在我们内；这正是爱所成就的。所以，祂就是那为爱的天主。最后，稍后他重复了同样的事，说\Quote{天主是爱}，随即加上：\Quote{那住在爱内的，就住在天主内，天主也住在他内}；在此之前他曾说：\Quote{我们所以知道我们住在祂内，祂也住在我们内，是因为祂将祂的圣神赐给了我们}。因此，在我们读到天主是爱的地方，所指的就是祂。所以，从圣父所发的天主圣神，当祂赐予人时，就激发人爱天主和近人，而祂自己就是爱。因为人没有从自己出发去爱天主的可能，除非出于天主；因此他稍后说：\Quote{我们爱祂，因为祂先爱了我们}。\Person{保禄}宗徒也说：\Quote{天主的爱借着所赐给我们的圣神，已倾注在我们心中了}。\par

% structure: p0055 source=h2 target=subsection reason=relative-heading-rank; base=h2

\subsection{第十八章——天主的恩赐没有比爱更高超的。}

32. 天主的恩赐没有比这更卓越的了。唯独它区分了永恒之国的子女与永恒丧亡的子女。其他恩赐也由圣神赐予；但若没有爱德，它们毫无益处。因此，除非圣神如此赐予每人，使他成为爱天主及爱近人的人，他就不会从左边被移到右边。圣神之所以特别被称为恩赐，也仅仅是因为爱德。没有这爱德的人，\BibleQuote{他若说人及天使的语言，就成了发声的锣，或发响的钹；他若有先知之恩，明白一切奥秘和各种知识，并有全备的信德，甚至能移山，他什么也不算；他若将一切财产施舍给人吃饭，并将身体交付火焚，为他毫无益处。}那么，没有它，那些如此伟大的恩惠也不能使任何人获得永生，这恩惠是多么的伟大啊！但是爱德本身——因为这是同一事物的两个名称——如果一个人拥有它，即使他不说舌音，没有先知之恩，不明白一切奥秘和知识，也没有将他的财物全施舍给穷人——或因为无可施舍，或因为某种需要阻碍——也没有将身体交付火焚，如果没有这种苦难的考验临到他——而将他带入天国，以致信德本身只有借着爱德才能有益，因为没有爱德的信德确实可以存在，却不能获益。因此保禄宗徒也说：\BibleQuote{在基督耶稣内，割损或不割损都毫无效力，唯有藉爱德工作的信德。}将他与甚至\BibleQuote{魔鬼也信，且战栗}的那种信德区分开来。因此，爱德既然出自天主并就是天主，尤其就是圣神，藉着祂，天主的爱德倾注在我们心中，藉着这爱德，整个天主圣三居住在我们内。因此，圣神虽是天主，却极正确地被称作天主的恩赐。而这恩赐，除了那引人归向天主的爱德之外，还能被恰当地理解为别的什么呢？因为没有它，任何其他天主的恩赐都不能引人归向天主。\par

% structure: p0057 source=h2 target=subsection reason=relative-heading-rank; base=h2

\subsection{第十九章.— 圣神在圣经中被称为天主的恩赐。藉着圣神的恩赐，意指那恩赐即是圣神。圣神特别被称为爱德，虽然天主圣三中并非只有圣神是爱德。}

33. 这也是需要证明的吗，即圣神在圣经中被称作天主的恩赐？如果有人也寻求这一点，我们在\BibleBook{若望}{福音}中找到了我们的主耶稣基督的话，祂说：\BibleQuote{谁若渴，到我这里来喝罢！信从我的人，就如经上所说的：从他的腹中要流出活水的江河。}而圣史接着又补充说：\BibleQuote{祂这话是指那信仰祂的人将要领受的圣神说的。}因此，\Person{保禄}宗徒也说：\BibleQuote{我们都饮于同一圣神。}那么，问题是：那水是否就是被称为天主的恩赐的圣神？但是，正如我们在此处发现这水就是圣神，我们同样在福音的其他地方发现这水被称为天主的恩赐。因为当同一主在井旁与撒玛黎雅妇人谈话时，祂曾对她说：\BibleQuote{给我水喝。}她回答说犹太人\BibleQuote{与撒玛黎雅人不来往。}耶稣回答她说：\BibleQuote{若是你知道天主的恩赐，并知道向你说\InnerQuote{给我水喝}的是谁，你或许早求了他，而他也早给了你活水。}那妇人对祂说：\BibleQuote{先生，你连打水器具也没有，而井又深，你从哪里得活水呢？等等？}耶稣回答她说：\BibleQuote{凡喝这水的，还要再渴；但谁若喝了我赐给他的水，他将永远不渴；并且我赐给他的水，将在他内成为涌到永生的水泉。}因此，既然这活水，正如圣史所给我们解释的，就是圣神，那么毫无疑问，圣神就是天主的恩赐，主在这里说：\BibleQuote{若是你知道天主的恩赐，并知道向你说\InnerQuote{给我水喝}的是谁，你或许早求了他，而他也早给了你活水。}因为一处经文所说\BibleQuote{从他的腹中要流出活水的江河}，在另一处则是\BibleQuote{将在他内成为涌到永生的水泉。}\par

34. 宗徒\Person{保禄}也说：\BibleQuote{我们每人蒙受恩宠，是按照基督恩赐的尺度。}然后，为表明他所说的基督恩赐是指圣神，他又接着说：\BibleQuote{因此经上说：\InnerQuote{祂上升至高天，带领了俘虏，施舍了恩赐给人。}}人人都知道，主\Person{耶稣}从死者中复活升天后，赐下了圣神，凡信祂的人都充满了圣神，并说起万国的方言。不要有人反对说，他说的是\Quote{恩赐}（复数），而不是\Quote{恩赐}（单数）；因为他引用了\BibleBook{}{圣咏}的一段话。\BibleBook{}{圣咏}中是这样写的：\BibleQuote{祢上升至高天，祢带领了俘虏，祢在人间接受了恩赐。}在许多抄本中，尤其是在希腊文抄本中，都是这样写的，我们根据希伯来文译出的也是如此。因此，宗徒像先知一样说\Quote{恩赐}（复数），而不是\Quote{恩赐}（单数）。先知说：\BibleQuote{祢在人间接受了恩赐}，宗徒更喜欢说：\BibleQuote{祂施舍了恩赐给人}；这样，从先知和宗徒的双重表述中，可以获得最圆满的意义，因为两者都具有天主圣言的权威。因为两者都是真实的：祂施舍给人，和祂在人内接受。祂施舍给人，如同头施舍给祂的肢体；祂自己施舍，也在人内接受，无疑是在祂自己的肢体内接受；祂因自己的肢体曾从天上呼喊：\BibleQuote{\Person{扫禄}，\Person{扫禄}，你为什么迫害我？}又论到祂的肢体说：\BibleQuote{你们对我这些最小兄弟中的一个所做的，就是对我做的。}因此，基督自己既从天上施舍，也在世上接受。再者，先知和宗徒之所以说\Quote{恩赐}（复数），是因为有许多恩赐，各别地分施给基督的所有肢体，这是通过那恩赐——即圣神——而实现的。因为并非每人都有全部恩赐，而是有些人有这些，有些人有那些；然而，众人都有那恩赐本身，即圣神，藉此才能将属于各人的恩赐分施给他。因为在别处，他列举了许多恩赐后说：\BibleQuote{这一切都是同一且同一圣神所运作的，祂按自己的意愿个别分施。}这句话也见于\BibleBook{}{希伯来书}，其中写道：\BibleQuote{天主也以征兆、奇迹、各种异能和圣神的恩赐来证明。}因此，当他在这里说了\BibleQuote{祂上升至高天，带领了俘虏，施舍了恩赐给人}之后，又接着说：\BibleQuote{但\InnerQuote{祂上升}是什么意思呢？岂不是说祂也曾下降到地下吗？那下降的，正是那上升到诸天之上，以充满万有的那一位。祂赐予这些人作宗徒，那些人作先知，有的作传福音者，有的作司牧和教师。}（我们明白这就是为什么说\Quote{恩赐}的原因；因为正如他在别处说的：\BibleQuote{众人岂都是宗徒？众人岂都是先知？}等等。）接着又说：\BibleQuote{为成全圣徒，以从事服务的工作，以建树基督的身体。}这就是那座\BibleBook{}{圣咏}所歌唱的在被掳之后建起的房屋；因为基督的房屋，即称为祂的教会的，是由那些从魔鬼手中被拯救出来的人建成的，他们原是被魔鬼掳去的。但基督自己亲自带领了这被掳的一群，祂战胜了魔鬼。为使那些将要成为这神圣首领之肢体的人不随魔鬼同受永罚，祂首先以正义的锁链束缚了魔鬼，然后以能力的锁链。因此，魔鬼本身被称为被掳的，那上升至高天、施舍恩赐给人、或在人内接受恩赐的基督，带领了他。\par

35. 宗徒\Person{伯多禄}，正如我们在那记载\BibleBook{}{宗徒大事录}的经典书卷中所读到的——当他讲论基督、犹太人心怀不安时说：\BibleQuote{诸位仁人弟兄，我们该作什么？请告诉我们}——他对他们说：\BibleQuote{你们悔改吧！你们每人要因主\Person{耶稣}基督的名受洗，为得罪之赦；并要领受圣神的恩赐。}我们在同一书中也读到，\Person{西满}术士愿给宗徒们银钱，以获得藉按手赐予圣神的权能。\Person{伯多禄}对他说：\BibleQuote{愿你的银钱与你一同丧亡！因为你以为天主的恩赐可以靠金钱买得。}在该书的另一处，当\Person{伯多禄}向\Person{科尔乃略}和与他在一起的人讲话、宣讲基督时，经上记载：\BibleQuote{\Person{伯多禄}还在讲这些话的时候，圣神降在所有听道的人身上。那些受割损的信徒，凡与\Person{伯多禄}同来的，都惊讶圣神的恩赐也倾注在外邦人身上，因为他们听见他们说各种方言，并颂扬天主。}后来，\Person{伯多禄}在耶路撒冷向弟兄们汇报他给未受割损的人施洗的事，因为圣神——为解决这个问题的结——在他们受洗之前已降临在他们身上，耶路撒冷的弟兄们听后感到不安。\Person{伯多禄}在其余的话之后说：\BibleQuote{我向他们开始讲话时，圣神就降在他们身上，正如当初降在我们身上一样。我就想起了主所说的话：\InnerQuote{\Person{若翰}固然用水施洗，但你们却要因圣神受洗。}所以，既然天主赐给了他们同样的恩赐，如同赐给了我们这些信主\Person{耶稣}基督的人一样，我是谁，能阻止天主把圣神赐给他们呢？}此外还有许多其他圣经的见证，一致证明圣神是天主的恩赐，就祂被赐予那些藉祂爱天主的人而言。但要收集所有见证，任务太繁重了。对于那些不满足于我们所引证的人来说，什么才算足够呢？\par

36. 当然，他们必须被警告，既然他们现在看到圣神被称为\OriginalTerm{天主的恩赐}{gift of God}，那么当他们听到\Quote{圣神的恩赐}时，就应当辨认出其中那种\OriginalTerm{言说方式}{mode of speech}，正如在\Quote{脱去肉体的身体}这句话中所发现的那样。因为正如肉体的身体无非就是肉体，同样，圣神的恩赐也无非就是圣神。因此，就其被赐予那些领受祂的人而言，祂是天主的恩赐。但祂本身是天主，即使祂没有被赐予任何人，因为在祂被赐予任何人之前，祂已与圣父和圣子\OriginalTerm{同为永恒}{co-eternal}。祂也不比圣父和圣子小，因为圣父和圣子赐予，而祂被赐予。因为祂作为天主的恩赐而被赐予，如此方式使得祂本身也作为天主而赐予自己。因为不能说祂不在自己的权能之下，关于祂曾说过：\BibleQuote{风随意向哪里吹}，并且宗徒正如我上面已经提到的说：\BibleQuote{这一切都是这同一位圣神所行的，祂随意分给各人。}我们在此没有那被赐予者的受造，也没有赐予者的管辖，而是被赐予者与赐予者之间的和谐。\par

37. 因此，如果圣经宣告\OriginalTerm{天主是爱}{God is love}，并且\OriginalTerm{爱出自天主}{love is of God}，且在我們身上工作，使我們住在天主内，天主也住在我們内，并且我们由此知道这事，因为祂赐给了我们祂的圣神，那么圣神本身就是天主，祂就是爱。其次，如果在天主的恩赐中没有比爱更大的，并且没有比圣神更大的天主的恩赐，那么还有什么比这更自然的结论：祂本身就是爱，祂被称为天主，又被称为出自天主？如果圣父爱圣子、圣子爱圣父的那种爱，不可言喻地展示了二者的\OriginalTerm{共融}{communion}，那么还有什么比这更合适的：那作为两者共有的圣神应该特别被称为爱？因为这样相信和这样理解更为稳妥：圣神并不是圣三中唯一的爱，但祂并非无故地特别被称为爱，原因我们已经阐述；正如在圣三中，祂也不是唯一的圣神或唯一的圣者，因为圣父是圣神，圣子也是圣神；圣父是圣者，圣子也是圣者——虔诚者对此毫不怀疑。然而，祂特别被称为圣神并非无因；因为祂是两者的共有者，所以特别被称为两者共有的那一位。否则，如果在圣三中只有圣神是爱，那么圣子显然就成为子，不仅出自圣父，也出自圣神。因为祂在无数地方被说成、被读成——天主圣父的\OriginalTerm{独生子}{only-begotten Son}；正如宗徒论及天主圣父的话也是真的：\BibleQuote{祂拯救了我们脱离黑暗的权势，把我们搬迁到祂所爱的子的国度。}祂没有说\BibleQuote{祂自己的子}。如果祂那样说，祂说得最真确，正如祂确实说得最真确，因为祂常这样说；但祂说：\BibleQuote{祂所爱的子。}因此，如果在圣三中天主内没有爱除了圣神，那么圣子也是圣神的子。而如果这极其荒谬，那么结论就是：圣神在那里面并不是唯一的爱，而是因我充分阐述的原因而被特别称为爱；并且\BibleQuote{祂所爱的子}这句话无非就是指祂自己的爱子——简言之，就是祂自己\OriginalTerm{本体的}{substance}子。因为圣父内的爱，在祂那\OriginalTerm{不可言喻的单纯本性}{ineffably simple nature}内，无非就是祂的本性和本体本身——正如我们已屡次说过，而且不怕屡次重复。因此，\BibleQuote{祂所爱的子}不是别的，正是那生于祂本体的那一位。\par

% structure: p0063 source=h2 target=subsection reason=relative-heading-rank; base=h2

\subsection{二十、驳斥\Person{欧诺米}，他声称天主子是子，不是出于祂的本性，而是出于祂的旨意。对以上所述之结语。}

38. 因此，欧诺弥乌斯的逻辑是可笑的，欧诺弥乌斯异端正是从他那里产生的。因为他不能理解，也不愿相信，天主的独生圣言——万物藉祂而受造——是本性上（即由圣父的本体所生）的天主圣子，便声称祂不是出于自己的本性或本体或本质的儿子，而是出于天主意志的儿子；这意思就是要断言，天主生圣子的意志是某种偶然[和可选]的属性——就像我们有时会愿意某件先前不愿意的事那样，仿佛正是由于这些事我们的本性才被看出是可变的——这是万万不可用来相信天主的。因为经上记载：\BibleQuote{人心多有计谋，唯有主的旨意永远立定}，这除了让我们理解或相信，正如天主是永恒的，祂的旨意也是永恒的，因此不可改变，如同祂自己一样。关于心思所说的，也完全可以真实地用于意志：人心中有许多意志，但主的意志永远立定。另一些人，为了避免说独生圣言是天主旨意或意志的儿子，便宣称圣言本身就是圣父的旨意或意志。但在我看来，说旨意之旨意、意志之意志，如本体之本体、智慧之智慧，更好；这样我们就不会陷入我们已经驳斥过的谬论，即说圣子使圣父成为有智慧或有意志的，如果圣父自己的本体中没有旨意或意志的话。确实，有人对异端者给出了一个机智的回答，那个异端者极为狡猾地问道：天主是愿意还是不愿意生圣子？如果回答说不愿意，那么最荒谬地得出天主是痛苦的；但如果回答愿意，他就会立即仿佛用一个不可战胜的理由得出他的目的，即圣子是意志的儿子，而不是本性的儿子。但那位（正统答复者）非常清醒地反过来问他：天主圣父是愿意还是不愿意作为天主？这样，如果他回答不愿意，那么痛苦就会随之而来，而相信天主痛苦是十足的疯狂；如果他说愿意，那么就可以回答他说：那么祂也是凭自己的意志成为天主的，而不是凭本性。那么，剩下的只有让他沉默，并看出他自己已被自己的问题束缚在一个无法解脱的捆绑中。但如果在天主圣三中也有谁可以特别被称为天主的意志，那么这名称，如同爱，更适合于圣神；因为爱不是别的，正是意志。\par

39. 我认为，按照圣经，本书中关于圣神的论述，对于已经知道圣神是天主、并非另有所体、也不比圣父和圣子小的信友们已经足够了——正如我们在前几卷中根据同一圣经所表明的真实的那样。我们也从天主所造的受造物进行了推理，并尽我们所能地警告那些在这样的主题上要求理性的人，要他们藉着所造之物，特别是按天主的肖像所造的理性或理智受造物，去看并理解祂那不可见的事，尽他们所能；透过这面镜子，可以说，他们可以在我们自己的记忆、理解和意志中，尽可能辨明那天主圣三。如果有人明智地认为这三者是天主性安排在他自己的心灵中，并用记忆去记忆，用理解去默观，用爱去拥抱，那么这心灵中那件何等伟大的事，藉此甚至那永恒不变的本性可以被忆起、被看见、被渴望，无疑这人就找到了那至高的天主圣三的肖像。他应当将整个生命归向于记住、看见、爱慕那至高的天主圣三，以便能忆起、默观、并以此为乐。但我也已足够地警告他，不要将这样由那天主圣三所造成、又因自己的过错而变坏的肖像，与那天主圣三本身如此比较，以致认为它在各方面都像那天主圣三，而应当在那种相似中，不管它是什么样的，也看出一种巨大的不相似。\par

% structure: p0066 source=h2 target=subsection reason=relative-heading-rank; base=h2

\subsection{第二十一章——论圣父与圣子被认为存在于我们记忆和理解中的肖像。论圣神在意志或爱中的肖像。}

40. 我无疑已尽我所能地努力，并不是要那事物被面对面看见，而是透过这肖像在隐谜中、在何等微小的程度上以推测，在我们的记忆和理解中暗示天主圣父和天主圣子：即生者的天主，祂以某种方式藉自己的同永恒圣言说出祂本体中所有的一切；以及天主圣言本身，祂在本质上既不多也不少，乃是那位不撒谎而真实地生了圣言的那一位；并且我将我们所知道的一切——即使我们不在思想它——都归于记忆，而将思维按某种特殊模式的形成归于理解。因为我们通常说，当我们通过思想发现某事为真时，我们就理解了它；而这一点正是我们再次留在记忆中的。但那是我们记忆更深隐藏的深处，在那里我们也首次发现了这一点，当我们想到它时，并且在那里生成了一个属于任何语言都不属于的内言——可以说是知识的知识、视觉的视觉，以及在[反思性]思维中出现的理解；这理解确实在记忆中原已存在，但隐藏在那里，尽管除非思维本身也有某种自己的记忆，否则它不会返回那些它在转向思考其他事物时留在记忆中的事物。\par

41. 但我在此关于圣神的谜中所显示的，并无任何似乎与祂相似的事物，唯有我们自己的意志、或爱、或情感，即一种更强烈的意志，因为我们本性所具有的意志，会随着不同对象的接近或出现而受其影响，被其吸引或排斥，从而产生各种不同的感受。那么，这又是什么呢？难道我们要说，我们的意志，当它正直时，会不知道渴望什么、回避什么吗？再者，若它知道，那么它无疑便拥有某种属于它自己的知识，这种知识不能没有记忆和理智。或者，我们要听信任何声称爱不知道它做了什么的人吗？而爱却不会做错事。因此，正如在那种原初的记忆中既存有理智和爱，我们在其中发现并储备着那些我们可以通过思想达至的事物，因为当我们通过思考发现我们既明白又爱某物时，我们也在那里找到这两样东西；而当我们并未思考它们时，它们也早已在那里；也正如在那种通过思想形成的理智中存有记忆和爱，这理智就是我们内在地说的真实的言语，不用任何民族的语言，当我们说出我们所知道的事物时；因为我们的思想的注视除非记起某物，否则不会返回任何事物，除非爱它，否则也不愿返回：同样地，爱——它结合了在记忆中产生的视觉和由之形成的思想的视觉，仿佛父母与后代——若没有对所渴望之物的知识，就不会知道如何正确地爱，而这种知识没有记忆和理智是不可能的。\par

% structure: p0069 source=h2 target=subsection reason=relative-heading-rank; base=h2

\subsection{第二十二章——我们在自身内所发现的天主圣三的肖像，与天主圣三本身之间的差异是何等巨大。}

42. 然而，既然这些存在于一个位格之内，如同人是位格，有人可能会对我们说：这三者——记忆、理智和爱——都是属于我的，而非它们自己的；它们所行的，并非为自身而行，而是为我而行，或者更确切地说，是我借着它们而行。因为是我借着记忆记忆，借着理智理解，借着爱去爱。当我将心灵之眼转向我的记忆，从而在心中说出我所知道的事物，一个真实的言语便由我的知识而生；这两者都是属于我的——知识无疑是属于我的，言语也是。因为是我知道，也是我在心中说出我所知道的事物。当我通过思想，发现我的记忆中存有我对某物的理解和爱（这些理解和爱在我思想之前也已存在）时，我在我自己的记忆中发现的是我自己的理智和我自己的爱，借此是我在理解，是我在爱，而不是那些事物本身在理解和爱。同样地，当我的思想记起并意愿返回那些它已留在记忆中的事物，去理解、观看并内在地说出它们时，那是我的记忆在记起，并且它是自己的，而非它的意志，借着那意志它意愿。就连我的爱本身，当它记起并理解它应当渴望什么、回避什么时，它是借着我的记忆而非它自己的记忆记起；它借着我的理智而非它自己的理智去理解它智慧地爱着的事物。简言之，借着这三者，是我在记忆，是我在理解，是我在爱；而我既不是记忆，也不是理智，也不是爱，而是拥有它们。因此，这些可以由一个拥有这三者、却不是这三者的单独位格说出。但在那至高的单纯本性——即天主——中，虽然只有一个天主，却有三个位格：圣父、圣子和圣神。\par

% structure: p0071 source=h2 target=subsection reason=relative-heading-rank; base=h2

\subsection{第二十三章——\Person{奥古斯定}进一步论述存在于人内的天主圣三与作为天主的圣三之间的差异。如今我们借着信德，通过镜子观看圣三，以便将来在应许的面对面中更清晰地看见。}

43. 因此，一个本身是三一的事物不同于另一个事物中三一的肖像；由于这个肖像，同时那三个东西所在者也称为肖像；就如木板和画在上面的画同时都称为肖像；但由于画在上面的画，木板也被称为肖像。但在那无可比拟地超越万物的至高的天主圣三中，存在着如此伟大的不可分性，以致多个人的三一不能被称作一个人，而在那里，却被称为且实际上是一个天主，并且那天主圣三并非在一个天主之内，而是它就是一个天主。同样，如同在人身上的那个肖像具有这三个东西但只是一个位格，天主圣三却不是这样；而是有三个位格，圣子的圣父、圣父的圣子、以及圣父和圣子双方的圣神。因为人的记忆（尤其是野兽所没有的那种记忆——\Emph{即}：理智所把握的事物并非通过身体感官进入该记忆的记忆）在这天主圣三的肖像中，按其微小的尺度，具有某种肖像相似于圣父，虽不可比拟地不平等，却仍有一些相似；同样，人的理智（它通过思想的目的被形成，当所认识之物被言说时，便产生一个不属于任何舌头的内心之言）在其巨大的差异中具有某种圣子的相似；而人的爱（它从知识出发，结合记忆与理智，仿佛共同属于父母与子女，由此被理解为既非父母亦非子女）在那天主圣三的肖像中具有某种圣神的相似，尽管极其不平等：然而，并非像在那天主圣三的肖像中，这三个东西不是一个人，而是属于一个人，同样在作为这肖像原型的天主圣三本身，这三个东西属于一个天主，但它们是同一个天主，并且是三个位格，而非一个。这确实是一件奇妙得不可言喻、或不可言喻地奇妙的事：虽然这天主圣三的肖像是一个位格，而至高的天主圣三本身是三个位格，但三位格的天主圣三却比这一个的肖像更不可分。因为那天主圣三，在天主性（或更恰当地说神性）的本性中，就是其所是，彼此永恒不变地相等；没有时间它不存在，或它不同于现在；也不会有时间它不存在或它不同于现在。但在那不完全的肖像中的这三个东西，虽然它们在空间上不分离（因为它们不是形体），但在现世生活中却在大小上彼此分离。因为它们中没有各自体积的事实，并不妨碍我们看到在一个人的记忆中理智更大，在另一个人中则相反；又在另一个人中，这两个都被爱的大小所超越；无论这两个彼此是否相等。因此，每一对都被单个超越，每个单个被一对超越，每个单个被另一个单个超越：较小的被较大的超越。当它们从一切软弱中被治愈，并彼此相等时，那因恩宠将不再改变的东西，仍然不能被等同于那因本性不能改变的东西，因为受造物不能被等同于造物主，并且当它从一切软弱中被治愈时，它将改变。\par

44. 但当那应许给我们面对面的视线来临时，我们将看到这天主圣三——不仅是非物质的，而且绝对不可分且真实不变——远比我们现在看到那自身就是其肖像的影像更清晰、更确定；然而，那些透过这镜子并在这谜语中（如同现世生活中允许看见的那样）看见的人，并不是那些在自己心中看到我们已整理并强调的事物的人；而是那些把这肖像当作肖像来看的人，以致他们能够以某种方式将所见者归向那原是肖像者，并通过猜测来看那他们通过注视肖像而看到者（因为他们还不能面对面地看）。因为宗徒并没有说\Quote{我们现在看见镜子}，而是说\Quote{我们现在是透过镜子观看}。\par

% structure: p0074 source=h2 target=subsection reason=relative-heading-rank; base=h2

\subsection{第二十四章——人类心智的软弱}

他们，那些以某种可能的方式看见自己的心智，并在其中看见我以多种方式尽力论述的那天主圣三，却仍不信或不明白这是天主的肖像，他们确实看见了镜子，但还没有透过镜子看见那位现在应透过镜子被看见者，以致他们甚至不知道他们所看见的镜子本身就是一面镜子，\Emph{即}一个肖像。如果他们知道这一点，也许他们会感到，那面镜子所属者应通过它被寻求，并暂且以某种方式被看见，由真诚的信德洁净他们的心，使那位现在透过镜子被看见者能够被面对面看见。如果他们轻视这洁净心灵的信德，那么他们通过理解有关人心智性质的最精微辩论又能成就什么呢？除非他们也因自己理解的见证而被定罪？他们之所以在理解上如此失败，几乎不能得出任何确定之事，岂不是因为他们陷入惩罚性的黑暗，并且被那压迫灵魂的腐朽身体所重负吗？除了罪的惩罚，还有什么过犯使他们遭受这邪恶？因此，被如此巨大邪恶的严重性所警告，他们应当追随那除免世罪的天主羔羊。\par

% structure: p0076 source=h2 target=subsection reason=relative-heading-rank; base=h2

\subsection{第二十五章——圣神为何不是受生，以及他如何从圣父和圣子发出，这一问题只有在享见福乐时才会明白。}

因为凡属于祂的人，即使智力远不如那些人，当他们在此生结束时脱离肉身，嫉妒的权势也无权扣留他们。因为那曾被他们宰杀、本无罪债的羔羊已战胜了他们；但祂首先以血的正义，而非以能力的权势战胜了他们。因此，他们脱离魔鬼的权势后，被圣天使扶持，藉着天主与人之间的中保——基督耶稣其人——摆脱一切邪恶。因为新旧两约圣经（既是预言基督的，也是宣告基督的）一致见证：天下没有赐下别的名字，使人赖以得救。当他们被净化一切腐朽的沾染后，被安置在和平的居所，直到他们再次取回自己的身体——他们自己的、现已不朽的身体——为装饰它们，而非成为负担。因为最良善、最智慧之造物主的旨意是：当人的灵魂虔诚地服从于天主时，应拥有一个幸福地服从的身体，而这幸福应永远持续。\par

45. 在那里，我们将毫无困难地看见真理，并充分享受它，最清楚、最确定。我们不再通过推理的理智探究什么，而是通过默观的理智辨别：为何圣神不是圣子，尽管祂由圣父所发。在那光明中，将没有探究的余地；但在此世，我凭经验本身感到如此困难——无疑，那些仔细而明智地阅读我所写内容的人也会感到同样困难——尽管在第二卷中我承诺将在别处谈论此事，但每当我想要藉着我们自身所是的受造形象来阐明它时，无论我的意思如何，适当的言辞总完全匮乏；甚至在我自己的意思中，我也感到我所达到的是努力而非成就。我确实在一个人（如人）身上找到了那至高圣三的形象，并特别在第九卷中想要藉着那受时间与变化支配的事物来说明并更清楚地展示三个位格的关系。但属于一个位格的三种事物，如人的目标所要求的，并不适用于那三个位格；这一点我们已在第十五卷中证明。\par

% structure: p0079 source=h2 target=subsection reason=relative-heading-rank; base=h2

\subsection{第26章——基督两次赐予圣神。圣神由圣父及圣子所发，超越时间，亦不可称为两者的圣子。}

此外，在那至高的天主圣三（即天主）内，没有时间的间隔，借以可以表明或至少探究：圣子是否先由圣父所生，然后圣神再由两者所发；因为圣经称祂为两者的圣神。正是祂，宗徒说：\Quote{但因为你们是儿子，天主派遣了祂圣子的圣神进入你们心中}；也正是祂，同一圣子说：\Quote{因为说话的不是你们，而是你们父的圣神在你们内说话}。天主圣言的许多其他见证也证明：在天主圣三内特别称为圣神的那一位，是由圣父和圣子所出的；关于祂，圣子自己又说：\Quote{我要从父那里派遣给你们}；在另一处又说：\Quote{父要因我的名派遣}。我们被教导祂由两者所发，因为圣子自己说祂由父所发。当祂从死者中复活，显现给门徒时，\Quote{祂向他们嘘了一口气，并说：你们领受圣神}，以此表明祂也由自己发出。而那正是\WorkTitle{福音}中所读到的\Quote{从祂身上发出的能力}，\Quote{治好了他们众人}。\par

46. 但在祂复活后，祂先在地上赐予圣神，后又从天上派遣祂，依我之见，其原因在于：\Quote{爱德藉此恩赐本身倾注在我们心中}，使我们按这两条诫命爱天主及近人，\Quote{全法律和先知都系于这两条诫命}。\Person{耶稣基督}为表明这一点，一次在地上赐给他们圣神，出于对近人的爱；第二次从天上赐下，出于对天主的爱。即使或许可以为圣神的这一双重恩赐提出其他理由，至少我们不应怀疑：当\Person{耶稣}向他们嘘气时，所赐的正是同一圣神；祂随即说道：\Quote{你们要去使万民成为门徒，因父及子及圣神之名给他们授洗}，在此天主圣三特别被推荐给我们。因此，在五旬节日，即主升天后的第十天，祂也是从天上被赐下的。那么，赐予圣神的那一位，怎会不是天主呢？不，赐予天主的祂何其伟大！因为祂的门徒中没有人赐予圣神，他们只祈求圣神降临在他们覆手的人身上：他们自己并不能赐予。教会如今在祂的治理者身上仍保守这习惯。最后，\Person{西满术士}向宗徒们献钱时，并没有说：\Quote{也赐给我这权能，使我可赐予}圣神；而是说：\Quote{使我按手在谁身上，谁就可以领受圣神。}因为圣经之前并没有说\Quote{西满看见宗徒们赐予圣神}，而是说：\Quote{西满看见宗徒们覆手，圣神就被赐下。}因此，主\Person{耶稣基督}自己不但作为天主赐予圣神，也作为人领受了圣神，所以祂被称为充满恩宠和圣神。在\WorkTitle{宗徒大事录}中，关于祂更明确地写道：\Quote{因为天主以圣神傅了祂。}当然不是用可见的油，而是用恩宠的恩赐，这恩赐由教会为领洗者傅油时所用的可见油膏所象征。当然，\Person{基督}并非在受洗时，当圣神如鸽子降在祂身上时，才被圣神傅油。因为那时祂屈尊预表祂的身体，即祂的教会，在教会中领洗者尤其领受圣神。但祂应被理解为，当天主圣言成为血肉时——即人性毫无任何先前善功的功绩，在童贞\Person{玛利亚}的胎中与天主圣言结合，以至于与之成为同一位格——那时祂已被那奥妙而不可见的傅油所傅。因此我们宣认祂是由圣神和童贞\Person{玛利亚}所生。因为相信祂在将近三十岁时才领受圣神是极其荒谬的：祂在那个年龄受\Person{若翰}的洗；但祂来受洗时既完全没有罪，也非没有圣神。因为如果关于祂的仆人及前驱\Person{若翰}本人，经上写道：\Quote{他还在母胎中就要充满圣神}，因为他虽由父亲所生，但在母腹成形时就领受了圣神；那么关于\Person{基督}的人性，祂的肉身的怀孕本身并非属血肉的，而是属神的，应当如何理解和相信呢？此外，祂的人性与神性两者也在关于祂所记载的这些话中显明：祂从圣父领受了圣神的恩许，并倾注了圣神；因为祂作为人领受，作为天主倾注。我们确实能按自己渺小的程度领受那恩赐，但绝不能将其倾注给别人；然而为做到这事，我们向天主呼求临于他们身上，此事是由天主完成的。\par

47. 那么，我们能够问：当圣子诞生时，圣神是否已经从圣父发出，还是尚未发出？而当圣子诞生时，圣神从二者发出，其中没有时间上的先后？就像我们在有时间的情形中所能问的那样：意志首先从人的心灵发出，以便寻求那当被找到时便可称为后裔的东西；那后裔一旦产生或诞生，那意志便达致圆满，安息于这终点，以致那在寻求时的欲望，在享受时便成为爱；这爱现在从二者发出，\Emph{即}从生育的心灵和被生的概念发出，如同从父母和子嗣发出。在这种情况下，绝对不可能提出这些问题，因为这里没有任何事是在时间中开始，继而在后续时间中完成的。所以，凡能理解圣子自圣父无时间地受生的人，也当理解圣神自二者无时间地发出。凡能理解圣子所说的：\BibleQuote{父怎样在自己内有生命，也赐给子在自己内有生命}——并非父将生命赐给早已存在却无生命的子，而是祂这样无时间地生了祂，以致父在生祂时所赐予的生命，与赐予生命的父的生命是永恒同等的——我说，凡能这样理解的人，也当理解：正如父在自己内有使圣神从祂发出的能力，祂也赐给子使同一圣神从祂发出，二者皆无时间；并且圣神被称为从父发出，同时要理解祂也从子发出，这是子从父所领受的特性。因为凡子从父所有的，当然祂从父也有圣神从祂发出。但不要让任何人认为其中有时间上的先后；因为这些事物根本不存在于时间中。那么，若称之为二者的子，岂非极其荒谬？因为正如从父受生，毫无本性变化，使子获得无时间起点的本质；同样，从二者发出，毫无本性变化，使圣神获得无时间起点的本质。我们虽然不说圣神是受生的，但我们也不敢因此说祂是非受生的，免得有人从这个词中怀疑在那天主圣三中有两位父，或者有两位不是从另一位而来的。因为只有圣父不是从另一位而来的，因此只有祂被称为非受生的——这并非在圣经中，而是在辩论者的惯用语中，他们对如此伟大的主题尽可能使用这样的语言。圣子由圣父受生；圣神主要从圣父发出，圣父毫无时间间隔地赐予这发出，但也是从二者[圣父和圣子]共同发出。然而，如果——这是所有健全心智所厌恶的——二者都\Emph{生了}祂，那么祂就会被称为圣父和圣子的子。因此，二者的圣神不是受生于二者，而是从二者发出。\par

% structure: p0083 source=h2 target=subsection reason=relative-heading-rank; base=h2

\subsection{第二十七章 —— 什么足以解决为何圣神不被说成受生，以及为何只有圣父是非受生的问题。那些不理解这些事的人应当做什么。}

48. 但因为在那永恒同等、无形、不可言喻、不可改变且不可分的天主圣三中，区分受生与发出极其困难，让我们暂时满足于将我们关于这个主题在一次讲道中向基督徒群众所说的话——说完后并写下——呈现给那些无法进一步深入的人。因为我藉着圣经的见证教导他们圣神从二者发出时，我继续说：\Quote{若圣神既从圣父又从圣子发出，那么为何圣子说：\BibleQuote{祂从父发出}？为什么，你以为呢？除非祂习惯于将祂自己的——即祂从祂所领受的——归于那一位，而祂自己也是从那一位而来的？由此，祂所说的：\BibleQuote{我的教训不是我的，而是那派遣我者的}又当如何理解？如果这里所理解的教训是祂的，而祂却说不是祂的，而是派遣祂者的，那么更当理解为圣神也从祂发出，因为祂说祂从父发出，而非说祂不从祂发出？当然，圣子从父那里得到祂的神性，因为祂是出自天主的 天主，祂也使圣神从祂发出；因此圣神从圣父本身领受了也应从圣子发出，正如祂从父发出一样。这里，也在某种程度上可以理解——就我们这样的人所能理解而言——为何圣神不说为生，而说为发出；因为如果祂也被称为子，那么祂当然会被称为二者的子，这是极其荒谬的，因为没有人是两个人的子，除非是父和母。但让我们远避在圣父天主和圣子天主之间有任何这样的事。因为就连人的儿子也不是同时从父母两者发出的；当他从父亲进入母亲时，他并非同时从母亲发出；当他从母亲发出进入这个世界的光时，他并非同时从父亲发出。但圣神不是从圣父进入圣子，再从圣子发出以圣化受造物，而是同时从二者发出；虽然圣父已赐予圣子，使圣神如同从祂自己发出一样，也从圣子发出。我们不能说圣神不是生命，而圣父是生命，圣子也是生命：因此，正如圣父在自己内有生命，也赐给圣子在自己内有生命；同样，祂也赐给祂生命从祂发出，如同从祂自己发出一样。} 我将此从那次讲道中移入本书，但我是对信徒讲的，而非对不信者。\par

49. 但如果他们不能注视这肖像，不能看出这些存在于他们心智中的事物是何等真实，然而这些事物并非三个到成为三个位格，而是全都属于一个作为一个人位格的人；那么，他们为何不相信在圣经中发现的、关于那至高的天主圣三的教导，反而坚持要求得到最清晰的理性证明，而这证明因人性的迟钝与软弱而无法被理解呢？当然，当他们坚定地相信圣经是最真实可靠的见证时，就让他们借着祈祷、寻求和善度生活而努力去理解，\Emph{即}，尽可能地看见那由信德所持守的事物。谁能禁止这事？谁不反而劝勉他们去做呢？但如果他们认为，因为他们盲目的心智无法分辨这些事物，就应当否认它们的存在，那么生来盲目的人也应当否认太阳的存在。光在黑暗中照耀；但如果黑暗不领悟它，就让他们先借着天主的恩赐被光照，成为信徒，并开始相对于不信者成为光；在这根基奠定之后，让他们被建立起来，去看见他们所相信的，以便有朝一日能够看见。因为有些事物是如此被相信，以至于完全无法被看见。基督不会第二次被看见在十字架上；但除非相信这已经完成并见过的事，以致不再期望将来发生和看见，否则就无法来到基督那里，如同无终地看见祂一样。然而，就借着理智以某种方式分辨那至高、不可言喻、无形体、不变的本性而言，人的心智之眼没有比在人的本性中——比其它动物更好的本性，也比人灵魂的其他部分——即心智本身更好的操练场所了，只要信德的规则引导它；心智被赋予了某种对无形事物的视力，并在一个更高更内在的位置上以荣誉的身份主持，身体的感官也将一切事物报告给它，以供判断；并且除了天主之外，没有比他更高的存在需要它服从，并由祂统治。\par

50. 但在现在所说的许多事中，没有一样我敢宣称自己说了配得上那至高圣三之不可言喻性的话，反而承认对祂的奇妙知识对我过于玄奥，我无法达到。哦，我的灵魂，你觉得自己在哪里？你卧在哪里？你站在哪里？直到那位赦免你一切罪愆者治愈你的一切软弱。你确知自己身处那家客店，就是撒玛黎雅人将那个被强盗打得半死的人带去的地方。然而你已经看见了许多真实的事，不是通过看见有色物体的眼睛，而是通过那祈祷者所求的眼睛，他说：\Quote{愿我的眼睛看见公平的事。}当然，你已看见了许多真实的事，并且借着那光你看见了它们，并将它们与那光区别开来。抬起你的眼睛注视那光本身，如果你能的话，就定睛在它上面。因为这样你将看到天主圣言的诞生与天主恩赐的出发有何不同，为此独生子没有说圣神是由圣父所生（否则祂就是祂的弟兄了），而是说祂由圣父所发。因此，既然两者的圣神是圣父与圣子某种同体的共融，祂并不被称为——我们决不可如此说——两者的儿子。但你无法将视线定在那里，以便清晰明白地分辨这事；我知道你不能。我说的是实话，我对自己说，我知道我不能做什么；然而那光本身向你展现了在你自己身上的这三样东西，在其中你可以认出那至高圣三的肖像，虽然你还不能以稳定的目光注视它。它向你展示：在你内有一个真实的话语，当它由你的知识所生时，\Emph{即}当我们说出我们所知道的事时；虽然我们既不发出也不思想任何具有各国语言意义的有声话语，但我们的思想是由我们所知的事物所形成；在思想者的心智之眼中，有一个与记忆所包含的思想相似的影像，而意志或爱作为第三位，将这两者如同父母与后代结合起来。凡能看见并分辨的人，就会看出这意志确实由思想发出（因为没有人会意愿他完全不知道是什么或是什么样的事物），但并非思想的影像；因此，在这可理解的事物中暗示了出生与出发的某种区别，因为通过思想来看，与欲求甚至以意志享受并不相同。你也能分辨这一点，虽然你尚未且不能以充分的言语表达你在身体相似物的云雾（它们不断在人类思想前后飘荡）中几乎看不到的东西。但那不是你的光也向你表明，这些无形体的身体相似物与真理不同，真理是通过摒弃它们而以理智沉思的。这些以及其他同样确定的事，那光已向你内在的眼睛显示了。那么，为何你不能以稳定的目光看见那光本身呢？除非是由于软弱。这软弱又是什么造成的呢？除非是由于不义。那么，谁治愈你的一切软弱呢？难道不是那赦免你一切不义者吗？因此，我现在将终于以祈祷而非辩论结束本书。\par

% structure: p0087 source=h2 target=subsection reason=relative-heading-rank; base=h2

\subsection{第二十八章——以祈祷结束本书，并为言辞繁多致歉。}

51. 上主、我们的天主，我们相信祢，圣父、圣子和圣神。因为真理不会说：\BibleQuote{去，因父、及子、及圣神之名给万民授洗}，除非祢是天主圣三。也不会，上主天主，命令我们因不是上主天主者的名受洗。天主的声音也不会说：\BibleQuote{以色列，请听：上主你的天主是唯一的天主}，除非祢是天主圣三，又是唯一的主天主。如果祢，天主，自己是父，自己是子，祢的圣言耶稣基督，以及圣神祢的恩赐，我们就不会在真理之书上读到：\BibleQuote{天主派遣了祂的圣子}；也不会，独生子，论到圣神说：\BibleQuote{父因我的名所要派遣的}；以及：\BibleQuote{我从父那里要派遣给你们的}。我以此信德准则指引我的目标，尽我所能，照祢使我所能，我寻求了祢，并渴望以理智看到我所相信的；我辩论辛劳颇多。上主我的天主，我唯一的希望，求祢垂听，免得我因疲惫而不愿寻求祢，\BibleQuote{但愿我常常热切寻求祢的面}。求祢赐予寻求的力量，因为是祢使我找到了祢，并赐予我更进一步找到祢的希望。我的力量和软弱都在祢眼前：求祢保全前者，治愈后者。我的知识和愚昧都在祢眼前：在祢为我开启之处，请接纳我进入；在祢关闭之处，请为我叩开。愿我记得祢，理解祢，爱慕祢。愿祢在我内增长这一切，直到祢完全更新我。我知道经上记载：\BibleQuote{多言多语难免犯罪}。但愿我讲话只为了宣讲祢的圣言并赞美祢！这样我不但能逃避犯罪，而且无论讲多少都能获得功绩。因为蒙祢祝福的人不会命令他信仰中的真子犯罪，他写信给他说：\BibleQuote{宣讲圣言：不论顺逆，务要坚持}。我们能说他讲了许多吗？他不只适时，而且连不适时都宣讲祢的圣言，因此并不算多，因为只讲了必要的。天主，求祢救我脱离那种在我灵魂内遭受的多言多语，它在祢眼中是可怜的，我逃到祢的仁慈庇护下；因为即使我沉默无言时，内心也在喋喋不休。如果我真的只想念祢所喜悦的事，我当然不会求祢救我脱离这种多言多语。但我的思绪很多，正如祢所知的，\BibleQuote{人的思念都是虚幻的}。求祢赐我不去赞同它们；即使它们有时令我喜悦，也要使我谴责它们，而不留在它们中，如同沉睡一般。也不让它们在我内得势，以致我的行为出自它们；但至少让我的意见、我的良心在祢的保护下免受它们侵扰。当那位智者在他那本现称\BibleBook{}{德训篇}的书中论到祢说：\BibleQuote{我们说了很多，仍然词不达意；总而言之，他就是一切}。因此，当我们来到祢面前时，这些我们所说的许多话（仍词不达意）都会止息；而祢，作为唯一者，将常作\BibleQuote{万有中的万有}。我们将无终地在合一中赞美祢，说同一句话，而我们也将在祢内成为一体。上主、唯一的天主，天主圣三，凡我在这些书中所说的属于祢的，愿属于祢的人承认；若有什么出于我自己的，愿祢和属于祢的人宽恕。阿们。\par

\SourceRecord{Translated by Arthur West Haddan.；Nicene and Post-Nicene Fathers, First Series；Vol. 3.；Edited by Philip Schaff.；Buffalo, NY: Christian Literature Publishing Co.,；1887.；<http://www.newadvent.org/fathers/130115.htm>.}
